\documentclass{report}

%%% Работа с языком
%\usepackage{cmap}					    % поиск в PDF				    % русские буквы в формулах
\usepackage[T2A]{fontenc}			% кодировка
\usepackage[utf8]{inputenc}             % кодировка исходного текста
\usepackage[english,ukrainian]{babel}	% локализация и переносы	    
\usepackage{indentfirst}

\usepackage[a4paper, top=25mm, bottom=25mm, left=30mm, right=30mm]{geometry}

\usepackage{amsmath,amsfonts,amssymb,amsthm,mathtools} % AMS

\usepackage{makecell}% Прикольная настройка ячеек в tabular 
%http://ctan.math.utah.edu/ctan/tex-archive/macros/latex/contrib/makecell/makecell-rus.pdf

\usepackage{diagbox}
% Разделение клеток tabular по диагонали

\usepackage{multicol} % Несколько колонок
\usepackage{wrapfig} % Картинки посреди текста
\usepackage{chngcntr} % для нумерации формул
\usepackage{enumitem} % чтоб можно было делать римскую нумерацию
\usepackage{mdwlist}
\usepackage{adjustbox}
\usepackage[psdextra]{hyperref}
\hypersetup{unicode=true}
\hypersetup{
    colorlinks=true,
    linkcolor=blue,
}
\setlist[enumerate]{nosep}

%%%% Теоремы, определения и т.д. и т.п.
\theoremstyle{plain} % Это стиль по умолчанию, его можно не переопределять.
\newtheorem{theorem}{Теорема}
\newtheorem{proposition}{Твердження}
\renewcommand{\qedsymbol}{$\blacktriangle$}
 
\theoremstyle{definition} % "Определение"
\newtheorem{definition}{Означення}[section]
\newtheorem*{definition*}{Означення}
\newtheorem*{example}{Приклад}

\theoremstyle{remark}
\newtheorem*{remark}{Зауваження}
\newtheorem*{exercise}{Вправа}

\counterwithout{equation}{chapter} % нумерация формул
\counterwithin*{equation}{section}

%%%% Картинки
\usepackage{graphicx}
\usepackage{tikz}
\usepackage{pgfplots}
\pgfplotsset{compat=1.17}

\usetikzlibrary{patterns}
\usetikzlibrary{shapes.geometric}
\usetikzlibrary{arrows.meta}
\usetikzlibrary{shapes.geometric}
\graphicspath{{pictures/}}
\DeclareGraphicsExtensions{.pdf,.png,.jpg}


\author{Каніовська І.Ю.}
\title{Конспект лекцій з теорії ймовірностей}
\date{}

\begin{document} 
    \maketitle
    \tableofcontents
    \chapter{Випадкові події}
        % !TEX root = ../main.tex

\section{Основні поняття теорії ймовірностей}
\subsection{Теоретико-множинний підхід до основних понять ТЙ}
\begin{definition}
    \emph{Стохастичним експериментом (випробуванням)} називається експеримент, 
    який можна повторювати неодноразово, зберігаючи певні умови, і результат якого 
    експерименту заздалегідь передбачити неможливо.
\end{definition}
\begin{definition}
    Будь-який результат СЕ називається \emph{подією}.
\end{definition}
\begin{example}
    \begin{enumerate}
        \item СЕ --- кидання кубика один раз, подія --- випало 6 очок.
        \item СЕ --- кидання кубика двічі, подія --- сума очок, що випала, дорівнює 6.
    \end{enumerate}
\end{example}

\noindent \textbf{Події бувають:}
\begin{enumerate}
    \item Випадкові --- можуть відбутися чи не відбутися при проведенні СЕ.
    \item Неможливі --- ніколи не відбуваються в даному СЕ.
    \item Вірогідні --- завжди відбуваються в даному СЕ.
\end{enumerate}

Будемо вважати, що кожному СЕ можна поставити у відповідність деяку множину, що
називається \emph{простором елементарних подій} $\Omega$. 
Під \emph{елементарними подіями} $\omega$ будемо розуміти єдині
логічно можливі події СЕ, що виключають одна одну.
\begin{example}
    \begin{enumerate}
        \item СЕ --- кидання кубика один раз. 
        $\Omega = \left\{\omega_1, \omega_2, ..., \omega_6\right\}$, 
        де події 

        $\omega_k = \left\{\text{випало } k \text{ очок}\right\}, k = 1,...,6$.
        \item СЕ --- кидання монетки до першої появи герба.
        $\Omega = \left\{\omega_1, \omega_2, ..., \omega_k, ...\right\}$, 
        
        де $\omega_i = \left\{\text{герб випав на }i\text{-тому киданні}\right\}, i\in \mathbb{N}$.
        \item СЕ --- зустріч двох осіб, що домовилися зустрітися протягом години.
        $x$ --- час приходу першої особи, $y$ --- час приходу другої, $0\leq x, y \leq 1$.
        $\Omega = \left\{ \left( x, y\right): 0\leq x, y \leq 1\right\}\subset \mathbb{R}^2$.
    \end{enumerate}
\end{example}
В подальшому випадкові події позначатимемо $A, B, C, ...$. 
\emph{Випадкова подія --- підмножина $\Omega$}. 
У прикладі з киданням кубика один раз $A = \left\{\text{випало }6\text{ очок}\right\} = \left\{ \omega_6\right\}$.
В загальному випадку маємо $A = \left\{\omega_{k1}, \omega_{k2}, ..., \omega_{kn}, ...\right\} \subset \Omega$.
Неможлива подія --- $\varnothing$, вірогідна --- $\Omega$. 

\begin{remark}
    Розглядаємо події в межах фіксованих СЕ та простору елементарних подій.
\end{remark}
\begin{enumerate}
    \item \emph{Включення} $A \subset B$ означає, що якщо відбулася подія $A$, то обов'язково відбудеться подія $B$.
    Наприклад, $A = \left\{\text{витягнуто даму пік}\right\}, B = \left\{\text{витягнуто карту чорної масті}\right\}$. 
    Очевидно, $A\subset B$.

    Рівність подій: $\left( A \subset B, B \subset A \right) \iff \left( A = B\right)$.
    \item \emph{Об'єднання подій} $A \cup B$ --- це подія, яка відбувається тоді, коли відбувається
    принаймні одна з подій $A$ чи $B$.

    Властивості: $A \cup A = A$,  $A \cup B = B \cup A$, $\left( A \subset B \right) \Rightarrow \left( A \cup B = B \right)$, 
    $A \cup \Omega = \Omega$, $A \cup \varnothing = A$, $\left( A \cup B \right) \cup C = A \cup \left( B \cup C \right)$.

    Операція узагальнюється на скінченну або зліченну кількість подій: $A = \bigcup\limits_{i=1}^{n \left( \infty \right)} A_i$.
    \item \emph{Перетин подій} $A \cap B$ --- це подія, яка відбувається тоді, коли $A$ і $B$ відбуваються одночасно.

    Властивості: $A \cap A = A$,  $A \cap B = B \cap A$, $\left( A \subset B \right) \Rightarrow \left( A \cap B = A \right)$, 
    $A \cap \Omega = A$, $A \cap \varnothing = \varnothing$, $\left( A \cap B \right) \cap C = A \cap \left( B \cap C \right)$,
    $A \cap B \subset A$, $A \cap B \subset B$, $\left( A \cup B \right) \cap C = \left( A \cap C \right) \cup \left( B \cap C \right)$,
    $\left( A \cap B \right) \cup C = \left( A \cup C \right) \cap \left( B \cup C \right)$.

    Операція узагальнюється на скінченну або зліченну кількість подій: $A = \bigcap\limits_{i=1}^{n \left( \infty \right)} A_i$.
\end{enumerate}
\begin{definition}
    Події $A$ та $B$ називається \emph{несумісними}, якщо вони не відбуваються одночасно: $A \cap B = \varnothing$.
    Узагальнення: події $A_1, A_2, ..., A_n, ...$ називаються \emph{попарно несумісними}, якщо $A_i \cap A_j = \varnothing$ для $i \neq j$.
\end{definition}
\begin{definition}
    Події $A_1, A_2, ..., A_n, ...$ утворюють \emph{повну групу подій}, якщо вони попарно несумісні 
    та $\bigcup\limits_{i=1}^{n \left( \infty \right)} A_i = \Omega$.
\end{definition}
\begin{enumerate}
    \setcounter{enumi}{3}
    \item \emph{Протилежна подія} $\overline{A}$ --- це подія, яка відбувається тоді, коли $A$ не відбувається.
    
    Властивості: $A \cup \overline{A} = \Omega$, $A \cap \overline{A} = \varnothing$, $\overline{\left( A \cup B \right)} = \overline{A} \cap \overline{B}$,
    $\overline{\left( A \cap B \right)} = \overline{A} \cup \overline{B}$.
    \item \emph{Різниця подій} $A \backslash B$ --- це подія $A \cap \overline{B}$. Для протилежної події маємо $\overline{A} =  \Omega \backslash A$.
\end{enumerate}

\begin{definition}
    Непорожня система підмножин $\mathcal{F}$ простору елементарних подій $\Omega$ утворює \emph{алгебру подій}, якщо:
    \begin{enumerate}
        \item $\Omega \in \mathcal{F}$;
        \item $\left( A, B \in \mathcal{F}\right) \Rightarrow \left( A \cup B \in \mathcal{F}\right)$;
        \item $\left( A \in \mathcal{F}\right) \Rightarrow \left( \overline{A} \in \mathcal{F}\right)$.
    \end{enumerate}
    \emph{Узагальнення:} якщо $\Omega$ містить зліченну кількість подій, то означення $\sigma$-алгебра отримаємо заміною умови
    2 на $\left(\forall n \in \mathbb{N}: A_n \in \mathcal{F} \right) \Rightarrow \left( \bigcup\limits_{n=1}^{\infty} A_n \in \mathcal{F}\right)$.
    
    Пара $\left\{\Omega, \mathcal{F}\right\}$ називається \emph{вимірним простором стохастичного експерименту}.
\end{definition}

\subsection{Класична модель ймовірності}
Дослідника завжди цікавить кількісна характеристика появи тієї чи іншої події.

Нехай $\Omega$ --- скінченний чи зліченний. 
Поставимо у відповідність кожній елементарній події $\omega_k$ число $p_k\geq 0$ так, що $\sum\limits_{k=0}^{ n \left( \infty\right)}p_k = 1$.
Тоді $\P(A) = \sum\limits_{\omega_k \in A} p_k$ --- кількісна характеристика, ймовірність події $A$.

\begin{example}
    \emph{Класична модель ймовірності}. Якщо простір елементарних подій $\Omega$ СЕ скінченний
    та всі $\omega_k$ рівноможливі, то такий СЕ називається \emph{класичним}.
    В такому випадку $p_1 = p_2 = ... = p_n = \frac{1}{n}$, де $n = \mathrm{card}(\Omega)$.
    
    \begin{equation}\label{eq:m_n}
        \P(A) = \sum\limits_{\omega_k \in A} p_k = \frac{m}{n} = \frac{\text{кількість елементарних подій в } A}{\text{загальна кількість елементарних подій}}
    \end{equation}
\end{example}
Ймовірності, що розраховуються за формулою \eqref{eq:m_n}, називаються \emph{класичними}.
\begin{example}
    \begin{enumerate}
        \item <<Задачі про вибір>> --- коли з великої кількості чогось вибирають певну кількість.
        Наприклад, з урни з 10 кульками, 3 чорними та 7 білими, навмання витягають 5 кульок.
        Обчислимо ймовірність події $A = \left\{ \text{серед них 2 чорних кульки}\right\}$:
        $\P(A) = \frac{C_3^2\cdot C_7^3}{C_{10}^5} = \frac{3\cdot 35}{252} = \frac{5}{12}$.
        \item <<Задачі про ліфт>>. 5 осіб одночасно зайшли в ліфт 11-поверхового будинку. 
        Яка ймовірність того, що вони всі вийдуть на різних поверхах, починаючи з другого?
        
        $\P(A) = \frac{A_{10}^5}{\widetilde{A}_{10}^5} = \frac{10\cdot 9\cdot 8\cdot 7 \cdot 6}{10^5} = 0.3024$. 
        Тут $A_n^k$ та $\widetilde{A}_n^k$ --- кількості розміщень без повторень та з повтореннями відповідно.
    \end{enumerate}
\end{example}

\noindent \textbf{Властивості класичної ймовірності:}
\begin{enumerate}
    \item $\forall A \in \mathcal{F}: 0 \leq \P(A) \leq 1$.
    \item $\P(\Omega) = 1$.
    \item $\P(\overline{A}) = 1 - \P(A)$.
    \item $\P(\varnothing) = 0$.
    \item $\P(A\cup B) = \P(A) + \P(B) - \P(A\cap B)$, для несумісних $A$ і $B$ $\P(A\cup B) = \P(A) + \P(B)$.
    \item $\left( A \subset B\right) \Rightarrow \left( \P(A) \leq \P(B)\right)$.
    \item Якщо $A_1, A_2, ..., A_n$ --- повна група подій СЕ, то $\P\left(\bigcup\limits_{i=1}^{n} A_i\right) = 1$.
\end{enumerate}

\subsection{Геометрична модель ймовірності}
\begin{example}
    Нехай точка кидається навмання на відрізок $\left[a; b\right]$. 
    Яка ймовірність її 
    потрапляння в $\left<\alpha; \beta\right> \subset  \left[a; b\right]$?
    Розглянемо подію $A = \left\{ 
        \text{точка потрапила в} \left<\alpha; \beta\right>
    \right\}$.
    \newline
    \hbox to \hsize{\hfil{
        \begin{tikzpicture}
            \draw [fill] (0, 0) circle [radius = 0.05];
            \node [below] at (0, 0) {$a$};
            \node [above] at (0, 0) { };
            \node [below] at (8, 0) {$b$};
            \draw [-{Straight Barb}] (6,0) to (2,0) {};
            \draw [-{Straight Barb}] (2,0) to (6,0) {};
            \node [below] at (2, 0) {$\alpha$};
            \node [below] at (6, 0) {$\beta$};
            \draw [fill] (8, 0) circle [radius = 0.05];
            \draw [thick] (0, 0) -- (8, 0);
        \end{tikzpicture}
    }\hfil}
    Логічно припустити, що $\P(A) = k\cdot l_{\left<\alpha; \beta\right>}$ для деякого $k > 0$.
    Тоді $\P(\Omega) = 1 = k\cdot l_{\left[a; b\right]}$. Таким чином, 
    $k = \frac{1}{l_{\left[a; b\right]}} = \frac{1}{b-a}$,
    тому $\P(A) = \frac{l_{\left<\alpha; \beta\right>}}{l_{\left[a; b\right]}}$.
\end{example}
Нехай простір елементарних подій інтерпретується як замкнена область в 
$ \mathbb{R} ^n$, а випадкові події --- її підмножини. В якості $\sigma$-алгебри 
подій $\mathcal{F}$ беремо підмножини, що мають міру $\mu$. Тоді в якості ймовірності 
деякої події $A$ беремо $\P(A) = \frac{\mu(A)}{\mu(\Omega)}$. 
Наприклад, в $\mathbb{R}^1$ беремо міру <<довжина>>, в $\mathbb{R}^2$ --- <<площа>>, а в $\mathbb{R}^3$ --- <<об'єм>>.

Ймовірності, що знаходяться таким чином називаються \emph{геометричними}, а сама модель 
називається \emph{геометричною моделлю ймовірності}.
Геометрична модель може використовуватись, 
коли $\Omega$ має геометричну інтерпретацію як замкнена область в $\mathbb{R}^n$,
а елементарні події --- рівноможливі.

\begin{example}[задача Бюффона]\label{buffon}
    На площині накреслені паралельні прямі на відстані $2a$, на них кидається 
    голка довжиною $2l,\; l \leq a $. Яка ймовірність того, що голка перетне 
    яку-небудь пряму?
    \newline
    \hbox to \hsize{\hfil{
        \begin{tikzpicture}[scale = 0.5]
            \draw [thick] (0, 0) -- (8, 0);
            \draw [thick] (0, 2) -- (8, 2);
            \draw [thick] (0, 4) -- (8, 4);
            \draw [thick] (0, 6) -- (8, 6);
            \draw [{Straight Barb}-{Straight Barb}] (7.5,4) to (7.5,6);
            \node [right] at (7.5, 5) {$2a$}; 
        \end{tikzpicture}
        \begin{tikzpicture}[scale = 1]
            \draw (0, 0) -- (4, 0);
            \draw (0, 2) -- (4, 2);
            \draw [thick] (0.5, -0.5)-- (2.5, 1.5);
            \draw [{Straight Barb}-{Straight Barb}] (3.5,0) to (3.5,2);
            \draw [{Straight Barb}-{Straight Barb}] (1.5,0) to (1.5,0.5);
            \node [right] at (3.5, 1) {$2a$};
            \node [right] at (1.5, 0.25) {$ x $};
            \draw [-{Straight Barb}](2.5, 0) to [out = 90, in = 0](1.75, 0.75);
            \node [right] at (2.28, 0.53) {$ \varphi $};
        \end{tikzpicture}
    }\hfil}
    Достатньо розглядати лише дві прямі. $\Omega = \left\{(\varphi, x)\in 
    \mathbb{R}^2: \varphi \in \left[0; \pi\right], x \in \left[0; a\right] 
    \right\}$
    \newline
    При такій побудові простору елементарних подій подія $A = \left\{\text{голка 
    перетне пряму}\right\} =$
    \newline
    $\left\{(\varphi, x)\in 
    \Omega: x\leq l\sin\varphi\right\}$
    \newline
    \hbox to \hsize{\hfil{
        \begin{tikzpicture}
            \draw [<->] (4, 0) -- (0, 0) -- (0, 2.5);
            \node [left] at (0, 2.5) {$x$};
            \node [below] at (4, 0) {$\varphi$};
            \draw [domain=0:pi,fill=gray] plot (\x, {1*sin(\x r)});
            \node [below] at (0, 0) {$0$};
            \node [below] at (pi, 0) {$\pi$};
            \draw [dashed] (0,1) to (pi/2,1);
            \node [left] at (0, 1) {$l$};
            \draw (0,2) -- (pi, 2) -- (pi, 0);
            \node [left] at (0, 2) {$a$};
            \node [below left] at (pi, 2) {$\Omega$};
            \node [above] at (pi/2, 0.25) {$A$};
        \end{tikzpicture}
    }\hfil}
    $\P(A) = \frac{S(A)}{S(\Omega)},\; S(\Omega) = \pi a,\;S(A) = \int\limits_0^\pi 
    l\sin\varphi\ d\varphi = l\left.(-\cos\varphi)\right|^\pi_0 = 2l$.
    Отже, $\P(A) = \frac{2l}{\pi a}$.
    \newline
    Якщо провести $n$ кидань голки, у $m$ з яких голка потрапить на пряму, то за допомогою приблизної рівності $\frac{m}{n} 
    \approx \frac{2l}{\pi a}$ можна приблизно обчислити число $\pi$. 
    Наприклад, якщо взяти $n=5000, m=2532, l/a = 0.8$, то отримаємо $\pi \approx 3.1596$.
\end{example}
\begin{example}[парадокс Бертрана]
    Нехай в колі радіусом $R$ навмання обирається хорда. Яка ймовірність того, 
    що її довжина буде більшою за довжину сторони правильного трикутника, 
    вписаного в це коло?
    \newline
    Існують три способи розв'язання цієї задачі, причому кожен з них дає різний результат.
    
    \textbf{Спосіб 1.}
    З міркувань симетрії обирається якийсь фіксований діаметр кола і розглядається
    всі перпендикулярні до нього хорди. Серед них і обирається навмання довільна хорда.
    Очевидно, що кожна хорда у цьому випадку може бути однозначно визначена своєю серединою,
    тобто кожній хорді можна взаємно однозначно поставити у відповідність координату її середини,
    якщо за початок відліку взяти якийсь з кінців фіксованого діаметру.
    \newline
    \hbox to \hsize{\hfil{
        \begin{tikzpicture}[scale = 1.5]
            \draw (0, 0) circle [radius = 1]; 
            \draw (0, 1) -- (-0.86602540378, -0.5); 
            \draw (0, 1) -- (0.86602540378, -0.5);
            \draw (0.86602540378, -0.5) -- (-0.86602540378, -0.5);
            \draw [->] (-1,0) to (1, 0);
            \node [left] at (-1, 0) {$0$};
            \node [right] at (1, 0) {$2R$};
            \draw [dashed] (0, -1) to (0, 1);
            \draw [dashed] (-0.3, -0.954) to (-0.3, 0.954);
            \draw [dashed] (-0.7, -0.714) to (-0.7, 0.714);
            \draw [dashed] (0.3, -0.954) to (0.3, 0.954);
            \draw [dashed] (0.7, -0.714) to (0.7, 0.714);
            \draw [ultra thick] (-0.5, 0) -- (0.5, 0);
            \draw [ultra thick] (-0.5, 0.05) -- (-0.5, -0.05);
            \draw [ultra thick] (0.5, 0.05) -- (0.5, -0.05);
            \node [below] at (-0.5, 0) {$\frac{R}{2}$};
            \node [below] at (0.5, 0) {$\frac{3R}{2}$};
        \end{tikzpicture}
    }\hfil}
    Таким чином, множина всіх значень координати середини хорди $\Omega = \left[0; 2R\right]$. 
    Множина, що відповідає події --- це відрізок $A = \left[\frac{R}{2}; \frac{3R}{2}\right]$.
    В якості міри візьмемо довжину. Тоді $\P(A) = \frac{l(A)}{l(\Omega)} = \frac{R}{2R}= \frac{1}{2}$.
    
    \textbf{Спосіб 2.}
    В цьому способі пропонується розглядати тільки хорди з одним закріпленим кінцем.
    Кожній хорді поставимо у відповідність ту частину дуги кола, яка потрапляє у кут $\varphi$,
    що утворює хорда з дотичною, проведеною через точку закріплення кінця всіх хорд.
    \newline
    \hbox to \hsize{\hfil{
        \begin{tikzpicture}[scale = 1.5]
            \draw (0, 0) circle [radius = 1];
            \draw (0, 1) -- (-0.86602540378, -0.5); 
            \draw (0, 1) -- (0.86602540378, -0.5);
            \draw (0.86602540378, -0.5) -- (-0.86602540378, -0.5);
            \draw (-1.2, 1) -- (1.2, 1);
            \draw [fill] (0, 1) circle [radius = 0.05];
            \draw [fill] (0.17, -0.985) circle [radius = 0.05];
            \draw [fill] (1, 0) circle [radius = 0.05];
            \draw [fill] (-0.967, -0.256) circle [radius = 0.05];
            \draw (0.16, 1) arc (0:-46:0.16);
            \draw (0.19, 1) arc (0:-46:0.19);
            %\node [above] at (0, 1);
            \draw [dashed] (0, 1) to (0.17, -0.985);
            \draw [dashed] (0, 1) to (1, 0);
            \draw [dashed] (0, 1) to (-0.967, -0.256);
            \node [below right] at (0.14, 1) {$_\varphi$};
            \draw [ultra thick](0.86602540378, -0.5) to [out = 242, in = 0](0, -1);
            \draw [ultra thick](0, -1) to [out = 180, in = 300](-0.86602540378, -0.5);
        \end{tikzpicture}
    }\hfil}
    Тоді множина $\Omega = \left[ 0; \pi\right]$, а $A = \left[ \frac{\pi}{3}; \frac{2\pi}{3}\right]$.
    Знову візьмемо за міру довжину і отримаємо $\P(A) = \frac{\pi /3}{\pi} = \frac{1}{3}$.
    
    \textbf{Спосіб 3.}
    Розглядаються всі хорди кола. Кожній з них взаємно однозначно ставиться у відповідність точка її середини $(x,y)$,
    якщо за початок координат прийняти центр кола.
    \newline
    \hbox to \hsize{\hfil{
        \begin{tikzpicture}[scale = 1.8]
            \draw (0, 0) circle [radius = 1];
            \draw (0, 0) circle [radius = 0.5];
            \draw (0, 0) [fill] circle [radius = 0.05];
            \draw (0, 1) -- (-0.86602540378, -0.5); 
            \draw (0, 1) -- (0.86602540378, -0.5);
            \draw (0.86602540378, -0.5) -- (-0.86602540378, -0.5);
            \draw (-0.737, 0.675) [fill] circle [radius = 0.05];
            \draw (0.878, 0.479) [fill] circle [radius = 0.05];
            \draw [dashed] (-0.737, 0.675) to (0.878, 0.479);
            \draw [dashed] (0, 0) to (0.0705, 0.577);
            \draw (0.0705, 0.577) [fill] circle [radius = 0.05];
        \end{tikzpicture}
    }\hfil}
    В такому випадку $\Omega = \left\{ (x,y)\in \mathbb{R}^2: x^2 + y^2 \leq R^2\right\}$,
    а $A = \left\{ (x,y)\in \Omega: x^2 + y^2 \leq \frac{R^2}{4}\right\}$.
    Тоді $\P(A) = \frac{\pi R^2/4}{\pi R^2} = \frac{1}{4}$.

    Жозеф Бертран був противником геометричного означення ймовірності.
    Він казав, що воно не дає можливості однозначно обчислити ймовірність
    однієї і тієї ж випадкової події та використовував вищенаведений приклад
    для доведення своєї правоти. Дійсно, було отримано три різних відповіді при
    розв'язанні однієї задачі. Так в чому ж справа? Насправді, помилка полягає
    у різних трактуваннях поняття <<навмання обрана хорда>>.
    В кожному з трьох способів це трактування було різним, що й стало причиною різних відповідей.
\end{example}

        % !TEX root = ../main.tex

\newlength\Radius
\setlength\Radius{2cm}

\section{Аксіоми теорії ймовірностей}
Розглянуті ймовірнісні схеми (класична, геометрична) мають свої недоліки і не дають можливості описати різноманіття стохастичних експериментів.
Сучасна теорія ймовірностей побудована на основі аксіом, які були запропоновані академіком 
\href{https://uk.wikipedia.org/wiki/%D0%9A%D0%BE%D0%BB%D0%BC%D0%BE%D0%B3%D0%BE%D1%80%D0%BE%D0%B2_%D0%90%D0%BD%D0%B4%D1%80%D1%96%D0%B9_%D0%9C%D0%B8%D0%BA%D0%BE%D0%BB%D0%B0%D0%B9%D0%BE%D0%B2%D0%B8%D1%87}{Андрієм Миколайовичем Колмогоровим} у 1929 році.
Класична та геометрична схеми виявляються частинними випадками цього аксіоматичного підходу.

\subsection{Система аксіом Колмогорова}\index{система аксіом Колмогорова}
\begin{enumerate}[label=\Roman*.]
    \item Побудова вимірного простору $\left\{ \Omega, \mathcal{F}\right\}$:
    \begin{enumerate}[label = \textbf{A\arabic*:}]
        \item Будь-якому стохастичному експерименту можна поставити у відповідність
        простір елементарних подій $\Omega$.\index{подія!простір елементраних подій}
        \item $\left(\forall n \in \mathbb{N}: A_n \in \mathcal{F} \right) \Rightarrow \left( \bigcup\limits_{n=1}^{\infty} A_n \in \mathcal{F}\right)$.
        \item $\left( A \in \mathcal{F}\right) \Rightarrow \left( \overline{A} \in \mathcal{F}\right)$.
        
        Дві останні аксіоми стосуються побудови $\sigma$-алгебри подій.
    \end{enumerate}
    \item Аксіоми ймовірності:
    \begin{enumerate}[label = \textbf{P\arabic*:}]
        \item $\forall A \in \mathcal{F}: \P(A)\geq 0$.
        \item $\P(\Omega) = 1$ --- аксіома нормування. \index{аксіома!нормування}
        \item $\forall A_1, A_2, ..., A_n, ... \in \mathcal{F},  A_i \cap A_j = \varnothing \text{ при } i \neq j: \P\left(\bigcup\limits_{n=1}^{\infty} A_n\right) = \sum\limits_{n=1}^{\infty} \P(A_n)$ ---
        аксіома зліченної адитивності. \index{аксіома!зліченної адитивності}
    \end{enumerate}
\end{enumerate}
\begin{remark}
    Іноді, у випадках <<простої>> алгебри, коли потребуємо у $\textbf{A2}$ лише $\bigcup\limits_{k=1}^n A_k \in \mathcal{F}$, достатньо
    аксіоми скінченної адитивності \textbf{P3$'$}: \index{аксіома!скінченної адитивності}
    $\forall A_1, A_2, ..., A_n \in \mathcal{F},  A_i \cap A_j  = \varnothing \text{ при } i \neq j: \P\left(\bigcup\limits_{k=1}^{n} A_k\right) = \sum\limits_{k=1}^{n} \P(A_k)$.
\end{remark}
\begin{definition}\index{ймовірність}
    Нормована міра $\P$, що введена на вимірному просторі $\left\{ \Omega, \mathcal{F}\right\}$,
    називається \emph{ймовірністю}. Всі події беруться з $\sigma$-алгебри подій.
    Трійка $\left\{ \Omega, \mathcal{F}, \P\right\}$ називається 
    \emph{ймовірнісним простором} стохастичного експерименту.
\end{definition}

Ця система аксіом несуперечна, бо існують стохастичні експерименти,
які задовольняють цим аксіомам, але неповна, бо в різних задачах
теорії ймовірностей розглядаються різні ймовірнісні простори.

\begin{example}
    Стохастичний експеримент полягає в киданні монети до першої появи герба. Введемо елементарні події $\omega_k = \left\{ 
        \text{герб випав при } k \text{-тому киданні}
    \right\}$. Простір елементарних подій буде зліченною множиною:
    $\Omega = \left\{\omega_k : k \in \mathbb{N} \right\}$. В якості $\sigma$-алгебри $\mathcal{F}$ візьмемо множину всіх підмножин $\Omega$.
    Оскільки ймовірність випадіння герба на <<чесній>> монеті рівна $\frac{1}{2}$ (як і ймовірність того, що герб не випаде), то логічно кожній елементарній події $\omega_k$ поставити 
    у відповідність число $p_k = \frac{1}{2^k}$. Також, для всіх $A \subset \Omega$ (або ж, в цьому випадку еквівалентно, $A \in \mathcal{F}$) покладемо
    $\P(A) = \sum\limits_{\omega_k \in A} p_k$.
\end{example}

\begin{exercise}
    Перевірити виконання усіх аксіом для цього прикладу, а також для класичної моделі ймовірності.
\end{exercise}

\begin{remark}
    Не завжди в якості $\sigma$-алгебри $\mathcal{F}$ можна обрати множину всіх підмножин $\Omega$. 
    У випадку класичної моделі ймовірності та у розглянутому вище прикладі
    це можна зробити, тому що, в деякому сенсі, $2^{\Omega}$ містить <<небагато>> подій: множина всіх підмножин скінченної множини також
    є скінченною, а множина всіх підмножин зліченної множини має потужність континууму (таку ж, як відрізок $[0; 1]$).
    У тому випадку, коли вже сам простір $\Omega$ має потужність континууму (як у випадку геометричної моделі ймовірності),
    множина всіх його підмножин має потужність, більшу за континуум, тому, в деякому сенсі, цих підмножин стає <<дуже багато>>.
    Для стохастичного експерименту, який полягає у киданні навмання точки на відрізок $[0; 1]$,
    можна навести \href{https://uk.wikipedia.org/wiki/%D0%9C%D0%BD%D0%BE%D0%B6%D0%B8%D0%BD%D0%B0_%D0%92%D1%96%D1%82%D0%B0%D0%BB%D1%96}{конкретний приклад}
    побудови підмножини $[0; 1]$, якій не можна надати \emph{жодного} значення ймовірності, навіть нульового.
    І хоча цей приклад виходить за рамки курсу, варто зауважити, що наводиться лише доведення того, що така множина існує:
    алгоритм побудови спирається на вибір елементів з незліченної множини, тому не може бути реалізований на будь-якому обчислювальному пристрої
    (можливість такого вибору ґрунтується на так званій <<аксіомі вибору>> з теорії множин).
    Через це для незліченних $\Omega$ в якості $\mathcal{F}$ треба розглядати такі набори підмножин, які все ще будуть досить великими,
    щоб включати в себе якомога більше подій, знання ймовірності яких може знадобитися, але не настільки великими, щоб це унеможливлювало введення ймовірності взагалі.
\end{remark}

\subsection{Властивості ймовірності, що випливають з аксіом}
\begin{enumerate}
    \item Якщо $A_1, A_2, ..., A_n, ... \in \mathcal{F}$ утворюють повну групу 
    подій, то $\P\left(\bigcup\limits_{k=1}^\infty A_k\right) = 1$.
    \begin{proof}
        Випливає з аксіоми \textbf{P2}.
    \end{proof}
    \item $\P(\overline{A}) = 1 - \P(A)$.
    \begin{proof}
        $A \cup \overline{A} = \Omega \Rightarrow 1 \overset{\text{\textbf{P2}}}{=} \P(\Omega) 
        = \P(A \cup \overline{A}) \overset{\text{\textbf{P3}}'}{=} \P(A) + \P(\overline{A})$.
    \end{proof}
    \item $\P(\varnothing) = 0$.
    \begin{proof}
        Наслідок із властивості 2 $(\varnothing = \overline{\Omega})$.
    \end{proof}
    \suspend{enumerate}
    \begin{remark}
        Якщо ймовірність події дорівнює нулю, то це не означає, що подія неможлива. Наприклад,
        якщо СЕ --- кидання точки на відрізок $[a; b]$, а подія $A$ полягає в тому, що точка потрапила в 
        певну точку $x \in [a; b]$, то $A$ не є неможливою, проте $\P(A) = 0$.
    \end{remark}
    \resume{enumerate}
    \item $A \subset B \Rightarrow \P(A) \leq \P(B)$.
    \begin{proof}
        $A \subset B \Rightarrow B = A \cup (B \setminus A) 
        \overset{\text{\textbf{P3}}'}{\Rightarrow} \P(B) = \P(A) + \P(B \setminus A) 
        \overset{\text{\textbf{P1}}}{\geq} \P(A)$.
    \end{proof}
    \item $\forall A \in \mathcal{F}: \P(A) \leq 1$.
    \begin{proof}
        Наслідок з властивості 4, $A \subset \Omega$ та \textbf{P2}.
    \end{proof}
    \item $A \subset B: \P(B \setminus A) = \P(B) - \P(A)$.
    \begin{proof}
        Наслідок з доведення властивості 4.
    \end{proof}
    \item $\forall A, B \in \mathcal{F}: \P(A \cup B) = \P(A) + \P(B) - \P(A \cap B)$.
    \begin{proof}
        $A \cup B = (A\setminus(A \cap B)) 
        \cup (A \cap B) 
        \cup (B\setminus(A \cap B))$ --- попарно несумісні події. 
        \newline
        З аксіоми \textbf{P3$'$}: $\P(A \cup B) = \P(A\setminus(A \cap B)) 
        + \P(A \cap B) + \P(B\setminus(A \cap B)) \overset{6}{=} \P(A) - \P(A \cap B) + \P(A \cap B)
        + \P(B) - \P(A \cap B) = \P(A) + \P(B) - \P(A \cap B)$.
    \end{proof}
    \item Узагальнення властивості 7 --- формула включення-виключення для ймовірностей: 
    
    \begin{math}
        \forall A_1, A_2, \dots, A_n \in \mathcal{F} 
    : \P\left(\bigcup\limits_{i=1}^n A_i\right) = \sum\limits_{i=1}^n \P(A_i) - \sum\limits_{i < j}^n \P(A_i \cap A_j)
    + \sum\limits_{i < j < k}^n \P(A_i \cap A_j \cap A_k) - ... + (-1)^{n-1}\P\left(\bigcap\limits_{i=1}^n A_i\right)
    \end{math}.
    \nopagebreak
    \suspend{enumerate}
    \begin{exercise}
        Довести властивість 8.
    \end{exercise}

\begin{example}[задача про неуважну секретарку]\index{задача!про неуважну секретарку}
    Секретарка поклала $n$ листів в $n$ чистих конвертів, заклеїла ці конверти і тільки 
    після цього написала адреси. Яка ймовірність того, що хоча б один з листів дійде 
    за призначенням?
    
    Введемо події
    $A_i = \left\{i\text{-тий лист дійшов за призначенням}\right\}, i = 
    1,...,n. \;\P(A_i) = \frac{1}{n} = \frac{(n-1)!}{n!}$.
    Тоді шукана подія
    $A = \left\{\text{хоча б один із листів дійшов за призначенням}\right\} = \bigcup\limits_{i=1}^n A_i$.
    Для застосування формули включення-виключення треба знайти відповідні ймовірності перетинів подій:
    $\P(A_i \cap A_j) = \frac{(n-2)!}{n!} = \frac{1}{n(n-1)}, i \neq j$,
    $\P(A_i \cap A_j \cap A_k) = \frac{1}{n(n-1)(n-2)}, i \neq j \neq k$, і так далі,
    $\P(A_1 \cap ... \cap A_n) = \frac{1}{n!}$. Формула для обчислення ймовірності перетину деяких $k$ подій --- це частка
    тих перестановок множини $\left\{1,...,n\right\}$ (тобто, усіх можливих варіантів того, куди можуть потрапити листи),
    в яких відповідні $k$ листів доходять за призначенням, а інші --- довільним чином. Зрозуміло, що усього варіантів
    обрати ці $k$ листів --- $C_n^k$. Тому отримуємо
    $\P(A) = n\cdot \frac{1}{n} - C_n^2 \cdot \frac{1}{n(n-1)}- \dots + (-1)^{n-1}\frac{1}{n!}
    = 1 - \frac{1}{2!} + \frac{1}{3!} - \dots +(-1)^{n-1} \cdot \frac{1}{n!} \approx 
    1 - e^{-1} \approx 0.63$.
\end{example}
\resume{enumerate}
    \item $\forall A_1, A_2, ..., A_n, ... \in \mathcal{F}: 
    \P\left(\bigcup\limits_{k=1}^\infty A_k\right) \leq \sum\limits_{k=1}^\infty \P(A_k)$
    \begin{proof}
        Введемо події $B_1 = A_1, B_2 = \overline{A_1} \cap A_2, 
        B_3 = \overline{A_1} \cap \overline{A_2} \cap A_3, ..., B_n = 
        \overline{A_1} \cap \overline{A_2} \cap ...$ 
        \newline
        $... \cap \overline{A_{n-1}} \cap A_n, ...$. Ці події попарно несумісні, 
        $\bigcup\limits_{k=1}^\infty B_k = \bigcup\limits_{k=1}^\infty A_k,\;B_k \subset A_k $ для всіх $k \in \mathbb{N}$.
        \newline
        $\P\left(\bigcup\limits_{k=1}^\infty A_k\right) = \P\left(\bigcup\limits_{k=1}^\infty B_k\right) 
        \overset{P3}{=} \sum\limits_{k=1}^\infty \P(B_k) \overset{4}{\leq} 
        \sum\limits_{k=1}^\infty \P(A_k)$.
    \end{proof}
    \item $\forall A_1, A_2, ..., A_n, ... \in \mathcal{F}: \P\left(\bigcap\limits_{k=1}^\infty
    A_k\right) = 1 - \P\left(\overline{\bigcap\limits_{k=1}^\infty A_k}\right) = 1 - \P\left(\bigcup\limits_{k=1}^\infty \overline{A_k}\right) \geq  
    1 - \sum\limits_{k=1}^\infty \P(\overline{A_k})$.
\end{enumerate}

\subsection{Теореми неперервності ймовірності}
\index{теорема!неперервності}
\begin{theorem}\label{th:1}
    Нехай є монотонно неспадна послідовність подій $A_1 \subset A_2 \subset ... \subset A_n ... \in \mathcal{F}$.
    Тоді $\P\left(\bigcup\limits_{n=1}^{\infty} A_n\right) = \lim\limits_{n\rightarrow \infty} \P(A_n)$.
\end{theorem}
\begin{proof}
    Зважаючи на монотонну неспадність послідовності, можна записати $\bigcup\limits_{n=1}^{\infty} A_n = A_1 \cup (A_2 \setminus A_1) \cup (A_3 \setminus A_2) \cup ... \cup (A_{n} \setminus A_{n-1}) \cup ... =
    A_1 \cup \left(\bigcup\limits_{n=1}^{\infty} (A_n \setminus A_{n-1})\right)$, причому $A_1$ та усі $A_n \setminus A_{n-1}$ є попарно несумісними.
    Тоді з аксіоми \textbf{P3}: $\P\left(\bigcup\limits_{n=1}^{\infty} A_n\right) = \P(A_1) + \sum\limits_{n=2}^{\infty} \P(A_{n} \setminus A_{n-1})$,
    і цей ряд є збіжним. Запишемо значення його часткової суми:

    \noindent $S_n = \P(A_1) + \sum\limits_{k=2}^{n} \P(A_{k} \setminus A_{k-1}) = \P(A_1) + \P(A_2) - \P(A_1) + \P(A_3) - \P(A_2) + ... + \P(A_n) - \P(A_{n-1}) = \P(A_n)$.
    Отже, $\P\left(\bigcup\limits_{n=1}^{\infty} A_n\right) = \lim\limits_{n\rightarrow \infty} S_n = \lim\limits_{n\rightarrow \infty} \P(A_n)$.
\end{proof}
\begin{remark}
    Застосовуючи природне позначення $\bigcup\limits_{n=1}^{\infty} A_k = \lim\limits_{n\rightarrow \infty} A_n$, твердження теореми \ref{th:1}
    можна записати як $\P\left( \lim\limits_{n\rightarrow \infty} A_n\right) = \lim\limits_{n\rightarrow \infty} \P(A_n)$, що пояснює назву теореми.
    Також, <<розширення>> подій $A_n$ до $A$ можна записати як $A_n \nearrow A$, тому твердження теореми можна ще записати як
    $\left(A_n \nearrow A\right) \Rightarrow \left(\P(A_n) \rightarrow \P(A)\right)$.
\end{remark}

\begin{theorem}\label{th:2}
    Нехай є монотонно спадна послідовність подій $A_1 \supset A_2 \supset ... \supset A_n ... \in \mathcal{F}$.
    Тоді $\P\left(\bigcap\limits_{n=1}^{\infty} A_n\right) = \lim\limits_{n\rightarrow \infty} \P(A_n)$.
\end{theorem}
\begin{proof}
    $\P\left(\bigcap\limits_{n=1}^{\infty} A_n\right) = 1 - \P\left(\bigcup\limits_{n=1}^{\infty} \overline{A_n}\right) \overset{\text{т. \refeq{th:1}}}{=} 1 -
    \lim\limits_{n\rightarrow \infty} \P(\overline{A_n}) = \lim\limits_{n\rightarrow \infty} (1-\P(\overline{A_n})) = \lim\limits_{n\rightarrow \infty} \P(A_n)$.
\end{proof}
\begin{remark}
    Аналогічно, застосовуючи позначення $\bigcap\limits_{n=1}^{\infty} A_k = \lim\limits_{n\rightarrow \infty} A_n$ та $A_n \searrow A$, твердження теореми
    \ref{th:2} можна записати як $\P\left( \lim\limits_{n\rightarrow \infty} A_n\right) = \lim\limits_{n\rightarrow \infty} \P(A_n)$ або 
    $\left(A_n \searrow A\right) \Rightarrow \left(\P(A_n) \rightarrow \P(A)\right)$.
\end{remark}

За допомогою теорем \ref{th:1} та \ref{th:2} можна довести, що замість аксіоми \textbf{P3} можна розглядати деяку іншу аксіому.
\begin{theorem}\index{аксіома!неперервності}
    Аксіома зліченної адитивності \emph{\textbf{P3}} еквівалентна одночасному виконанню аксіом неперервності
    \emph{\textbf{P4}} та скінченної адитивності \emph{\textbf{P3$'$}}: $$\forall n \in \mathbb{N}: B_n \in \mathcal{F}, B_1 \supset B_2 \supset ... \supset B_n \supset ..., \bigcap\limits_{n=1}^{\infty} B_n = \varnothing \Rightarrow \lim_{n\rightarrow \infty} \P(B_n) = 0$$
    \begin{proof}
        За теоремою \refeq{th:2} маємо наслідок \textbf{P3} $\Rightarrow$ \textbf{P4}, \textbf{P3$'$}.

        Доведемо \textbf{P4}, \textbf{P3$'$} $\Rightarrow$ \textbf{P3}. Для подій $A_1, A_2, ... , A_n, ... \in \mathcal{F}, A_i \cap A_j = \varnothing \text{ при } i \neq j$
        введемо події $B_n = \bigcup\limits_{k=n+1}^{\infty} A_k$. Для них виконується $B_1 \supset B_2 \supset ... \supset B_n \supset ...$ та $\bigcap\limits_{n=1}^{\infty} B_n = \varnothing$,
        тому з аксіоми \textbf{P4} $\lim\limits_{n\rightarrow \infty} \P(B_n) = 0$.

        З іншого боку, з аксіоми \textbf{P3$'$} маємо таку рівність:
        $$\P\left(\bigcup\limits_{k=1}^{\infty} A_k\right) = \P\left(\bigcup\limits_{k=1}^{n} A_k\right) + \P\left(\bigcup\limits_{k=n+1}^{\infty} A_k\right) = \sum\limits_{k=1}^{n} \P(A_k) + \P(B_n)$$
        Перейшовши у правій частині до границі при $n \rightarrow \infty$, отримаємо $\P\left(\bigcup\limits_{k=1}^{\infty} A_k\right) = \sum\limits_{k=1}^{\infty} \P(A_k)$, 
        що є твердженням аксіоми зліченної адитивності.
    \end{proof}
\end{theorem}

        % !TEX root = ../main.tex
\section{Умовна ймовірність та її застосування}

\subsection{Поняття умовної ймовірності}
Припустимо, що спостерігається деякий експеримент, що описується класичною моделлю, а для події $B$ $\P(B)>0$.
\emph{Умовна ймовірність} $\P(A/B)$ --- це ймовірність події $A$ за умови, що відбулась подія $B$.
Наприклад, експеримент --- витягання карт з колоди, $A=\left\{\text{витягнуто даму пік}\right\}$, $B=\left\{\text{витягнуто карту чорної масті}\right\}$.

Оскільки експеримент описується класичною моделлю, то можемо позначити $m_B$ та $m_{A\cap B}$ кількості елементарних подій, що сприяють появами подій $B$ та $A \cap B$ відповідно.
Тоді $\P(A/B) = \frac{m_{A\cap B}}{m_B} = \frac{m_{A\cap B}/n}{m_B/n} = \frac{\P(A\cap B)}{\P(B)}$, де $n = \mathrm{card}(\Omega)$.

Для довільних ймовірнісних просторів ця формула вводиться як означення умовної ймовірності.
\begin{definition}
    Якщо подія $B$ має додатну ймовірність $\P(B)>0$, то \emph{умовна ймовірність} події $A$ за умови, що відбулась подія $B$,
    обчислюється за формулою 
    \begin{equation}\label{eq:cond_prob}
        \P(A/B) = \frac{\P(A\cap B)}{\P(B)}
    \end{equation}
\end{definition}
\noindent \textbf{Властивості умовної ймовірності:}
\begin{enumerate}
    \item $\P(A/B) \geq 0$.
    \item $\P(\Omega /B) = \P(B/B) = 1$.
    \item $ \forall A_1, A_2, ... , A_n, ... \in \mathcal{F}: A_i \cap A_j = \varnothing \text{ при } i \neq j : \\
    \P\left(\left(\bigcup\limits_{n=1}^{\infty} A_n\right)/B\right) = \frac{\P\left(\left(\bigcup\limits_{n=1}^{\infty} A_n\right)\cap B\right)}{\P(B)} = \frac{\P\left(\bigcup\limits_{n=1}^{\infty} (A_n\cap B)\right)}{\P(B)} \overset{\text{\textbf{P3}}}{=} \frac{\sum\limits_{n=1}^{\infty} \P(A_n \cap B)}{\P(B)} = \sum\limits_{n=1}^{\infty} \P(A_n/B)$.
\end{enumerate}
\vspace{1em}
Таким чином, для умовної ймовірності виконуються аксіоми \textbf{P1}, \textbf{P2} та \textbf{P3}. 
\begin{exercise}
    Нехай $\left\{ \Omega, \mathcal{F}, \P\right\}$ --- деякий ймовірнісний простір, а $B\in\mathcal{F}$ --- деяка
    подія з $\P(B) > 0$. Перевірити, що для простору елементарних подій $\Omega_B$, що складається
    з усіх елементарних подій, які містить $B$, $\mathcal{F}_B = \left\{ A\cap B : A \in \mathcal{F} \right\}$ є $\sigma$-алгеброю.
    Разом з вже доведеними властивостями умовної ймовірності $\P(\cdot / B)$ це означає, що $\left\{ \Omega_B, \mathcal{F}_B, \P(\cdot / B)\right\}$
    також є ймовірнісним простором, в якому $B$ є вірогідною подією.
\end{exercise}

\subsection{Незалежність подій}
\begin{definition}
    Дві події $A$ та $B$ називаються \emph{незалежними}, якщо
    \begin{equation}\label{eq:indep_events}
        \P(A/B) = \P(A) \text{ або } \P(B/A) = \P(B)
    \end{equation}
    Це означає, що на появу однієї події не впливає поява іншої.
\end{definition}
\noindent \textbf{Властивості незалежних подій:}
\begin{enumerate}
    \item Якщо події $A$ та $B$ незалежні, то $\P(A\cap B) = \P(A)\cdot \P(B)$. Це випливає з \eqref{eq:cond_prob} та \eqref{eq:indep_events}.
    \begin{remark}
        Несумісні події залежні, якщо їх ймовірності не нульові.
    \end{remark}

    \item Якщо події $A$ та $B$ незалежні, то пари $A$ та $\overline{B}$, 
    $\overline{A}$ та $B$, $\overline{A}$ та $\overline{B}$ --- теж незалежні.
    \begin{proof}
        Доведемо спочатку для пари $A$ та $\overline{B}$. $\P(A\cap \overline{B}) = \P(A\setminus (A\cap B)) = \P(A) - \P(A\cap B) = \P(A) - \P(A)\cdot \P(B) = \P(A)\cdot (1 - \P(B)) = \P(A)\cdot \P(\overline{B})$.
        Отже, $A$ та $\overline{B}$ --- незалежні. Доведення для інших пар аналогічне.
    \end{proof}
    \item Якщо $A$ та $B$, $A$ та $C$ незалежні, а $B$ та $C$ несумісні, то $A$ та $B\cup C$ --- незалежні події.
    \begin{proof}
        $\P(A\cap (B \cup C)) = \P((A \cap B) \cup (A \cap C)) = \P(A\cap B) + \P(A\cap C) = \P(A)\cdot \P(B) + \P(A)\cdot \P(C) = \P(A) \cdot (\P(B) + \P(C)) = \P(A)\cdot \P(B\cup C).$
    \end{proof}
\end{enumerate}

\begin{definition}
    Події $A_1, A_2, ..., A_n \in \mathcal{F}$ називаються \emph{незалежними у сукупності}, якщо
    \begin{equation}\label{eq:indep}
        \forall i_1, i_2, ..., i_r \in \{1,...,n\}: \P\left(\bigcap\limits_{k=1}^r A_{i_k}\right) = \prod\limits_{k=1}^r \P\left(A_{i_k}\right)
    \end{equation}
    Зокрема:
    \nopagebreak
    \begin{enumerate}
        %\nopagebreak
        \item $\forall i, j \in {1,..,n}, i\neq j: \P(A_i \cap A_j) = \P(A_i)\cdot \P(A_j)$ --- попарна незалежність.
        \item $\P\left(\bigcap\limits_{k=1}^n A_k\right) = \prod\limits_{k=1}^n \P(A_k)$.
    \end{enumerate}
\end{definition}
\begin{remark}
    З незалежності подій у сукупності випливає попарна незалежність, 
    але наслідку в іншу сторону, взагалі кажучи, немає.
\end{remark}
\begin{example}[приклад Бернштейна]
    Дитина кидає на підлогу тетраедр, три грані якого розфарбовані відповідно зеленим, синім та червоним кольором,
    а четверта --- усіма трьома кольорами. Введемо події $\text{З}, \text{С}, \text{Ч}$, які означають, що випала грань,
    на якій є зелений, синій або червоний колір відповідно.

    $\P(\text{З}) = \P(\text{С}) = \P(\text{Ч}) = \frac{1}{2}$, 
    $\P(\text{З} \cap \text{С}) = \P(\text{З} \cap \text{Ч}) = \P(\text{С} \cap \text{Ч}) = \frac{1}{4}$,
    тому ці події є попарно незалежними. Але $\P(\text{З} \cap \text{С} \cap \text{Ч}) = \frac{1}{4} \neq \frac{1}{8}$,
    тому незалежності у сукупності немає.
\end{example}

\subsection{Теореми додавання та множення}
\noindent\textbf{Теорема множення.} 
З формули \eqref{eq:cond_prob} випливає, що
\begin{equation}\label{eq:mult_for_2}
    \P(A\cap B) = \P(A)\cdot \P(B/A) = \P(B) \cdot \P(A/B)
\end{equation}
Для трьох подій: $\P(A \cap B \cap C) = \P(A\cap B) \cdot \P(C/(A\cap B)) = \P(A)\cdot \P(B/A) \cdot \P(C/(A\cap B))$.
За методом математичної індукції неважко довести, що
\begin{equation}\label{eq:mult_for_n}
    \P\left( \bigcap\limits_{k=1}^{n} A_k\right) = \P\left(A_1\right) \cdot \P\left(A_2/A_1\right) \cdot \P\left( A_3 / \left( A_1 \cap A_2\right)\right) \cdot ... \cdot \P\left( A_n / \left( \bigcap\limits_{k=1}^{n-1} A_k\right) \right)
\end{equation}

Якщо $A_1, A_2, ..., A_n \in \mathcal{F}$ --- незалежні у сукупності, то, як було зазначено, з \eqref{eq:indep} випливає, що
формула \eqref{eq:mult_for_n} спрощується до $\P\left(\bigcap\limits_{k=1}^n A_k\right) = \prod\limits_{k=1}^n \P(A_k) = \prod\limits_{k=1}^n \left(1-\P(\overline{A_k})\right)$.

\begin{example}
    На 10 картках написано по одній букві так, що можна скласти слово <<математика>>.
    Дитина навмання витягує без повернення по одній картці. Яка ймовірність того, що після витягання 4 карток
    утвориться слово <<мама>>?

    Треба обчислити ймовірність $\P(M^{(1)} \cap A^{(2)} \cap M^{(3)} \cap A^{(4)})$,
    де верхній індекс означає крок, на якому витягнуто літеру.
    $\P(M^{(1)}) = \frac{2}{10}$, $\P(A^{(2)} / M^{(1)}) = \frac{3}{9}$,
    $\P(M^{(3)}/(M^{(1)} \cap A^{(2)})) = \frac{1}{8}$, $\P(A^{(4)} / (M^{(1)} \cap A^{(2)} \cap M^{(3)})) = \frac{2}{7}$.
    Тому за формулою \eqref{eq:mult_for_n} $\P(M^{(1)} \cap A^{(2)} \cap M^{(3)} \cap A^{(4)}) = \frac{2 \cdot 3 \cdot 1 \cdot 2}{10 \cdot 9 \cdot 8 \cdot 7}$.
\end{example}

\noindent\textbf{Теорема додавання.} Нехай $A_1, A_2, ..., A_n \in \mathcal{F}$ --- незалежні у сукупності.
Тоді з рівностей $\P\left(\bigcup\limits_{k=1}^n A_k\right) = 1 - \P\left(\overline{\bigcup\limits_{k=1}^n A_k}\right) = 1 - \P\left(\bigcap\limits_{k=1}^n \overline{A_k}\right)$
випливає формула
\begin{equation}
    \P\left(\bigcup\limits_{k=1}^n A_k\right) = 1 - \prod\limits_{k=1}^n \P(\overline{A_k})
\end{equation}
Зокрема, для двох незалежних подій $A$ і $B$: $\P(A\bigcup B) = 1 - \P(\overline{A}) \cdot \P(\overline{B})$.

\subsection{Формула повної ймовірності}
Розглядаємо деякий ймовірнісний простір $\left\{ \Omega, \mathcal{F}, \P\right\}$.

Нехай $H_1, H_2, ..., H_n \in \mathcal{F}$ (може бути й зліченна кількість) --- деяка повна група подій,
які називаємо \emph{гіпотезами (припущеннями)}, а подія $A$ відбувається разом з якоюсь гіпотезою. Тоді має місце
\emph{формула повної ймовірності}:
\begin{equation}\label{eq:total_prob}
    \P\left( A \right) = \sum\limits_{k=1}^n \P(H_k)\cdot \P(A/H_k)
\end{equation}
\begin{proof}
    $\P(A) = \P(A \cap \Omega) = \P\left(A \cap \left(\bigcup\limits_{k=1}^n H_k\right)\right) = \P\left(\bigcup\limits_{k=1}^n (A\cap H_k)\right) = \sum\limits_{k=1}^n \P(A\cap H_k) = \sum\limits_{k=1}^n \P(H_k)\cdot \P(A/H_k)$.
\end{proof}
У випадку зліченної кількості гіпотез формула \eqref{eq:total_prob} доводиться аналогічно та перетворюється
на $\P\left( A \right) = \sum\limits_{k=1}^{\infty} \P(H_k)\cdot \P(A/H_k)$.

\begin{example}
    \begin{enumerate}
        \item У магазин постачають $80\%$ телефонів з Китаю, $15\%$ з В'єтнаму та $5\%$ з Кореї,
        причому бракованих відповідно $1\%$, $0.1\%$ та $0.01\%$.
        Знайти ймовірність події $A = \left\{ \text{куплений телефон буде бракованим}\right\}$.

        За умовою введемо гіпотези $H_1 = \left\{ \text{телефон з Китаю}\right\}$,
        $H_2 = \left\{ \text{телефон з В'єтнаму}\right\}$ та $H_3 = \left\{ \text{телефон з Кореї}\right\}$.
        $\P(H_1) = 0.8$, $\P(H_2) = 0.15$, $\P(H_3) = 0.05$,
        $\P(A/H_1) = 0.01$, $\P(A/H_2) = 0.001$, $\P(A/H_3) = 0.0001$.
        Тоді за формулою \eqref{eq:total_prob} $\P(A) = 0.8\cdot 0.01 + 0.15\cdot 0.001 + 0.05\cdot 0.0001 = 0.008155$.
        \item Серед $N$ екзаменаційних білетів $n$ <<щасливих>>, $n<N$.
        У якого студента ймовірність витягнути <<щасливий>> білет більша --- у того, хто тягне першим, чи того, хто тягне другим?

        Позначимо через $A$ і $B$ події, що <<щасливий>> білет витягнув перший та другий студент відповідно. За умовою $\P(A) = \frac{n}{N}$.
        Щоб скористатися формулою \eqref{eq:total_prob} для обчислення $\P(B)$, введемо гіпотези $H_1$ = $A$ та $H_2$ = $\overline{A}$.
        $\P(H_1) = \frac{n}{N}$, $\P(H_2) = \frac{N-n}{N}$.
        \\ $\P(B) = \P(H_1)\cdot \P(B/H_1) + \P(H_2)\cdot \P(B/H_2) = \frac{n}{N}\cdot \frac{n-1}{N-1} + \frac{N-n}{N}\cdot \frac{n}{N-1} = 
        \frac{n^2 - n + N\cdot n - n^2}{N\cdot(N-1)} = \frac{n\cdot (N-1)}{N\cdot(N-1)} = \frac{n}{N}$.
    \end{enumerate}
\end{example}

\subsection{Формула Баєса}
Як і раніше, нехай $H_1, H_2, ..., H_n \in \mathcal{F}$ --- повна група подій деякого СЕ, які називаємо гіпотезами, причому
перед проведенням експерименту відомі їх \emph{апріорні} ймовірності $\P(H_1), \P(H_2), ..., \P(H_n)$.
В результаті проведення експерименту відбулась деяка подія $A$.
Постає питання: чому рівні \emph{апостеріорні} ймовірності $\P(H_1/A), \P(H_2/A), ..., \P(H_n/A)$?
Тобто, чому дорівнюють ймовірності, що мала місце кожна з гіпотез за умови, що подія $A$ відбулась? На це питання дає відповідь \emph{формула Баєса}:
\begin{equation}\label{eq:bayes}
    \P(H_i/A) = \frac{\P(H_i \cap A)}{\P(A)} = \frac{\P(H_i) \cdot \P(A/H_i)}{\sum\limits_{k=1}^n \P(H_k)\cdot \P(A/H_k)}, i = 1,...,n
\end{equation}

\begin{example}
    В групі з 10 осіб троє вчаться на <<5>>, четверо на <<4>>, двоє на <<3>> та один на <<2>>.
    Викладач підготував на екзамен 20 питань. Студенти, які вчаться на <<5>>, знають відповіді на всі,
    на <<4>> --- на 16, на <<3>> --- на 10, на <<2>> --- на 5.
    Екзаменаційний білет містить 3 питання. Деякий студент відповів правильно на всі 3.
    Яка ймовірність того, що він вчиться на <<2>>?

    Введемо гіпотези $H_1 = \left\{ \text{студент вчиться на <<5>>}\right\}$, $H_2 = \left\{ \text{студент вчиться на <<4>>}\right\}$,
    $H_3 = \left\{ \text{студент вчиться на <<3>>}\right\}$, $H_4 = \left\{ \text{студент вчиться на <<2>>}\right\}$.
    За умовою $\P(H_1) = \frac{3}{10}$, $\P(H_2) = \frac{4}{10}$, $\P(H_3) = \frac{2}{10}$, $\P(H_4) = \frac{1}{10}$. 
    Позначимо $A = \left\{ \text{студент відповів на всі три питання}\right\}$. 
    Тоді $\P(A/H_1) = 1$, $\P(A/H_2) = \frac{16 \cdot 15 \cdot 14}{20 \cdot 19 \cdot 18}$, 
    $\P(A/H_3) = \frac{10 \cdot 9 \cdot 8}{20 \cdot 19 \cdot 18}$, 
    $\P(A/H_4) = \frac{5 \cdot 4 \cdot 3}{20 \cdot 19 \cdot 18}$.
    За формулою \eqref{eq:bayes} шукана ймовірність $\P(H_4/A) = \frac{\P(H_4) \cdot \P(A/H_4)}{\sum\limits_{k=1}^4 \P(H_k)\cdot \P(A/H_k)} =
    \frac{0.1 \cdot 5 \cdot 4 \cdot 3}{0.3 \cdot 20 \cdot 19 \cdot 18 + 0.4 \cdot 16 \cdot 15 \cdot 14 + 0.2 \cdot 10 \cdot 9 \cdot 8 + 0.1 \cdot 5 \cdot 4 \cdot 3}
     = \frac{60}{35460} = \frac{1}{591}$.

\end{example}
        % !TEX root = ../main.tex

\section{Незалежні випробування. Схема Бернуллі.}
\subsection{Поліноміальна схема}
Припустимо, що події $A_1, A_2, ..., A_k$ утворюють повну групу подій деякого 
СЕ, причому відомі ймовірності $\P(A_i) = p_i, \sum\limits_{i=1}^k p_i = 1$.
Експеримент проводиться $n$ разів, в кожному з яких може відбутись одна з 
подій $A_i$. Можливими результатами $n$ раз проведеного експерименту будуть події
$A_{i_1}^{(1)} \cap A_{i_2}^{(2)} \cap ... \cap A_{i_n}^{(n)}$,
де $A_{i_r}^{(r)}, r = 1,...,n, i_r = 1, ..., k$ --- один із можливих результатів $r$-того випробування.
\begin{definition}\index{незалежні випробування}
    Випробування, що проводяться, називаються \emph{незалежними}, якщо 
    $$\P\left(A_{i_1}^{(1)} \cap A_{i_2}^{(2)} \cap ... \cap A_{i_n}^{(n)}\right) = 
    \P\left(A_{i_1}^{(1)}\right) \cdot \P\left(A_{i_2}^{(2)}\right) \cdot ... \cdot \P\left(A_{i_n}^{(n)}\right)$$
\end{definition}
Прикладами незалежних випробувань є багаторазове підкидання монети або грального кубика,
проведення декількох тиражів лотереї з незмінними правилами тощо.

Основна задача --- знайти ймовірність того, що подія $A_1$ відбудеться $m_1$ разів, 
подія $A_2$ відбудеться $m_2$ разів і т.д., подія $A_k$ відбудеться $m_k$ разів, причому $\sum\limits_{j=1}^k m_j = n$. Такі ймовірності будемо позначати як $\P_n(m_1, m_2, ..., m_k)$.
Одним з можливих варіантів послідовності появи цих подій є поява події $A_1$ у перших $m_1$ випробуваннях,
$A_2$ --- у наступних $m_2$ і так далі, $A_k$ у останніх $m_k$. Знайдемо відповідну ймовірність:
$$\P\left(\underbrace{A_1 \cap ... \cap A_1}_{\text{перші }m_1} 
\cap \underbrace{ A_2 \cap ... \cap A_2}_{\text{наступні }m_2}
\cap ... \cap  \underbrace{A_k \cap ... \cap A_k}_{\text{останні }m_k}\right)
= p_1^{m_1} p_2^{m_2} ... p_k^{m_k}$$

Це лише один із можливих варіантів послідовності появи подій в серії випробувань. Всього таких варіантів 
рівно $C_n^{m_1, m_2, ..., m_k} = \frac{n!}{m_1!m_2! ... m_k!}$. Таким чином отримуємо:
\begin{equation}\label{eq:poly_prob}
    \P_n(m_1, ..., m_k) = \frac{n!}{m_1!m_2! ... m_k!} p_1^{m_1} p_2^{m_2} ... p_k^{m_k}
\end{equation}
\begin{definition}\index{схема ймовірності!поліноміальна}
    Ймовірності, що обчислюємо за формулою \eqref{eq:poly_prob}, називаються \emph{поліноміальними}, а 
    сама схема --- \emph{поліноміальною схемою ймовірності}. 
\end{definition}
\begin{remark} Назва <<поліноміальна>> пов'язана з тим, що $\frac{n!}{m_1!m_2! ... m_k!}$
    є коефіцієнтом при $p_1^{m_1} p_2^{m_2} ... p_k^{m_k}$ в розкладі полінома $(p_1 + p_2 + ... + p_k)^n$.
    Звідси ж можна отримати, що
    $$1 = (p_1 + p_2 + ... + p_k)^n = \sum\limits_{m_1 + ... + m_k = n} \P_n(m_1, ..., m_k)$$
\end{remark}
\begin{example}
    Студент ІПСА за рівнем підготовки з ймовірністю $0.3$ вважається слабким студентом, 
    з ймовірністю $0.5$ вважається середнім студентом та 
    з ймовірністю $0.2$ --- сильним студентом. Яка ймовірність того, що з 6 навмання 
    обраних студентів кількість слабких та сильних буде однаковою?

    Розглядаємо подію $A = \left\{\text{кількість слабких рівна кількості сильних}\right\}$. 
    Позначимо ймовірності того, що студент має певний рівень підготовки, таким чином:
    $p_1 = 0.3$, $p_2 = 0.5$, $p_3 = 0.2$.
    Треба знайти $\P(A)$.
    Скористаємось поліноміальною схемою:
    $\P(A) = \P_6(3, 0, 3) + \P_6(2, 2, 2) + \P_6(1, 4, 1) + \P_6(0, 6, 0) = 
    \frac{6!}{3!3!}\cdot(0.3)^3 \cdot(0.2)^3 + \frac{6!}{2!2!2!}\cdot(0.3)^2 \cdot(0.5)^2 \cdot(0.2)^2 + 
    \frac{6!}{1!4!1!}\cdot0.3 \cdot(0.5)^4 \cdot 0.2 + \frac{6!}{6!} \cdot(0.5)^6 = 0.213445$.
\end{example}
\subsection{Схема Бернуллі}
\begin{definition}\index{схема ймовірності!біноміальна (Бернуллі)}
    Поліноміальна схема, в кожному випробуванні якої є тільки дві події $A_1$ та $A_2$, 
    називається \emph{біноміальною схемою} або \emph{схемою Бернуллі}.
\end{definition}

Якщо $A_1$ та $A_2$ утворюють повну групу подій, то $A_2 = \overline{A_1}$.
Позначимо $A_1 = A$ --- <<успіх>>, $A_2 = \overline{A}$ --- <<невдача>>, 
$\P(A) = p$, $\P(\overline{A}) = 1-p = q$.

Яка ймовірність того, що в $n$ випробуваннях <<успіх>> з'явиться $m$ разів?
На це питання дає відповідь \emph{формула Бернуллі}, яку можна отримати з \eqref{eq:poly_prob}:
\index{формула!Бернуллі}
\begin{equation}\label{eq:bern_prob}
    \P_n(m) = \frac{n!}{m!(n-m)!}p^mq^{n-m} = 
    C_n^m p^m q^{n-m}
\end{equation}
Ймовірності, що обчислюються за цією формулою, називаються \emph{біноміальними}.

\begin{example}\label{ex:sportsman}
    Спортсмен 5 разів стріляє по мішені. Ймовірність влучення при кожному пострілі --- 0.2. 
    Яка ймовірність того, що він влучив не менше 3 разів?

    $\P\left\{\text{влучив} \geq  \text{3 разів} \right\} = \P_5(3) + \P_5(4) + \P_5(5) = 
    C_5^3 (0.2)^3 (0.8)^2 + C_5^4 (0.2)^4 0.8 + C_5^5 (0.2)^5 = 0.05792$.
\end{example}
\subsection{Найбільш імовірна кількість успіхів в схемі Бернуллі}
\begin{definition}
    Натуральне число $m_0$, при якому $\P_n(m)$ набуває найбільшого значення, 
    називається \emph{найбільш імовірною кількістю успіхів}.
\end{definition}

Знайдемо $m_0$:
$$\begin{cases}
    \P_n(m_0 - 1) \leq \P_n(m_0)\\
    \P_n(m_0 + 1) \leq \P_n(m_0)
\end{cases}
\Leftrightarrow 
\begin{cases}
    \frac{n!}{(m_0 - 1)!(n-m_0+1)!} p^{m_0 - 1} q^{n - m_0 + 1} 
    \leq 
    \frac{n!}{m_0!(n-m_0)!} p^{m_0} q^{n - m_0}\\
    \frac{n!}{(m_0 + 1)!(n-m_0-1)!} p^{m_0 + 1} q^{n - m_0 - 1} 
    \leq 
    \frac{n!}{m_0!(n-m_0)!} p^{m_0} q^{n - m_0}
\end{cases}
\Leftrightarrow
$$

$$
\Leftrightarrow
\begin{cases}
    \frac{q}{n-m_0+1} \leq \frac{p}{m_0} \\
    \frac{p}{m_0 + 1} \leq \frac{q}{n-m_0}
\end{cases}
\Leftrightarrow
\begin{cases}
    m_0 \leq np+p\\
    m_0 \geq np-q
\end{cases}
$$
\hbox to \hsize{\hfil{
    \begin{tikzpicture}[scale = 0.8]
        \node [below] at (2, 0) {$np-q$};
        \node [below] at (6, 0) {$np+p$};
        \draw [fill] (4, 0) circle [radius = 0.08];
        \draw [fill] (2, 0) circle [radius = 0.05];
        \draw [fill] (6, 0) circle [radius = 0.05];
        \node [below] at (4, 0) {$m_0$};
        \draw [-{Straight Barb}] [thick] (0, 0) -- (8, 0);
        \draw [thick] (2, 0) to (2, 0.5);
        \draw [thick] (6, 0) to (6, 0.5);
        \draw [{Straight Barb}-{Straight Barb}] [thick] (2, 0.5) to (6, 0.5);
        \node [above] at (4, 0.5) {1};
    \end{tikzpicture}
}\hfil}

Таким чином, отримуємо два варіанти:
\begin{enumerate}
    \item Якщо $np+p$ не ціле, то $m_0 = \left[np+p\right]$.
    \item Якщо $np+p$ ціле, то існують два найбільш ймовірних числа $m_0^{(1)} = np+p$ та $m_0^{(2)} = np-q$.
\end{enumerate}

\begin{remark}
    Фактично, під час знаходження $m_0$ було доведено, що $\P_n(m)$ спочатку зростають за $m$ до певного значення, 
    а потім --- спадають, причому максимальне значення досягається або при єдиному значенні $m$,
    або при двох сусідніх.
\end{remark}

\begin{example}
    Знайти найбільш імовірну кількість успіхів для попереднього \hyperref[ex:sportsman]{прикладу}.

    Оскільки $np + p = 1.2$, то $m_0 = \left[np + p\right] = \left[1.2\right] = 1$.
    Ймовірність лише одного влучення буде найбільшою, рівною $0.4096$.
\end{example}

\subsection{Асимптотичні наближення формули Бернуллі}
Формула Бернуллі \eqref{eq:bern_prob} не є зручною для великих значень $m$ та $n$ через складність обчислень, особливо для малих значень ймовірності успіху $p$. Існують формули, якими можна 
скористатись для наближення результатів формули Бернуллі.
\begin{proposition*}[асимптотична формула Пуассона]\index{формула!Пуассона, асимптотична}
    Нехай $n$ --- достатньо велике натуральне число, а p --- достатньо мале, щоб 
    $a = np \in (1, 20)$. Тоді для наближення формули Бернуллі можна використати 
    асимптотичну формулу Пуассона:
    \begin{equation}\label{eq:poiss_aprox}
        \P_n(m) \approx \frac{(np)^m}{m!}e^{-np}
    \end{equation}
\end{proposition*}
\begin{proof}[Доведення (нестроге)] Розпишемо $\P_n(m)$:
    \begin{gather*}
        \P_n(m) = C_n^m p^m (1-p)^{n-m} = \frac{n(n-1)...(n-m+1)}{m!} \left(\frac{a}{n}\right)^m 
        \left(1-\frac{a}{n}\right)^{n-m} = \\
        \frac{a^m}{m!}\cdot 1 \left(1 - \frac{1}{n}\right)\left(1 - \frac{2}{n}\right)
        ...\left(1 - \frac{m-1}{n}\right)\left(1-\frac{a}{n}\right)^n \left(1-\frac{a}{n}\right)^{-m} 
        \to \frac{a^m}{m!}e^{-a}, n\to\infty
    \end{gather*}
    Отже, ця наближена рівність є наслідком існування границі $\P_n(m)$ при $n\to\infty$. 
\end{proof}
\begin{proposition*}[формула Муавра-Лапласа]\index{формула!Муавра-Лапласа, асимптотична}
    Нехай $npq \geq 20$. Тоді для наближення формули Бернуллі можна використати 
    формулу Муавра-Лапласа:
    \begin{equation}\label{eq:laplace_approx}
        \P_n(m) \approx \frac{1}{\sqrt{2\pi npq}} e^{-\frac{(m-np)^2}{2npq}}
    \end{equation}
\end{proposition*}
Детальне доведення та пояснення обох тверджень буде в темі <<Граничні теореми теорії ймовірностей>> (ст. \pageref{binom_theorems}).

\begin{example}
    \begin{enumerate}
        \item Словник має 1500 сторінок. Ймовірність друкарської помилки на одній сторінці дорівнює
        0.001. Знайти ймовірність того, що в словнику буде рівно 3 помилки.

        За умовою $p=0.001$, $n=1500$, тому $np = 1.5$ і можна застосувати формулу Пуассона \eqref{eq:poiss_aprox}:
        $\P_{1500}(3) \approx \frac{1.5^3}{3!} e^{-1.5} \approx 0.12551$.
        Якщо ж скористатися точною формулою \eqref{eq:bern_prob}, отримаємо
        $\P_{1500}(3) \approx 0.12554$.
        \item Завод у середньому виготовляє 80\% виробів першого сорту та 20\% --- другого.
        Знайти ймовірність того, що серед 400 навмання відібраних виробів, які виготовлені на цьому заводі,
        84 вироби будуть другого сорту.

        За умовою $n=400$, $m=84$, $p=0.2$, $q=0.8$, тому $npq = 64$ і можна застосувати формулу Муавра-Лапласа \eqref{eq:laplace_approx}:
        $\P_{400}(84) \approx \frac{1}{\sqrt{2\pi \cdot 64}} \cdot \exp\left\{
            -\frac{(84 - 400\cdot0.2)^2}{2\cdot 64}\right\} \approx 0.044$. 
            Якщо ж скористатися точною формулою \eqref{eq:bern_prob}, отримаємо
        $\P_{400}(84) \approx 0.04323$. Зауважимо, що застосування в цьому випадку формули Пуассона \eqref{eq:poiss_aprox}
        дало б відповідь, далеку від правильної: приблизно $0.0394$.
    \end{enumerate}
\end{example}
    \chapter{Випадкові величини}
        % !TEX root = ../main.tex

\section{Поняття випадкової величини та її задання}
\begin{definition}
    \emph{Випадковою величиною} називається дійснозначна вимірна функція, що здійснює 
    відображення з $\Omega$ в $\mathbb{R}$, тобто $\xi = \xi(\omega): \Omega 
    \rightarrow \mathbb{R}$.
\end{definition}
\begin{remark}
    Вимірність функції $\xi(\omega)$ означає, що 
    \begin{equation}
        \forall x \in \mathbb{R}: 
        \left\{ \omega \in \Omega\; |\; \xi(\omega) < x\right\} \in \mathcal{F}
    \end{equation} В курсі 
    функціонального аналізу доводиться, що якщо виконується (1), то
    \begin{equation}
        \left.\begin{aligned}
            \{\omega\;|&\;\xi(\omega) > x\}\\
            \{\omega\;|&\;\xi(\omega) \leq x\}\\
            \{\omega\;|&\;\xi(\omega) \geq x\}\\
            \{\omega\;|&\;a < \xi(\omega) < b\}
        \end{aligned}\right\} \in \mathcal{F}
    \end{equation}
\end{remark}

Всього існують два види випадкових величин --- дискретні та неперервні.

\subsection{Дискретні випадкові величини та їх задання}
\begin{definition}
    Випадкова величина $\xi = \xi(\omega)$ називається 
    \emph{дискретною випадковою величиною} (ДВВ), якщо вона набуває скінченну або зліченну 
    кількість значень.
\end{definition}
Для задання ДВВ крім знання значень випадкової величини необхідно знати ймовірності, 
з якими ці значення приймаються.
$$p_i = P\left\{\omega: \xi(\omega) = x_i\right\} = P(\xi = x_i)$$
$$\sum_{i=1}^\infty p_i = 1$$
\begin{definition}
    \emph{Законом розподілу (імовірностей)} ДВВ називається співвідношення, яке вказує, 
    які ця випадкова величина приймає значення і з якими імовірностями.

    Закон розподілу ДВВ записується у вигляді \emph{ряду розподілу}:

    \begin{tabular}{c|c|c|c|c|c}
        $\xi$ & $x_1$ & $x_2$ & ... & $x_n$ & ... \\
        \hline
        $p$ & $p_1$ & $p_2$ & ... & $p_n$ & ...
    \end{tabular}
    \hspace{40pt}
    $x_1 < x_2 < ... < x_n < ...,\;\; \sum\limits_i p_i = 1$
\end{definition}
\begin{example}
    $\xi$ задає кількість влучень при чотирьох пострілах з імовірністю влучення 
    $p = \frac{1}{2}$. Скласти закон розподілу цієї випадкової величини.

    $P(\xi = k) = C_4^k \left(\frac{1}{2}\right)^k \left(\frac{1}{2}\right)^{4-k} = 
    C_4^k \left(\frac{1}{2}\right)^4 = \frac{C_4^k}{2^4}$

    \begin{tabular}{c|c|c|c|c|c}
        $\xi$ & 0 & 1 & 2 & 3 & 4 \\
        \hline
        $p$ & $\frac{1}{2^4}$ & $\frac{4}{2^4}$ & $\frac{6}{2^4}$ & $\frac{5}{2^4}$ & 
        $\frac{1}{2^4}$
    \end{tabular}
\end{example}
        % !TEX root = ../main.tex
\section{Числові характеристики випадкових величин}
Закон розподілу випадкової величини в будь-якому вигляді (ряд розподілу для ДВВ, щільність розподілу для НВВ, функція розподілу в загальному випадку)
повністю задає випадкову величину. Однак в деяких прикладних задачах достатньо мати лише деяке сумарне уявлення про деякі характерні невипадкові риси розподілу.

\subsection{Математичне сподівання}
\begin{definition}
    \emph{Математичним сподіванням} $\E\xi$ випадкової величини називається значення
    інтеграла Стілтьєса
    \begin{equation}\label{eq:e_xi}
        \E\xi = \int\limits_{-\infty}^{+\infty} x dF_\xi(x) = \begin{cases}
            \sum\limits_{k=1}^{n(\infty)} x_k \P\left\{\xi = x_k\right\}, & \xi \text{ --- ДВВ} \\
            \int\limits_{-\infty}^{+\infty} x f_\xi(x)dx, & \xi \text{ --- НВВ}
        \end{cases}
    \end{equation}
\end{definition}
Інтеграл або ряд \eqref{eq:e_xi} має збігатися \emph{абсолютно}, інакше кажуть, що
випадкова величина не має математичного сподівання. Ця вимога пояснюється тим, що, наприклад, у випадку ДВВ зі зліченною множиною значень не має бути різниці,
в якому порядку додавати всі значення.
\begin{example}
    НВВ, розподілена за законом Коші з щільністю $f_\xi(x) = \frac{1}{\pi (1+x^2)}$
    не має математичного сподівання, бо інтеграл $\frac{1}{\pi}\int\limits_{-\infty}^{+\infty} \frac{x}{1+x^2}dx$ 
    розбіжний.
\end{example}
Математичне сподівання характеризує середнє значення випадкової величини. Наприклад,
якщо ДВВ набуває скінченну кількість значень з однаковими ймовірностями, то математичним сподіванням є просто середнє арифметичне усіх цих значень.
Фізична інтерпретація --- центр мас системи точок $\left\{x_1, x_2, ..., x_n,...\right\}$ з масами $\left\{p_1, p_2, ..., p_n, ...\right\}$ (ДВВ) або стрижня,
розподіл маси в якому задано функцією щільності (НВВ).

\vspace{0.5em}
\noindent \textbf{Властивості математичного сподівання:}
\begin{enumerate}
    \item Математичне сподівання константи --- сама константа, оскільки
    її можна інтерпретувати як ДВВ, що приймає єдине значення з ймовірністю $1$:
    $c = const, \E c = c$.
    \item $\E \left(c\cdot\xi\right) = c\cdot \E\xi$ --- це властивість рядів та інтегралів.
    \item $\E\left( \xi_1 + \xi_2\right) = \E\xi_1 + \E\xi_2$.
    \begin{proof}[Доведення для ДВВ]
        $\P\left\{\xi_1 = x_k\right\} = p_k$, $\P\left\{\xi_2 = y_j\right\} = p_j$, $\P\left\{\xi_1 + \xi_2 = x_k + y_j\right\} = p_{kj}$.
        В цих позначеннях
        $\E\left( \xi_1 + \xi_2\right) = \sum\limits_k \sum\limits_j (x_k+y_j) p_{kj} =
        \sum\limits_k x_k \sum\limits_j p_{kj} + \sum\limits_j y_j \sum\limits_k p_{kj} = \sum\limits_k x_k p_k + \sum\limits_j y_j p_j = \E\xi_1 + \E\xi_2$.
    \end{proof}
\suspend{enumerate}
\begin{definition}
    Дві випадкові величини $\xi_1$ та $\xi_2$, задані на одному ймовірнісному просторі, називаються \emph{незалежними}, якщо
    події $A=\left\{\omega : \; \xi_1(\omega) < x\right\}$ та
    $B=\left\{\omega : \; \xi_2(\omega) < y\right\}$ є незалежними для всіх $x, y \in \mathbb{R}$.
    В цьому випадку $\P\left\{\xi_1 < x, \xi_2 < y\right\} = \P(A \cap B) = \P(A)\cdot \P(B) = F_{\xi_1}(x)\cdot F_{\xi_2}(y)$.
\end{definition}
\resume{enumerate}
    \item Якщо $\xi_1$ та $\xi_2$ незалежні, то $\E\xi_1\xi_2 = \E\xi_1 \cdot \E\xi_2$.
    \begin{proof}[Доведення для ДВВ]
        Позначимо $p_{kj} = \P\left\{\xi_1 = x_k, \xi_2 = y_j\right\}$.
        
        $\E\xi_1\xi_2 = \sum\limits_k \sum\limits_j x_k y_j p_{kj} = \sum\limits_k \sum\limits_j x_k y_j p_k p_j = \left( \sum\limits_k x_k p_k\right)\cdot \left( \sum\limits_j y_j p_j\right) = \E\xi_1 \cdot \E\xi_2$.
    \end{proof}
\end{enumerate}
\begin{remark}
    Властивості 3 та 4 для НВВ буде доведено в темі <<Функції випадкових аргументів>> (ст. \pageref{proof:expectation}).
\end{remark}

\begin{example}
    \begin{enumerate}
        \item Обчислити математичне сподівання двох ДВВ:
        
        \begin{center}
            \begin{tabular}{c c}
                \begin{tabular}{|c|c|c|}
                    \hline
                    $\xi_1$ & $-1$ & $1$ \\ 
                    \hline
                    $p$ & $1/2$ & $1/2$ \\
                    \hline
                \end{tabular} &
                \begin{tabular}{|c|c|c|}
                    \hline
                    $\xi_2$ & $-100$ & $100$ \\ 
                    \hline
                    $p$ & $1/2$ & $1/2$ \\
                    \hline
                \end{tabular}
            \end{tabular}
        \end{center}
        $\E\xi_1 = -1\cdot \frac{1}{2} + 1\cdot \frac{1}{2} = 0, \E\xi_2 = -100\cdot \frac{1}{2} + 100\cdot \frac{1}{2} = 0$.
        
        Цей приклад показує, що попри однакове значення математичного сподівання, можливі значення цих ДВВ знаходяться на різній відстані від нього,
        тому знання математичного сподівання ще не дає повного уявлення про закон розподілу випадкової величини.        
        \item Обчислити математичне сподівання НВВ, розподіленої за законом Сімпсона.

        \begin{tabular}{c c}
            \begin{tikzpicture}[baseline={(current bounding box.center)}, yscale=1]
                \pgfmathsetmacro{\a}{0.7}
                \draw [->] (-2, 0) -- (2, 0);
                \draw [->] (0, -0.1) -- (0, 1.7);
                \draw [ultra thick] (-2, 0) -- (-\a, 0);
                \draw [ultra thick] (\a, 0) -- (2, 0);
                \draw [ultra thick] (-\a, 0) -- (0, 1/\a);
                \draw [ultra thick] (\a, 0) -- (0, 1/\a);
                \node [below] at (2, 0) {$x$};
                \node [left] at (0, 1.6) {$f_\xi(x)$};
                \node [below] at (\a, 0) {$a$};
                \node [below] at (-\a, 0) {$-a$};
                \node [right] at (0, 1/\a) {$\frac{1}{a}$};
            \end{tikzpicture} &
            $f_\xi(x) = \begin{cases}
                \frac{1}{a} \left(1 - \frac{|x|}{a}\right), & |x| \leq a \\
                0, & |x| > a
            \end{cases} \;(a>0)$
        \end{tabular}
        
        $\E\xi = \int\limits_{-\infty}^{+\infty} x f_\xi(x)dx = \int\limits_{-a}^a \frac{x}{a}\left(1 - \frac{|x|}{a}\right)dx = 0$, 
        оскільки інтегрується непарна функція по симетричному проміжку.
    \end{enumerate}
\end{example}

\subsection{Дисперсія}
Як було показано, випадкові величини можуть мати однакові математичні сподівання, але відхилення значень цих випадкових величин
від нього може бути різним. Через це треба ввести ще одну числову характеристику,
що описуватиме це.
\begin{definition}
    \emph{Дисперсією} випадкової величини називається 
    \begin{equation}\label{eq:d_xi}
        \D\xi = \E\left(\xi-\E\xi\right)^2 = \int\limits_{-\infty}^{+\infty} \left(x-\E\xi\right)^2 dF_\xi(x) = \begin{cases}
            \sum\limits_{k=1}^{n(\infty)} \left(x_k-\E\xi\right)^2 \P\left\{\xi = x_k\right\}, & \xi \text{ --- ДВВ} \\
            \int\limits_{-\infty}^{+\infty} \left(x-\E\xi\right)^2 f_\xi(x)dx, & \xi \text{ --- НВВ}
        \end{cases}
    \end{equation}
\end{definition}
Дисперсія --- характеристика розсіювання випадкової величини навколо свого математичного сподівання.
Фізична інтерпретація --- момент інерції маси системи точок або стрижня (як і у випадку інтерпретації математичного сподівання)
відносно свого центру мас.
\begin{remark}
    Дисперсія визначена тоді і тільки тоді, коли математичне сподівання існує та є скінченним.
    Оскільки в \eqref{eq:d_xi} і члени ряду, і підінтегральна функція є невід'ємними, то
    дисперсія може бути або скінченною, або рівною $+\infty$.
\end{remark}

\begin{example}
    \begin{enumerate}
        \item Обчислити дисперсії двох ДВВ:
        
        \begin{center}
            \begin{tabular}{c c}
                \begin{tabular}{|c|c|c|}
                    \hline
                    $\xi_1$ & $-1$ & $1$ \\ 
                    \hline
                    $p$ & $1/2$ & $1/2$ \\
                    \hline
                \end{tabular} &
                \begin{tabular}{|c|c|c|}
                    \hline
                    $\xi_2$ & $-100$ & $100$ \\ 
                    \hline
                    $p$ & $1/2$ & $1/2$ \\
                    \hline
                \end{tabular}
            \end{tabular}
        \end{center}
        
        $\D\xi_1 = (-1)^2\cdot \frac{1}{2} + 1^2\cdot \frac{1}{2} = 1, \D\xi_2 = (-100)^2\cdot \frac{1}{2} + 100^2\cdot \frac{1}{2} = 10000$.
        Тут вже видно різницю між двома випадковими величинами: значення $\xi_2$ розташовані далеко від математичного сподівання і тому дисперсія більша, ніж у $\xi_1$.
        \item Обчислити дисперсію НВВ, розподіленої за законом Сімпсона.
        
        Відповідна щільність розподілу має вигляд $f_\xi(x) = \begin{cases}
            \frac{1}{a} \left(1 - \frac{|x|}{a}\right), & |x| \leq a \\
            0, & |x| > a
        \end{cases}$.
        Математичне сподівання вже було обчислено, воно рівне 0. Тому $\D\xi = \int\limits_{-\infty}^{+\infty} \left(x-\E\xi\right)^2 f_\xi(x)dx =
        \int\limits_{-a}^a \frac{x^2}{a}\left(1 - \frac{|x|}{a}\right)dx = \frac{2}{a}\int\limits_0^a {x^2}\left(1 - \frac{x}{a}\right)dx =
        \frac{2}{a}\cdot \left( \frac{x^3}{3} - \frac{x^4}{4a}\right)\Bigr\vert_0^a = \frac{2}{a} \cdot \left( \frac{a^3}{3} - \frac{a^3}{4}\right) = \frac{a^2}{6}$.
    \end{enumerate}
\end{example}

\noindent \textbf{Властивості дисперсії:}
\begin{enumerate}
    \item $\D\xi \geq 0$, $\sqrt{\D\xi} = \sigma_\xi$ --- \emph{середньоквадратичне (стандартне) відхилення}.
    \item $\D\xi = 0 \Leftrightarrow \xi = const$.
    \begin{proof}
        В одну сторону: $\D c = \E(c - \E c)^2 = \E(c - c)^2 = 0$. В іншу: якщо $\xi$ --- ДВВ,
        то з $\D\xi = 0$ випливає, що $\left(x_k-\E\xi\right)^2 = 0$ для всіх $k$,
        тому для всіх $k$ $x_k = \E \xi$ --- $\xi$ приймає лише одну значення.
        Для НВВ аналогічно випливає, що $\left(x-\E\xi\right)^2 f_\xi(x) \equiv 0$. Оскільки $\left(x-\E\xi\right)^2$
        рівне $0$ лише в точці $x=\E\xi$, то $f_\xi(x)$ може бути відмінним від 0 лише в цій точці. Але тоді $f_\xi(x)$ не може бути щільністю,
        тому що площа під кривою розподілу має бути рівною 1. Це свідчить про дискретність $\xi$,
        і тому за вже доведеним $\xi$ приймає лише одне значення.
    \end{proof}
    \item Для обчислення дисперсії більш зручною є наступна формула:

    $\D\xi = \E(\xi^2 - 2\xi \E\xi +(\E\xi)^2) = \E\xi^2 - 2(\E\xi)^2 + (\E\xi)^2 = \E\xi^2 - (\E\xi)^2$. 
    \item Для сталої $c$: $\D(c\cdot \xi) = c^2\cdot \D\xi$, оскільки $\D(c\cdot \xi) = \E\left(c\cdot\xi-\E(c\cdot\xi\right))^2 = \E\left(c\cdot\xi - c\cdot \E\xi\right)^2 = c^2 \cdot \E\left(\xi-\E\xi\right)^2 = c^2 \cdot \D\xi$.
    \item Для сталої $c$: $\D(\xi + c) = \E(\xi + c - \E(\xi + c))^2 = \E(\xi + c - \E\xi - \E c) = \E(\xi - \E\xi)^2 = \D\xi$.
\suspend{enumerate}
\begin{definition}
    \emph{Центрованою} випадковою величиною, що відповідає випадковій величині $\xi$, називається 
    випадкова величина $\mathring{\xi} = \xi - \E\xi$. Для неї $\E\mathring{\xi} = 0$ та 
    $\E\mathring{\xi^2} = \D\xi$.
\end{definition}
\resume{enumerate}
    \item Якщо випадкові величини $\xi_1$ та $\xi_2$ --- незалежні, то
    $\D\left(\xi_1 + \xi_2\right) = \D\xi_1 + \D\xi_2$.
    \begin{proof}
        $\D\left(\xi_1 + \xi_2\right) = \E((\xi_1 + \xi_2) - \E(\xi_1 + \xi_2))^2 = \E(\mathring{\xi}_1 + \mathring{\xi}_2)^2 =
        \E\mathring{\xi}_1^2 + 2\E\mathring{\xi}_1\mathring{\xi}_2 + \E\mathring{\xi}_2^2 = 
        \E\mathring{\xi}_1^2 + 2\E\mathring{\xi}_1 \E\mathring{\xi}_2 + \E\mathring{\xi}_2^2 =
        \E\mathring{\xi}_1^2 + \E\mathring{\xi}_2^2 = \D\xi_1 + \D\xi_2$.
    \end{proof}
\end{enumerate}

\subsection{Моменти випадкової величини}
\begin{definition}
    \emph{Початковим моментом k-того порядку} ($k \in \mathbb{N}$) ВВ $\xi$ називається
    математичне сподівання $k$-того степеня $\xi$:
    \begin{equation}\label{eq:e_alpha_k}
        \alpha_k = \E\xi^k = \int\limits_{-\infty}^{+\infty} x^k dF_\xi(x) = \begin{cases}
            \sum\limits_{m=1}^{n(\infty)} x_m^k \P\left\{\xi = x_m\right\}, & \xi \text{ --- ДВВ} \\
            \int\limits_{-\infty}^{+\infty} x^k f_\xi(x)dx, & \xi \text{ --- НВВ}
        \end{cases}
    \end{equation}
\end{definition}
\begin{definition}
    \emph{Центральним моментом k-того порядку} ($k \in \mathbb{N}$) ВВ $\xi$ називається
    \begin{equation}\label{eq:d_beta_k}
        \beta_k = \E\mathring{\xi}^k = \E\left(\xi-\E\xi\right)^k = \int\limits_{-\infty}^{+\infty} \left(x-\E\xi\right)^k dF_\xi(x) = \begin{cases}
            \sum\limits_{m=1}^{n(\infty)} \left(x_m-\E\xi\right)^k \P\left\{\xi = x_m\right\}, & \xi \text{ --- ДВВ} \\
            \int\limits_{-\infty}^{+\infty} \left(x-\E\xi\right)^k f_\xi(x)dx, & \xi \text{ --- НВВ}
        \end{cases}
    \end{equation}
\end{definition}
Математичне сподівання є початковим моментом першого порядку, а дисперсія --- центральним моментом другого порядку.
Узагальненням формули для дисперсії
$\D\xi = \beta_2 = \E\xi^2 - (\E\xi)^2 = \alpha_2 - \alpha_1^2$ є наступна формула,
що дозволяє виразити будь-який центральний момент через початкові моменти:
\begin{equation}
    \beta_k = \E\left(\xi-\E\xi\right)^k = \E\left( \sum\limits_{j=0}^k C_k^j \xi^j (-1)^{k-j} (\E\xi)^{k-j}\right) =
\sum\limits_{j=0}^k (-1)^{k-j} C_k^j \E\xi^j (\E\xi)^{k-j}
\end{equation}
\begin{exercise}
    Виразити через початкові моменти $\beta_3$ та $\beta_4$.
\end{exercise}
Розглядаються також \emph{абсолютні початкові моменти} $\E\vert \xi \vert ^k$, \emph{абсолютні центральні моменти} $\E\vert \xi - \E\xi \vert ^k$
та \emph{факторіальні моменти} $\gamma_k = \E\left( \xi (\xi - 1) ... (\xi - k + 1)\right)$.
\begin{remark}
    Якщо за змістом задачі випадкова величина має якісь одиниці вимірювання (наприклад, метри чи кілограми),
    то моменти $k$-того порядку мають відповідний порядок одиниць вимірювання.
    Наприклад, якщо деяка випадкова величина задає похибку вимірювання в міліметрах,
    то математичне сподівання теж буде у міліметрах, а дисперсія --- в квадратних міліметрах
    (але середньоквадратичне відхилення --- просто у міліметрах).
\end{remark}


\subsection{Мода та медіана випадкової величини}
\begin{definition}
    \emph{Модою} $\mathrm{Mo}\xi$ називається абсциса точки максимуму щільності 
    розподілу у випадку НВВ та значення, ймовірність 
    появи якого є найбільшим, у випадку ДВВ.

    \emph{Унімодальний закон} --- такий, що має лише одну моду. 
    \emph{Полімодальний закон} --- такий, що має декілька мод.
    \emph{Антимодальний закон} --- такий, що не має моди.
\end{definition}
\begin{definition}
    \emph{Медіаною} $\mathrm{Me}\xi$ НВВ $\xi$ називається така точка $x_0$, для якої 
    $\P\left\{\xi < x_0\right\} = \P\left\{\xi \geq x_0\right\} 
    = \frac{1}{2} \Leftrightarrow F_\xi(x_0) = \frac{1}{2}$.
    Для ДВВ $\xi$ медіаною є така точка $x_0$, для якої одночасно
    виконуються нерівності $\P\left\{\xi < x_0\right\} \geq \frac{1}{2}$ та
    $\P\left\{\xi \geq x_0\right\} \geq \frac{1}{2}$.
    Медіана є окремим випадком \emph{квантиля}.
\end{definition}
\begin{definition}
    Точка $x_0$ називається \emph{квантилем q-го порядку}, якщо $F_\xi(x_0) = q$ 
    для НВВ $\xi$ або $F_\xi(x_0) \geq q$ та $F_\xi(x_0) \geq 1-q$ для ДВВ.
\end{definition}
\begin{example}
    Знайти моду та медіану НВВ, розподіленої за законом Релея.
    
    \begin{tabular}{c c}
        \begin{tikzpicture}[baseline={(current bounding box.center)}, yscale=2]
            \pgfmathsetmacro{\s}{1}
            \draw [->] (-0.5, 0) -- (5, 0);
            \draw [->] (0, -0.1) -- (0, 1);
            \node [below] at (5, 0) {$x$};
            \node [below left] at (0, 1) {$f_\xi(x)$};
            \draw [domain=0:5, smooth, variable = \x, ultra thick] plot ({\x}, {((\x/(\s^2)) * exp(-(\x)^2/(2*\s^2))});
            \draw [ultra thick] (-0.5, 0) -- (0, 0);
        \end{tikzpicture} &
        $f_\xi(x) = \begin{cases}
            \frac{x}{\sigma^2} e^{-\frac{x^2}{2\sigma^2}}, & x \geq 0 \\
            0, & x < 0
        \end{cases} (\sigma > 0)$ 
    \end{tabular}

    Для визначення моди знайдемо максимум $f_\xi(x)$ при $x\geq 0$.
    $f'_\xi(x) = \frac{\sigma^2 - x^2}{\sigma^4} e^{-\frac{x^2}{2\sigma^2}} = 0$ при $x=\sigma$.
    $f''_\xi(x) = \left( -\frac{3x}{\sigma^4} + \frac{x^3}{\sigma^6}\right) e^{-\frac{x^2}{2\sigma^2}}$,
    $f''_\xi(\sigma) = -\frac{2}{\sigma^3} e^{-\frac{1}{2}} < 0$. Отже, в точці $x=\sigma$ дійсно максимум $f_\xi(x)$, тому $\mathrm{Mo}\xi = \sigma$.

    Медіану знайдемо з рівності $\P\left\{\xi < x_0\right\} = 0.5$. $\P\left\{\xi < x_0\right\} = \int\limits_{-\infty}^{x_0}f_\xi(x)dx = 
    \int\limits_0^{x_0} \frac{x}{\sigma^2} e^{-\frac{x^2}{2\sigma^2}} dx = 1 - e^{-\frac{x_0^2}{2\sigma^2}}$.
    З рівняння $1 - e^{-\frac{x_0^2}{2\sigma^2}} = 0.5$ знаходимо $x_0 = \mathrm{Me}\xi = \sigma \sqrt{2\ln{2}}$.
\end{example}

\subsection{Асиметрія та ексцес випадкової величини}
\begin{definition}
    \emph{Асиметрією} випадкової величини $\mathrm{As}\xi$ називається безрозмірна 
    числова характеристика, що дорівнює $\frac{\beta_3}{\sigma_\xi^3} = 
    \frac{\E(\xi - \E\xi)^3}{(\D\xi)^{3/2}}$. 
    Ця характеристика показує порушення чи наявність симетрії кривої розподілу відносно математичного сподівання.
\end{definition}
\begin{center}
    \begin{tabular}{c c c}
        \begin{tikzpicture}[yscale = 1.5]
            \draw [->] (-0.5, 0) -- (3.8, 0);
            \draw [->] (0, -0.1) -- (0, 1);
            \draw [domain=-0.5:3.7, smooth, variable = \x, ultra thick] plot ({\x}, {0.797884560803 * exp(-2*(\x-1)^2)});
            \draw [dashed] (1, 0) -- (1, 0.797884560803);
            \node [below] at (1, 0) {$\E\xi$};
        \end{tikzpicture} &
        \begin{tikzpicture}[yscale = 1.5]
            \draw [->] (-0.5, 0) -- (3.2, 0);
            \draw [->] (0, -0.1) -- (0, 1);
            \draw [ultra thick] (-0.5, 0) -- (0, 0);
            \draw [ultra thick] (3, 0) -- (3.1, 0);
            \draw [domain=0:3, smooth, variable = \x, ultra thick] plot ({\x}, {10 * (\x/3)^4 * (1-(\x/3))});
            \draw [dashed] (2.1428, 0) -- (2.1428, 0.744);
            \node [below] at (2.1428, 0) {$\E\xi$};
        \end{tikzpicture} &
        \begin{tikzpicture}[yscale = 1.5]
            \draw [->] (-0.5, 0) -- (3.2, 0);
            \draw [->] (0, -0.1) -- (0, 1);
            \draw [ultra thick] (-0.5, 0) -- (0, 0);
            \draw [domain=0:3, smooth, variable = \x, ultra thick] plot ({\x}, {10 * (\x/3) * (1-(\x/3))^4});
            \draw [dashed] (0.8571, 0) -- (0.8571, 0.744);
            \node [below] at (0.8571, 0) {$\E\xi$};
        \end{tikzpicture} \\
        $As\xi = 0$ & $As\xi < 0$ & $As\xi > 0$ 
    \end{tabular}
\end{center}

\begin{definition}
    \emph{Ексцесом} випадкової величини $\mathrm{Ex}\xi$ називається безрозмірна 
    числова характеристика, що дорівнює $\mathrm{Ex}\xi = \frac{\beta_4}{\sigma_\xi^4} - 3 = 
    \frac{\E(\xi - \E\xi)^4}{(\D\xi)^{2}} - 3$.
    Ця характеристика показує, наскільки швидко крива розподілу 
    прямує до точки максимуму.
\end{definition}
\begin{center}
    \begin{tikzpicture}[yscale = 2]
        \draw [->] (-1, 0) -- (7, 0);
        \draw [->] (0, -0.1) -- (0, 1);
        \draw [domain=-1:7, smooth, variable = \x, ultra thick] plot ({\x}, {0.398942280401 * exp(-(\x-3)^2/2)});
        \draw [domain=-1:7, smooth, variable = \x, ultra thick] plot ({\x}, {0.797884560803 * exp(-2*(\x-3)^2)});
        \draw [domain=-1:7, smooth, variable = \x, ultra thick] plot ({\x}, {0.199471140201 * exp(-(\x-3)^2/8)});
        \draw [->] (2, 0.7) -- (2.4, 0.6);
        \node [left] at (2, 0.7) {${Ex}\xi > 0$};
        \draw [->] (3.7, 0.7) -- (3, 0.45);
        \node [right] at (3.7, 0.7) {${Ex}\xi = 0$};
        \draw [->] (5, 0.4) -- (3, 0.25);
        \node [right] at (5, 0.4) {${Ex}\xi < 0$};
    \end{tikzpicture}
\end{center}

\subsection{Генератриса (твірна функція) ДВВ}
\begin{definition}
    Нехай $\xi$ --- ДВВ, що приймає цілі невід'ємні значення. 
    \emph{Генератрисою (твірною функцією)} цієї ДВВ називається функція комплексного аргументу
    \begin{equation}\label{eq:gen_func}
        G_\xi(z) = \sum\limits_{k=0}^{\infty} \P\left\{\xi = k\right\} z^k
    \end{equation}
\end{definition}
\noindent \textbf{Властивості генератриси:}
\begin{enumerate}
    \item Відповідний ряд рівномірно збігається принаймні в колі $|z|\leq 1$.
    \begin{proof}
        Випливає з того, що $\left| \sum\limits_{k=0}^{\infty} \P\left\{\xi = k\right\} z^k \right| \leq \sum\limits_{k=0}^{\infty} \P\left\{\xi = k\right\} = 1$.
    \end{proof}
    \item Генератриса --- аналітична функція в колі $|z|\leq 1$ (як наслідок властивості 1).
    \item За генератрисою можна відновити розподіл $\xi$:
    $\P\left\{\xi = k\right\} = \frac{1}{k!}\cdot G_\xi^{(k)}(0)$.
    \item За допомогою генератриси можна знайти факторіальні моменти $\xi$.
    \begin{proof}
        За означенням факторіальний момент $\gamma_m = \E\left( \xi (\xi - 1) ... (\xi - m + 1)\right)$.
        Для ДВВ, що розглядаються, він рівний $\sum\limits_{k=0}^{\infty} k(k-1)(k-2)...(k-m+1) \P\left\{\xi = k\right\}$.
        З іншого боку, $G'_\xi(z) = \sum\limits_{k=0}^{\infty} k \P\left\{\xi = k\right\} z^{k-1}$,
        $G''_\xi(z) = \sum\limits_{k=0}^{\infty} k(k-1) \P\left\{\xi = k\right\} z^{k-2}$ і так далі,
        $G^{(m)}_\xi(z) = \sum\limits_{k=0}^{\infty} k(k-1)(k-2)...(k-m+1) \P\left\{\xi = k\right\} z^{k-m}$.
        Отже, $\gamma_m = G^{(m)}_\xi(1)$.
    \end{proof}
\end{enumerate}
Як наслідок властивості 4, можна також знайти математичне сподівання та дисперсію.
Факторіальний та початковий моменти першого порядку збігаються, тому $\E\xi = G'_\xi(1)$.
Дисперсію знайдемо за допомогою другого факторіального моменту: $\gamma_2 = \E\left( \xi (\xi - 1)\right) = \E\xi^2 - \E\xi = G''_\xi(1)$,
звідки $\E\xi^2 = G''_\xi(1) + G'_\xi(1)$. 
Отже, $\D\xi = \E\xi^2 - (\E\xi)^2 = G''_\xi(1) + G'_\xi(1) - \left( G'_\xi(1)\right)^2$.
\begin{exercise}
    Виразити через похідні генератриси центральні моменти 3-го та 4-го порядків.
\end{exercise}
        % !TEX root = ../main.tex
\section{Деякі закони розподілу випадкових величин}

\subsection{Біноміальний розподіл}
\noindent\textbf{Означення:}
    ДВВ $\xi$ розподілена за \emph{біноміальним законом}, 
    якщо набуває значень $0,1,...,n$ з ймовірностями \begin{equation}
        \P\left\{\xi = k\right\} = C_n^k p^k q^{n-k}, q = 1 - p
    \end{equation}
\textbf{Коротке позначення:} $\xi \sim \mathrm{Bin}(n, p)$.
    $n$ і $p$ --- параметри закону, $n\in \mathbb{N}$, $p\in (0;1)$.

Окремим випадком біноміального закону є \emph{розподіл Бернуллі}: $\xi \sim \mathrm{Bin}(1, p)$.
За законом Бернуллі розподілені випадкові величини-індикатори $\mathds{1}_A = \begin{cases}
    0, & \text{подія }A\text{ не відбулась}\\ 1, & \text{подія }A\text{ відбулась}
\end{cases}$ при $\P(A) = p$.

\noindent\textbf{Ряд розподілу:}
\begin{center}
    \begin{tabular}{|c|c|c|c|c|c|c|}
        \hline
        $\xi$ & $0$ & $1$ & $2$ & $...$ & $n-1$ & $n$ \\
        \hline
        $p$ & $q^n$ & $C_n^1 pq^{n-1}$ & $C_n^2 p^2 q^{n-2}$ & $...$ & $C_n^{n-1}p^{n-1}q$ & $p^n$ \\
        \hline
    \end{tabular}
\end{center}

\noindent\textbf{Полігон розподілу:} приклад для $n=5$, $p=0.6$.
\begin{center}
    \begin{tikzpicture}
        \pgfmathsetmacro{\s}{5};
        \pgfmathsetmacro{\p}{0.6};
        \pgfmathsetmacro{\q}{1-\p};
        \pgfmathsetmacro{\n}{5};
        \draw [->] (-1,0) -- (\n+1, 0);
        \draw [->] (0, -0.1*\s) -- (0, 0.4*\s);
        \draw [thick] (0,\q^\n*\s) -- (1, {\n*\p*\q^(\n-1)*\s}) -- (2, {(\n*(\n-1)/2)*\p^2*\q^(\n-2)*\s});
        \draw [thick] (2, {(\n*(\n-1)/2)*\p^2*\q^(\n-2)*\s}) -- (3, {(\n*(\n-1)*(\n-2)/6)*\p^3*\q^(\n-3)*\s});
        \draw [thick] (3, {(\n*(\n-1)*(\n-2)/6)*\p^3*\q^(\n-3)*\s}) -- (4, {\n*\p^4*\q^(\n-4)*\s});
        \draw [thick] (4, {\n*\p^4*\q^(\n-4)*\s}) -- (5, \p^\n*\s);
        \draw [fill] (0,\q^\n) circle [radius = 0.05];
        \draw [fill] (1, {\n*\p*\q^(\n-1)*\s}) circle [radius = 0.05];
        \draw [fill] (2, {(\n*(\n-1)/2)*\p^2*\q^(\n-2)*\s}) circle [radius = 0.05];
        \draw [fill] (3, {(\n*(\n-1)*(\n-2)/6)*\p^3*\q^(\n-3)*\s}) circle [radius = 0.05];
        \draw [fill] (4, {\n*\p^4*\q^(\n-4)*\s}) circle [radius = 0.05];
        \draw [fill] (5, \p^\n*\s) circle [radius = 0.05];
        \node [below left] at (0, 0) {0};
        \foreach \k in {1,...,\n}:
            \node [below] at (\k, 0) {\k};
        \node [below] at (\n+1, 0) {$x$};
        \node [left] at (0, 0.4*\s) {$p$};
    \end{tikzpicture}
\end{center}

\noindent\textbf{Функція розподілу:} приклад для $n=5$, $p=0.6$.
\begin{center}
    \begin{tabular}{c c}
        $
            F_\xi(x) = \begin{cases}
                0, & x \leq 0 \\
                q^n, & 0 < x \leq 1 \\
                q^n + npq^{n-1}, & 1 < x \leq 2 \\
                \dots \\
                1, & x > n
            \end{cases}
        $ &
        \begin{tikzpicture}[baseline={(current bounding box.center)}, yscale=2.5]
            \pgfmathsetmacro{\p}{0.6};
            \pgfmathsetmacro{\q}{1-\p};
            \pgfmathsetmacro{\n}{5};
            \draw [->] (-1,0) -- (\n+1, 0);
            \draw [->] (0, -0.1) -- (0, 1.2);
            \draw [ultra thick] (-1, 0) -- (0,0);
            \draw [ultra thick] [<-] (0,\q^\n) -- (1, \q^\n);
            \draw [ultra thick] [<-] (1, {\q^\n + \n*\p*\q^(\n-1)}) -- (2, {\q^\n + \n*\p*\q^(\n-1)});
            \draw [ultra thick] [<-] (2, {\q^\n + \n*\p*\q^(\n-1) + (\n*(\n-1)/2)*\p^2*\q^(\n-2)}) -- (3, {\q^\n + \n*\p*\q^(\n-1) + (\n*(\n-1)/2)*\p^2*\q^(\n-2)});
            \draw [ultra thick] [<-] (3, {1 - (\q^\n + \n*\p*\q^(\n-1) + (\n*(\n-1)/2)*\p^2*\q^(\n-2))}) -- (4, {1 - (\q^\n + \n*\p*\q^(\n-1) + (\n*(\n-1)/2)*\p^2*\q^(\n-2))});
            \draw [ultra thick] [<-] (4, {1 - (\q^\n + \n*\p*\q^(\n-1))}) -- (5, {1 - (\q^\n + \n*\p*\q^(\n-1))});
            \draw [ultra thick] [<-] (5, 1) -- (6, 1);
            \draw [dashed] (1, 0) -- (1, {\q^\n + \n*\p*\q^(\n-1)});
            \draw [dashed] (2, 0) -- (2, {\q^\n + \n*\p*\q^(\n-1) + (\n*(\n-1)/2)*\p^2*\q^(\n-2)});
            \draw [dashed] (3, 0) -- (3, {1 - (\q^\n + \n*\p*\q^(\n-1) + (\n*(\n-1)/2)*\p^2*\q^(\n-2))});
            \draw [dashed] (4, 0) -- (4, {1 - (\q^\n + \n*\p*\q^(\n-1))});
            \draw [dashed] (5, 0) -- (5, 1);
            \node [below left] at (0, 0) {0};
            \foreach \k in {1,...,\n}:
                \node [below] at (\k, 0) {\k};
            \draw [dashed] (0, 1) -- (\n, 1);
            \node [left] at (0, 1) {1};
            % \node [right] [align=center] at (3.5, 0.2) {Приклад для \\ $n = \n, p = \p$};
            \node [below] at (\n+1, 0) {$x$};
            \node [left] at (0, 1.2) {$F_\xi(x)$};
        \end{tikzpicture}
    \end{tabular}
\end{center}

Для дослідження числових характеристик скористаємося генератрисою розподілу:

$G_\xi(z) = \sum\limits_{k=0}^{n} \P\left\{\xi = k\right\} z^k = \sum\limits_{k=0}^{n} C_n^k p^k q^{n-k} z^k = (pz+q)^n$,
звідки 
$G'_\xi(z) = np(pz+q)^{n-1}$, $G''_\xi(z) = n(n-1)p^2(pz+q)^{n-2}$.
Тому 
$G'_\xi(1) = np(p+q)^{n-1} = np$, $G''_\xi(1) = n(n-1)p^2 = n^2p^2 - np^2$, $G''_\xi(1) + G'_\xi(1) - \left( G'_\xi(1)\right)^2 = n^2p^2 - np^2 + np - n^2p^2 = np(1-p) = npq$.
Отже, знайдено значення математичного сподівання та дисперсії.

\begin{samepage}
\noindent\textbf{Числові характеристики:}
\begin{enumerate}
    \item $\E\xi = np$.
    \item $\D\xi = npq$, $\sigma_\xi = \sqrt{npq}$.
    \item $\Mo\xi = \begin{cases}
        \left[np+p\right], & \text{якщо } np+p \text{ не ціле}\\
        np+p, np-q, & \text{якщо } np+p \text{ ціле}
    \end{cases}$ --- як найбільш ймовірна кількість успіхів у схемі Бернуллі.
    \item $\Me\xi$ --- одне зі значень $\left[np\right] - 1$, $\left[np\right]$, $\left[np\right] + 1$.
\end{enumerate}
\end{samepage}

\noindent\textbf{Застосування:} якщо проводиться $n$ незалежних випробувань з ймовірністю успіху $p$, 
то $\xi \sim \mathrm{Bin}(n, p)$ задає кількість успіхів.

\subsection{Геометричний розподіл}
\noindent\textbf{Означення:}
    ДВВ $\xi$ розподілена за \emph{геометричним законом}, 
    якщо набуває значень $1,2,3,..$ з ймовірностями \begin{equation}
        \P\left\{\xi = k\right\} = pq^{k-1}, q = 1 - p
    \end{equation}
    \textbf{Коротке позначення:} $\xi \sim \mathrm{Geom}(p)$.
    $p$ --- параметр закону, $p\in (0;1)$.    

\noindent\textbf{Ряд розподілу:}
\begin{center}
    \begin{tabular}{|c|c|c|c|c|c|c|}
        \hline
        $\xi$ & $1$ & $2$ & $3$ & $...$ & $k$ & $...$ \\
        \hline
        $p$ & $p$ & $pq$ & $pq^2$ & $...$ & $pq^{k-1}$ & $...$ \\
        \hline
    \end{tabular}
\end{center}

\noindent\textbf{Полігон розподілу:} приклад для $p=0.5$.
\begin{center}
    \begin{tikzpicture}[baseline={(current bounding box.center)}]
        \pgfmathsetmacro{\s}{4};
        \pgfmathsetmacro{\p}{0.5};
        \pgfmathsetmacro{\q}{1-\p};
        \pgfmathsetmacro{\n}{5};
        \draw [->] (-0.3,0) -- (\n+1, 0);
        \draw [->] (0, -0.1*\s) -- (0, 0.6*\s);
        \foreach \k in {1,...,\n}:
            \draw [thick] (\k, {\p*\q^(\k - 1)*\s}) -- (\k+1, {\p*\q^(\k)*\s});
        \foreach \k in {1,...,\n}:
            \draw [fill] (\k, {\p*\q^(\k - 1)*\s}) circle [radius = 0.05];
        \node [below left] at (0, 0) {0};
        \foreach \k in {1,...,\n}:
            \node [below] at (\k, 0) {\k};
        \node [below] at (\n+1, 0) {$x$};
        \node [left] at (0, 0.6*\s) {$p$};
    \end{tikzpicture}
\end{center}

\noindent\textbf{Функція розподілу:} приклад для $p=0.5$.

\begin{tabular}{c c}
    $
        F_\xi(x) = \begin{cases}
            0, & x \leq 1 \\
            p, & 1 < x \leq 2 \\
            p+pq, & 2 < x \leq 3 \\
            \dots \\
            \sum\limits_{m=0}^{k-1}pq^m = 1-q^k, & k < x \leq k+1 \\
            \dots
        \end{cases}
    $ &
    \begin{tikzpicture}[baseline={(current bounding box.center)}, yscale=2.5, xscale=0.88]
        \pgfmathsetmacro{\p}{0.5};
        \pgfmathsetmacro{\q}{1-\p};
        \pgfmathsetmacro{\n}{5};
        \draw [->] (-0.3,0) -- (\n+1, 0);
        \draw [->] (0, -0.1) -- (0, 1.2);
        \draw [ultra thick] (-0.3, 0) -- (1,0);
        \foreach \k in {1,...,\n}:
            \draw [ultra thick] [<-] (\k, 1-\q^\k) -- (\k+1, 1-\q^\k);
        \node [below left] at (0, 0) {0};
        \foreach \k in {1,...,\n}:
            \node [below] at (\k, 0) {\k};
        \draw [dashed] (0, 1) -- (\n+1, 1);
        \node [left] at (0, 1) {1};
        \node [below] at (\n+1, 0) {$x$};
        \node [left] at (0, 1.2) {$F_\xi(x)$};
    \end{tikzpicture}
\end{tabular}

Для дослідження числових характеристик скористаємося генератрисою розподілу:

$G_\xi(z) = \sum\limits_{k=1}^{\infty} \P\left\{\xi = k\right\} z^k = \sum\limits_{k=1}^{\infty} pq^{k-1} z^k = \frac{pz}{1-qz}$, якщо $\left| qz\right|<1$, звідки
$G'_\xi(z) = \frac{p}{(1-qz)^2}$, $G''_\xi(z) = \frac{2pq}{(1-qz)^3}$. Тому
$G'_\xi(1) = \frac{1}{p}$, $G''_\xi(1) = \frac{2q}{p^2}$, $G''_\xi(1) + G'_\xi(1) - \left( G'_\xi(1)\right)^2 = \frac{2q}{p^2} + \frac{1}{p} - \frac{1}{p^2} = \frac{2q+p-1}{p^2} = \frac{q}{p^2}$.
Отже, знайдено значення математичного сподівання та дисперсії.

\noindent\textbf{Числові характеристики:}
\begin{enumerate}
    \item $\E\xi = \frac{1}{p}$.
    \item $\D\xi = \frac{q}{p^2}$, $\sigma_\xi = \frac{\sqrt{q}}{p}$.
    \item $\Mo\xi = 1$.
    \item $\Me\xi = \left[ \frac{-1}{\log_2(1-p)}\right]$.
\end{enumerate}

\noindent\textbf{Застосування:} якщо проводяться незалежні випробування з ймовірністю успіху $p$ до першого успішного,
    то $\xi$, що задає кількість проведених випробувань, має розподіл $\mathrm{Geom}(p)$.

\begin{exercise}
    Записати ряд розподілу, функцію розподілу, генетратрису та обчислити
    числові характеристики для іншого означення геометричного розподілу, 
    де $\xi$ приймає значення $0,1,2,...$ з ймовірностями $\P\left\{\xi = k\right\} = pq^k$.
\end{exercise}
Геометричний розподіл \textbf{<<не має пам'яті>>}: $\P\left\{\xi = n+m / \xi \geq n\right\} = \P\left\{\xi = m\right\}$.
Це означає, що кількість минулих <<невдач>> не впливає на кількість майбутніх <<невдач>>.
\begin{exercise}
    Довести цю властивість і те, що геометричний розподіл --- 
    єдиний дискретний розподіл, що <<не має пам'яті>>.
\end{exercise}
До геометричного закону можна звести \emph{закон розподілу Паскаля} $\xi \sim \mathrm{Pas}(a)$,
що задається $\P\left\{\xi = k\right\} = \frac{a^k}{(1+a)^{k+1}}, k = 0,1,2,...$ для $a>0$,
заміною $p=\frac{1}{a+1}$. Для нього $\E\xi = a$, $\D\xi = a + a^2$.

\subsection{Розподіл Пуассона}
\noindent\textbf{Означення:}
    ДВВ $\xi$ розподілена за \emph{законом Пуассона}, 
    якщо набуває значень $0,1,2,..$ з ймовірностями \begin{equation}
        \P\left\{\xi = k\right\} = \frac{a^k}{k!}e^{-a}
    \end{equation}
    \textbf{Коротке позначення:} $\xi \sim \mathrm{Poiss}(a)$.
    $a$ --- параметр закону, $a > 0$.

\noindent\textbf{Ряд розподілу:}
\begin{center}
    \begin{tabular}{|c|c|c|c|c|c|c|}
        \hline
        $\xi$ & $0$ & $1$ & $2$ & $...$ & $k$ & $...$ \\
        \hline
        $p$ & $e^{-a}$ & $ae^{-a}$ & $\frac{a^2}{2}e^{-a}$ & $...$ & $\frac{a^k}{k!}e^{-a}$ & $...$\\
        \hline
    \end{tabular}
\end{center}

\noindent\textbf{Полігон розподілу:} приклад для $a=2$.
\begin{center}
    \begin{tikzpicture}[baseline={(current bounding box.center)}]
        \pgfmathsetmacro{\s}{5};
        \pgfmathsetmacro{\a}{2};
        \pgfmathsetmacro{\n}{5};
        \draw [->] (-0.3,0) -- (\n+1, 0);
        \draw [->] (0, -0.1*\s) -- (0, 0.3*\s);
        \foreach \k in {0,...,\n}:
            \draw [thick] (\k, {(\a)^(\k)/factorial(\k)*e^(-\a)*\s}) -- (\k+1, {(\a)^(\k+1)/factorial(\k+1)*e^(-\a)*\s});
        \foreach \k in {0,...,\n}:
            \draw [fill] (\k, {(\a)^(\k)/factorial(\k)*e^(-\a)*\s}) circle [radius = 0.05];
        \node [below left] at (0, 0) {0};
        \foreach \k in {1,...,\n}:
            \node [below] at (\k, 0) {\k};
        \node [below] at (\n+1, 0) {$x$};
        \node [left] at (0, 0.3*\s) {$p$};
    \end{tikzpicture}
\end{center}

\noindent\textbf{Функція розподілу:} приклад для $a=2$.

\begin{tabular}{c c}
    $
        F_\xi(x) = \begin{cases}
            0, & x \leq 0 \\
            e^{-a}, & 0 < x \leq 1 \\
            e^{-a}+ae^{-a}, & 1 < x \leq 2 \\
            \dots \\
            \sum\limits_{m=0}^{k-1}\frac{a^m}{m!}e^{-a}& k-1 < x \leq k \\
            \dots
        \end{cases}
    $ &
    \begin{tikzpicture}[baseline={(current bounding box.center)}, yscale=2.5, xscale=0.88]
        \pgfmathsetmacro{\a}{2};
        \pgfmathsetmacro{\n}{5};
        \draw [->] (-0.3,0) -- (\n+1, 0);
        \draw [->] (0, -0.1) -- (0, 1.2);
        \draw [ultra thick] (-0.3, 0) -- (0,0);
        \draw [ultra thick] [<-] (0, {e^(-\a)}) -- (1, {e^(-\a)});
        \draw [ultra thick] [<-] (1, {e^(-\a)*(1 + \a)}) -- (2, {e^(-\a)*(1 + \a)});
        \draw [ultra thick] [<-] (2, {e^(-\a)*(1 + \a + \a^2/2)}) -- (3, {e^(-\a)*(1 + \a + \a^2/2)});
        \draw [ultra thick] [<-] (3, {e^(-\a)*(1 + \a + \a^2/2 + \a^3/6)}) -- (4, {e^(-\a)*(1 + \a + \a^2/2 + \a^3/6)});
        \draw [ultra thick] [<-] (4, {e^(-\a)*(1 + \a + \a^2/2 + \a^3/6 + \a^4/24)}) -- (5, {e^(-\a)*(1 + \a + \a^2/2 + \a^3/6 + \a^4/24)});
        \draw [ultra thick] [<-] (5, {e^(-\a)*(1 + \a + \a^2/2 + \a^3/6 + \a^4/24 + \a^5/120)}) -- (6, {e^(-\a)*(1 + \a + \a^2/2 + \a^3/6 + \a^4/24 + \a^5/120)});
        \node [below left] at (0, 0) {0};
        \foreach \k in {1,...,\n}:
            \node [below] at (\k, 0) {\k};
        \draw [dashed] (0, 1) -- (\n+1, 1);
        \node [left] at (0, 1) {1};
        \node [below] at (\n+1, 0) {$x$};
        \node [left] at (0, 1.2) {$F_\xi(x)$};
    \end{tikzpicture}
\end{tabular}

Для дослідження числових характеристик скористаємося генератрисою розподілу:

$G_\xi(z) = \sum\limits_{k=0}^{\infty} \P\left\{\xi = k\right\} z^k = e^{-a} \sum\limits_{k=0}^{\infty} \frac{a^k}{k!}z^k = e^{-a(1-z)}$, звідки
$G'_\xi(z) = ae^{az-a}$, $G''_\xi(z) = a^2e^{az-a}$. Тому
$G'_\xi(1) = a$, $G''_\xi(1) = a^2$, $G''_\xi(1) + G'_\xi(1) - \left( G'_\xi(1)\right)^2 = a^2 + a - a^2 = a$.
Отже, знайдено значення математичного сподівання та дисперсії.

\noindent\textbf{Числові характеристики:}
\begin{enumerate}
    \item $\E\xi = a$.
    \item $\D\xi = a$, $\sigma_\xi = \sqrt{a}$.
    \item $\Mo\xi = \left[a\right]$.
    \item $\Me\xi \approx \left[a + 1/3 - 0.02/a\right]$.
\end{enumerate}

\noindent\textbf{Застосування:} як вже було показано в асимптотичній формулі Пуассона \eqref{eq:poiss_aprox},
цим розподілом можна наблизити біноміальний, якщо ймовірність появи події дуже мала. Також цей розподіл
має кількість подій в \hyperref[poiss_proc]{потоці Пуассона}.

\subsection{Потік Пуассона}\label{poiss_proc}
\emph{Потоком} називають послідовність подій, які наступають в певні моменти часу одна за одною.
Потік характеризується випадковою величиною $\xi(t)$ --- 
кількістю подій, що наступили протягом проміжку часу $\left[ 0; t\right)$.
Ймовірності $\P\left\{\xi(t) = m\right\}$ позначатимемо $p_m(t)$.

Накладемо на потік такі вимоги:
\begin{enumerate}
    \item \emph{Стаціонарність (однорідність)} --- кількість подій, що наступають
    за певний проміжок часу, залежить лише від довжини проміжку і не залежить
    від того, де цей проміжок розташований на часовій осі.
    \item \emph{Відсутність післядії} --- якщо проміжки часу не перетинаються, то кількості подій,
    які за ці проміжки відбулися, є незалежними подіями.
    \item \emph{Ординарність} --- події наступають поодинці, тобто
    $\P\left\{\xi(t+\Delta t) - \xi(t) = 1\right\} = \lambda\cdot\Delta t + o(\Delta t)$
    ($\lambda > 0$ --- параметр інтенсивності) і 
    $\P\left\{\xi(t+\Delta t) - \xi(t) \geq 2\right\} = o(\Delta t)$.
\end{enumerate}
\begin{definition}
    Потік, що має перелічені властивості, називається \emph{потоком Пуассона}.
\end{definition}
\begin{remark}
    Насправді, умова $\P\left\{\xi(t+\Delta t) - \xi(t) = 1\right\} = \lambda\cdot\Delta t + o(\Delta t)$
    є надлишковою і її можна вивести з інших умов, але для спрощення приймемо її як частину означення.
\end{remark}

Отримаємо явний вираз для ймовірностей $p_m(t)$ у потоці Пуассона.
Позначимо $\tilde{p}_k(t+\Delta t) = \P\left\{\xi(t+\Delta t) - \xi(t) = k\right\}$.
Тоді $p_m(t+\Delta t) = p_m(t)\cdot \tilde{p}_0(t+\Delta t) + p_{m-1}(t)\cdot \tilde{p}_1(t+\Delta t) +
p_{m-2}(t)\cdot \tilde{p}_2(t+\Delta t) + ... + p_0(t)\cdot \tilde{p}_m(t+\Delta t)$ для $m \geq 1$,
для $m=0$ $p_0(t+\Delta t) = p_0(t)\cdot\tilde{p}_0(t+\Delta t)$. Застосуємо ординарність:

$p_0(t+\Delta t) = p_0(t) \cdot (1 - \lambda \cdot \Delta t + o(\Delta t))$

$p_m(t+\Delta t) = p_m(t) \cdot (1 - \lambda \cdot \Delta t + o(\Delta t)) + p_{m-1}(t) \cdot (\lambda \cdot \Delta t + o(\Delta t)) + o(\Delta t)$

\noindent Розкриємо дужки:

$p_0(t+\Delta t) - p_0(t) = -\lambda \cdot \Delta t \cdot p_0(t) + o(\Delta t)$
\nopagebreak

$p_m(t+\Delta t) - p_m(t) = -\lambda \cdot \Delta t \cdot p_m(t) + \lambda\cdot\Delta t \cdot p_{m-1}(t) + o(\Delta t)$

\noindent Поділимо на $\Delta t$ та перейдемо до границі при $\Delta t \rightarrow 0$.
Отримаємо систему диференціальних рівнянь:

$\begin{cases}
    p'_0(t) = - \lambda p_0(t) \\
    p'_m(t) = - \lambda p_m(t) + \lambda p_{m-1}(t), \; m \geq 1
\end{cases}$

\noindent Щоб розв'язати цю систему, введемо генератрису ймовірностей $p_m(t)$:

$G_\xi(z, t) = \sum\limits_{m=0}^{\infty} \P\left\{\xi(t) = m\right\} z^m = \sum\limits_{m=0}^{\infty} p_m(t) z^m$ 

\noindent Помножимо кожне рівняння отриманої системи на $z$ у відповідному степені та складемо:

$\sum\limits_{m=0}^{\infty} p'_m(t) z^m = -\lambda \sum\limits_{m=0}^{\infty} p_m(t) z^m + \lambda \sum\limits_{m=1}^{\infty} p_{m-1}(t) z^m = 
-\lambda \sum\limits_{m=0}^{\infty} p_m(t) z^m + \lambda z \sum\limits_{m=0}^{\infty} p_{m}(t) z^m$

\noindent Отримали диференціальне рівняння для генератриси:

$\frac{\partial G_\xi}{\partial t} = - \lambda G_\xi + \lambda z G_\xi = -\lambda(1-z) G_\xi$

\noindent Його розв'язком буде $G_\xi(z, t) = C(z) \cdot e^{-\lambda(1-z)t}$.
Оскільки $G_\xi(z, 0) = 1$, то $C(z) = 1$.

Отже, генератриса рівна $G_\xi(z, t) = e^{-\lambda(1-z)t}$ --- це генератриса закону Пуассона з параметром $a=\lambda t$.

Таким чином, $p_m(t) = \frac{(\lambda t)^m}{m!}e^{-\lambda t}$. 
Тепер зрозумілою стає назва параметру $\lambda$ (<<інтенсивність>>): $\E\xi = \lambda t$, $\lambda = \frac{\E\xi}{t}$ --- середня кількість подій за одиницю часу.

\begin{example}
    На АТС за годину в середньому за годину надходить 180 дзвінків. Яка ймовірність того,
    що за 2 хвилини надійде більше 4 дзвінків, якщо вважати, що моменти надходження дзвінків
    утворюють потік Пуассона?

    За умовою можна виписати інтенсивність надходження дзвінків: $\lambda = \frac{180}{60} = 3$ дзвінків за хвилину,
    тому кількість дзвінків за 2 хвилини матиме розподіл $\mathrm{Poiss}(3\cdot 2)$.
    Спочатку можна знайти ймовірність того, що надійде 4 або менше дзвінків:
    $\left(1 + 6 + \frac{6^2}{2!} + \frac{6^3}{3!} + \frac{6^4}{4!}\right)e^{-6} \approx 0.285$.
    Отже, шукана ймовірність приблизно дорівнює $1 - 0.285 = 0.715$.

\end{example}

\subsection{Рівномірний розподіл}
\noindent\textbf{Означення:}
    НВВ $\xi$ розподілена за \emph{рівномірним законом}, 
    якщо її щільність розподілу має вигляд \begin{equation}
        f_\xi(x) = \begin{cases}
            \frac{1}{b-a}, & x \in \left<a; b\right> \\
            0, & x \notin \left<a; b\right>
        \end{cases}
    \end{equation}
\textbf{Коротке позначення:} $\xi \sim \mathrm{U}\left<a; b\right>$.
    $a$ і $b$ --- параметри закону, $a<b$, $a, b \in \mathbb{R}$.

Розподіл $\mathrm{U}\left<0; 1\right>$ називається \emph{стандартним рівномірним розподілом}.

\begin{remark}
    Насправді, не важливо, що саме розуміти під проміжком $\left<a; b\right>$
    ($[a; b]$, $(a, b)$, $[a, b)$ чи $(a, b]$), оскільки $\P\left\{\xi = a \right\} = \P\left\{\xi = b \right\} = 0$,
    бо $\xi$ є неперервною.
\end{remark}

\begin{samepage}
    \noindent \textbf{Крива розподілу:}
    \begin{center}
        \begin{tikzpicture}[yscale = 2]
            \pgfmathsetmacro{\a}{1};
            \pgfmathsetmacro{\b}{3};
            \draw [->] (-2, 0) -- (5, 0);
            \draw [->] (0, -0.5) -- (0, 1);
            \draw [ultra thick] (\a, {1/(\b -\a)}) -- (\b, {1/(\b -\a)});
            \draw [dashed] (\a, 0) -- (\a, {1/(\b -\a)});
            \draw [dashed] (\b, 0) -- (\b, {1/(\b -\a)});
            \draw [ultra thick] [->] (-2, 0) -- (\a, 0);
            \draw [ultra thick] [<-] (\b, 0) -- (4.9, 0);
            \node [below] at (\a, -0.03) {$a$};
            \node [below] at (\b, 0) {$b$};
            \draw [dashed] (0, {1/(\b -\a)}) -- (\a, {1/(\b -\a)});
            \node [left] at (0, {1/(\b -\a)}) {$\frac{1}{b-a}$};
            \node [below] at (5, 0) {$x$};
            \node [left] at (0, 0.9) {$f_\xi(x)$};
        \end{tikzpicture}
    \end{center}
\end{samepage}

\noindent \textbf{Функція розподілу:}
\begin{center}
    \begin{tabular}{c c}
        $
            F_\xi(x) = \begin{cases}
                0, & x \leq a \\
                \frac{x-a}{b-a}, & a< x \leq b \\
                1, & x > b
            \end{cases}
        $ &
        \begin{tikzpicture}[baseline={(current bounding box.center)}, yscale=1.5]
            \pgfmathsetmacro{\a}{1};
        \pgfmathsetmacro{\b}{3};
        \draw [->] (-2, 0) -- (5, 0);
        \draw [->] (0, -0.5) -- (0, 1.2);
        \draw [ultra thick] (-2, 0) -- (\a, 0);
        \draw [domain=\a:\b, smooth, variable = \x, ultra thick] plot ({\x}, {(\x-\a)/(\b-\a)});
        \draw [ultra thick] (\b, 1) -- (5, 1);
        \node [below] at (\a, -0.03) {$a$};
        \node [below] at (\b, 0) {$b$};
        \node [below] at (5, 0) {$x$};
        \node [left] at (0, 1.2) {$F_\xi(x)$};
        \draw [dashed] (0, 1) -- (\b, 1);
        \draw [dashed] (\b, 0) -- (\b, 1);
        \node [left] at (0, 1) {$1$};
        %\node [right] [align=center] at (2.5, 0.3) {Приклад для \\ $a = \a, b = \b$};
        \end{tikzpicture}
    \end{tabular}
\end{center}

\noindent\textbf{Числові характеристики:}
\begin{enumerate}
    \item $\E\xi = \int\limits_{-\infty}^{+\infty} x f_\xi(x)dx = \int\limits_{a}^{b} \frac{x}{b-a}dx = \frac{a+b}{2}$.
    \item $\D\xi = \E\xi^2 - (\E\xi)^2 = \int\limits_{a}^{b} \frac{x^2}{b-a}dx - \frac{(a+b)^2}{4} = \frac{(b-a)^2}{12}$, $\sigma_\xi = \frac{b-a}{2\sqrt{3}}$.
    \item $\Mo\xi$ --- будь-яка точка з $\left<a; b\right>$.
    \item $\Me\xi = \frac{a+b}{2}$.
    \item $\As\xi = 0$.
    \item $\Ex\xi = -\frac{6}{5}$.
\end{enumerate}
% \begin{exercise}
%     Обчислити $\Me\xi$, $\As\xi$ та $\Ex\xi$ та перевірити наведені значення.
% \end{exercise}

\noindent\textbf{Застосування:} отримання вибірок з інших законів розподілу;
похибка округлення до найближчої поділки приладу при ціні поділки $a$ має розподіл $\mathrm{U}\left<-\frac{a}{2}; \frac{a}{2}\right>$.

\subsection{Експоненційний (показниковий) розподіл}
\noindent\textbf{Означення:}
    НВВ $\xi$ розподілена за \emph{експоненційним законом}, 
    якщо її щільність розподілу має вигляд \begin{equation}
        f_\xi(x) = \begin{cases}
            \lambda e^{-\lambda x}, & x \geq 0 \\
            0, & x < 0
        \end{cases}
    \end{equation}
\textbf{Коротке позначення:} $\xi \sim \mathrm{Exp}(\lambda)$.
    $\lambda > 0$ --- параметр закону.

\noindent \textbf{Крива розподілу:}
\begin{center}
    \begin{tikzpicture}[yscale = 2]
        \pgfmathsetmacro{\l}{1};
        \draw [->] (-2, 0) -- (5, 0);
        \draw [->] (0, -0.5) -- (0, 1.2);
        \draw [ultra thick] [->] (-2, 0) -- (0, 0);
        \draw [domain=0:5, smooth, variable = \x, ultra thick] plot ({\x}, {\l*e^(-\l*\x)});
        \node [below] at (5, 0) {$x$};
        \node [left] at (0, 1.2) {$f_\xi(x)$};
        \node [below right] at (0, 0) {$0$};
        \node [left] at (0, \l) {$\lambda$};
    \end{tikzpicture}
\end{center}

\noindent \textbf{Функція розподілу:}
\begin{center}
    \begin{tabular}{c c}
        $
            F_\xi(x) = \begin{cases}
                \lambda \int\limits_0^x e^{-\lambda t} dt = 1 - e^{-\lambda x}, & x> 0 \\
                0, & x \leq 0
            \end{cases}
            $ &
        \begin{tikzpicture}[baseline={(current bounding box.center)}, yscale=2]
            \pgfmathsetmacro{\l}{1};
            \draw [->] (-2, 0) -- (5, 0);
            \draw [->] (0, -0.5) -- (0, 1.2);
            \draw [ultra thick] (-2, 0) -- (0, 0);
            \draw [domain=0:5, smooth, variable = \x, ultra thick] plot ({\x}, {1 - e^(-\l*\x)});
            \node [below] at (5, 0) {$x$};
            \node [left] at (0, 1.2) {$F_\xi(x)$};
            \node [below right] at (0, 0) {$0$};
            \node [left] at (0, 1) {$1$};
            \draw [dashed] (0, 1) -- (5, 1);
        \end{tikzpicture}
    \end{tabular}
\end{center}

Для знаходження основних числових характеристик знайдемо всі \emph{початкові моменти} експоненційного розподілу:
\begin{gather*}
    \E\xi^k = \int\limits_{-\infty}^{+\infty} x^k f_\xi(x)dx = \lambda \int\limits_{0}^{+\infty} x^k e^{-\lambda x}dx = \left[\; \lambda x = t \;\right]=
\frac{\lambda}{\lambda^k \cdot \lambda} \int\limits_{0}^{+\infty} t^k e^{-t} dt = \frac{\Gamma(k+1)}{\lambda^k} = \frac{k!}{\lambda^k}
\end{gather*}

\noindent\textbf{Числові характеристики:}
\begin{enumerate}
    \item $\E\xi = \frac{1}{\lambda}$.
    \item $\D\xi = \E\xi^2 - (\E\xi)^2 = \frac{2}{\lambda^2} - \frac{1}{\lambda^2} = \frac{1}{\lambda^2}$, $\sigma_\xi = \frac{1}{\lambda}$.
    \item $\Mo\xi = 0$.
    \item $\Me\xi = \frac{\ln2}{\lambda}$.
    \item $\As\xi = 2$ --- не залежить від $\lambda$, бо $\beta_3 = \alpha_3 - 3\alpha_2 \alpha_1 + 2 \alpha_1^3 = \frac{2}{\lambda^3}$.
    \item $\Ex\xi = 6$ --- теж не залежить від $\lambda$, бо $\beta_4 = \alpha_4 - 4\alpha_3 \alpha_1 + 6 \alpha_2 \alpha_1^2 - 3\alpha_1^4 = \frac{9}{\lambda^4}$.
\end{enumerate}

\begin{exercise}
    Дослідити характеристики <<зсунутого>> експоненційного розподілу
    зі щільністю
    $\mathrm{Exp}(\lambda, x_0)$ з щільністю $f_\xi(x) = \begin{cases}
        \lambda e^{-\lambda (x-x_0)}, & x \geq x_0 \\
        0, & x < x_0
    \end{cases}$.
\end{exercise}

Як і геометричний розподіл, експоненційний \textbf{<<не має пам'яті>>}.
Доведемо, що якщо проміжок часу $T$, розподілений за експоненційним законом,
тягнувся час $\tau$, то проміжок $T_1 = T - \tau$ розподілений так само --- 
за експоненційним законом з тим самим параметром:
\begin{gather*}
    F_{T_1}(t) = \P\left\{T_1 < t / T \geq \tau\right\} = 
    \frac{\P\left\{(T_1 < t) \cap (T \geq \tau)\right\}}{\P\left\{ T \geq \tau \right\}} =
    \frac{\P{\left\{\tau \leq T < t +\tau\right\}}}{1 - \P\left\{ T < \tau \right\}} = \\
    \frac{F_T(t+\tau) - F_T(\tau)}{1-F_T(\tau)} = 
    \frac{1-e^{-\lambda (t+\tau)} - (1 - e^{-\lambda \tau})}{1-(1-e^{-\lambda \tau})} =
    \frac{e^{-\lambda \tau}(1-e^{-\lambda t})}{e^{-\lambda \tau}} = 1-e^{-\lambda t} = F_T(t)
\end{gather*}

\noindent\textbf{Застосування:}
\begin{enumerate}
    \item Час між двома сусідніми подіями в потоці Пуассона має експоненційний розподіл.
    \begin{exercise}
        Довести цю властивість.
    \end{exercise}
    \item Як правило, час безвідмовної роботи приладу має експоненційний розподіл.
    
    Нехай в момент часу $t=0$ увімкнули прилад. Припустимо, що умовна ймовірність виходу з ладу
    приладу в інтервалі часу $\left[ t; t+\Delta t\right)$ за умови, що до часу $t$ він працював,
    пропорційна $\Delta t$ і рівна $\lambda\cdot \Delta t + o(\Delta t)$. Нехай $\xi$ --- час безвідмовної роботи,
    покладемо $Q(t) = \P\left\{\xi \geq t\right\}$.
    $Q(t+\Delta t) = \P\left\{\xi \geq t + \Delta t\right\} = \P\left\{\xi \geq t\right\} \cdot \P\left\{\xi \geq t+\Delta t / \xi \geq t\right\} = 
    Q(t) \cdot (1 - \lambda\cdot \Delta t + o(\Delta t))$.
    $\frac{Q(t+\Delta t) - Q(t)}{\Delta t} = -\lambda Q(t) + \frac{o(\Delta t)}{\Delta t} \cdot Q(t)$,
    при $\Delta t \rightarrow 0$ отримаємо $Q'(t) = -\lambda Q(t)$.
    Розв'язком цього диференціального рівняння буде $Q(t) = C\cdot e^{-\lambda t}$.
    Оскільки вважаємо, що в момент часу $t=0$ прилад працював, то $Q(0) = 1$ і $C=1$.

    Тоді $F_\xi(t) = \P\left\{ \xi < t\right\} = 1 - \P\left\{ \xi \geq t\right\} = 1 - Q(t) = \begin{cases}
        1 - e^{-\lambda t}, & t \geq 0 \\
        0, & t < 0
    \end{cases}$. Отже, $\xi \sim \mathrm{Exp}(\lambda)$.
\end{enumerate}

\begin{example}
    Яка ймовірність того, що людина проживе 100 років? Нехай $\xi$ --- час життя людини.
    Оскільки $\P\left\{\xi = 100\right\} = 0$, знайдемо $\P\left\{ \xi \geq 100\right\}$.
    $\P\left\{ \xi \geq 100\right\}$.
    Оскільки $\lambda = \frac{1}{\E\xi}$, то, прийнявши $\E\xi = 70$ (середня тривалість життя),
    отримаємо $\P\left\{ \xi \geq 100\right\} = e^{-100/70} \approx 0.24$.
\end{example}

\subsection{Гамма-розподіл}
\noindent\textbf{Означення:}
    НВВ $\xi$ розподілена за \emph{гамма-законом}, 
    якщо її щільність розподілу має вигляд 
    \begin{equation}
        f_\xi(x) = \begin{cases}
            \frac{\beta^\alpha}{\Gamma(\alpha)} x^{\alpha-1} e^{-\beta x}, & x \geq 0 \\
            0, & x < 0
        \end{cases}
    \end{equation}
\textbf{Коротке позначення:} $\xi \sim {\Gamma}(\alpha, \beta)$.
    $\alpha >0, \beta > 0$ --- параметри закону.
\begin{remark}
    Експоненційний розподіл $\mathrm{Exp}(\lambda)$ є частинним випадком гамма-розподілу, а саме ${\Gamma}(1, \lambda)$.
\end{remark}
\begin{samepage}
    \noindent \textbf{Крива розподілу:} приклад для $\alpha = 2, \beta=1$.
    \begin{center}
        \begin{tikzpicture}[yscale = 7, xscale = 0.5, baseline={(current bounding box.center)}]
            \pgfmathsetmacro{\a}{2};
            \pgfmathsetmacro{\b}{1};
            \pgfmathsetmacro{\s}{2};
            \draw [->] (-2, 0) -- (20, 0);
            \draw [->] (0, -0.05) -- (0, 0.4);
            \draw [ultra thick] (-2, 0) -- (0, 0);
            \draw [domain=0:19.7, smooth, variable = \x, ultra thick, samples=100] plot ({\x}, {\b^(\a)*(\x/\s)^(\a-1) * e^(-\x*\b/\s)});
            \node [below] at (20, 0) {$x$};
            \node [left] at (0, 0.4) {$f_\xi(x)$};
            \node [below left] at (0, 0) {$0$};
        \end{tikzpicture} 
    \end{center}
\end{samepage}

% \begin{center}
%     \begin{tikzpicture}[yscale = 5, xscale = 0.5, baseline={(current bounding box.center)}]
%         \pgfmathsetmacro{\a}{1};
%         \pgfmathsetmacro{\b}{0.5};
%         \draw [->] (-2, 0) -- (20, 0);
%         \draw [->] (0, -0.05) -- (0, 0.6);
%         \draw [ultra thick] [->] (-2, 0) -- (0, 0);
%         \draw [domain=0:19.7, smooth, variable = \x, ultra thick] plot ({\x}, {\b^(\a)*\x^(\a-1) * e^(-\x*\b)});
%         \node [below] at (20, 0) {$x$};
%         \node [left] at (0, 0.6) {$f_\xi(x)$};
%         \node [below right] at (0, 0) {$0$};
%         \node at (18, 0.375) {$\alpha = 1$};
%         \node at (18, 0.285) {$\beta = 1/2$};
%     \end{tikzpicture} 
% \end{center}

\noindent \textbf{Функція розподілу:}
\nopagebreak

\begin{tabular}{c c}
    \begin{tabular}{c}
        $
        F_\xi(x) = \begin{cases}
            0, & x \leq 0 \\
            \int\limits_0^x f_\xi(t)dt = \frac{\gamma(\alpha, \beta x)}{\Gamma(\alpha)}, & x>0 
        \end{cases}$
        \\
        де $\gamma(s, x) = \int\limits_0^x t^{s-1} e^{-t} dt$
    \end{tabular}
    &
    \begin{tikzpicture}[baseline={(current bounding box.center)}, yscale=0.9]
        \pgfmathsetmacro{\a}{2};
        \pgfmathsetmacro{\b}{3};
        \draw [->] (-2, 0) -- (5, 0);
        \draw [->] (0, -0.5) -- (0, 2.7);
        \draw [ultra thick] (-2, 0) -- (0, 0);
        \draw [dashed] (0, 2.3) -- (5, 2.3);
        \node [below] at (5, 0) {$x$};
        \node [left] at (0, 2.7) {$F_\xi(x)$};
        \node [left] at (0, 2.3) {$1$};
        \node [below right] at (0, 0) {$0$};
        \begin{axis}[width=0.4\textwidth,height=0.25\textwidth, samples=100, domain = 0:4, restrict y to domain = 0:1, axis x line=bottom,axis y line=left, clip=false, axis line style={draw=none}, ticks=none]
        \addplot [ultra thick] plot ({\x}, {(\x^\a) * (
            1/\a -
            \x/(1+\a) +
            \x^2/(2*(2+\a)) -
            \x^3/(6*(3+\a)) +
            \x^4/(24*(4+\a)) -
            \x^5/(120*(5+\a)) +
            \x^6/(720*(6+\a)) -
            \x^7/(5040*(7+\a)) +
            \x^8/(40320*(8+\a)) -
            \x^9/(362880*(9+\a)) +
            \x^10/(3628800*(10+\a)) -
            \x^11/(39916800*(11+\a))
        )});
        \end{axis}
    \end{tikzpicture}
\end{tabular}

Знайдемо всі \emph{початкові моменти} гамма-розподілу.
\begin{gather*}
    \E\xi^k = \int\limits_{-\infty}^{+\infty} x^k f_\xi(x)dx =
    \frac{\beta^\alpha}{\Gamma(\alpha)} \int\limits_0^{+\infty} x^k x^{\alpha-1} e^{-\beta x} dx =
    \left[\; \beta x = t \;\right] = \\
    = \frac{\beta^\alpha}{\Gamma(\alpha)} \cdot \frac{1}{\beta^k\cdot \beta^{\alpha}} \int\limits_0^{+\infty} t^{k+\alpha-1} e^{-t} dt =
    \frac{\Gamma(\alpha+k)}{\Gamma(\alpha) \cdot \beta^k}
\end{gather*}
Зокрема, $\E\xi = \frac{\Gamma(\alpha+1)}{\Gamma(\alpha) \cdot \beta} = \frac{\alpha}{\beta}$ та 
$\E\xi^2 = \frac{\Gamma(\alpha+2)}{\Gamma(\alpha) \cdot \beta^2}  = \frac{(\alpha+1)\alpha}{\beta^2}$.

\noindent\textbf{Числові характеристики:}
\begin{enumerate}
    \item $\E\xi = \frac{\alpha}{\beta}$.
    \item $\D\xi = \E\xi^2 - (\E\xi)^2 = \frac{(\alpha+1)\alpha}{\beta^2} - \frac{\alpha^2}{\beta^2} = \frac{\alpha}{\beta^2}$, $\sigma_\xi = \frac{\sqrt{\alpha}}{\beta}$.
    \item $\Mo\xi = \frac{\alpha-1}{\beta}$.
    \item $\Me\xi$ --- немає виразу у замкненій формі.
    \item $\As\xi = \frac{2}{\sqrt{\alpha}}$.
    \item $\Ex\xi = \frac{6}{\alpha}$.
\end{enumerate}

\noindent\textbf{Застосування:} гамма-розподіл застосовується для моделювання
складних потоків подій, в економіці, теорії масового обслуговування, логістиці.
У випадку натурального параметру $\alpha = k \in \mathbb{N}$ за законом
$\Gamma(k, \lambda)$ розподілений час очікування появи $k$-тої події
в процесі Пуассона з інтенсивністю $\lambda$. Таку версію гамма-розподілу іноді називають \emph{розподілом Ерланга}.
\begin{exercise}
    Нехай $T$ --- час очікування появи $k$-тої події
    в потоці Пуассона з інтенсивністю $\lambda$. Довести $T \sim \Gamma(k, \lambda)$.
\end{exercise} % https://towardsdatascience.com/gamma-distribution-intuition-derivation-and-examples-55f407423840

\subsection{Гауссівський (нормальний) розподіл}
\noindent\textbf{Означення:}
    НВВ $\xi$ розподілена за \emph{нормальним законом}, 
    якщо її щільність розподілу має вигляд 
    \begin{equation}
        f_\xi(x) = \frac{1}{\sqrt{2\pi}\sigma} e^{-\frac{(x-a)^2}{2\sigma^2}}
    \end{equation}
\textbf{Коротке позначення:} $\xi \sim \mathrm{N}(a, \sigma)$ або 
    $\xi \sim \mathrm{N}(a, \sigma^2)$.
    $a \in \mathbb{R}, \sigma > 0$ --- параметри закону.

    Розподіл $\mathrm{N}(0, 1)$ називають \emph{стандартним гауссівським розподілом}.

\begin{samepage}
    \noindent \textbf{Крива розподілу:}
    \begin{center}
        \begin{tikzpicture}[yscale = 4, xscale = 2]
            \pgfmathsetmacro{\a}{2.5};
            \pgfmathsetmacro{\s}{1}
            \draw [->] (-1, 0) -- (5, 0);
            \draw [->] (0, -0.2) -- (0, 0.5);
            \draw [dashed] (\a, 0) -- (\a, 0.398942280401/\s);
            \draw [dashed] (\a-\s, 0) -- (\a-\s, 0.241970724519/\s);
            \draw [dashed] (\a+\s, 0) -- (\a+\s, 0.241970724519/\s);
            \draw [domain=-1:5, smooth, variable = \x, ultra thick] plot ({\x}, {0.398942280401/\s * e^(-(\x-\a)^2/(2*\s^2))});
            \node [below] at (5, 0) {$x$};
            \node [below] at (\a, -0.025) {$a$};
            \node [below] at (\a-\s, 0) {$a - \sigma$};
            \node [below] at (\a+\s, 0) {$a + \sigma$};
            \node [below right] at (0, 0) {$0$};
            \node [left] at (0, 0.5) {$f_\xi(x)$};
        \end{tikzpicture}
    \end{center}
\end{samepage}

\noindent \textbf{Функція розподілу:}
$F_\xi(x) = \int\limits_{-\infty}^{x} f_\xi(t) dt = 
\frac{1}{\sqrt{2\pi}\sigma} \int\limits_{-\infty}^{x} 
e^{-\frac{(t-a)^2}{2\sigma^2}} dt$ не виражається в елементарних 
функціях.
Для зручного використання гауссівського закону в задачах введемо 
спеціальну функцію, \emph{функцію Лапласа} $\Phi(x)$, значення якої занесено до \hyperref[table:laplace]{таблиці}:
\begin{gather*}
    \Phi(x) = \frac{1}{\sqrt{2\pi}} 
    \int\limits_{0}^{x} e^{-\frac{t^2}{2}} dt
\end{gather*}
\begin{center}
    \begin{tabular}{c c}
        \begin{tikzpicture}[baseline={(current bounding box.center)}, yscale=3, 
            scale = 1]
            \fill [lightgray, domain=0:1, smooth, variable = \x] plot ({\x}, 
            {
                (0.3989422804) * e^(- (\x * \x / 2))
            }) -- (1, 0) -- (0, 0) -- (0, 0.3989422804);
            \draw [->] (-3, 0) -- (3, 0);
            \node [below] at (3, 0) {$t$};
            \draw [->] (0, -0.2) -- (0, 0.7);
            \draw [domain=-3:3, smooth, variable = \x, ultra thick] plot ({\x}, 
            {
                (0.3989422804) * e^(- (\x * \x / 2))
            });
            \node [below] at (1, 0) {$x$};
            \draw [dashed] (1, 0) -- (1, 0.25);
            \draw [->, thick] (1.5, 0.4) -- (0.7, 0.2);
            \node [below left] at (3.1, 0.5) {$S = \Phi(x)$};
            \draw [->] (-0.8, 0.5) -- (-0.495, 0.355);
            \node [left] at (-0.8, 0.5) {$\frac{1}{\sqrt{2\pi}}e^{-\frac{t^2}{2}}$};
        \end{tikzpicture} &
        \begin{tikzpicture}[baseline={(current bounding box.center)}, yscale=2, scale = 1.2]
            \draw [->] (-2, 0) -- (2, 0);
            \draw [->] (0, -0.6) -- (0, 0.6);
            \draw [domain=-2:2, smooth, variable = \x, ultra thick] plot ({\x}, 
            {
                (0.564189583547756) * (
                (\x / 1.414213562373095) 
                - ((\x / 1.414213562373095)^3 / 3) 
                + ((\x / 1.414213562373095)^5 / 10) 
                - ((\x / 1.414213562373095)^7 / 42)
                + ((\x / 1.414213562373095)^9 / 216) 
                - ((\x / 1.414213562373095)^11 / 1320)
                + ((\x / 1.414213562373095)^13 / 9360)
                )
            });
            \node [below] at (2, 0) {$x$};
            \node [above left] at (0, 0.5) {$\Phi(x)$};
            \node [below right] at (0, 0) {$0$};
            \draw [dashed] (-2, -0.5) -- (2, -0.5);
            \draw [dashed] (-2, 0.5) -- (2, 0.5);
            \node [below left] at (0, 0.5) {$0.5$};
            \node [above left] at (0, -0.5) {$-0.5$};
        \end{tikzpicture}
    \end{tabular}
\end{center}

\noindent
\emph{Властивості функції Лапласа: }
    \begin{enumerate}
        \item $\Phi(-x) = -\Phi(x)$ --- непарна.
        \item $\Phi(0) = 0$.
        \item $y = \pm \;0.5$ --- горизонтальні асимптоти
        \item $\forall x \geq 5: \Phi(x) \approx 0.5 $.
        \item $\forall x \leq -5: \Phi(x) \approx -0.5 $.
    \end{enumerate}

\noindent Виразимо функцію розподілу гауссівського закону через функцію Лапласа:
\begin{equation}
    F_\xi (x) = \frac{1}{\sqrt{2\pi}\sigma} \int\limits_{-\infty}^{x} 
    e^{-\frac{(t-a)^2}{2\sigma^2}} dt = 
    \left[\frac{t-a}{\sigma} = z\right] = 
    \frac{1}{\sqrt{2\pi}} \int\limits_{-\infty}^{\frac{x-a}{\sigma}} 
    e^{-\frac{z^2}{2}} dz = \frac{1}{2} + 
    \Phi\left(\frac{x-a}{\sigma}\right)
\end{equation}

Виявляється, що параметри $a$ та $\sigma$ мають простий ймовірнісний сенс.
Знайдемо спочатку математичне сподівання:
\begin{gather*}
    \E\xi = 
    \frac{1}{\sqrt{2\pi}\sigma}
    \int\limits_{-\infty}^{+\infty}x e^{-\frac{(x-a)^2}{2\sigma^2}} dx 
    = \left[ \frac{x-a}{\sigma \sqrt{2}} = t, \; dx = \sigma\sqrt{2}dt \right] = 
    \frac{1}{\sqrt{\pi}}\int\limits_{-\infty}^{+\infty} 
    (a + \sigma\sqrt{2}t)e^{-t^2} dt = \\
    =\frac{a}{\sqrt{\pi}}
    \underbrace{\int\limits_{-\infty}^{+\infty}e^{-t^2}dt}_{ = \sqrt{\pi}} \; + \; 
    \frac{\sigma \sqrt{2}}{\sqrt{\pi}} 
    \underbrace{\int\limits_{-\infty}^{+\infty} t e^{-t^2}dt}_{= 0}
    = a
\end{gather*}
Знайдемо тепер всі центральні моменти:
\begin{gather*}
    \beta_k = \E(\xi - \E\xi)^k = 
    \frac{1}{\sigma\sqrt{2\pi}}
    \int\limits_{-\infty}^{+\infty}(x-a)^k 
    e^{-\frac{(x-a)^2}{2\sigma^2}} dx = 
    \left[ \frac{x-a}{\sigma \sqrt{2}} = t\right] 
    =
    \frac{\sigma^k 2^{k/2}}{\sqrt{\pi}}
    \int\limits_{-\infty}^{+\infty}t^k 
    e^{-t^2} dt
\end{gather*}
Якщо $k = 2l + 1$, то $\beta_k = 0$, оскільки під інтегралом непарна функція.
Зокрема, $\beta_3 = 0$, тому $\As\xi = 0$.
Візьмемо $k = 2l$:
\begin{gather*}
    \beta_{2l} = 
    \frac{\sigma^{2l} 2^{l}}{\sqrt{\pi}}
    \int\limits_{-\infty}^{+\infty}t^{2l} 
    e^{-t^2} dt = \left[ t^2 = z, \; dt = \frac{1}{2} z^{-\frac{1}{2}}dz\right] = 
    2\frac{\sigma^{2l}2^l}{\sqrt{\pi}}
    \int\limits_{0}^{+\infty}z^l e^{-z}\cdot \frac{1}{2} z^{-\frac{1}{2}}
    dz = \\
    = \frac{\sigma^{2l}2^l}{\sqrt{\pi}}
    \int\limits_{0}^{+\infty}z^{l-\frac{1}{2}} e^{-z}
    dz = 
    \frac{\sigma^{2l}2^l}{\sqrt{\pi}}\cdot\Gamma\left(l+\frac{1}{2}\right) = 
    \frac{\sigma^{2l}2^l}{\sqrt{\pi}} \cdot \frac{(2l-1)!!}{2^l} \sqrt{\pi} = \sigma^{2l} (2l-1)!!
\end{gather*}
Зокрема, $\beta_2 = \D\xi = \frac{2 \sigma^2}{\sqrt{\pi}}\cdot\Gamma\left(\frac{3}{2}\right) = \sigma^2$,
$\beta_4 = \frac{\sigma^4\cdot2^2}{\sqrt{\pi}}\cdot \Gamma\left(\frac{5}{2}\right) = 3\sigma^4$.

\noindent \textbf{Числові характеристики:}
\begin{enumerate}
    \item $\E\xi = \Mo\xi = \Me\xi = a$.
    \item $\D\xi = \sigma^2$, $\sigma_\xi = \sigma$.
    \item $\As\xi = 0$.
    \item $\Ex\xi = \frac{\beta_4}{\sigma_\xi^4} - 3 = 
    \frac{3\sigma^4}{\sigma^4} - 3 = 0$. 
\end{enumerate}
\begin{remark}
    Доданок $-3$ у формулі ексцесу був введений саме для того, щоб нормальний розподіл мав нульовий ексцес,
    тому додатній ексцес означає, що крива розподілу біля точки максимуму є крутішою за криву гауссівського розподілу,
    а від'ємний --- що більш пласкою.
\end{remark}

\noindent\textbf{Застосування:} гауссівський закон є найважливішим в теорії ймовірностей,
оскільки може описати розподіл багатьох величин, що зустрічаються в природі.
Цьому сприяє, наприклад, симетричність розподілу відносно математичного
сподівання (значення якого також є модою і медіаною, тобто <<середнім>> у всіх можливих в цьому випадку сенсах) та наведене нижче \emph{<<правило $3\sigma$>>},
яке означає, що майже всі значення гауссівської випадкової величини
лежать на відстані не більше трьох середньоквадратичних 
відхилень від її математичного сподівання.
Ці властивості природно очікувати від деяких <<загальних>> характеристик, 
як-от зріст дорослої людини: інтуїтивно зрозуміло, що є деякий <<середній>> зріст (наприклад, 170 см),
а майже всі інші можливі значення знаходяться в деякому діапазоні навколо нього (наприклад, 155-185 см), окрім досить малої кількості
винятково малих чи великих. Важливість гауссівського розподілу ще пояснюється тим,
що за деяких умов він може наближено замінити розподіл суми достатньо великої кількості випадкових величин, 
навіть якщо вони мають різні закони розподілу --- про це йтиме мова у розділі \hyperref[ch:limit_theorems]{<<Граничні теореми теорії ймовірностей>>}.

\vspace{0.5em}
\noindent \textbf{Робочі формули для розв'язання задач:}
\begin{enumerate}
    \item $\P\left\{\xi \in \left<\alpha; \beta\right>\right\} 
    = F_\xi(\beta) - F_\xi(\alpha) = \Phi\left(\frac{\beta - a}{\sigma}\right)
    - \Phi\left(\frac{\alpha - a}{\sigma}\right)$.
    \item $\P\left\{\left| \xi - a \right| < \varepsilon\right\} = 
    \P\left\{\xi \in \left(a - \varepsilon, a + \varepsilon\right)\right\} = 
    \Phi\left(\frac{a + \varepsilon - a}{\sigma}\right)
    - \Phi\left(\frac{a - \varepsilon - a}{\sigma}\right) = 
    2\Phi(\frac{\varepsilon}{\sigma})$.
    \item \emph{<<Правило $3\sigma$>>}: 
    $\P\left\{\left| \xi - a \right| < 3\sigma\right\} = 
    \P\left\{\xi \in \left( a-3\sigma; a+3\sigma\right)\right\} = 2\Phi\left(3\right)
    \approx 0.9973$.
\end{enumerate}
    \chapter{Випадкові вектори}
        % !TEX root = ../main.tex
\section{Поняття випадкового вектора}
В багатьох прикладних задачах розглядають не одну випадкову величину,
а деяку їх систему. Наприклад,
стан деякого приладу характеризується станом його компонентів,
які в кожен момент можуть вийти з ладу, причому відмова одного компоненту 
може пришвидшити відмову інших.
Системою випадкових величин також можна описати оцінки
навмання вибраного випускника університету, і в цьому випадку теж
можуть виникати деякі залежності між оцінками з різних дисциплін.
Ці приклади показують, що дослідження
систем випадкових величин є складнішим за дослідження окремих величин,
тому що необхідно враховувати, як ці величини можуть бути пов'язані між собою.

\begin{definition}\index{випадковий вектор}
    Нехай $\left\{\Omega, \mathcal{F}, \P \right\}$ --- ймовірнісний просторі
    деякого стохастичного експерименту.
    \emph{Випадковим вектором (системою випадкових величин)}
    називається вимірна функція $\vec{\xi}(\omega) = 
    \left(\xi_1(\omega), \xi_2(\omega), ... , \xi_n(\omega)\right)^{T}: \Omega \rightarrow \mathbb{R}^n$, де
    під вимірністю мається на увазі, що 
    \begin{equation*}
        \forall \vec{x} \in \mathbb{R}^n: 
        A=\left\{\omega \in \Omega: \xi_1(\omega) < x_1, 
                                    \xi_2(\omega) < x_2,
                                    ... ,
                                    \xi_n(\omega) < x_n\right\}
        \in \mathcal{F}
    \end{equation*}
\end{definition}
\begin{remark}
    Можна довести інтуїтивно зрозумілий факт: координати випадкового вектора є випадковими величинами в сенсі означення 
    \ref{def:random_variable}, і навпаки --- функція з $\Omega$ в $\mathbb{R}^n$, координатами якої є випадкові величини,
    є випадковим вектором в сенсі наведеного вище означення.
\end{remark}
Випадковий вектор можна трактувати як випадкову точку в $\mathbb{R}^n$.
Наприклад, при $n = 2$ $\vec{\xi} = (\xi_1, \xi_2)^T$ --- випадкова точка на площині. 

\subsection{Сумісна функція розподілу}
Як у і випадку випадкових величин, випадковий вектор теж характеризується своєю функцією розподілу.
Спочатку розглянемо лише двовимірний випадок, оскільки його простіше інтерпретувати геометрично.
\begin{definition} \index{функція розподілу!сумісна}
    \emph{Сумісною функцією розподілу} двовимірного випадкового вектора $\vec{\xi}$ 
    називається $F_{\vec{\xi}}(x, y) = \P\left\{\omega: \xi_1 < x, \xi_2 < y\right\}$, $x, y \in \mathbb{R}$.
\end{definition}
Значення $F_{\vec{\xi}}(x, y)$ рівне ймовірності 
    потрапляння точки в нескінченний квадрант з вершиною в точці $(x, y)$.
\begin{center}
    \begin{tikzpicture}[baseline={(current bounding box.north)}]
        \fill [lightgray] (2, 2) rectangle (-0.55, -0.51);
        \draw [->] (-0.5, 0) -- (3, 0);
        \draw [->] (0, -0.5) -- (0, 3);
        \draw [fill] (2, 2) circle [radius=0.05];
        \node [above right] at (2, 2) {$(x, y)$};
        \draw [dashed] (-0.5, 2) -- (2, 2) -- (2, -0.5);
        \node [below] at (3, 0) {$x$};
        \node [above right] at (0.2, 0.2) {$F_{\vec{\xi}}(x, y)$};
        \node [left] at (0, 3) {$y$};
    \end{tikzpicture}
\end{center}

\noindent\textbf{Властивості сумісної функції розподілу:}
\begin{enumerate}
    \item Область визначення --- $\mathbb{R}^2$, область значень --- $\left<0; 1\right>$.
    \item Монотонно неспадна по кожному з аргументів.
    \begin{proof}
        Для $x_1 > x_2$ позначимо
        $A = \left\{\xi_1 < x_2, \xi_2 < y \right\}$, 
        $B = \left\{x_2 \leq \xi_1 < x_1,\xi_2<y\right\}$, 
        $C = \left\{\xi_1 < x_1, \xi_2 < y\right\}$.
        $A$ і $B$ --- несумісні,
        $C = A \cup B$, тому $\P(C) = \P(A) + \P(B)$.
        Оскільки 
        $\P(C) = F_{\vec{\xi}}(x_1, y)$,  
        $\P(A) = F_{\vec{\xi}}(x_2, y)$, то 
        $F_{\vec{\xi}}(x_1, y) = F_{\vec{\xi}}(x_2, y) + \P(B)
        \Rightarrow F_{\vec{\xi}}(x_1, y) \geq  F_{\vec{\xi}}(x_2, y)$. 
        Доведення для другого аргумента аналогічне.
    \end{proof}
    \item $\lim\limits_{x \to -\infty} F_{\vec{\xi}}(x, y) = 
           \lim\limits_{y \to -\infty} F_{\vec{\xi}}(x, y) = 
           \lim\limits_{x,y \to -\infty} F_{\vec{\xi}}(x, y) = 0$.
    \begin{proof}
        Як і в одновимірному випадку (властивість \refeq{cdf:limit}, ст. \pageref{cdf:limit}),
        доведення ґрунтується на теоремі про неперервність 
        ймовірності (\ref{th:2}).
        Для будь-якої послідовності $x_n \to -\infty$ маємо
        $\lim\limits_{n \to \infty} F_{\vec{\xi}}(x_n, y) = 
        \lim\limits_{n \to \infty} \P\left\{\xi_1<x_n,\xi_2<y\right\} 
        = \P(\varnothing \cap \{\xi_2 < y\})
        = \P(\varnothing) = 0$.
        Інші випадки доводяться аналогічно.
    \end{proof}
    \item \emph{<<Умови узгодженості>>}\index{умова!узгодженості (умови)}:
    $\lim\limits_{x \to +\infty} F_{\vec{\xi}}(x, y) = F_{\xi_2}(y)$, 
    $\lim\limits_{y \to +\infty} F_{\vec{\xi}}(x, y) = F_{\xi_1}(x)$.
    \begin{proof}
        Доведення ґрунтується на теоремі про неперервність 
        ймовірності (\ref{th:1}). Для будь-якої
        послідовності $x_n \to +\infty$ маємо
        $\lim\limits_{n \to \infty} F_{\vec{\xi}}(x_n, y) = 
        \lim\limits_{n \to \infty} 
        \P\left\{\xi_1 < x_n, \xi_2 < y\right\} = 
        \P\left(\Omega \cap \left\{\xi_2<y\right\}\right) = 
        \P\left\{\xi_2<y\right\} = F_{\xi_2}(y)$. 
        Для іншого аргумента доведення аналогічне.
    \end{proof}
    \item $\lim\limits_{x,y \to +\infty} F_{\vec{\xi}}(x, y) = 1$.
    \begin{proof}
        Випливає з умов узгодженості.
    \end{proof}
    \item Функція розподілу є неперервною зліва по кожному з аргументів: 
    
    $\lim\limits_{x \to x_0 - 0} F_{\vec{\xi}}(x, y) = F_{\vec{\xi}}(x_0, y)$, 
    $\lim\limits_{y \to y_0 - 0} F_{\vec{\xi}}(x, y) = F_{\vec{\xi}}(x, y_0)$.
    \begin{proof}
        Доведення аналогічне одновимірному випадку (властивість \refeq{cdf:left_con}, 
        с. \pageref{cdf:left_con}).
    \end{proof}
    \item Для $\Pi = [a; b)\times [c; d)$ $\P\left\{\vec{\xi} \in \Pi \right\} = F(b,d) + F(a,c) - F(b,c) - F(a,d)$.
    \begin{proof}
        Нехай $A = \left\{\xi_1 < a, \xi_2 < c\right\}$,
        $B = \left\{\xi_1 \in [a;b), \xi_2 < c\right\}$,
        $C = \left\{\xi_1 < a, \xi_2 \in [c;d)\right\}$,

        \begin{tabular}{c p{8.4cm}}
            \begin{tikzpicture}[baseline={(current bounding box.north)}]
                \draw [->] (-1, 0) -- (3, 0);
                \draw [->] (0, -1) -- (0, 2);
                \draw [very thick] (0.5,1.5) -- (0.5,0.5) -- (2.5,0.5);
                \draw [dashed, very thick] (2.5,0.5) -- (2.5,1.5) -- (0.5,1.5);
                \fill [lightgray] (0.5,0.5) rectangle (2.5,1.5);
                \draw [dashed] (0.5, 0.5) -- (-1, 0.5);
                \draw [dashed] (0.5, 0.5) -- (0.5, -1);
                \draw [dashed] (0.5, 1.5) -- (-1, 1.5);
                \draw [dashed] (2.5, 0.5) -- (2.5, -1);
                \node [below right] at (0.5, 0) {$a$};
                \node [below left] at (2.5, 0) {$b$};
                \node [above left] at (0, 0.5) {$c$};
                \node [below left] at (0, 1.5) {$d$};
                \node [below] at (3, 0) {$x$};
                \node [left] at (0, 2) {$y$};
                \node [below left] at (2.5, 1.5) {$\Pi$};
            \end{tikzpicture} & 
            $D = \left\{\vec{\xi} \in \Pi\right\}$ та 
            $E = \left\{\xi_1 < b, \xi_2 < d\right\}$.
            Оскільки
            $E = A \cup B \cup C \cup D$, причому $A$, $B$, $C$, $D$ попарно несумісні, то
            $\P(E) = \P(A) + \P(B) + \P(C) + \P(D)$.
            $\P(E) = F(b, d)$, $\P(A) = F(a, c)$, $\P(B) = F(b,c) - F(a,c)$,
            $\P(C) = F(a, d) - F(a, c)$. 
            
            Отже, $\P\left\{\vec{\xi} \in \Pi\right\} = F(b,d) + F(a,c) - F(b,c) - F(a,d)$.
        \end{tabular}
    \end{proof}
\end{enumerate}

Означення сумісної функції розподілу та всі властивості узагальнюються і на випадкові вектори довільної розмірності.
\begin{definition} \index{функція розподілу!сумісна}
    \emph{Сумісною функцією розподілу} $n$-вимірного випадкового вектора $\vec{\xi}$ 
    називається $F_{\vec{\xi}}(\vec{x}) = \P\left\{\omega: \xi_1 < x_1, \xi_2 < x_2, ..., \xi_n < x_n\right\}$, $\vec{x} \in \mathbb{R}^n$.
    Функції розподілу координат цього випадкового вектора називаються \emph{маргінальними функціями розподілу} \index{функція розподілу!маргінальна}.
\end{definition}
\begin{definition}\index{випадкова величина!незалежні величини!в сукупності}
    Випадкові величини $\xi_1, \xi_2, ..., \xi_n$ називаються
    \emph{незалежними (у сукупності)}, якщо для всіх
    $x_1, x_2, ... , x_n$
    події $\left\{\xi_k < x_k\right\}, k=1,...,n$ є незалежними в сукупності.
\end{definition}

\noindent\textbf{Властивості сумісної функції розподілу:}
\begin{enumerate}
    \item Область визначення --- $\mathbb{R}^n$, область значень --- $\left<0; 1\right>$.
    \item Є монотонно неспадною за кожним аргументом.
    \item Є неперервною зліва за кожним аргументом.
    \item $\forall k = 1,...,n: \underset{x_k \to -\infty}{\lim} F_{\vec{\xi}}(\vec{x}) = 0$.
    \item $\lim\limits_{x_1, ..., x_n \rightarrow +\infty} F_{\vec{\xi}}(\vec{x}) = 1$.
    \item \emph{<<Умови узгодженості>>} \index{умова!узгодженості (умови)}: 
        $$\forall k = 1,...,n: \lim\limits_{\forall i\neq k : x_i \to + \infty} F_{\vec{\xi}}(\vec{x}) = F_{\xi_k}(x_k)$$
        $$\forall I = \left\{i_1, i_2, ..., i_k\right\} \in \{1,...,n\}: 
        \lim\limits_{\forall i \notin I : x_i \to + \infty} F_{\vec{\xi}}(\vec{x}) =
        F_{\xi_{i_1} \xi_{i_2} ... \xi_{i_k}}(x_{i_1}, x_{i_2}, ..., x_{x_k})$$.
    \item Якщо координати $\xi_1, \xi_2, ..., \xi_n$ є незалежними у сукупності, то 
    $F_{\vec{\xi}}(\vec{x}) = F_{\xi_1}(x_1) \cdot ... \cdot F_{\xi_n}(x_n)$.
\end{enumerate}
\begin{exercise}
    Довести ці властивості.
\end{exercise}

\subsection{Дискретні випадкові вектори}
\begin{definition}\index{випадковий вектор!дискретний}
    Випадковий вектор називається \emph{дискретним}, якщо всі його координати --- 
    дискретні випадкові величини.
\end{definition}

Закон розподілу двовимірного дискретного випадкового вектора задається 
\emph{таблицею розподілу}:\index{таблиця розподілу}

\begin{center}
\begin{tabular}{c c}
    \begin{tabular}{|c|c|c|c|c|}
        \hline
        \diagbox{$\xi_2$}{$\xi_1$} & $x_1$ & $x_2$ & ... & $x_n$ \\
        \hline
        $y_1$ & $p_{11}$ & $p_{21}$ & ... & $p_{n1}$ \\
        \hline
        $y_2$ & $p_{12}$ & $p_{22}$ & ... & $p_{n2}$ \\
        \hline
        ... & ... & ... & $p_{ij}$ & ... \\
        \hline
        $y_m$ & $p_{1m}$ & $p_{2m}$ & ... & $p_{nm}$ \\
        \hline
    \end{tabular} 
    &
    \begin{tikzpicture}[baseline={(current bounding box.center)}]
        \draw [->] (-0.5, 0) -- (5, 0); 
        \draw [->] (0, -0.5) -- (0, 2.5);
        %\draw [fill] (1, 1) circle [radius = 0.05];
        \foreach \i in {1,...,2}:
            %\draw [dashed] (\i, 0) -- (\i, 1.5);
            \foreach \j in {1,...,2}:
                \draw [fill] (\i, {\j/2}) circle [radius = 0.05];
        \foreach \k in {1,...,2} {
            \draw [fill] (4, {\k/2}) circle [radius = 0.05];
            \draw [fill] (\k, 2) circle [radius = 0.05];
            \draw [dashed] (\k, 0) -- (\k, 2);
            \node [below] at (\k, 0) {$x_\k$};
            \node [left] at (0, {\k/2}) {$y_\k$};
            \draw [dashed] (0, {\k/2}) -- (4, {\k/2});
        }
        \draw [fill] (4, 2) circle [radius = 0.05];
        \draw [dashed] (4, 0) -- (4, 2);
        \node [below] at (3, -0.1) {$\ldots$};
        \node [below] at (4, 0) {$x_n$};
        \node [below] at (5, 0) {$x$};
        \draw [dashed] (0, 2) -- (4, 2);
        \node [left] at (0, 1.6) {$\vdots$};
        \node [left] at (0, 2) {$y_m$};
        \node [left] at (0, 2.5) {$y$};
    \end{tikzpicture} \\
    $p_{ij} = \P\left\{\xi_1 = x_i, \xi_2 = y_j\right\}$, $\sum\limits_{i,j} p_{ij} = 1$ & \\
\end{tabular}
\end{center}

З таблиці розподілу можна обчислити ряди розподілу $\xi_1$ та $\xi_2$:\index{ряд розподілу}
\begin{center}
    \begin{tabular}{c c}
        \begin{tabular}{|c|c|c|c|c|}
            \hline
            $\xi_1$ & $x_1$ & $x_2$ & $...$ & $x_n$ \\
            \hline
            $p$ & $\sum\limits_{k=1}^m p_{1k}$ & $\sum\limits_{k=1}^m p_{2k}$ & $...$ & $\sum\limits_{k=1}^m p_{nk}$ \\
            \hline
        \end{tabular} &
        \begin{tabular}{|c|c|c|c|c|}
            \hline
            $\xi_2$ & $y_1$ & $y_2$ & $...$ & $y_m$ \\
            \hline
            $p$ & $\sum\limits_{k=1}^n p_{k1}$ & $\sum\limits_{k=1}^n p_{k2}$ & $...$ & $\sum\limits_{k=1}^n p_{km}$ \\
            \hline
        \end{tabular}
    \end{tabular}
\end{center}

\begin{remark}
    Обчислити таблицю розподілу, знаючи ряди розподілу координат, можна лише у випадку їх незалежності.
\end{remark}

Побудову функції розподілу для дискретного випадкового вектора простіше розглянути на прикладі.
\begin{example}
    Нехай двовимірний випадковий вектор  задано таблицею розподілу:
    \begin{center}
        \begin{tabular}{c c}
            \begin{tabular}{|c|c|c|}
                \hline
                \diagbox{$\xi_2$}{$\xi_1$} & $1$ & $2$ \\
                \hline
                $0$ & $0.3$ & $0.4$ \\
                \hline
                $1$ & $0.2$ & $0.1$ \\
                \hline
            \end{tabular}
            &
            \begin{tikzpicture}[baseline={(current bounding box.center)}]
                \draw [->] (-0.5, 0) -- (3, 0); 
                \draw [->] (0, -0.5) -- (0, 2);
                \draw [fill] (1, 0) circle [radius = 0.05];
                \draw [fill] (1, 1) circle [radius = 0.05];
                \draw [fill] (2, 0) circle [radius = 0.05];
                \draw [fill] (2, 1) circle [radius = 0.05];
                \draw [dashed] (0, 1) -- (2, 1);
                \draw [dashed] (1, 1) -- (1, 0);
                \draw [dashed] (2, 1) -- (2, 0);
                \node [below] at (1, 0) {$1$};
                \node [below] at (2, 0) {$2$};
                \node [left] at (0, 1) {$1$};
                \node [below left] at (0, 0) {$0$};
                \node [below] at (3, 0) {$t$};
                \node [left] at (0, 2) {$s$};
            \end{tikzpicture}
        \end{tabular}
    \end{center}
\end{example}
Значення функції розподілу також зручно звести в таблицю:
\begin{center}
    \begin{tabular}{|c|c|c|c|}
        \hline
        \diagbox[height=2em, width=6em]{$y$}{$x$} & $x\leq1$ & $1<x\leq2$ & $x> 2$ \\
        \hline
        $y\leq0$ & $0$ & $0$ & $0$ \\
        \hline
        $0<y\leq 1$ & $0$ & $0.3$ & $0.7$ \\
        \hline
        $y>1$ & $0$ & $0.5$ & $1$ \\
        \hline
    \end{tabular}
\end{center}
Суть побудови цієї таблиці --- у підрахунку кількості точок, що потрапляють у
квадрант $\left\{t < x, s < y \right\}$ при різних значеннях
$x$ та $y$, та обчисленні відповідних ймовірностей. Зауважимо,
що з умов узгодженості випливає, що
останній рядок таблиці --- значення $F_{\xi_1} (x)$, 
а останній стовпчик --- значення $F_{\xi_2} (y)$.
        % !TEX root = ../main.tex
\section{Неперервні випадкові вектори}

\begin{definition}
    Вимірна функція $\xi(\omega): \Omega \rightarrow \mathbb{R}^n$ називається 
    \emph{неперервним випадковим вектором}, якщо координати $\xi_i$ є 
    неперервними випадковими величинами для $i = 1,...,n$.
\end{definition}
Функція розподілу $F_{\vec{\xi}}\left(\vec{x}\right) = 
P\left\{\omega:\xi_1(\omega)<x_1,...,\xi_n(\omega)<x_n\right\}$ є неперервною по всім аргументам.

\subsection{Щільність розподілу двовимірного неперервного випадкового вектора}
\begin{definition}
    \emph{Щільністю розподілу} двовимірного випадкового вектора називається 
    подвійна границя
    \begin{equation}
        f_{\vec{\xi}}(x, y) = \lim_{\substack{\Delta x \to 0 \\ 
        \Delta y \to 0}} \frac{P\left\{\vec{\xi} \in \Pi\right\}}
        {\Delta x \Delta y}
    \end{equation}
    де $\Pi = \left[x; x+\Delta x\right] \times \left[y; y+\Delta y\right]$.
\end{definition}
\begin{remark}
    В означенні 1 замість можна брати будь-яку обмежену замкнену область (компакт) 
    $D$ і тоді
    \begin{equation*}
        f_{\vec{\xi}}(x, y) = \lim_{S(D) \to 0} 
        \frac{P\left\{\vec{\xi} \in D\right\}}
        {S(D)}
    \end{equation*}
\end{remark}

\noindent\textbf{Зв’язок щільності розподілу з сумісною функцією розподілу: }
\begin{gather*}
    \lim_{\substack{\Delta x \to 0 \\ 
\Delta y \to 0}} \frac{P\left\{\vec{\xi} \in \Pi\right\}}
{\Delta x \Delta y} = 
\lim_{\substack{\Delta x \to 0 \\ \Delta y \to 0}} 
\left(
    \frac{F_{\vec{\xi}}(x+\Delta x, y+\Delta y) - F_{\vec{\xi}}(x, y+\Delta y)}
    {\Delta x \Delta y}
    -
    \frac{F_{\vec{\xi}}(x+\Delta x, y) - F_{\vec{\xi}}(x, y)}
    {\Delta x \Delta y}
\right) = \\
= \left[\text{формула Лагранжа про скінченні прирости}, \theta_1, \theta_2 \in (0, 1)\right] = \\
= \lim_{\substack{\Delta x \to 0 \\ \Delta y \to 0}} 
\left(
    \frac{\frac{\partial F}{\partial x}(x + \theta_1 \Delta x, y + \Delta y)\Delta x}
    {\Delta x \Delta y}
    -
    \frac{\frac{\partial F}{\partial x}(x + \theta_2 \Delta x, y)\Delta x}
    {\Delta x \Delta y}
\right) =  \\
= \left[\theta_3, \theta_4 \in (0, 1)\right] 
= \lim_{\substack{\Delta x \to 0 \\ \Delta y \to 0}}
\frac{\frac{\partial^2 F}{\partial x \partial y}(x + \theta_3 \Delta x, 
y + \theta_4 \Delta y)\Delta y}
{\Delta y} = \frac{\partial^2 F_{\vec{\xi}}}{\partial x \partial y}(x, y)\end{gather*}

Таким чином, якщо існує неперервна $ \frac{\partial^2 F_{\vec{\xi}}}{\partial x \partial y}$, то
\begin{equation}\label{eq:dens_r2}
    f_{\vec{\xi}}(x, y) = \frac{\partial^2 F_{\vec{\xi}}}{\partial x \partial y}(x, y)
\end{equation}

\begin{definition}
    Поверхня, що є графіком щільності двовимірного випадкового вектора називається 
    \emph{поверхнею розподілу}.
\end{definition}
\begin{definition}
    Лінії, де $f_\xi(x, y) = const$, називаються \emph{лініями рівних ймовірностей}.
\end{definition}

\noindent\textbf{Властивості щільності: }

\begin{enumerate}
    \item $f_{\vec{\xi}}(x, y) \geq 0$.
    \item З формули \refeq{eq:dens_r2}: $F_{\vec{\xi}}(x, y) 
    = \int\limits_{-\infty}^x \int\limits_{-\infty}^y f_{\vec{\xi}}(t, s) 
    dt ds $.
    \item Умова нормування: 
    $\int\limits_{-\infty}^{+\infty} \int\limits_{-\infty}^{\infty} f_{\vec{\xi}}(t, s) 
    dt ds = 1$ --- об'єм під поверхнею розподілу 
    дорівнює 1.
    \item $F_{\xi_1}(x) = \int\limits_{-\infty}^{x} \int\limits_{-\infty}^{+\infty} 
    f_{\vec{\xi}}(t, s) dt ds$.

    $F_{\xi_2}(y) = \int\limits_{-\infty}^{+\infty} \int\limits_{-\infty}^{y} 
    f_{\vec{\xi}}(t, s) dt ds$.
    \item Щільності розподілу окремих координат називаються 
    \emph{маргінальними щільностями}.

    $f_{\xi_1}(x) = F'_{\xi_1}(x) 
    = \int\limits_{-\infty}^{+\infty} f_{\vec{\xi}}(x, y) dy $.

    $f_{\xi_2}(y) = F'_{\xi_2}(y) 
    = \int\limits_{-\infty}^{+\infty} f_{\vec{\xi}}(x, y) dx $.
    
    \item D --- замкнена обмежена область в $\mathbb{R}^2$.
    $P\left\{ \vec{\xi} \in D \right\} = \iint\limits_D f_{\vec{\xi}}(x, y) 
    dx dy$.
\end{enumerate}

\subsection{Рівномірний закон розподілу на площині}
\begin{definition}
    D --- замкнена обмежена область в $\mathbb{R}^2$. 
    Вектор $\vec{\xi} = (\xi_1,\xi_2)^T$ називається 
    \emph{рівномірно розподіленим в області D}, якщо 
    \begin{equation*}
        f_{\vec{\xi}}(x, y) = 
        \begin{cases}
            \frac{1}{S(D)},&(x, y) \in D \\
            0,&(x, y) \notin D
        \end{cases}
    \end{equation*}
\end{definition}
\begin{example}
    \begin{tabular}{c m{4cm}}
        $f_{\vec{\xi}}(x, y) = 
        \begin{cases}
            2, & (x, y) \in D \\
            0, & (x, y) \notin D
        \end{cases}
        $
        &
        \begin{tikzpicture}[scale = 1.5]
            \draw [->] (0, -0.5) -- (0, 1.5);
            \draw [->] (-0.5, 0) -- (1.5, 0);
            \draw (1, 0) -- (0, 1);
            \fill [lightgray] (0, 0) -- (1, 0) -- (0, 1);
            \node [below] at (1.5, 0) {$x$};
            \node [left] at (0, 1.5) {$y$};
            \node [above right] at (0.15, 0.15) {$D$};
            \node [right] at (0.6, 0.6) {$y = 1 - x$};
        \end{tikzpicture}
    \end{tabular}
    
    \begin{enumerate}
        \item $f_{\xi_1}(x) = \int\limits_{-\infty}^{+\infty} f_{\vec{\xi}}(x, y) dy = 
        \begin{cases}
            0 , &x\leq0 \text{ або } x>1\\
            \int\limits_0^{1-x} 2 dy = 2(1-x), & 0 < x \leq 1 
        \end{cases}$

        $f_{\xi_2}(y) = \int\limits_{-\infty}^{+\infty} f_{\vec{\xi}}(x, y) dx = 
        \begin{cases}
            0 , &y\leq0 \text{ або } y>1\\
            \int\limits_0^{1-y} 2 dx = 2(1-y), & 0 < y \leq 1 
        \end{cases}$
        \item Функції розподілу координат: 
        
        $F_{\xi_1}(x) = \int\limits_{-\infty}^{x} f_{\xi_1}(t) dt = \begin{cases}
            0, & x\leq 0 \\
            2\int\limits_0^x (1-t) dt = 2(x-\frac{x^2}{2}) = 2x - x^2, & 0<x\leq 1 \\
            1, x>1
        \end{cases}$
        
        $F_{\xi_2}(y) = \int\limits_{-\infty}^{y} f_{\xi_2}(t) dt = \begin{cases}
            0, & y\leq 0 \\
            2y - y^2, & 0<y\leq 1 \\
            1, y>1
        \end{cases}$
        \item Сумісна функція розподілу: 
        
        $F_{\vec{\xi}}(x, y) = 
        \int\limits_{-\infty}^{x} \int\limits_{-\infty}^{y} 
        f_{\vec{\xi}}(t, \tau) dt d\tau$

        \begin{enumerate}[label=(\Roman*)]
            \item $(x, y) \in D_0$
            
            $D_0 = \left\{(x, y):\; x \leq 0 \lor y \leq 0\right\}$

            $F_{\vec{\xi}}(x, y) = 0$

            \item $(x, y) \in D_1$
            
            $D_1 = \left\{(x, y):\; x \in \left[0, 1\right],
            y \in \left[0, 1-x\right]\right\}$

            $F_{\vec{\xi}}(x, y) = 2xy$

            \item $(x, y) \in D_2$
            
            $D_2 = \left\{(x, y):\; x \in \left[0, 1\right],
            y \in \left[1-x, 1\right]\right\}$

            $F_{\vec{\xi}}(x, y) = 2(xy - \frac{1}{2}(x-1+y)(y-1+x))$

            \item $(x, y) \in D_3$
            
            $D_3 = \left\{(x, y):\; x > 1,
            y \in \left[0, 1\right]\right\}$

            $F_{\vec{\xi}}(x, y) = 2\frac{1+1-y}{2}y = y(2-y)$

            \item $(x, y) \in D_4$
            
            $D_4 = \left\{(x, y):\; y > 1,
            x \in \left[0, 1\right]\right\}$

            $F_{\vec{\xi}}(x, y) = 2\frac{1+1-x}{2}x = x(2-x)$

            \item $(x, y) \in D_5$
            
            $D_5 = \left\{(x, y):\; x > 1 \land y > 1\right\}$

            $F_{\vec{\xi}}(x, y) = 1$
        \end{enumerate}
        \begin{tikzpicture}[scale = 1.5]
            \fill [orange] (-0.5, -0.5) -- (-0.5, 1.5) -- (0, 1.5) --
                           (0, 0) -- (1.5, 0) -- (1.5, -0.5) -- (-0.5, -0.5);
            \fill [lightgray] (0, 0) -- (1, 0) -- (0, 1);
            \fill [pink] (0, 1) -- (1, 0) -- (1, 1);
            \fill [red] (0, 1) rectangle (1, 1.5);
            \fill [yellow] (1, 0) rectangle (1.5, 1);
            \fill [lime] (1, 1) rectangle (1.5, 1.5);
            \draw [->] (0, -0.5) -- (0, 1.5);
            \draw [->] (-0.5, 0) -- (1.5, 0);
            \draw (1, 0) -- (0, 1);
            \node [right] at (1.5, 0) {$t$};
            \node [above] at (0, 1.5) {$\tau$};
            \node [below left] at (0, 0) {$D_0$};
            \node [above right] at (0.15, 0.15) {$D_1$};
            \node [above right] at (0.5, 0.5) {$D_2$};
            \node at (0.5, 1.25) {$D_4$};
            \node at (1.25, 0.5) {$D_3$};
            \node at (1.25, 1.25) {$D_5$};
        \end{tikzpicture}
    \end{enumerate}
\end{example}

\begin{remark}
    Перевірити правильність побудови функції розподілу можна за допомогою 
    умови узгодженості, перевірки точок стику та неперервності на лініях стику.
\end{remark}

\subsection{Щільність розподілу n-вимірного неперервного випадкового 
вектора}
\begin{definition}
    $\vec{\xi} = \left(\xi_1, \xi_2, ..., \xi_n\right)^T,\; \Pi = 
    \left[x_1; x_1+\Delta x_1\right] \times ... \times 
    \left[x_n; x_n+\Delta x_n\right]$
    \emph{Щільністю розподілу} n-вимірного неперервного випадкового 
    вектора називається границя
    \begin{equation*}
        f_{\vec{\xi}} (\vec{x}) = \lim_{\substack{\Delta x_i \to 0 \\
        \forall i = \overline{1,n}}} 
        \frac{P\left\{\vec{\xi} \in \Pi\right\}}{\Delta x_1...\Delta x_n}
    \end{equation*}

    Якщо існує і неперервна $\frac{\partial^n F_{\vec{\xi}}(\vec{x})}
    {\partial x_1 ... \partial x_n}$, то
    \begin{equation*}
        f_{\vec{\xi}} (\vec{x}) = \frac{\partial^n F_{\vec{\xi}}}
        {\partial x_1 ... \partial x_n}(\vec{x})
    \end{equation*}
\end{definition}

\noindent \textbf{Властивості}
\begin{enumerate}
    \item $f_{\vec{\xi}}(\vec{x}) \geq 0$
    \item $\int\limits_{-\infty}^{+\infty}...\int\limits_{-\infty}^{+\infty}
    f_{\vec{\xi}} (\vec{x})dx_1...dx_n = 1$
    \item $F_{\vec{\xi}}(\vec{x}) = 
    \int\limits_{-\infty}^{x_1}...\int\limits_{-\infty}^{x_n}
    f_{\vec{\xi}} (\vec{t})dt_1...dt_n$
    \item $F_{\xi_1\xi_2...\xi_k}(x_1, ..., x_k) = 
    \int\limits_{-\infty}^{x_1}...\int\limits_{-\infty}^{x_k}
    \int\limits_{-\infty}^{+\infty}...\int\limits_{-\infty}^{+\infty}
    f_{\vec{\xi}} (\vec{t})dt_1...dt_n$

    $F_{\xi_i}(x_i) = \int\limits_{-\infty}^{+\infty}...
    \int\limits_{-\infty}^{x_i}...\int\limits_{-\infty}^{+\infty}
    f_{\vec{\xi}} (\vec{t})dt_1...dt_n$

    \item $f_{\xi_1...\xi_k}(x_1, ..., x_k) = 
    \underbrace{
        \int\limits_{-\infty}^{+\infty} 
        ... 
        \int\limits_{-\infty}^{+\infty}
    }_{n-k} f_{\vec{\xi}}(\vec{x}) dx_{k+1}..dx_n$
    \item D - замкнена обмежена область в $\mathbb{R}^n$ (компакт).
    
    $P\left\{\vec{\xi} \in D\right\} = {\int\int...\int}_D f_{\vec{\xi}}(\vec{x})
    dx_1 ... dx_n$
\end{enumerate}
\begin{exercise}
    Довести властивості.
\end{exercise}

\subsection{Незалежні випадкові величини}
\begin{definition}
    $\xi_1$ та $\xi_2$ називаються незалежними, якщо 
    $\forall x, y \in \mathbb{R}$:  
    $\left\{\omega\;:\;\xi_1 < x\right\}$ та $\left\{\omega\;:\;\xi_2 < y\right\}$ ---
    незалежні (як події).
\end{definition}
\begin{remark}
    Якщо $\xi_1$, $\xi_2$ - незалежні координати випадкового вектора $\vec{\xi}$, 
    то $F_{\vec{\xi}}(x, y) = 
    P\left(\left\{\xi_1 < x\right\} \cap \left\{\xi_2 < y\right\}\right) = 
    F_{\xi_1}(x)F_{\xi_2}(y)$.

    Якщо ж вектор ще й є неперервним випадковим вектором та існує неперервна 
    друга часткова похідна, то
    $f_{\vec{\xi}}(x, y) = \frac{\partial^2 F_{\vec{\xi}}(x, y)}
    {\partial x \partial y} = \frac{\partial^2(F_{\xi_1}(x)F_{\xi_2}(y))}{\partial x \partial y} 
    = F_{\xi_1}^\prime (x)F_{\xi_2}^\prime (y) = 
    f_{\xi_1}(x)f_{\xi_2}(y)$.
\end{remark}
        % !TEX root = ../main.tex
\section{Числові характеристики випадкових векторів}

\subsection{Математичне сподівання випадкових векторів}
\begin{definition}
    \emph{Математичним сподіванням $E{\vec{\xi}}$} 
    випадкового вектора $\vec{\xi}$ називається вектор 
    $\left(E{\xi_1}, ..., E{\xi_n}\right)^T$.
\end{definition}
\begin{remark}
    Математичне сподівання двовимірного випадкового вектора 
    називається \emph{центром розсіювання}.
\end{remark}

\noindent \textbf{Способи знаходження:}
\begin{enumerate}
    \item Знайти закони розподілу окремих координат, а далі --- математичні сподівання цих координат.
    \item Знайти математичні сподівання координат одразу.
    Для $\xi = (\xi_1, \xi_2)^T$:
    
    $E\xi_1 = \begin{cases}
        \sum\limits_i \sum\limits_j x_i p_{ij}, & \vec{\xi} \text{ --- ДВВ} \\
        \int\limits_{-\infty}^{+\infty} \int\limits_{-\infty}^{+\infty} x f_{\vec{\xi}}(x,y) dx dy, & \vec{\xi} \text{ --- НВВ}
    \end{cases}$,
    $E\xi_2 = \begin{cases}
        \sum\limits_j \sum\limits_i y_j p_{ij}, & \vec{\xi} \text{ --- ДВВ} \\
        \int\limits_{-\infty}^{+\infty} \int\limits_{-\infty}^{+\infty} y f_{\vec{\xi}}(x,y) dx dy, & \vec{\xi} \text{ --- НВВ}
    \end{cases}$.
    
    В загальному випадку для НВВ $E\xi_k = \int\limits_{-\infty}^{+\infty}...\int\limits_{-\infty}^{+\infty} x_k f_{\vec{\xi}}(\vec{x}) dx_1 ... dx_n$.
\end{enumerate}

\subsection{Мішані початкові та центральні 
            моменти випадкових векторів}
\begin{definition}
    \emph{Мішаним початковим моментом} порядку 
    $k_1+k_2+...+k_n$ 
    $\left(k_i \in \mathbb{N}\right)$
    випадкового вектора 
    $\vec{\xi} = \left(\xi_1, ..., \xi_n\right)^T$
    називається число
    \begin{equation*}
        \alpha_{k_1+k_2+...+k_n} = 
        E{\xi_1^{k_1}...\xi_n^{k_n}}
    \end{equation*}
\end{definition}
\begin{remark}
    $E\xi_1 = \alpha_{1+0+...+0}, E\xi_2 = \alpha_{0+1+...+0},
    ..., E\xi_n = \alpha_{0+0+...+1}$.
\end{remark}
\begin{definition}
    \emph{Мішаним центральним моментом} порядку 
    $k_1+k_2+...+k_n$
    $\left(k_i \in \mathbb{N}\right)$ 
    випадкового вектора
    $\vec{\xi} = \left(\xi_1, ..., \xi_n\right)^T$
    називається число
    \begin{equation*}
        \beta_{k_1+k_2+...+k_n} = 
        E{\mathring{\xi}_1^{k_1}
        ...
        \mathring{\xi}_n^{k_n}},
        \mathring{\xi_k} = \xi_k - E\xi_k
    \end{equation*}
\end{definition}
\begin{remark}
    Всі центральні моменти 1-го порядку --- 
    нульові.
\end{remark}
Центральні моменти порядку $0+...+0+2+0+...+0$ --- дисперсії відповідних координат:
    $\forall k = \overline{1,n}$: 
    $\beta_{0+...+0+\underset{k}{2}
    +0+...+0} = E\mathring{\xi}_k^2 = 
    D\xi_k$

Дисперсії координат задають розсіювання 
вздовж відповідних координатних осей.

\begin{definition}
    \emph{Кореляційним моментом} або 
    \emph{коваріацією} випадкових величин 
    $\xi_i$ та $\xi_j$ називається змішаний 
    центральний момент 
    порядку
    $0+...+0+\underset{i}{1}+0+...+0+\underset{j}{1}+0+....+0$:
    \begin{equation*}
        {cov}(\xi_i,\xi_j) = K\xi_i\xi_j = \beta_{0+...+0+\underset{i}{1}+0+...
        +0+\underset{j}{1}+0+...+0}
        =
        E(\xi_i-E\xi_i)(\xi_j-E\xi_j)
    \end{equation*}
\end{definition}
\begin{remark}
    ${cov}(\xi_k, \xi_k) = E(\xi_k-E\xi_k)^2 = D\xi_k$.
\end{remark}
\begin{definition}
   \emph{Кореляційною матрицею} випадкового 
   вектора $\vec{\xi}$ називається матриця, у яку зібрано усі 
   кореляційні моменти випадкового вектора
   \begin{equation*}
       K = 
       \begin{pmatrix}
           D\xi_1 & K\xi_1\xi_2 & \cdots & K\xi_1\xi_n \\
           K\xi_1\xi_2 & D\xi_2 & \cdots & K\xi_2\xi_n \\
           \vdots & \vdots & \ddots & \vdots \\
           K\xi_1\xi_n & K\xi_2\xi_n & \cdots & D\xi_n
       \end{pmatrix}
   \end{equation*} 
\end{definition}
\begin{remark}
    $\mathring{\vec{\xi}} = \left(\mathring{\xi_1}, ..., \mathring{\xi_n}
    \right)^T$.
    Тоді $K = E\,\mathring{\vec{\xi}}\,\mathring{\vec{\xi}}^T$.
\end{remark}

\noindent \textbf{Властивості кореляційного моменту випадкових 
величин:}
\begin{enumerate}
    \item $K\xi_i\xi_j = K\xi_j\xi_i$.
    \begin{proof}
        Випливає з означення кореляційного моменту.
    \end{proof}
    \item $\eta_1 = a_1\xi_1 + b_1$, $\eta_2 = a_2\xi_2 + b_2$.
    Тоді $K\eta_1\eta_2 = a_1a_2K\xi_1\xi_2$.
    \begin{proof}
        $K\eta_1\eta_2 = E(\eta_1 - E\eta_1)(\eta_2 - E\eta_2) = 
        E(a_1\xi_1 + b_1 - E(a_1\xi_1 + b_1))
        (a_2\xi_2 + b_2 - E(a_2\xi_2 + b_2)) = 
        Ea_1(\xi_1 - E\xi_1)a_2(\xi_2 - E\xi_2) = 
        a_1a_2E(\xi_1 - E\xi_1)(\xi_2 - E\xi_2) = 
        a_1a_2K\xi_1\xi_2$.
    \end{proof}

    \item Зручна формула для обчислення кореляційного моменту.
    \begin{equation}
        K\xi_1\xi_2 = E\xi_1\xi_2 - E\xi_1 E\xi_2
    \end{equation}
    \begin{proof}
        $K\xi_1\xi_2 = E(\xi_1 - E\xi_1)(\xi_2 - E\xi_2) = 
        E(\xi_1\xi_2 - \xi_1E\xi_2 - \xi_2E\xi_1 + E\xi_1E\xi_2) = 
        E\xi_1\xi_2 - E\xi_1E\xi_2 - E\xi_2E\xi_1 + E\xi_1\xi_2 = 
        E\xi_1\xi_2 - E\xi_1E\xi_2$
    \end{proof}
    
    \item $D(a\xi_1 \pm b\xi_2) = a^2D\xi_1 \pm 2abK\xi_1\xi_2 + 
    b^2D\xi_2$.

    \begin{proof}
        $D(a\xi_1 \pm b\xi_2) = 
        E(a\xi_1 \pm b\xi_2 - E(a\xi_1 \pm b\xi_2))^2 =
        E(a(\xi_1-E\xi_1)\pm b(\xi_2 - E\xi_2))^2 = 
        a^2 E\mathring{\xi}_1^2 \pm 2abE\mathring{\xi}_1 E\mathring{\xi} _2 
        + b^2 E\mathring{\xi}_2^2 = 
        a^2D\xi_1 \pm 2abK\xi_1\xi_2 + b^2D\xi_2$.
    \end{proof}

    \item $\left|K\xi_1\xi_2\right| \leq \sigma_{\xi_1} \sigma_{\xi_2}$.
    \begin{proof}
        Два способи: 
        \begin{enumerate}[label=\Roman*.]
            \item $D(x\xi_1 \pm \xi_2) = x^2 D\xi_1 \pm 2xK\xi_1\xi_2 + D\xi_2 \geq 0 
            \Rightarrow 4(K\xi_1\xi_2)^2 - 4D\xi_1D\xi_2 \leq 0 
            \Rightarrow $

            $\left|K\xi_1\xi_2\right| \leq \sigma_{\xi_1} \sigma_{\xi_2}$.

            \item Введемо центровані нормовані випадкові величини $\xi_H = \frac{\xi - E\xi}
            {\sigma_\xi}$, $E\xi_H =0$, $D\xi_H = 1$.

            $E(\xi_{1H} \pm \xi_{2H})^2 \geq 0 \Leftrightarrow 
            E(\xi^2_{1H}\pm2\xi_{1H}\xi_{2H}+\xi^2_{2H}) \geq 0 
            \Leftrightarrow 2 \pm 
            2\frac{K\xi_1\xi_2}{\sigma_{\xi_1}\sigma_{\xi_2}} \geq 0 
            \Leftrightarrow $

            $\left|K\xi_1\xi_2\right| \leq \sigma_{\xi_1} \sigma_{\xi_2}$.
        \end{enumerate}
    \end{proof}
    \item Якщо $\xi_1$ та $\xi_2$ --- незалежні, то $K\xi_1\xi_2 = 0$.
    \begin{proof}
        Нехай $\xi_1$ та $\xi_2$ --- незалежні.

        $K\xi_1\xi_2 = E\xi_1\xi_2 - E\xi_1E\xi_2 = 
        \iint\limits_{\mathbb{R}^2} xy f_{\vec{\xi}}(x,y)dx dy 
        - E\xi_1E\xi_2
        = \int\limits_{-\infty}^{+\infty}xf_{\xi_1}(x)dx \cdot
        \int\limits_{-\infty}^{+\infty}yf_{\xi_2}(y)dy
        - E\xi_1E\xi_2 = 
        E\xi_1E\xi_2 - E\xi_1E\xi_2 = 0$.
    \end{proof}

    Обернене твердження, взагалі кажучи, не має місця. Тому випадкові 
    величини, для яких $K\xi_1\xi_2 = 0$, називаються 
    \emph{некорельованими}.
\end{enumerate}

\begin{example}(некорельованість $\nRightarrow$ незалежність)
    
    \begin{tabular}{c p{7cm}}
        \begin{tikzpicture}[baseline={(current bounding box.center)}]
            \draw [fill = black!10] (0, 0) circle [radius = 1];
            \draw [->] (-1.5, 0) -- (1.5, 0);
            \draw [->] (0, -1.5) -- (0, 1.5);
            \node [below right] at (0.95, 0) {$R$};
            \node [below] at (1.5, 0) {$x$};
            \node [left] at (0, 1.5) {$y$};
        \end{tikzpicture} &
        $S = \left\{(x,y) \in \mathbb{R}^2 : x^2 + y^2 \leq R^2\right\}$, $\vec{\xi} \sim \mathrm{U}(S)$.

        $f_{\vec{\xi}}(x, y) = 
            \begin{cases}
            \frac{1}{\pi R^2} & (x, y) \in S \\
            0 & (x, y) \notin S
        \end{cases}$
    \end{tabular}

\begin{tabular}{c c}
    $f_{\xi_1}(x) = 
        \begin{cases}
            \int\limits_{-\sqrt{R^2 - x^2}}^{\sqrt{R^2 - x^2}} \frac{1}{\pi R^2} dy
            = \frac{2\sqrt{R^2-x^2}}{\pi R^2},& |x| \leq R \\
            0, & |x| > R
        \end{cases}
    $,
    &
    $f_{\xi_2}(y) = 
    \begin{cases}
        \frac{2\sqrt{R^2-y^2}}{\pi R^2},& |y| \leq R \\
        0, & |y| > R
    \end{cases}
    $
\end{tabular}

$E\xi_1 = E\xi_2 = \int\limits_{-R}^R t\cdot\frac{2\sqrt{R^2-t^2}}{\pi R^2} dt = 0$, бо інтегрується непарна функція по симетричному проміжку.

$K\xi_1\xi_2 = \iint\limits_S xy \frac{1}{\pi r^2} dx dy - E\xi_1E\xi_2 =
\frac{1}{\pi R^2}\int\limits_0^{2\pi} d\varphi \int\limits_0^R r^2 \cos \varphi \sin \varphi r dr = 
\frac{1}{2\pi R^2}\int\limits_0^{2\pi} \sin 2\varphi d\varphi \int\limits_0^R r^3 dr = 0$.
Отже, $f_{\vec{\xi}}(x,y) \neq f_{\xi_1}(x)f_{\xi_2}(y)$ і координати вектора залежні, але $K\xi_1\xi_2 = 0$.
\end{example}

\subsection{Коваріація як скалярний добуток випадкових величин}
Зафіксуємо ймовірнісний простір $\left\{\Omega, \mathcal{F}, P\right\}$.
Позначимо $\mathcal{L}^2(\Omega) = \left\{\xi : \Omega \rightarrow \mathbb{R} \;|\; E\xi^2 < +\infty\right\}$.
Оскільки $D\xi = E\xi^2 - (E\xi)^2 \geq 0$, то для $\xi \in \mathcal{L}^2(\Omega)$ $|E\xi| < +\infty$.
Таким чином, для будь-яких $\xi, \eta \in \mathcal{L}^2(\Omega)$ 
існує ${cov}(\xi,\eta)$, бо $|{cov}(\xi,\eta)| \leq \sqrt{D\xi} \cdot \sqrt{D\eta}$.

Для коваріації маємо 
${cov}(a\xi_1 + b\xi_2,\eta) = E(a\xi_1+b\xi_2)\eta - E(a\xi_1 + b\xi_2)E\eta = 
a E\xi_1\eta + b E\xi_2\eta - a E\xi_1 E\eta - b E\xi_2 E\eta = a\cdot{cov}(\xi_1,\eta) + b\cdot{cov}(\xi_2,\eta)$.
Також було доведено ${cov}(\xi, \eta) = {cov}(\eta, \xi)$ та ${cov}(\xi, \xi) \geq 0$.
Виконуються всі умови скалярного добутку, окрім ${cov}(\xi, \xi) = 0 \Rightarrow \xi = 0$.
Ця умова буде виконуватися, якщо розглядати інший простір випадкових величин:
$\mathcal{L}_0^2(\Omega) = \left\{\xi : \Omega \rightarrow \mathbb{R} \;|\; E\xi = 0, E\xi^2 < +\infty\right\}$.
\begin{remark}
    Якщо $\xi \in \mathcal{L}^2(\Omega)$, то $\mathring{\xi} \in \mathcal{L}_0^2(\Omega)$.
\end{remark}
\begin{exercise}
    Перевірити, що $\mathcal{L}_0^2(\Omega)$ є лінійним простором.
\end{exercise}

\noindent \textbf{Висновок:} ${cov}(\xi, \eta)$ задає \emph{скалярний добуток} на $\mathcal{L}_0^2(\Omega)$.

Таким чином, властивість $|{cov}(\xi,\eta)| \leq \sqrt{D\xi} \cdot \sqrt{D\eta}$ ---
це нерівність Коші-Буняковського, а кореляційна матриця $K$ --- це матриця Грама системи
випадкових величин $\left\{\xi_1, \xi_2, ..., \xi_n\right\}$.
З курсу лінійної алгебри відомо, що якщо ці випадкові величини лінійно незалежні, то матриця $K$ є додатно визначеною, 
і невід'ємно визначеною в загальному випадку.

\subsection{Коефіцієнт кореляції}
\begin{definition}
    \emph{Коефіцієнтом кореляції} називають безрозмірну числову характеристику
    \begin{equation*}
        \rho_{\xi_1\xi_2} = \frac{K\xi_1\xi_2}{\sigma_{\xi_1}\sigma_{\xi_2}}
    \end{equation*}
\end{definition}

\begin{definition}
    Всі коефіцієнти кореляції збираються в 
    \emph{нормовану кореляційну матрицю}:

    \begin{equation*}
        R = 
        \begin{pmatrix}
            1 & \rho_{\xi_1\xi_2} & \cdots & \rho_{\xi_1\xi_n} \\
            \rho_{\xi_1\xi_2} & 1 & \cdots & \rho_{\xi_2\xi_n} \\
            \vdots & \vdots & \ddots & \vdots \\
            \rho_{\xi_1\xi_n} & \rho_{\xi_2\xi_n} & \cdots & 1
        \end{pmatrix}
    \end{equation*}

\end{definition}

\noindent \textbf{Властивості коефіцієнту кореляції:}
\begin{enumerate}
    \item Для некорельованих випадкових величин $\xi_1$ та $\xi_2$ 
    $\rho_{\xi_1\xi_2} = 0$.
    \begin{proof}
        Випливає з визначення коефіцієнта кореляції.
    \end{proof}
    \item $\left|\rho_{\xi_1\xi_2}\right| \leq 1$.
    \begin{proof}
        Випливає з властивості кореляційного моменту.
    \end{proof}
    \item $\left|\rho_{\xi_1\xi_2}\right| = 1 \Leftrightarrow 
    \exists \; a, b \in \mathbb{R}: \xi_2 = a\xi_1 + b$.
    \begin{proof}
        Нехай 
        $\xi_2 = a\xi_1 + b$. $K\xi_1\xi_2 = 
        aD\xi_1$
        , 
        $\sigma_{\xi_2} = \sqrt{D\xi_2} = \sqrt{a^2D\xi_1} = 
        |a|\sqrt{D\xi_1}$.

        Тоді $\rho_{\xi_1\xi_2} = \frac{aD\xi_1}{|a|(\sqrt{D\xi_1})^2} 
        \Rightarrow
        \left|\rho\xi_1\xi_2\right| = 1$.

        Нехай $\left|\rho\xi_1\xi_2\right| = 1$.
        $E\left(\frac{\xi_1 - E\xi_1}{\sigma_{\xi_1}} 
        \pm \frac{\xi_2 - E\xi_2}{\sigma_{\xi_2}}\right)^2 = 
        2(1\pm\rho_{\xi_1\xi_2}) = 0$.

        $E\xi^2 = 0$ $D\xi = -(E\xi)^2 \geq 0 \Rightarrow E\xi = 0
        \Rightarrow \frac{\xi_1 - E\xi_1}{\sigma_{\xi_1}} 
        \pm \frac{\xi_2 - E\xi_2}{\sigma_{\xi_2}} = 0 \Rightarrow 
        $ лінійний зв'язок.
    \end{proof}
\end{enumerate}
        % !TEX root = ../main.tex

\section{Умовні закони розподілу}
З поняттям умовної ймовірності пов'язано поняття умовного закону розподілу випадкової величини.
Розглядаємо випадок $n = 2$, $\vec{\xi} = \left(\xi_1, \xi_2\right)^T$.

\begin{definition}
    \emph{Умовним законом розподілу} $\xi_1$ ($\xi_2$) називається закон розподілу 
    $\xi_1$ ($\xi_2$) за умови того, що $\xi_2$ ($\xi_1$) набула відповідного значення (для дискретного випадку) 
    або потрапила в деякий проміжок (для неперервного).
\end{definition}

\begin{definition}
    Універсальним умовним законом розподілу ВВ є \emph{умовна функція 
    розподілу}:
    \begin{equation*}
        F_{\xi_1}(x/y) = \P\left\{{\xi_1 < x}/{\xi_2 < y}\right\} = 
        \frac{\P\left\{\xi_1 < x, \xi_2 < y\right\}}
        {\P\left\{\xi_2 < y\right\}} = 
        \frac{F_{\vec{\xi}}(x, y)}{F_{\xi_2}(y)}, F_{\xi_2}(y/x) = \frac{F_{\vec{\xi}}(x, y)}{F_{\xi_1}(x)}
    \end{equation*}
\end{definition}
\begin{remark}
    З означення випливає, що 
    $F_{\vec{\xi}}(x,y) = F_{\xi_1}(x)F_{\xi_2}(y/x)$
    та 
    $F_{\vec{\xi}}(x,y) = F_{\xi_2}(y)F_{\xi_1}(x/y)$.
\end{remark}

\subsection{Умовний закон розподілу ДВВ}
Для дискретного випадкового вектора умовні розподіли координат також можна задати умовними ймовірностями
\begin{gather}
    \label{eq:discr_cond_x}
    \P\left\{\xi_1 = x_i / \xi_2 = y_j\right\} = 
    \frac{\P\left\{\xi_1 = x_i , \xi_2 = y_j\right\}}
    {\P\left\{\xi_2 = y_j\right\}} \\
    \label{eq:discr_cond_y}
    \P\left\{\xi_2 = y_j / \xi_1 = x_i\right\} = 
    \frac{\P\left\{\xi_1 = x_i , \xi_2 = y_j\right\}}
    {\P\left\{\xi_1 = x_i\right\}}
\end{gather}
Всі ці ймовірності вносяться до \emph{умовного ряду розподілу}.

\begin{example}\label{ex:discr_cond_distr}
    Дискретний випадковий вектор має закон розподілу:
    \begin{center}
        \begin{tabular}{ccc}
            \begin{tabular}{|c|c|c|c|}
                \hline
                \diagbox{$\xi_2$}{$\xi_1$} & $-1$ & $0$ & $1$\\
                \hline
                0 & $0.1$ & $0.2$ & $0.1$ \\
                \hline
                $1$ & $0.2$ & $0.3$ & $0.1$ \\
                \hline
            \end{tabular}
            &
            \begin{tabular}{|c|c|c|c|}
                \hline
                $\xi_1$ & $-1$ & $0$ & $1$ \\
                \hline
                $p$ & $0.3$ & $0.5$ & $0.2$ \\
                \hline
            \end{tabular}
            &
            \begin{tabular}{|c|c|c|}
                \hline
                $\xi_2$ & $0$ & $1$ \\
                \hline
                $p$ & $0.4$ & $0.6$ \\
                \hline
            \end{tabular}
        \end{tabular}
    \end{center}
    За формулами \eqref{eq:discr_cond_x} та \eqref{eq:discr_cond_y} умовні ряди розподілу мають вигляд
    \begin{center}
        \begin{tabular}{c c}
            \begin{tabular}{|c|c|c|c|}
                \hline
                $\xi_1$ & $-1$ & $0$ & $1$ \\
                \hline
                $\P\left\{\xi_1 / \xi_2 = 0\right\}$ & $1/4$ 
                & $2/4$ & $1/4$ \\
                \hline
                $\P\left\{\xi_1 / \xi_2 = 1\right\}$ & $2/6$ 
                & $3/6$ & $1/6$ \\
                \hline
            \end{tabular}
            &
            \begin{tabular}{|c|c|c|}
                \hline
                $\xi_2$ & $0$ & $1$ \\
                \hline
                $\P\left\{\xi_2 / \xi_1 = -1\right\}$ & $1/3$ 
                & $2/3$ \\
                \hline
                $\P\left\{\xi_2 / \xi_1 = 0\right\}$ & $2/5$ 
                & $3/5$\\
                \hline
                $\P\left\{\xi_2 / \xi_1 = 1\right\}$ & $1/2$ 
                & $1/2$\\
                \hline
            \end{tabular}
        \end{tabular}
    \end{center}  
\end{example}

\subsection{Умовне математичне сподівання ДВВ}
\begin{definition}
    \emph{Умовним математичним сподіванням} випадкової величини $\xi_1$ 
    є математичне сподівання цієї випадкової величини за умови, що 
    $\xi_2$ набула певного значення.
    $$\E(\xi_1 / \xi_2 = y_j) =
\sum\limits_{i=1}^{n(\infty)}x_i 
\P\left\{\xi_1 = x_i / \xi_2 = y_j\right\}, 
\E(\xi_2 / \xi_1 = x_i) = \sum\limits_{j=1}^{n(\infty)}y_j 
\P\left\{\xi_2 = y_j / \xi_1 = x_i\right\}$$
\end{definition}

\begin{remark}
    Якщо $\xi_1$ та $\xi_2$ незалежні, то $\E(\xi_1 / \xi_2 = y_j) 
    = \E\xi_1$, $\E(\xi_2 / \xi_1 = x_i) = \E\xi_2$.
\end{remark}

Умовні математичні сподівання координат дискретного випадкового вектора є дискретними випадковими величинами, оскільки 
приймають декілька значень з певними ймовірностями, тому можна скласти 
їх закон розподілу.
\begin{example}\label{ex:discr_cond_distr_2}
    Знайдемо умовні математичні сподівання для величин з прикладу \ref{ex:discr_cond_distr}.
    З умовного ряду розподілу $\xi_1$ маємо $\E(\xi_1 / \xi_2 = 0) = -\frac{1}{4} + \frac{1}{4}$ та
    $\E(\xi_1 / \xi_2 = 1) = -\frac{2}{6} + \frac{1}{6} = -\frac{1}{6}$. Для побудови ряду розподілу $\E(\xi_1 / \xi_2)$
    залишилося зауважити, що $\P\left\{\xi_2 = 0\right\} = 0.4$ та $\P\left\{\xi_2 = 1\right\} = 0.6$.
    Аналогічно будується ряд розподілу $\E(\xi_2 / \xi_1)$:
    \begin{center}
        \begin{tabular}{c c}
            \begin{tabular}{|c|c|c|}
                \hline
                $\E(\xi_1 / \xi_2)$ & $-1/6$ & $0$ \\
                \hline
                $p$ & $0.6$ & $0.4$ \\
                \hline
            \end{tabular}
            &
            \begin{tabular}{|c|c|c|c|}
                \hline
                $\E(\xi_2 / \xi_1)$ & $1/2$ & $3/5$ & $2/3$ \\
                \hline
                $p$ & $0.2$ & $0.5$ & $0.3$ \\
                \hline
            \end{tabular}
        \end{tabular}
    \end{center}
\end{example}

\subsection{Умовний закон розподілу НВВ}
Спробуємо визначити умовну щільність розподілу за аналогією до звичайної:
\begin{gather*}
    \P\left\{x \leq \xi_1 < x + \Delta x / y \leq \xi_2 < y + \Delta y\right\} = 
    \frac{
        \P\left\{\vec{\xi} \in \left[x; x+\Delta x\right) \times \left[y; y+\Delta y\right)\right\}
    }{
        \P\left\{y \leq \xi_2 < y + \Delta y \right\}
    }
    = \\ =
    \frac{f_{\vec{\xi}}(x,y)\Delta x \Delta y + 
    o\left(\sqrt{\Delta x^2 + \Delta y^2}\right)}{f_{\xi_2}(y)\Delta y + o(\Delta y)} = 
    \frac{f_{\vec{\xi}}(x,y)\Delta x  + 
    o\left(\sqrt{\Delta x^2 + \Delta y^2}\right)/\Delta y}{f_{\xi_2}(y) + o(1)}
\end{gather*}
Цей вираз приводить до наступного означення:
\begin{definition}
    \emph{Умовною щільністю розподілу} називається функція 
    вигляду:
    \begin{equation}
        f_{\xi_1}(x/y) = \frac{f_{\vec{\xi}}(x, y)}{f_{\xi_2}(y)}, f_{\xi_2}(y/x) = \frac{f_{\vec{\xi}}(x, y)}{f_{\xi_1}(x)}
    \end{equation}
\end{definition}

\begin{remark}
    Графік умовної щільності розподілу можна інтерпретувати як лінію перетину 
    поверхні розподілу та площини $y = y_{\text{знач.}}$ або
    $x = x_{\text{знач.}}$ відповідно, нормовану на одиничну площу під нею.
\end{remark}

\noindent \textbf{Властивості умовної щільності:}
\begin{enumerate}
    \item $f_{\xi_1}(x / y) \geq 0, f_{\xi_2}(y / x) \geq 0$.
    \item $f_{\xi_1}(x/y) = \frac{f_{\vec{\xi}}(x,y)}
    {\int\limits_{-\infty}^{+\infty}f_{\vec{\xi}}(x,y)dx},
    f_{\xi_2}(y/x) = \frac{f_{\vec{\xi}}(x,y)}
    {\int\limits_{-\infty}^{+\infty}f_{\vec{\xi}}(x,y)dy}$.
    \item $\int\limits_{-\infty}^{+\infty} f_{\xi_1}(x / y) dx = 
    \int\limits_{-\infty}^{+\infty}\frac{f_{\vec{\xi}}(x,y)}
    {\int\limits_{-\infty}^{+\infty}f_{\vec{\xi}}(x,y)dx}dx = 1,
    \int\limits_{-\infty}^{+\infty} f_{\xi_2}(y / x) dy = 
    \int\limits_{-\infty}^{+\infty}\frac{f_{\vec{\xi}}(x,y)}
    {\int\limits_{-\infty}^{+\infty}f_{\vec{\xi}}(x,y)dy}dy = 1$.
\end{enumerate}
Властивості умовної щільності аналогічні властивостям звичайної (безумовної) щільності.

\subsection{Умовне математичне сподівання НВВ}
\emph{Умовні математичні сподівання} координат НВВ визначаються через умовні щільності:

\begin{equation}\label{eq:cond_dens}
    \E(\xi_1 / \xi_2 = y) = \int\limits_{-\infty}^{+\infty} xf_{\xi_1}(x/y)dx, \;
    \E(\xi_2 / \xi_1 = x) = \int\limits_{-\infty}^{+\infty} yf_{\xi_2}(y/x)dy
\end{equation}

\begin{definition}
    Графіки функцій $\varphi(y) = \E(\xi_1 / \xi_2 = y)$ та 
    $\psi(x) = \E(\xi_2 / \xi_1 = x)$ називаються \emph{лініями регресії}
    $\xi_1$ на $\xi_2$ та $\xi_2$ на $\xi_1$ відповідно.
\end{definition}
\begin{example}\label{ex:cont_cond_prob}
    Знайти умовні розподіли та математичні сподівання координат вектора $\vec{\xi}$, що рівномірно розподілений в трикутнику $D$ з вершинами в точках
    $(0, 0)$, $(1, 0)$ та $(1, 1)$.
    \begin{center}
        \begin{tabular}{c m{4cm}}
            $f_{\vec{\xi}}(x, y) = 
            \begin{cases}
                2, & (x, y) \in D \\
                0, & (x, y) \notin D
            \end{cases}
            $
            &
            \begin{tikzpicture}[scale = 1.5]
                \draw [->] (0, -0.5) -- (0, 1.5);
                \draw [->] (-0.5, 0) -- (1.5, 0);
                \draw (0, 0) -- (1, 0) -- (1, 1) -- cycle;
                \draw [dashed] (0, 1) -- (1, 1);
                \fill [lightgray] (0, 0) -- (1, 0) -- (1, 1);
                \node [below] at (1.5, 0) {$x$};
                \node [left] at (0, 1.5) {$y$};
                \node [above right] at (0.45, 0.15) {$D$};
                \node [right] at (0, 0.75) {$y = x$};
                \node [below] at (1, 0) {$1$};
                \node [left] at (0, 1) {$1$};
                \node [below left] at (0, 0) {$0$};  
            \end{tikzpicture}
        \end{tabular}
    \end{center}
    Знайдемо маргінальні щільності розподілу:
    \begin{gather*}
        f_{\xi_1}(x) = \begin{cases}
            \int\limits_0^x 2 dy = 2x, & x \in \left[ 0; 1\right] \\
            0, & x \notin \left[ 0; 1\right]
        \end{cases}, 
        \;
        f_{\xi_2}(y) = \begin{cases}
            \int\limits_y^1 2 dx = 2(1-y), & y \in \left[ 0; 1\right] \\
            0, & y \notin \left[ 0; 1\right]
        \end{cases}
    \end{gather*}
    Запишемо умовні щільності за формулами \eqref{eq:cond_dens}: 
    \begin{gather*}
        f_{\xi_1}(x/y) = \frac{f_{\vec{\xi}}(x, y)}{f_{\xi_2}(y)} = \begin{cases}
            \frac{1}{1-y}, & x \in \left[ y; 1\right], \\
            & y \in \left[ 0; 1\right)\\
            0, & \text{інакше}
        \end{cases}\;
        f_{\xi_2}(y/x) = \frac{f_{\vec{\xi}}(x, y)}{f_{\xi_1}(x)} = \begin{cases}
            \frac{1}{x}, & y \in \left[ 0; x\right], \\
            & x \in \left( 0; 1\right] \\
            0, & \text{інакше}
        \end{cases}
    \end{gather*}
    Отже, умовні розподіли обох координат є рівномірними з параметрами $\left< y, 1\right>$ та $\left<0, x\right>$ відповідно,
    тому 
    $\E(\xi_1 / \xi_2 = y) = \frac{1+y}{2}$, $y \in \left[ 0; 1\right)$ та
    $\E(\xi_2 / \xi_1 = x) = \frac{x}{2}$, $x \in \left( 0; 1\right]$.
\end{example}
\subsection{Формули повного математичного сподівання та дисперсії}
\begin{proposition*}[формула повного математичного сподівання]
    Оскільки умовні математичні сподівання є випадковими величинами, то можна знайти їх власне математичне сподівання,
    причому 
    \begin{equation}
        \E(\E(\xi_1 / \xi_2)) = \E\xi_1, \E(\E(\xi_2 / \xi_1)) = \E\xi_2
    \end{equation}
\end{proposition*}
\begin{proof} Достатньо довести лише одну з формул.
    Дискретний випадок:
    $$\E(\E(\xi_1 / \xi_2)) = \sum\limits_{j = 1}^m 
    \left(
        \sum\limits_{i=1}^n x_i \P\{\xi_1 = x_i / \xi_2 = y_j\}
    \right) \P\{\xi_2 = y_j\} = 
    \sum\limits_{i=1}^n x_i \P\{\xi_1 = x_i\} = \E\xi_1$$
    Неперервний випадок:
    \begin{gather*}
        \E(\E(\xi_1 / \xi_2)) = \int\limits_{-\infty}^{+\infty} 
        \left(
            \int\limits_{-\infty}^{+\infty} x f_{\xi_1}(x / y) dx
        \right) f_{\xi_2}(y) dy = 
        \int\limits_{-\infty}^{+\infty} x \left(\int\limits_{-\infty}^{+\infty} 
        \underbrace{f_{\xi_1}(x / y) f_{\xi_2}(y)}_{f_{\vec{\xi}}(x, y)}
        dy\right) dx = \\
        = \int\limits_{-\infty}^{+\infty} x f_{\xi_1}(x) dx = \E\xi_1
    \end{gather*}
\end{proof}

%\begin{remark}
%    В англомовній літературі також використовується назва «Adam's law».
%\end{remark}
\begin{example}
    Застосуємо формулу повного математичного сподівання до попередніх прикладів.

    \noindent Приклади \ref{ex:discr_cond_distr} та \ref{ex:discr_cond_distr_2}:
        $\E\xi_1 = -0.3 + 0.2 = -0.1$, $\E(\E(\xi_1 / \xi_2)) = -\frac{1}{6}
        \cdot0.6 = -0.1$.
        $\E\xi_2 = 0.6$, $\E(\E(\xi_2 / \xi_1)) = 0.3\cdot\frac{2}{3} + 
        \frac{1}{2}\cdot\frac{3}{5} + \frac{1}{2} \cdot\frac{2}{10} = 
        0.6$.

    \noindent Приклад \ref{ex:cont_cond_prob}:
        $\E\xi_1 = \int\limits_0^1 2x^2 dx = \frac{2}{3}$, $\E(\E(\xi_1 / \xi_2)) = \int\limits_0^1 \frac{1+y}{2} \cdot 2(1-y) dy = \int\limits_0^1 (1-y^2) dy = \frac{2}{3}$.
        $\E\xi_2 = \int\limits_0^1 2y(1-y) dy = \frac{1}{3}$, $\E(\E(\xi_2 / \xi_1)) = \int\limits_0^1 \frac{x}{2} \cdot 2x dx = \int\limits_0^1 x^2 dx = \frac{1}{3}$. 
\end{example}


Крім умовного математичного сподівання розглядають
\emph{умовну дисперсію} --- міру розсіювання однієї випадкової величини 
за умови того, що інша набула певного значення.

\begin{definition}
    \emph{Умовною дисперсією} дискретної випадкової величини $\xi_1$ 
    є дисперсія цієї випадкової величини за умови, що 
    $\xi_2$ набула певного значення.
    $$\D(\xi_1 / \xi_2 = y_j) = 
    \E((\xi_1 - \E(\xi_1 / \xi_2 = y_j))^2/\xi_2 = y_j)$$
\end{definition}

\begin{proposition*}[формула повної дисперсії]
    Нехай $\D\xi_1 < +\infty$. Тоді
    \begin{equation}
        \D\xi_1 = \E(\D(\xi_1/\xi_2)) + \D(\E(\xi_1 / \xi_2))
    \end{equation}
\end{proposition*}

\begin{proof}
    $\D\xi_1 = \E\xi_1^2 - (\E \xi_1)^2 = 
    \E(\E(\xi_1^2 / \xi_2)) - (\E(\E(\xi_1/\xi_2)))^2 = 
    \E(\D(\xi_1/\xi_2) + (\E(\xi_1/\xi_2))^2) - (\E(\E(\xi_1/\xi_2)))^2 = 
    \E(\D(\xi_1 / \xi_2)) + (\E(\E(\xi_1/\xi_2)^2) - 
    (\E(\E(\xi_1 / \xi_2)))^2) = \E(\D(\xi_1/\xi_2)) + \D(\E(\xi_1/\xi_2))$.
\end{proof}


%\begin{remark}
%    В англомовній літературі також використовується назва \Eve's law».
%\end{remark}

\subsection{Випадок незалежних координат}
Необхідною і достатньою умовою незалежності координат, як вже було розглянуто, є 
$F_{\vec{\xi}}(x, y) = F_{\xi_1}(x)\cdot F_{\xi_2}(y)$. З цього випливає,
що у випадку незалежності координат всі умовні характеристики дорівнюють безумовним: в першу чергу, 
$F_{\xi_1}(x/y) = F_{\xi_1}(x)$, $F_{\xi_2}(y/x) = F_{\xi_2}(y)$.
Зокрема, для ДВВ:
\begin{enumerate}
    \item $\P\{\xi_1 = x_i, \xi_2 = y_j\} = \P\{\xi_1 = x_i\}\cdot \P\{\xi_2 = y_j\}$.
    \item $\P\{\xi_1 = x_i / \xi_2 = y_j\} = \P\{\xi_1 = x_i\}$, $\P\{\xi_2 = y_j / \xi_1 = x_i\} = \P\{\xi_2 = y_j\}$.
    \item $\E(\xi_1 / \xi_2 = y_j) = \E\xi_1$, $\E(\xi_2 / \xi_1 = x_i) = \E\xi_2$.
\end{enumerate}
Для НВВ:
\begin{enumerate}
    \item $f_{\vec{\xi}}(x, y) = f_{\xi_1}(x)\cdot f_{\xi_2}(y)$.
    \item $f_{\xi_1}(x/y) = f_{\xi_1}(x)$, $f_{\xi_2}(y/x) = f_{\xi_2}(y)$.
    \item $\E(\xi_1 / \xi_2 = y) = \E\xi_1$, $\E(\xi_2 / \xi_1 = x) = \E\xi_2$.
\end{enumerate}
    \chapter{Характеристичні функції. Гауссівські випадкові вектори}
        % !TEX root = ../main.tex
\section{Характеристичні функції випадкових величин}
\subsection{Поняття характеристичної функції}
\begin{definition}
    \emph{Характеристичною функцією} випадкової величини $\xi$
    називається комплекснозначна функція $\chi_\xi(t) = Ee^{it\xi}, t\in \mathbb{R}$.
    \begin{equation}\label{eq:char_func}
        \chi_\xi(t) = \int\limits_{-\infty}^{+\infty} e^{itx} dF_\xi(x) = \begin{cases}
            \sum\limits_{k=1}^{n(\infty)} e^{itx_k} P\left\{\xi = x_k\right\}, & \xi \text{ --- ДВВ} \\
            \int\limits_{-\infty}^{+\infty} e^{itx} f_\xi(x)dx, & \xi \text{ --- НВВ}
        \end{cases}
    \end{equation}
\end{definition}
Інтеграл $\int\limits_{-\infty}^{+\infty} e^{itx} f_\xi(x)dx$ у курсі
гармонічного аналізу називається \emph{перетворенням Фур'є} функції $f_\xi(x)$.
Отже, у випадку НВВ $\xi$ $\chi_\xi(t)$ --- перетворення Фур'є щільності.
Також з курсу гармонічного аналізу відомо, що за допомогою \emph{оберненого перетворення Фур'є}
можна відновити $f_\xi(x)$ за $\chi_\xi(t)$: $f_\xi(x) = \frac{1}{2\pi}\int\limits_{-\infty}^{+\infty} e^{-itx} \chi_\xi(t)dt$.

\subsection{Властивості характеристичної функції}
\begin{enumerate}
    \item $\chi_\xi(0) = 1$, оскільки $\int\limits_{-\infty}^{+\infty} f_\xi(x)dx = 1$.
    
    $\left|\chi_\xi(t)\right| \leq 1$, оскільки $\left|\int\limits_{-\infty}^{+\infty} e^{itx} dF_\xi(x)\right| \leq \int\limits_{-\infty}^{+\infty} \left|e^{itx}\right| dF_\xi(x) = \int\limits_{-\infty}^{+\infty} dF_\xi(x) = 1$.
    \item Характеристична функція є рівномірно неперервною.
    \begin{proof}
        $\left| \chi_\xi(t+h) - \chi_\xi(t) \right| = 
        \left| \int\limits_{-\infty}^{+\infty} \left( e^{i(t+h)x} - e^{itx}\right) dF_\xi(x) \right| \leq \\
        \leq \int\limits_{-\infty}^{+\infty} \left|e^{i(t+h)x} - e^{itx}\right| dF_\xi(x) = 
        \int\limits_{-\infty}^{+\infty} \left|e^{itx}\right|\left| e^{ihx} - 1\right| dF_\xi(x) = 
        \int\limits_{-\infty}^{+\infty} \left|e^{ihx} - 1\right| dF_\xi(x)$.

        $e^{ihx} - 1 = (\cos(hx) - 1) + i\sin(hx) = -2\sin^2(\frac{hx}{2}) + 2i\sin(\frac{hx}{2})\cos(\frac{hx}{2})$, 
        
        $\left| e^{ihx} - 1\right| = 
        \sqrt{4\sin^4(\frac{hx}{2}) + 4\sin^2(\frac{hx}{2})\cos^2(\frac{hx}{2})} = 2\left| \sin(\frac{hx}{2})\right|$.

        $\int\limits_{-\infty}^{+\infty} \left|e^{ihx} - 1\right| dF_\xi(x) = 2 \int\limits_{-\infty}^{+\infty} \left| \sin(\frac{hx}{2})\right| dF_\xi(x) =
        \\\overset{A>0}{=} 2\cdot \left( \int\limits_{-\infty}^{-A}\left| \sin(\frac{hx}{2})\right| dF_\xi(x) + 
        \int\limits_{-A}^{A}\left| \sin(\frac{hx}{2})\right| dF_\xi(x) + 
        \int\limits_{A}^{+\infty}\left| \sin(\frac{hx}{2})\right| dF_\xi(x)\right)
        $.

        Інтеграл $\int\limits_{-A}^{A}\left| \sin(\frac{hx}{2})\right| dF_\xi(x)$
        можна зробити як завгодно малим за рахунок вибору $h$,
        а інші два --- за рахунок вибору $A$.
    \end{proof}
    \item $\eta = a\xi + b$ --- афінне перетворення $\xi$, $a, b \in \mathbb{R}$.

    $\chi_\eta(t) = Ee^{it\eta} = Ee^{it(a\xi + b)} = e^{itb}\cdot Ee^{ita\xi} = e^{itb}\cdot \chi_\xi(at)$.
    \item Якщо $\xi_1, ..., \xi_n$ --- незалежні у сукупності, то
    $\chi_{\sum_{k=1}^n {\xi_k}} (t) = \prod\limits_{k=1}^n \chi_{\xi_k}(t)$.
    \begin{proof}
        $\chi_{\sum_{k=1}^n \xi_k}(t) = E\left( e^{it\sum_{k=1}^n \xi_k}\right) =
        E\left( \prod\limits_{k=1}^n e^{it\xi_k}\right) = \\
        = \left[e^{it\xi_k} \text{ --- теж незалежні у сукупності}\right] = 
        \prod\limits_{k=1}^n Ee^{it\xi_k} = \prod\limits_{k=1}^n \chi_{\xi_k}(t)$.
    \end{proof}
    \item За допомогою характеристичної функції можна знайти початкові моменти будь-якого порядку: $\chi_\xi^{(k)}(t) = \int\limits_{-\infty}^{+\infty} (ix)^k e^{itx} dF_\xi(x)$,
    причому $\chi_\xi^{(k)}(0) = \int\limits_{-\infty}^{+\infty} (ix)^k dF_\xi(x) = i^k \cdot E\xi^k$.
    
    Отже, $E\xi = -i \chi_\xi'(0)$, $D\xi = E\xi^2 - \left( E\xi\right)^2 = - \chi_\xi''(0) + \left( \chi_\xi'(0) \right)^2$.
    \item Зв'язок характеристичної функції та генератриси ДВВ.

    $\varphi_\xi(z) = \sum\limits_{k=0}^{\infty} P\left\{\xi = k\right\} z^k$,
    $\chi_\xi(t) = \sum\limits_{k=0}^{\infty} e^{itk} P\left\{\xi = k\right\}$. Отже, $\chi_\xi(t) = \varphi_\xi(e^{it})$.
    \item $\overline{\chi_\xi(t)} = E e^{-it\xi} = \chi_\xi(-t) = \chi_{-\xi} (t)$. 
    \item Для того, щоб характеристична функція була дійснозначною,
    необхідно і достатньо, щоб розподіл ВВ був симетричним відносно 0.
    \begin{proof}
        Розглянемо випадок НВВ. Нехай розподіл є симетричним відносно 0, тоді $f_\xi(x)$ --- парна.
        $\chi_\xi(t) = \int\limits_{-\infty}^{+\infty} e^{itx} f_\xi(x)dx = 
        \int\limits_{-\infty}^{+\infty} \cos(tx) f_\xi(x)dx +
        i\int\limits_{-\infty}^{+\infty} \sin(tx) f_\xi(x)dx = 
        \int\limits_{-\infty}^{+\infty} \cos(tx) f_\xi(x)dx = {Re}\chi_\xi(t)$, 
        інтеграл з синусом рівний 0, бо інтегрується непарна функція по симетричному проміжку.

        Нехай $\chi_\xi(t)$ --- дійснозначна, тоді $\chi_\xi(t) = \overline{\chi_\xi(t)} = \chi_\xi(-t)$.
        Отже, $\chi_\xi(t)$ --- парна, тоді з оберненого перетворення Фур'є $f_\xi(x)$ --- теж парна.
    \end{proof}
\end{enumerate}

\subsection{Необхідні умови того, що функція є характеристичною}
Нехай $\chi(t)$ --- деяка комплекснозначна функція дійсного аргументу.
Якщо вона є характеристичною функцією деякої ВВ, то для неї має виконуватися:
\begin{enumerate}
    \item $\chi(0) = 1$, $\left| \chi(t)\right| \leq 1$.
    \item $\chi(t)$ --- рівномірно неперервна.
    \item $\overline{\chi(t)} = \chi(-t)$.
\end{enumerate}
\begin{example}
    \begin{enumerate}
        \item $\chi(t) = \cos(t)$ --- може бути характеристичною.
        Оскільки $\cos(t) = \frac{1}{2}e^{it} + \frac{1}{2}e^{-it}$, то відповідна ВВ --- дискретна: 
        \begin{tabular}{|c|c|c|}
            \hline
            $\xi$ & $-1$ & $1$ \\
            \hline
            $p$ & $1/2$ & $1/2$ \\
            \hline
        \end{tabular}
        \item $\chi(t) = \sin(t)$ --- не може бути характеристичною, бо $\chi(0) = 0 \neq 1$.
        \item $\chi(t) = \cos^2(t)$ --- може бути характеристичною.
        $\cos^2(t) = \frac{1}{2} + \frac{1}{2}\cos(2t) = \frac{1}{2} + \frac{1}{4}e^{2it} + \frac{1}{4}e^{-2it}$, тому
        відповідна ВВ --- дискретна:
        \begin{tabular}{|c|c|c|c|}
            \hline
            $\xi$ & $-2$ & $0$ & $2$ \\
            \hline
            $p$ & $1/4$ & $1/2$ & $1/4$ \\
            \hline
        \end{tabular}
        \item $\chi(t) = \frac{\alpha^2}{\alpha^2 + t^2}$ --- може бути характеристичною. 
        Знайдемо щільність розподілу відповідної НВВ за формулою $f_\xi(x) = \frac{1}{2\pi}\int\limits_{-\infty}^{+\infty} \frac{\alpha^2 \cdot e^{-itx}}{\alpha^2 + t^2} dt$.

        При $x<0$: $\int\limits_{-\infty}^{+\infty} \frac{e^{-itx}}{\alpha^2 + t^2} dt = 2\pi i \cdot \underset{z=\alpha i}{\mathrm{Res}} \frac{e^{-izx}}{\alpha^2 + z^2} = 2\pi i \cdot \frac{e^{-i\cdot \alpha i\cdot x}}{2\alpha i} = \pi \cdot \frac{e^{\alpha x}}{\alpha}$.

        При $x>0$ аналогічно отримуємо $\pi \cdot \frac{e^{-\alpha x}}{\alpha}$.
        Остаточно $f_\xi(x) = \frac{\alpha}{2}e^{-\alpha |x|}$ --- це щільність розподілу Лапласа.
    \end{enumerate}
\end{example}

\begin{exercise}
    Нехай $\chi(t)$ --- характеристична функція деякого розподілу. Чи можуть бути характеристичними функції $\overline{\chi}$, $\chi^2$, $|\chi|^2$, ${Re}\chi$, $|\chi|$, ${Im}\chi$? 
\end{exercise}

Необхідні та достатні умови того, що функція є характеристичною, дає
\emph{теорема Бохнера-Хінчина}: $\chi(t)$ має бути рівномірно неперервною, $\chi(0) = 1$ та 
невід'ємно визначеною: $$
    \forall \; t_1, t_2, ..., t_n \in \mathbb{R}, c_1, c_2, ..., c_n \in \mathbb{C}: 
    \sum\limits_{k=1}^n {\sum\limits_{m=1}^n \chi(t_k - t_m) c_k \overline{c_m}} \geq 0
$$

\subsection{Характеристичні функції деяких розподілів}
\begin{enumerate}
    \item $\xi \sim \mathrm{Bin}(n,p)$ --- біноміальний розподіл.
    $\varphi_\xi(z) = (pz+q)^n \Rightarrow \chi_\xi(t) = \left( p e^{it} + q\right)^n$.
    \item $\xi \sim \mathrm{Geom}(p)$ --- геометричний розподіл.
    $\varphi_\xi(z) = \frac{pz}{1-qz} \Rightarrow \chi_\xi(t) = \frac{p e^{it}}{1-q e^{it}}$.
    \item $\xi \sim \mathrm{Poiss}(a)$ --- розподіл Пуассона.
    $\varphi_\xi(z) = e^{a(z-1)} \Rightarrow \chi_\xi(t) = e^{a(e^{it}-1)}$.
    \item $\xi \sim \mathrm{U}\left< a; b\right>$ --- рівномірний розподіл.

    $\chi_\xi(t) = \frac{1}{b-a} \int\limits_a^b e^{itx} dx = \frac{e^{itb} - e^{ita}}{it(b-a)}$.
    Зокрема, при $\xi \sim {U}\left< -a; a\right>$ $\chi_\xi(t) = \begin{cases}
        \frac{\sin(at)}{at}, & t \neq 0 \\
        1, & t = 0
    \end{cases}$.
    \item $\xi \sim \mathrm{Exp}(\lambda)$ --- експоненційний розподіл.
    
    $\chi_\xi(t) = \lambda \int\limits_0^{+\infty} e^{itx} e^{-\lambda x} dx = \lambda \int\limits_0^{+\infty} e^{-(\lambda-it)x} dx = \frac{\lambda}{\lambda - it}$.
    \item $\xi \sim \mathrm{N}(a, \sigma^2)$ --- нормальний розподіл. Розглянемо стандартний розподіл $\eta \sim \mathrm{N}(0, 1)$,
    тоді $\xi = a + \sigma\cdot\eta$.

    $\chi_\eta(t) = \frac{1}{\sqrt{2\pi}} \int\limits_{-\infty}^{+\infty} e^{itx} e^{-x^2/2} dx = 
    \frac{1}{\sqrt{2\pi}} \int\limits_{-\infty}^{+\infty} e^{-x^2/2 + itx} dx = 
    \frac{1}{\sqrt{2\pi}} \int\limits_{-\infty}^{+\infty} e^{-\frac{1}{2}(x^2 - 2itx + i^2t^2) + \frac{1}{2}i^2t^2} dx = 
    \frac{1}{\sqrt{2\pi}} e^{-\frac{t^2}{2}} \int\limits_{-\infty}^{+\infty} e^{-\frac{(x-it)^2}{2}} dx = 
    \left[ x-it = u, dx = du\right] = e^{-\frac{t^2}{2}}$.
    Отже, $\chi_\xi(t) = e^{iat - \frac{\sigma^2 t^2}{2}}$.
    \item $\xi \sim \Gamma(\alpha, \beta)$ --- гамма-розподіл.
    
    $\chi_\xi(t) = \frac{\beta^\alpha}{\Gamma(\alpha)} \int\limits_0^{+\infty} e^{itx} x^{\alpha-1} e^{-\beta x} dx =
    \frac{\beta^\alpha}{\Gamma(\alpha)} \int\limits_0^{+\infty} x^{\alpha-1} e^{-(\beta-it) x} dx = \left[ (\beta-it) x = u, dx = \frac{du}{\beta-it}\right] = 
    \frac{\beta^\alpha}{\Gamma(\alpha)} \cdot \frac{1}{(\beta-it)^\alpha} \int\limits_0^{+\infty} u^{\alpha-1} e^{-u} du = 
    \frac{\beta^\alpha}{\Gamma(\alpha)} \cdot \frac{1}{(\beta-it)^\alpha} \cdot \Gamma(\alpha) = (1-\frac{it}{\beta})^{-\alpha}$.
\end{enumerate}

\begin{definition}
    \emph{Композицією законів розподілу} називається складання закону розподілу суми незалежних випадкових величин.
    
    Закон розподілу називається \emph{стійким по відношенню до операції додавання},
    якщо закон розподілу суми незалежних випадкових величин, розподілених за цим законом (в загальному випадку з різними параметрами),
    є таким самим.
\end{definition}
\begin{example}
    \begin{enumerate}
        \item Перевірити стійкість нормального розподілу. Нехай $\xi_1 \sim \mathrm{N}(a_1, \sigma_1), ..., \xi_n \sim \mathrm{N}(a_n, \sigma_n)$ ---
        незалежні у сукупності. 
    
        Знайдемо характеристичну функцію суми: $\chi_{\xi_1 + ... + \xi_n}(t) = \prod\limits_{k=1}^n \chi_{\xi_k}(t)=
        \prod\limits_{k=1}^n e^{ia_k t - \frac{\sigma_k^2 t^2}{2}} = 
        \exp\left\{i\cdot \left( \sum\limits_{k=1}^n a_k\right)\cdot t - \frac{t^2}{2}\cdot\left( \sum\limits_{k=1}^n \sigma_k^2\right)\right\}$ ---
        це характеристична функція нормального розподілу. Отже, $\xi_1 + ... + \xi_n \sim \mathrm{N}\left(\sum\limits_{k=1}^n a_k,  \sqrt{\sum\limits_{k=1}^n \sigma_k^2}\right)$.
        \item Перевірити стійкість біноміального розподілу. Нехай $\xi_1 \sim \mathrm{Bin}(n_1, p_1)$, $\xi_2 \sim \mathrm{Bin}(n_2, p_2)$ --- незалежні.
        $\chi_{\xi_1 + \xi_2} (t) = \left( p_1 e^{it} + q_1\right)^{n_1} \cdot \left( p_2 e^{it} + q_2\right)^{n_2} \overset{?}{=} \left(P e^{it} + Q\right)^N$.

        Взагалі кажучи, біноміальний розподіл не є стійким по відношенню до  операції додавання.
        Його називають \emph{умовно стійким} при $p_1 = p_2 = p$, тоді
        $\chi_{\xi_1 + \xi_2} (t) = \left( p e^{it} + q\right)^{n_1 + n_2}$ і $\xi_1 + \xi_2 \sim \mathrm{Bin}(n_1+n_2, p)$.
    \end{enumerate}
\end{example}

\begin{exercise}
    Перевірити стійкість гамма-розподілу $\Gamma(\alpha, \beta)$.
\end{exercise}

\section{Характеристичні функції випадкових векторів}
\subsection{Поняття характеристичної функції випадкового вектору}
\begin{definition}
    \emph{Характеристичною функцією} випадкового вектору $\vec{\xi} = \left( \xi_1, \xi_2, ..., \xi_n\right)^T$ 
    називається комплекснозначна функція $n$ дійсних аргументів $\chi_{\vec{\xi}}(\vec{t}) = Ee^{i(\vec{\xi}, \vec{t})}$, $\vec{t} \in \mathbb{R}^n$.
    \begin{equation}\label{eq:char_func_vect}
        \chi_{\vec{\xi}}(\vec{t}) = \int\limits_{\mathbb{R}^n} e^{i(\vec{x}, \vec{t})} dF_\xi(x) = \begin{cases}
            \sum\limits_{k_1,k_2, ..., k_n} e^{i \sum\limits_{j=1}^n x_{k_j}t_j} P\left\{\xi_1 = x_{k_1}, ..., \xi_n = x_{k_n}\right\}, & \vec{\xi} \text{ --- ДВВ} \\
            \int\limits_{\mathbb{R}^n} e^{i(\vec{x}, \vec{t})} f_{\vec{\xi}}(x_1, x_2, ..., x_n)dx_1 dx_2 ... dx_n, & \vec{\xi} \text{ --- НВВ}
        \end{cases}
    \end{equation}
\end{definition}

\begin{example} Обчислення характеристичної функції для дискретного та неперервного випадкових векторів
    \begin{enumerate}
        \item \begin{tabular}{c c}
            \begin{tabular}{|c|c|c|}
                \hline
                \diagbox{$\xi_2$}{$\xi_1$} & $1$ & $2$ \\
                \hline
                $0$ & $0.1$ & $0.3$ \\
                \hline
                $1$ & $0.2$ & $0.4$ \\
                \hline
            \end{tabular} 
        \end{tabular}

            $\chi_{\vec{\xi}}(t_1, t_2) = 0.1 e^{i(1\cdot t_1 + 0\cdot t_2)} +
            0.3 e^{i(2\cdot t_1 + 0\cdot t_2)} + 0.2 e^{i(1\cdot t_1 + 1\cdot t_2)} +
            0.4 e^{i(2\cdot t_1 + 1\cdot t_2)}$.
        \item $\vec{\xi} \sim {U}(D)$, $D = \left\{(x;y) \in \mathbb{R}^2:|x|\leq 1, |y|\leq 1\right\}$.
        
        $\chi_{\vec{\xi}}(t_1, t_2) = \frac{1}{4} \iint\limits_{D} e^{i(t_1x + t_2y)} dx dy = \frac{1}{4}\cdot\left( \int\limits_{-1}^1 e^{it_1x} dx\right) 
        \cdot\left( \int\limits_{-1}^1 e^{it_2y} dy\right) = \frac{e^{it_1}-e^{-it_1}}{2it_1}\cdot\frac{e^{it_2}-e^{-it_2}}{2it_2} = \begin{cases}
            \frac{\sin(t_1)}{t_1} \cdot \frac{\sin(t_2)}{t_2}, & t_1 \neq 0, t_2 \neq 0 \\
            \frac{\sin(t_1)}{t_1}, & t_2 = 0 \\
            \frac{\sin(t_2)}{t_2}, & t_1 = 0 \\
            1, & t_1 = t_2 = 0
        \end{cases}$
    \end{enumerate}
\end{example}

\subsection{Властивості характеристичних функцій векторів}
\begin{enumerate}
    \item Характеристична функція існує для будь-якого розподілу випадкового вектора, бо \\
    $\left| \chi_{\vec{\xi}}(\vec{t})\right| \leq 1$, $\chi_{\vec{\xi}}(\vec{0}) = 1$.
    \item Характеристична функція рівномірно неперервна по кожному з аргументів.
    \item За характеристичною функцією вектора можна знайти 
    характеристичну функцію будь-якої підсистеми цього вектора.
    $\vec{\xi} = (\xi_1, \xi_2, ..., \xi_n)$, $\vec{\eta} = (\xi_1, \xi_2, ..., \xi_k)$, $k<n$.

    $\chi_{\vec{\eta}}(t_1, t_2, ..., t_k) = \chi_{\vec{\xi}}(t_1, t_2, ..., t_k, 0, ..., 0)$.
    \item Якщо координати $\vec{\xi}$ незалежні у сукупності, то $\chi_{\vec{\xi}}(\vec{t}) = \prod\limits_{k=1}^n \chi_{\xi_k}(t_k)$.
    \begin{proof}
        $\chi_{\vec{\xi}}(\vec{t}) = E\left( e^{i(\xi_1 t_1 + ... + \xi_n t_n)}\right) = 
        E\left( e^{i\xi_1 t_1} ... e^{i\xi_n t_n}\right) = \left[ \text{незалежні у сукупності}\right] =
        E\left( e^{i\xi_1 t_1}\right) \cdot ... \cdot E\left( e^{i\xi_n t_n}\right) = \prod\limits_{k=1}^n \chi_{\xi_k}(t_k)$.
    \end{proof}
    \item $\vec{\eta} = A\vec{\xi} + \vec{b}$ --- афінне перетворення $\vec{\xi}$.

    $\chi_{\vec{\eta}}(\vec{t}) = E e^{i(\vec{\eta}, \vec{t})} = E e^{i(A\vec{\xi} + \vec{b}, \vec{t})}=
    E e^{i(\vec{b}, \vec{t})} \cdot E e^{(A\vec{\xi}, \vec{t})} = e^{i(\vec{b}, \vec{t})} \cdot E e^{(\vec{\xi}, A^{*}\vec{t})} = e^{i(\vec{b}, \vec{t})} \cdot \chi_{\vec{\xi}}(A^{*}\vec{t})$.
    \item Характеристична функція вектора дозволяє знайти початкові моменти будь-якого порядку:
    \begin{gather*}
        E\xi_1^{k_1} \xi_2^{k_2} ... \xi_n^{k_n} = \frac{1}{i^{k_1 + k_2 + ... + k_n}} \cdot \frac{\partial^{k_1 + k_2 + ... + k_n}\chi_{\vec{\xi}}(\vec{t})}{\partial t_1^{k_1} \partial t_2^{k_2} ... \partial t_n^{k_n}} \Biggr \vert_{\vec{t} = \vec{0}}, 
        \; E\xi_j = \frac{1}{i} \cdot\frac{\partial \chi_{\vec{\xi}}(\vec{t})}{\partial t_j} \Biggr \vert_{\vec{t} = \vec{0}} \\
        {cov}(\xi_k, \xi_j) = E\xi_k\xi_j - E\xi_k\cdot E\xi_j = \frac{1}{i^2} \frac{\partial^2 \chi_{\vec{\xi}}(\vec{t})}{\partial t_k \partial t_j} \Biggr \vert_{\vec{t} = \vec{0}} - 
        \left( \frac{1}{i} \cdot\frac{\partial \chi_{\vec{\xi}}(\vec{t})}{\partial t_k} \Biggr \vert_{\vec{t} = \vec{0}}\right) \cdot \left( \frac{1}{i} \cdot\frac{\partial \chi_{\vec{\xi}}(\vec{t})}{\partial t_j} \Biggr \vert_{\vec{t} = \vec{0}}\right) = \\
        = -\frac{\partial^2 \chi_{\vec{\xi}}(\vec{t})}{\partial t_k \partial t_j} \Biggr \vert_{\vec{t} = \vec{0}} + 
        \left(\frac{\partial \chi_{\vec{\xi}}(\vec{t})}{\partial t_k} \Biggr \vert_{\vec{t} = \vec{0}}\right) \cdot \left(\frac{\partial \chi_{\vec{\xi}}(\vec{t})}{\partial t_j} \Biggr \vert_{\vec{\vec{t}} = \vec{0}}\right)\end{gather*}
\end{enumerate}

Необхідні та достатні умови того, що комплекснозначна функція $n$ дійсних змінних є характеристичною, дає
\emph{теорема Бохнера-Хінчина}: $\chi(\vec{t})$ має бути рівномірно неперервною по кожній змінній, $\chi(\vec{0}) = 1$ та 
невід'ємно визначеною: $$
    \forall \; \vec{t^1}, \vec{t^2}, ..., \vec{t^n} \in \mathbb{R}^n, c_1, c_2, ..., c_n \in \mathbb{C}: 
    \sum\limits_{k=1}^n {\sum\limits_{m=1}^n \chi(\vec{t^k} - \vec{t^m}) c_k \overline{c_m}} \geq 0$$ 
        % !TEX root = ../main.tex

\section{Багатовимірний нормальний розподіл}
\subsection{Виведення характеристичної функції та щільності}

Розглянемо випадковий вектор $\vec{\xi} = (\xi_1, \xi_2, ..., \xi_n)^T$, $\xi_k \sim \mathrm{N}(a_k^o, \sigma_k)$ для $k=1,...,n$, координати незалежні у сукупності.
За властивостями характеристичної функції та щільності $\chi_{\vec{\xi}}(\vec{t}) = \prod\limits_{k=1}^n \chi_{\xi_k}(t_k)$,
$f_{\vec{\xi}}(\vec{x}) = \prod\limits_{k=1}^n f_{\xi_k}(x_k)$.

\begin{gather*}
    \chi_{\vec{\xi}}(\vec{t}) = \prod\limits_{k=1}^n e^{i a_k^o t_k - \frac{\sigma_k^2 t_k^2}{2}} = \exp\left\{i(\vec{a^o}, \vec{t}) - \frac{1}{2}\sum\limits_{k=1}^n \sigma_k^2 t_k^2\right\}
    \\
    f_{\vec{\xi}}(\vec{x}) = \prod\limits_{k=1}^n \frac{1}{\sqrt{2\pi}\sigma_k} e^{-\frac{(x_k-a_k^o)^2}{2\sigma_k^2}} = \frac{1}{(2\pi)^{n/2}\prod_{k=1}^n \sigma_k} \exp \left\{ -\frac{1}{2} \sum_{k=1}^n \frac{(x_k-a_k^o)^2}{\sigma_k^2}\right\}
\end{gather*}

Введемо матрицю $D = \begin{pmatrix}
    \sigma_1^2 & 0 & \cdots & 0 \\
    0 & \sigma_2^2 & \cdots & 0 \\
    \vdots & \vdots & \ddots & \vdots \\
    0 & 0 & \cdots & \sigma_n^2
\end{pmatrix}$, $D^{-1} = \begin{pmatrix}
    1/\sigma_1^2 & 0 & \cdots & 0 \\
    0 & 1/\sigma_2^2 & \cdots & 0 \\
    \vdots & \vdots & \ddots & \vdots \\
    0 & 0 & \cdots & 1/\sigma_n^2
\end{pmatrix}$.

\noindent Тоді отримаємо
\begin{gather}
    \chi_{\vec{\xi}}(\vec{t}) = \exp\left\{i(\vec{a^o}, \vec{t}) - \frac{1}{2}(D\vec{t}, \vec{t})\right\}
    \\
    f_{\vec{\xi}}(\vec{x}) = \frac{1}{(2\pi)^{n/2} \sqrt{{\det{D}}}} \exp \left\{ -\frac{1}{2} \left( D^{-1}(\vec{x} - \vec{a^o}), (\vec{x} - \vec{a^o})\right)\right\}
\end{gather}
Причому в такому випадку $D$ --- кореляційна матриця.

Тепер розглянемо вектор $\vec{\eta} = A\vec{\xi}$, де $A$ --- невироджена матриця. За властивістю 
$\chi_{\vec{\eta}}(\vec{t}) = \chi_{\vec{\xi}}(A^{*}\vec{t})$. Маємо
\begin{equation}
    \begin{gathered}
        \chi_{\vec{\eta}}(\vec{t}) = \exp\left\{i(\vec{a^o}, A^{*}\vec{t}) - \frac{1}{2}(DA^{*}\vec{t}, A^{*}\vec{t})\right\} =
    \exp\left\{i(A\vec{a^o}, \vec{t}) - \frac{1}{2}(ADA^{*}\vec{t}, \vec{t})\right\} = \\
    = \exp\left\{i(\vec{a}, \vec{t}) - \frac{1}{2}(K\vec{t}, \vec{t})\right\}
    \end{gathered}
\end{equation}
Доведемо, що $K = ADA^{*}$ --- симетрична й додатно визначена, тому її можна вважати кореляційною матрицею.
\begin{enumerate}
    \item $K^{*} = \left( ADA^{*}\right)^{*} = \left( A^{*}\right)^{*}DA^{*} = ADA^{*} = K$.
    \item $\forall \;\vec{u} \in \mathbb{R}^n$ $\left( K \vec{u}, \vec{u}\right) = \left( ADA^{*} \vec{u}, \vec{u}\right) =
    \left( DA^{*} \vec{u}, A^{*}\vec{u}\right) = \left[ A^{*}\vec{u} = \vec{v}\right] = \left( D\vec{v}, \vec{v}\right) > 0$ для $\vec{u}\neq \vec{0}$.
\end{enumerate}

Нехай в $n$-вимірному евклідовому просторі задано вектор $\vec{a}$ та симетричну додатно визначену матрицю $K$.
Існує ортогональне перетворення $U$, яке дає можливість записати $K=UDU^{*}$, причому на діагоналі $D$ стоять строго додатні
власні числа $K$. Тоді функцію вигляду $\exp\left\{i(\vec{a}, \vec{t}) - \frac{1}{2}(K\vec{t}, \vec{t})\right\}$ 
можна розглядати як характеристичну функцію випадкового вектора $\vec{\eta} = U \vec{\xi}$, де $\vec{\xi}$ ---
вектор, координати якого незалежні у сукупності та мають нормальний розподіл.

\begin{definition}
    $n$\emph{-вимірним гауссівським вектором} $\vec{\xi}$ називається випадковий вектор,
    характеристична функція якого має вигляд $\chi_{\vec{\xi}}(\vec{t}) = \exp\left\{i(\vec{a}, \vec{t}) - \frac{1}{2}(K\vec{t}, \vec{t})\right\}$,
    де $K$ --- кореляційна матриця, а $\vec{a} = \left( E\xi_1, E\xi_2, ..., E\xi_n\right)^T$ --- центр розсіювання.
\end{definition}
\noindent \textbf{Позначення:} $\vec{\xi} \sim \mathrm{N}( \vec{a}, K)$. $\mathrm{N}( \vec{0}, {I})$ --- стандартний нормальний розподіл.
\begin{exercise}
    Перевірити, що $\vec{a}$ дійсно є центром розсіювання.
\end{exercise}
Знайдемо щільність сумісну розподілу такого вектору:
\begin{gather*}
    P\left\{\vec{\eta} \in C \subset \mathbb{R}^n\right\} = P\left\{U\vec{\xi}\in C\right\} = 
    P\left\{\vec{\xi}\in U^{-1}C\right\} = 
    \int\limits_{U^{-1}C} f_{\vec{\xi}}(\vec{x}) d\vec{x} = \\
    \int\limits_{U^{-1}C} \frac{1}{(2\pi)^{n/2} \sqrt{{\det{D}}}} \exp \left\{ -\frac{1}{2} \left( D^{-1}(\vec{x} - \vec{a^o}), (\vec{x} - \vec{a^o})\right)\right\} d\vec{x} =
    \left[ \vec{x} - \vec{a^o} = U^{-1}(\vec{y} - \vec{a}) = \right. \\ \left. = U^{*}(\vec{y} - \vec{a}), d\vec{x} = d\vec{y}\right] = 
    \frac{1}{(2\pi)^{n/2} \sqrt{{\det{D}}}} \int\limits_{C} \exp \left\{ -\frac{1}{2} \left( D^{-1}(U^{*}\vec{y} - U^{*}\vec{a}), (U^{*}\vec{y} - U^{*}\vec{a})\right)\right\} d\vec{y} = \\
    = \left[ UD^{-1}U^{*} = K^{-1}\right] =
    \frac{1}{(2\pi)^{n/2} \sqrt{{\det{K}}}} \int\limits_{C} \exp \left\{ -\frac{1}{2} \left( K^{-1}(\vec{y} - \vec{a}), (\vec{y} - \vec{a})\right)\right\} d\vec{y}
\end{gather*}
Оскільки $P\left\{\vec{\eta} \in C\right\} = \int\limits_{C} f_{\vec{\eta}}(\vec{y}) d\vec{y}$,
отримуємо \textbf{щільність розподілу} гауссівського вектора $\vec{\xi} \sim \mathrm{N}(\vec{a}, K)$:
\begin{equation}
    f_{\vec{\xi}}(\vec{x}) = \frac{1}{(2\pi)^{n/2} \sqrt{{\det{K}}}} \exp \left\{ -\frac{1}{2} \left( K^{-1}(\vec{x} - \vec{a}), (\vec{x} - \vec{a})\right)\right\}
\end{equation}

\subsection{Властивості гауссівських векторів}
\begin{enumerate}
    \item Якщо $\vec{\xi}$ --- гауссівський вектор, то всі його координати гауссівські,
    а будь-яка підсистема теж є гауссівським вектором.
    \begin{proof}
        $\chi_{\vec{\xi}}(0,0,...,t_j, ..., 0) = \chi_{\xi_j}(t_j) = 
        e^{i a_j t_j - \frac{1}{2}\sigma_j^2 t_j^2} \Rightarrow \xi_j \sim \mathrm{N}(a_j, \sigma_j^2)$.

        Аналогічно для будь-якої підсистеми $\vec{\eta} = (\xi_{k_1}, \xi_{k_2}, ..., \xi_{k_n})^T$.
    \end{proof}
    \begin{remark}
        Обернене твердження, взагалі кажучи, не є вірним.

        Розглянемо випадковий вектор із щільністю
        $$ f_{\vec{\xi}}(x, y) = \frac{1}{2\pi} \left( 
            \left( \sqrt{2} e^{-x^2/2} - e^{-x^2}\right)e^{-y^2} + 
            \left( \sqrt{2} e^{-y^2/2} - e^{-y^2}\right)e^{-x^2}\right)$$

        Очевидно, це не щільність нормального розподілу. Знайдемо щільності розподілу координат $\xi_1$ та $\xi_2$.
        \begin{gather*}
            f_{\xi_1}(x) = \int\limits_{-\infty}^{+\infty}f_{\vec{\xi}}(x, y) dy =
            \frac{\sqrt{2}}{2\pi} e^{-x^2/2}\int\limits_{-\infty}^{+\infty}e^{-y^2}dy -
            \frac{1}{2\pi}e^{-x^2}\int\limits_{-\infty}^{+\infty}e^{-y^2}dy \; +\\
            + \; \frac{\sqrt{2}}{2\pi}e^{-x^2}\int\limits_{-\infty}^{+\infty} e^{-y^2/2}dy - 
            \frac{1}{2\pi} e^{-x^2}\int\limits_{-\infty}^{+\infty}e^{-y^2}dy =
            \frac{\sqrt{2}}{2\pi} e^{-x^2/2}\cdot \sqrt{\pi} - \frac{1}{2\pi}e^{-x^2} \cdot \sqrt{\pi} \; + \\
            + \; \frac{\sqrt{2}}{2\pi}e^{-x^2} \cdot \sqrt{2\pi} - \frac{1}{2\pi}e^{-x^2}\cdot \sqrt{\pi} = 
            \frac{1}{\sqrt{2\pi}}e^{-x^2/2} - \frac{1}{2\sqrt{\pi}}e^{-x^2} + \frac{1}{\sqrt{\pi}}e^{-x^2} - \frac{1}{2\sqrt{\pi}}e^{-x^2} \; = \\ 
            \\ = \; \frac{1}{\sqrt{2\pi}}e^{-x^2/2} \Rightarrow \xi_1 \sim \mathrm{N}(0, 1).
        \end{gather*}
        Аналогічно $\xi_2 \sim \mathrm{N}(0, 1)$. Отже, координати вектора мають нормальний розподіл, а сам вектор --- ні.
    \end{remark}
    \item Для гауссівського випадкового вектора поняття незалежності та некорельованості координат є еквівалентними.
    \begin{proof}
        Було доведено, що з незалежності координат випливає ії некорельованість.

        Якщо координати некорельовані, то кореляційна матриця $K = \mathrm{diag}(\sigma_1^2, \sigma_2^2, ..., \sigma_n^2)$.
        Тоді $\left( K \vec{t}, \vec{t}\right) = \sum_{k=1}^n \sigma_k^2 t_k^2$.
        Тоді $\chi_{\vec{\xi}}(\vec{t}) = \exp\left\{i(\vec{a}, \vec{t}) - \frac{1}{2}\left( K \vec{t}, \vec{t}\right)\right\} = \prod\limits_{k=1}^n e^{i a_k t_k - \frac{1}{2}\sigma_k^2 t_k^2} = \prod\limits_{k=1}^n \chi_{\xi_k}(t_k)$
        і координати незалежні.
    \end{proof}
    \item В $n$-вимірному евклідовому просторі $\mathbb{R}^n$ завжди можна перейти до
    ортонормованого базису з власних векторів матриці $K$, в якому $K$ приймає діагональний вигляд.
    Отже, в базисі з власних векторів матриці $K$ координати відповідного гауссівського вектора стають незалежними.
    \item \emph{Афінне перетворення гауссівських векторів}.

    $\vec{\xi} \sim \mathrm{N}(\vec{a}, K)$. $A : \mathbb{R}^n \rightarrow \mathbb{R}^m$ --- лінійний оператор, заданий матрицею, $\vec{b} \in \mathbb{R}^m$, $\vec{\eta} = A\vec{\xi} + \vec{b}$.
    За властивістю характеристичної функції $\chi_{\vec{\eta}}(\vec{t}) = e^{i(\vec{b}, \vec{t})}\cdot\chi_{\vec{\xi}}(A^{*}\vec{t}) =
    \exp\left\{i(\vec{b}, \vec{t})\right\}\cdot \exp\left\{i(\vec{a}, A^{*}\vec{t}) - \frac{1}{2}\left( K A^{*}\vec{t}, A^{*}\vec{t}\right)\right\} =
    \exp\left\{i(A\vec{a} + \vec{b}, \vec{t}) - \frac{1}{2}\left( A K A^{*}\vec{t}, \vec{t}\right)\right\}$. 
    
    Отже, $\vec{\eta} \sim \mathrm{N}(A\vec{a} + \vec{b}, A K A^{*})$.
\end{enumerate}
\begin{example}
    \begin{enumerate}
        \item Задано вектор $\vec{\xi} \sim \mathrm{N}(\vec{a}, K)$. Знайти розподіл $\eta = \xi_1 + \xi_2 + ... + \xi_n$.

        $\eta = \begin{pmatrix}
            1 & 1 & \cdots & 1
        \end{pmatrix}\cdot
        \begin{pmatrix}
            \xi_1 & \xi_2 & \cdots & \xi_n
        \end{pmatrix}^T = A\vec{\xi}$.
        $E\eta = A\vec{a} = E\xi_1 + E\xi_2 + ... + E\xi_n = \sum\limits_{k=1}^n E\xi_k$.
        $D\eta = A K A^{*} = \begin{pmatrix}
            1 & 1 & \cdots & 1
        \end{pmatrix}
        \cdot K \cdot \begin{pmatrix}
            1 & 1 & \cdots & 1
        \end{pmatrix}^T$. Якщо $K = \mathrm{diag}(\sigma_1^2, \sigma_2^2, ..., \sigma_n^2)$, то $D\eta = \sum\limits_{k=1}^n D\xi_k$,
        а в загальному випадку у цій сумі ще будуть доданки вигляду $2K\xi_i\xi_j$.
        \item $\vec{\xi} = \begin{pmatrix}
            \xi_1 \\ \xi_2
        \end{pmatrix} \sim \mathrm{N}\left( \begin{pmatrix}
            0 \\ 1
        \end{pmatrix}, \begin{pmatrix}
            1 & 2 \\
            2 & 6
        \end{pmatrix}\right) = \mathrm{N}(\vec{a}, K)$. Знайти розподіл, характеристичну функцію, щільність розподілу 
        $\vec{\eta} = (3\xi_1-2\xi_2+1, \xi_1+3\xi_2)^T$ та розподіл координати $\eta_1$.

        Запишемо $\vec{\eta}$ у вигляді $\vec{\eta} = A\vec{\xi} + \vec{b}$: $\vec{\eta} = \begin{pmatrix}
            3 & -2 \\
            1 & 3 
        \end{pmatrix}\begin{pmatrix}
            \xi_1 \\ \xi_2
        \end{pmatrix} + \begin{pmatrix}
            1 \\ 0 
        \end{pmatrix}$. Тоді $\vec{\eta} \sim \mathrm{N}(\vec{d}, B)$,
        $\vec{d} = A\vec{a} + \vec{b} = \begin{pmatrix}
            3 & -2 \\
            1 & 3 \\
        \end{pmatrix}\begin{pmatrix}
            0 \\ 1
        \end{pmatrix} + \begin{pmatrix}
            1 \\ 0 
        \end{pmatrix} = \begin{pmatrix}
            -2 \\ 3 
        \end{pmatrix} + \begin{pmatrix}
            1 \\ 0
        \end{pmatrix} = \begin{pmatrix}
            -1 \\ 3
        \end{pmatrix} = \begin{pmatrix}
            E\eta_1 \\ E\eta_2
        \end{pmatrix}$, $B=A K A^{*} = \begin{pmatrix}
            3 & -2 \\
            1 & 3 \\
        \end{pmatrix}\begin{pmatrix}
            1 & 2 \\
            2 & 6
        \end{pmatrix}\begin{pmatrix}
            3 & 1 \\
            -2 & 3
        \end{pmatrix} = \begin{pmatrix}
            9 & -19 \\
            -19 & 67
        \end{pmatrix}$.

        $\chi_{\vec{\eta}}(t_1, t_2) = \exp\left\{i(\vec{d}, \vec{t}) - \frac{1}{2}\left( B\vec{t},\vec{t}\right)\right\} = 
        \exp\left\{i(-1 t_1 + 3 t_2) - \frac{1}{2}(9t_1^2 - 38 t_1 t_2 + 67 t_2^2)\right\}$.

        $B^{-1} = \begin{pmatrix}
            67/142 & 19/142 \\
            19/142 & 9/142
        \end{pmatrix}$, $\det{B} = 142$.

        $f_{\vec{\eta}}(x_1, x_2) = \frac{1}{2\pi \cdot \sqrt{142}} \cdot \exp\left\{-\frac{1}{2}\left( B^{-1}(\vec{x} - \vec{d}), (\vec{x} - \vec{d}))\right)\right\} = $

        $ = \frac{1}{2\pi \cdot \sqrt{142}} \cdot \exp\left\{-\frac{1}{2}\left( \frac{67}{142}(x_1+1)^2 + \frac{38}{142}(x_1+1)(x_2-3) + \frac{9}{142}(x_2-3)^2\right)\right\}$.

        Знайдемо розподіл $\eta_1$ за допомогою характеристичної функції: 
        $\chi_{\eta_1}(t) = \chi_{\vec{\eta}}(t, 0) = $

        $= \exp\left\{-i t - \frac{1}{2}\cdot 9t^2\right\} \Rightarrow
        \eta_1 \sim \mathrm{N}(-1, 9)$.
    \end{enumerate}
\end{example}

\subsection{Вироджені гауссівські вектори}
Розглянемо випадок $n$-вимірного нормального розподілу $\vec{\xi} \sim \mathrm{N}(\vec{a}, K)$ з виродженою матрицею $K$,
у якої $\mathrm{rang}K = m < n$. Існує ортогональне перетворення $U$ таке, що $K = UDU^{*}$,
$D = \mathrm{diag}(\lambda_1, \lambda_2, ..., \lambda_m, 0, ..., 0)$. 
Знайдемо розподіл вектора $\vec{\eta} = U^{*}\vec{\xi}$ за властивістю характеристичної функції:
\begin{gather*}
    \chi_{\vec{\eta}} (\vec{t}) = \chi_{\vec{\xi}} (U^{**}\vec{t}) = \chi_{\vec{\xi}} (U\vec{t}) = 
\exp\left\{i(\vec{a}, U\vec{t}) - \frac{1}{2}(KU\vec{t}, U\vec{t})\right\} = 
\exp\left\{i(U^{*}\vec{a}, \vec{t}) - \frac{1}{2}(U^{*}KU\vec{t}, \vec{t})\right\} = \\
= \left[ U^{*}\vec{a} = \vec{b}, U^{*}KU = D\right] = 
\exp\left\{i(\vec{b}, \vec{t}) - \frac{1}{2}(D\vec{t}, \vec{t})\right\} = 
\exp\left\{i\sum\limits_{k=1}^n b_k t_k - \frac{1}{2}\sum\limits_{k=1}^m \lambda_k t_k^2\right\} = \\
= \exp\left\{i\sum\limits_{k=m+1}^n b_k t_k\right\} \cdot
\exp\left\{i\sum\limits_{k=1}^m b_k t_k - \frac{1}{2}\sum\limits_{k=1}^m \lambda_k t_k^2\right\}
\end{gather*}

Перший множник --- це характеристична функція сталого вектору, а другий --- характеристична функція гауссівського вектору меншої розмірності.

Позначимо $\vec{c} = (0, ..., 0, b_{m+1}, ..., b_{n})^T$,
$\vec{d} = (b_1, ..., b_m, 0, ..., 0)^T$,
$\vec{\zeta} = (\zeta_1, ..., \zeta_m, 0, ..., 0)^T$, $\zeta_k \sim {N}(0, 1)$ та незалежні у сукупності,
$\Lambda = \mathrm{diag}(\sqrt{\lambda_1}, \sqrt{\lambda_2}, ..., \sqrt{\lambda_m}, 1, ..., 1)$.

Тоді $\vec{\eta} = \vec{c} + \left(\Lambda \vec{\zeta} + \vec{d} \right)$. 
Оскільки $\vec{\eta} = U^{*}\vec{\xi}$, то $\vec{\xi} = U\vec{\eta} = U\Lambda \vec{\zeta} + U(\vec{c} + \vec{d}) = U\Lambda \vec{\zeta} + \vec{a}$.
Позначимо $U\Lambda = L$. Остаточно: $\vec{\xi} = L\vec{\zeta} + \vec{a}$.

Що тепер можна сказати про розподіл такого вектора? Оскільки перші $m$ координат $\vec{\zeta}$ можуть приймати довільні дійсні значення, то
можливі значення вектора $L\vec{\zeta}$ --- це точки з лінійної оболонки векторів $\left\{u_1, u_2, ..., u_m\right\}$ (це власні вектори, що відповідають ненульовим власним числам $K$).
Оскільки власні вектори, що відповідають нульовому власному числу, утворюють базис $\mathrm{Ker}K$, то лінійна оболонка $\left\{u_1, u_2, ..., u_m\right\}$ --- це $\mathrm{Im}K$.


\textbf{Висновок:} $n$-вимірний гауссівський вектор з виродженою кореляційною матрицею, що має ранг $m<n$, можна представити
у вигляді суми свого вектора математичного сподівання та вектора, перші $m$ координат є незалежними у сукупності та мають стандартний нормальний розподіл, 
а інші координати нульові, помноженого на невироджену матрицю. 

Усі значення $n$-вимірного гауссівського вектора $\vec{\xi} \sim \mathrm{N}(\vec{a}, K)$ з виродженою кореляційною матрицею, що має ранг $m<n$,
зосереджені на $m$-вимірному підпросторі $\mathbb{R}^n$ --- $\mathrm{Im}K + \vec{a}$.
\subsection{Нормальний розподіл на площині} 
Розглянемо двовимірний гауссівський вектор $\vec{\xi} = (\xi_1, \xi_2)^T \sim \mathrm{N}(\vec{a}, K)$.

Нехай $\vec{a} = \begin{pmatrix}
    a_1 \\ a_2
\end{pmatrix}$, $K = \begin{pmatrix}
    \sigma_1^2 & \rho \sigma_1 \sigma_2 \\
    \rho \sigma_1 \sigma_2 & \sigma_2^2
\end{pmatrix}$, $\det{K} = \sigma_1^2 \sigma_2^2 - \rho^2 \sigma_1^2 \sigma_2^2 = (1-\rho^2)\sigma_1^2\sigma_2^2$, 
$K^{-1} = \frac{1}{(1-\rho^2)\sigma_1^2\sigma_2^2}\begin{pmatrix}
    \sigma_2^2 & -\rho \sigma_1 \sigma_2 \\
    -\rho \sigma_1 \sigma_2 & \sigma_1^2
\end{pmatrix}$.
Окремо запишемо щільність:
\begin{equation}
    \begin{gathered}
        f_{\vec{\xi}}(x,y) = \frac{1}{2\pi\sigma_1\sigma_2\sqrt{1-\rho^2}} \exp\left\{-\frac{1}{2}\cdot
        \frac{\left( \sigma_2^2(x-a_1)^2 - 2\rho\sigma_1\sigma_2(x-a_1)(y-a_2) + \sigma_1^2(y-a_2)^2\right)}{\sigma_1^2\sigma_2^2(1-\rho^2)}
        \right\} = \\
        = \frac{1}{2\pi\sigma_1\sigma_2\sqrt{1-\rho^2}} \exp\left\{
            -\frac{1}{2(1-\rho^2)}\cdot
            \left(\frac{(x-a_1)^2}{\sigma_1^2} -
            2\rho\frac{x-a_1}{\sigma_1}\frac{y-a_2}{\sigma_2} +
            \frac{(y-a_2)^2}{\sigma_2^2}
            \right) 
        \right\}
    \end{gathered}
\end{equation}
При $\rho = 0$ отримуємо $f_{\vec{\xi}}(x,y) = \frac{1}{\sqrt{2\pi}\sigma_1} e^{-\frac{(x-a_1)^2}{2\sigma_1^2}} \cdot \frac{1}{\sqrt{2\pi}\sigma_2} e^{-\frac{(y-a_2)^2}{2\sigma_2^2}} = f_{\xi_1}(x)\cdot f_{\xi_2}(y)$.

Поверхня, утворена графіком щільності,
називається \emph{поверхнею} або \emph{палаткою Гаусса}.
 
\hbox to \hsize{\hfil{
    \includegraphics[scale=0.45]{bivar}
}\hfil}
Лінії рівня цієї поверхні задають \emph{еліпси розсіювання}:
$$f_{\vec{\xi}}(x,y) = C \Leftrightarrow \frac{(x-a_1)^2}{\sigma_1^2} -
2\rho\frac{x-a_1}{\sigma_1}\frac{y-a_2}{\sigma_2} +
\frac{(y-a_2)^2}{\sigma_2^2} = \lambda^2$$

\begin{tabular}{c p{7.8cm}}
\begin{tikzpicture}[baseline={(current bounding box.north)}]
    \draw [->] (-0.5, 0) -- (5, 0);
    \draw [->] (0, -0.5) -- (0, 3);
    \draw (2, 1.5) circle [x radius=1, y radius=0.5, rotate=40];
    \draw (2, 1.5) circle [x radius=1.3, y radius=0.7, rotate=40];
    \draw (2, 1.5) circle [x radius=1.5, y radius=0.9, rotate=40];
    \draw [fill] (2, 1.5) circle [radius = 0.05];
    \draw [<-] (3.5, 2.824) -- (0.3, 0);
    \draw [<-] (1, 2.64) -- (3, 0.38);
    \draw [dashed] (2, 1.5) -- (0, 1.5);
    \draw [dashed] (2, 1.5) -- (2, 0);
    \node [below] at (5, 0) {$x$};
    \node [left] at (0, 2.8) {$y$};
    \node [below] at (2, 0) {$a_1$};
    \node [left] at (0, 1.5) {$a_2$};
    \draw (0.6, 0) arc (0:40:0.3);
    \node [above right] at (0.55, 0) {$\alpha$};
\end{tikzpicture}&
$\mathrm{tg}2\alpha = \frac{2\rho\sigma_1\sigma_2}{\sigma_1^2 - \sigma_2^2}$\newline
Осі еліпсів називаються \emph{головними осями розсіювання}.\newline
Якщо $\rho = 0$, то головні осі розсіювання паралельні осям координат.
\end{tabular}

Позначимо $E_\lambda = \left\{(x;y) \in \mathbb{R}^2: \frac{(x-a_1)^2}{\sigma_1^2} -
2\rho\frac{x-a_1}{\sigma_1}\frac{y-a_2}{\sigma_2} +
\frac{(y-a_2)^2}{\sigma_2^2} \leq \lambda^2\right\}$.
Знайдемо ймовірність потрапляння в цю область у випадку $\rho = 0$:

\begin{gather*}
    P\left\{\vec{\xi} \in E_\lambda\right\} = \iint\limits_{E_\lambda} f_{\vec{\xi}}(x,y) dx dy = 
    \frac{1}{2\pi\sigma_1\sigma_2} \iint\limits_{E_\lambda} \exp\left\{-\frac{1}{2}\left( 
        \frac{(x-a_1)^2}{\sigma_1^2} + 
        \frac{(y-a_2)^2}{\sigma_2^2}
    \right)\right\} dx dy = \\
    \left[ \begin{gathered}
        x-a_1 = \lambda\sigma_1 r \cos\varphi \\ 
        y-a_2 = \lambda\sigma_2 r \sin\varphi \\
        |\mathcal{J}| = \lambda^2 \sigma_1 \sigma_2 r
    \end{gathered}\right] = 
    \frac{\lambda^2}{2\pi} \int\limits_0^{2\pi}  d\varphi
    \int\limits_0^1 r e^{-\frac{\lambda^2r^2}{2}} dr = 
    \int\limits_0^1 e^{-\frac{\lambda^2r^2}{2}} d\left(\frac{\lambda^2r^2}{2} \right) = 
    1 - e^{-\frac{\lambda^2}{2}}
\end{gather*}

\begin{exercise}
    Довести, що у випадку $\rho > 0$ $P\left\{\vec{\xi} \in E_\lambda\right\} = 
    1 - e^{-\frac{\lambda^2}{2(1-\rho^2)}}$.
\end{exercise}

\begin{example}
    $\vec{\xi} \sim \mathrm{N}\left( \begin{pmatrix}
        -1 \\ 1
    \end{pmatrix}, \begin{pmatrix}
        1 & 2 \\
        2 & 16
    \end{pmatrix}\right)$. Знайти рівняння еліпса розсіювання, для якого \\
    $P\left\{\vec{\xi} \in E_\lambda\right\} = 0.93$. З кореляційної матриці $\rho = 0.5$.
    Розв'яжемо відносно $\lambda^2$ рівняння $1 - e^{-\frac{\lambda^2}{2(1-0.5^2)}} = 0.93
    \Leftrightarrow e^{-\frac{\lambda^2}{1.5}} = 0.07 \Rightarrow \lambda^2 \approx 4$. Тому рівняння еліпса
    $\frac{(x+1)^2}{1} - \frac{x+1}{1} \frac{y-1}{4} + \frac{(y-1)^2}{16} = 4$ або 
    $\frac{(x+1)^2}{4} - \frac{(x+1)(y-1)}{16} + \frac{(y-1)^2}{64} = 1$.
\end{example}

Розглянемо гауссівський вектор $\vec{\xi} = (\xi_1, \xi_2)$ з незалежними координатами.
Для нього $F_{\vec{\xi}}(x,y) = F_{\xi_1}(x) F_{\xi_2}(y) = 
\left(\frac{1}{2} + \Phi\left( \frac{x-a_1}{\sigma_1}\right) \right)
\left(\frac{1}{2} + \Phi\left( \frac{y-a_2}{\sigma_2}\right) \right)$.
Тоді ймовірність потрапляння в прямокутник $\Pi = \left\{(x;y)\in\mathbb{R}^2 : \alpha \leq x < \beta, \gamma \leq y < \delta\right\}$
дорівнює $$\left( \Phi\left( \frac{\beta-a_1}{\sigma_1}\right) - \Phi\left( \frac{\alpha-a_1}{\sigma_1}\right)\right) \cdot
\left( \Phi\left( \frac{\delta-a_2}{\sigma_2}\right) - \Phi\left( \frac{\gamma-a_2}{\sigma_2}\right)\right)$$
\begin{exercise}
    Довести формулу для $P\left\{\vec{\xi}\in\Pi\right\}$.
\end{exercise}
Для гауссівського вектору з незалежними координатами також має місце <<правило $3\sigma$>>:
\begin{gather*}
    P\left\{\vec{\xi} \in (a_1-3\sigma_1; a_1+3\sigma_1)\times(a_2-3\sigma_2; a_2+3\sigma_2)\right\} = \\
    = P\left\{\xi_1 \in(a_1-3\sigma_1; a_1+3\sigma_1)\right\}\cdot P\left\{\xi_2 \in(a_2-3\sigma_2; a_2+3\sigma_2)\right\}
    \approx 0.9973^2 \approx 0.9946
\end{gather*}

\subsection{Колове розсіювання}
\begin{definition}
    Двовимірний гауссівський вектор має \emph{колове розсіювання}, якщо
    $\vec{a} = \vec{0}, \; K = \begin{pmatrix}
        \sigma^2 & 0 \\
        0 & \sigma^2
    \end{pmatrix}$. В цьому випадку $f_{\vec{\xi}} = \frac{1}{2\pi\sigma^2} \exp\left\{-\frac{x^2+y^2}{2\sigma^2}\right\}$, 
    а еліпси розсіювання стають колами.
\end{definition}
Знайдемо ймовірність потрапляння такого $\vec{\xi}$ в коло $K_R = \left\{(x; y) \in \mathbb{R}^2 : x^2 + y^2 \leq R^2\right\}$:
\begin{gather*}
    P\left\{\vec{\xi} \in K_R\right\} = \frac{1}{2\pi\sigma^2} \iint\limits_{K_R} e^{-\frac{x^2+y^2}{2\sigma^2}} dx dy = 
    \left[ \begin{gathered}
        x = \sigma r\cos\varphi \\
        y = \sigma r\sin\varphi \\
        |\mathcal{J}| = \sigma^2 r
    \end{gathered}\right] =
    \frac{1}{2\pi} \int\limits_0^{2\pi} d\varphi \int\limits_0^{R/\sigma} r e^{-\frac{r^2}{2}} dr = \\
    = \int\limits_0^{R/\sigma} e^{-\frac{r^2}{2}} d \left( \frac{r^2}{2}\right) = 
    1 - e^{-\frac{R^2}{2\sigma^2}}
\end{gather*}
Таким чином, тепер відомий розподіл норми $\vec{\xi}$. $\eta = \Vert \vec{\xi} \Vert = \sqrt{\xi_1^2 + \xi_2^2}$ --- випадкова величина,
для якої $F_{\eta}(R) = P\left\{\Vert \vec{\xi} \Vert < R\right\} = \begin{cases}
    1 - e^{-\frac{R^2}{2\sigma^2}}, & R > 0 \\
    0, & R \leq 0
\end{cases}$ --- це \emph{розподіл Релея}. Ця властивість переноситься й на довільну скінченну розмірність.

\begin{exercise}
    Знайти основні числові характеристики розподілу Релея.
\end{exercise}

\subsection{Умовний гауссівський розподіл на площині}
Знайдемо одну з умовних щільностей розподілу:
\begin{gather*}
    f_{\xi_1}(x/\xi_2 = y) = \frac{f_{\vec{\xi}}(x,y)}{f_{\xi_2}(y)} = 
    \frac{\frac{1}{2\pi\sigma_1\sigma_2\sqrt{1-\rho^2}} \exp\left\{
        -\frac{1}{2(1-\rho^2)}\cdot
        \left(\frac{(x-a_1)^2}{\sigma_1^2} -
        2\rho\frac{x-a_1}{\sigma_1}\frac{y-a_2}{\sigma_2} +
        \frac{(y-a_2)^2}{\sigma_2^2}
        \right)\right\} }{
            \frac{1}{\sqrt{2\pi}\sigma_2} \exp\left\{-\frac{1}{2\sigma_2^2}(y-a_2)^2\right\}
        } = \\
        \frac{1}{\sqrt{2\pi}\sigma_1\sqrt{1-\rho^2}} \exp\left\{
            -\frac{1}{2(1-\rho^2)}
        \left(\frac{(x-a_1)^2}{\sigma_1^2} -
        2\rho\frac{(x-a_1)(y-a_2)}{\sigma_1\sigma_2} +
        \frac{(y-a_2)^2}{\sigma_2^2}
        \right) +
        \frac{(y-a_2)^2}{2\sigma_2^2}
        \right\} = \\
        = \left [ \frac{(y-a_2)^2}{\sigma_2^2} - \frac{(1-\rho^2)(y-a_2)^2}{\sigma_2^2} = \rho^2 \frac{(y-a_2)^2}{\sigma_2^2}\right]=\\
        = \frac{1}{\sqrt{2\pi}\sigma_1\sqrt{1-\rho^2}} \exp\left\{-\frac{1}{2(1-\rho^2)}\left( \frac{x-a_1}{\sigma_1} - \rho\cdot\frac{y-a_2}{\sigma_2}\right)^2\right\} = \\
        = \frac{1}{\sqrt{2\pi}\sigma_1\sqrt{1-\rho^2}} \exp\left\{-\frac{1}{2\sigma_1^2(1-\rho^2)}\left( x- \left( a_1 + \frac{\rho\sigma_1(y-a_2)}{\sigma_2}\right)\right)^2\right\}
\end{gather*}
Аналогічно $f_{\xi_2}(y/\xi_1 = x) = \frac{1}{\sqrt{2\pi}\sigma_2\sqrt{1-\rho^2}} \exp\left\{-\frac{1}{2\sigma_2^2(1-\rho^2)}\left( y- \left( a_2 + \frac{\rho\sigma_2(x-a_1)}{\sigma_1}\right)\right)^2\right\}$.

\noindentОтже, обидві умовні щільності є щільностями нормального розподілу.

$E\left( \xi_1 / \xi_2 = y\right) = a_1 + \frac{\rho\sigma_1(y-a_2)}{\sigma_2}$, 
$E\left( \xi_2 / \xi_1 = x\right) = a_2 + \frac{\rho\sigma_2(x-a_1)}{\sigma_1}$. Лінії регресії --- прямі,

\noindent
$\rho\frac{\sigma_1}{\sigma_2} = \frac{K\xi_1\xi_2}{D\xi_2}$ та $\rho\frac{\sigma_2}{\sigma_1}  = \frac{K\xi_1\xi_2}{D\xi_1}$ --- відповідні кутові коефіцієнти прямих регресії.

$D\left( \xi_1 / \xi_2 = y\right) = \sigma_1^2(1-\rho^2)$, 
$D\left( \xi_2 / \xi_1 = x\right) = \sigma_2^2(1-\rho^2)$. Умовні дисперсії є сталими, ця властивість називається \emph{гомоскедастичністю}.
    \chapter{Функції випадкових аргументів}
        % !TEX root = ../main.tex

\section{Функції одного випадкового аргументу}
При розв'язанні задач іноді виникає необхідність визначити закон розподілу невипадкової функції від випадкового аргумента,
розподіл якого є відомим.
\subsection{Функції від дискретного випадкового аргументу}
Нехай $\xi$ --- дискретна випадкова величина з рядом розподілу 
\begin{center}
    \begin{tabular}{|c|c|c|c|c|c|}
        \hline
        $\xi$ & $x_1$ & $x_2$ & $...$ & $x_n$ & $...$ \\
        \hline
        $p$ & $p_1$ & $p_2$ & $...$ & $p_n$ & $...$ \\
        \hline
    \end{tabular}    
\end{center}
а $\varphi$ --- деяка вимірна числова функція, 
область визначення якої містить можливі значення $\xi$.
Можна розглядати випадкову величину $\eta = \varphi(\xi)$.

Очевидно, що $\eta$ --- теж ДВВ, що приймає значення $y_k = \varphi(x_k)$, 
$k = 1, 2, ...$ з ймовірностями $\P\{\eta = y_k\} = \P\{\xi = x_k : \varphi(x_k) = y_k\}$.
Якщо існує лише одне $\varphi^{-1} (y_k) = x_k$, то $\P\{\eta = y_k\} = \P\{\xi = x_k\} = p_k$.
Однак, можливо, що $y_k = \varphi(x_{k_1}) = \varphi(x_{k_2}) = ... = \varphi(x_{k_m})$.
Тоді $\P\{\eta = y_k\} = p_{k_1} + p_{k_2} + ... + p_{k_m}$.

\begin{example}
    Нехай $\xi$ має ряд розподілу
    \begin{center}
        \begin{tabular}{|c|c|c|c|c|c|}
            \hline
            $\xi$ & $-2$ & $-1$ & $0$ & $1$ & $2$ \\
            \hline
            $p$ & $1/9$ & $2/9$ & $1/9$ & $2/9$ & $3/9$ \\
            \hline
        \end{tabular}
    \end{center}
    Знайти закон розподілу $\eta = \xi^2$.
    
    З ряду розподілу $\xi$ видно, що $\xi^2$ приймає значення $0, 1, 4$ з
    ймовірностями $1/9$, $2/9 + 2/9$ та $1/9 + 3/9$.
    Отже, маємо ряд розподілу $\eta$: 
    \begin{center}
        \begin{tabular}{|c|c|c|c|}
            \hline
            $\eta$ & $0$ & $1$ & $4$ \\
            \hline
            $p$ & $1/9$ & $4/9$ & $4/9$ \\
            \hline
        \end{tabular}
    \end{center}
\end{example}

З побудови ряду розподілу функції від дискретного
випадкового аргументу випливає, що $\E(\varphi(\xi))^k = \sum\limits_{i=1}^{n(\infty)} \varphi^k(x_i) \P\{\xi = x_i\}$.

\subsection{Функції від неперервного випадкового аргументу}
Нехай $\xi$ --- неперервна випадкова величина, $f_\xi(x)$ --- її щільність розподілу, а $\varphi$ --- монотонна неперервна числова функція,
для якої обернена функція $\varphi^{-1}$ є диференційовною майже скрізь.
Розглянемо два випадки.

\begin{enumerate}
    \item $\eta = \varphi(\xi)$, де $\varphi$ --- монотонно зростаюча.

    \begin{tabular}{c p{8.8cm}}
        \begin{tikzpicture}[xscale = 0.7, yscale = 0.3, baseline={(current bounding box.north)}]
            \draw [->] (-3, 0) -- (3, 0);
            \draw [->] (0, -0.5) -- (0, 5.5);
            \draw [domain=-3:3, smooth, variable = \x, ultra thick] plot ({\x}, {e^(\x/2});
            \node [below] at (3, 0) {$x$};
            \node [left] at (0, 5.5) {$y$};
        \end{tikzpicture} &
        $F_\eta (y) = \P\left\{ \eta < y\right\} = \P\left\{ \varphi(\xi) < y\right\}$.

        Нехай $y = \varphi(x)$, тоді з монотонності $\varphi$ маємо рівність подій $\left\{ \varphi(\xi) < y\right\} = \left\{ \xi < x\right\}$.
        Тому $F_\eta (y) = F_\xi (x) = \int\limits_{-\infty}^x f_\xi(t) dt$. $\varphi$ має обернену, $x = \varphi^{-1} (y)$.
    \end{tabular}

    Отже, $F_\eta (y) = \int\limits_{-\infty}^{\varphi^{-1} (y)} f_\xi(t) dt$,
    $f_\eta(y) = F^\prime_\eta (y) = f_\xi\left(\varphi^{-1} (y)\right) \cdot \left(\varphi^{-1} (y) \right)^{\prime}$.
    \item $\eta = \varphi(\xi)$, де $\varphi$ --- монотонно спадна.
    
    \begin{tabular}{c p{8.8cm}}
        \begin{tikzpicture}[xscale = 0.7, yscale = 0.3, baseline={(current bounding box.north)}]
            \draw [->] (-3, 0) -- (3, 0);
            \draw [->] (0, -0.5) -- (0, 5.5);
            \draw [domain=-3:3, smooth, variable = \x, ultra thick] plot ({\x}, {e^(-\x/2});
            \node [below] at (3, 0) {$x$};
            \node [left] at (0, 5.5) {$y$};
        \end{tikzpicture} &
        $F_\eta (y) = \P\left\{ \eta < y\right\} = \P\left\{ \varphi(\xi) < y\right\}$.

        Нехай $y = \varphi(x)$, тоді з монотонності $\varphi$ маємо рівність подій $\left\{ \varphi(\xi) < y\right\} = \left\{ \xi > x\right\}$.
        Тому $F_\eta (y) = 1 - F_\xi (x) = 1 - \int\limits_{-\infty}^x f_\xi(t) dt$. $\varphi$ має обернену, $x = \varphi^{-1} (y)$.
    \end{tabular}

    Отже, $F_\eta (y) = 1 - \int\limits_{-\infty}^{\varphi^{-1} (y)} f_\xi(t) dt$,
    $f_\eta(y) = F^\prime_\eta (y) = - f_\xi\left(\varphi^{-1} (y)\right) \cdot \left(\varphi^{-1} (y) \right)^{\prime}$.
\end{enumerate}

\noindent Остаточно, для монотонної неперервної $\varphi$,
для якої обернена функція $\varphi^{-1}$ є диференційовною майже скрізь, має місце формула:
\begin{equation}\label{eq:phi_xi}
    \eta = \varphi(\xi), f_\eta(y) = f_\xi\left(\varphi^{-1} (y)\right) \cdot \left|\left(\varphi^{-1} (y) \right)^{\prime}\right|
\end{equation}
При цьому обов'язково треба зазначити допустимі значення $y$.

\begin{example}
    Нехай $\xi \sim \mathrm{U}(-\frac{\pi}{2}; \frac{\pi}{2})$. Знайти розподіл $\eta = \sin \xi$.

    \noindent $f_\xi(x) = \begin{cases}
        \frac{1}{\pi}, & |x| < \frac{\pi}{2} \\
        0, & |x| \geq \frac{\pi}{2}
    \end{cases}$, $\varphi(x) = \sin x$, $\varphi^{-1} (y) = \arcsin(y)$, $\left(\varphi^{-1} (y)\right)^{\prime} = \frac{1}{\sqrt{1-y^2}}$.
    $\eta$ може приймати значення від $-1$ до $1$. 
    Отже, $f_\eta(y) = \begin{cases}
        \frac{1}{\pi\sqrt{1-y^2}}, & |y| < 1 \\
        0, & |y| \geq 1
    \end{cases}$ --- це <<закон арксинуса>>.
\end{example}

\begin{example}
    Нехай $\xi \sim \mathrm{N}(a, \sigma^2)$. Знайти розподіл $\eta = e^\xi$.

    В цьому випадку $\varphi(x) = e^x$, $\varphi^{-1}(y) = \ln y$, $\left(\varphi^{-1} (y)\right)^{\prime} = \frac{1}{y}$.
    $\eta$ може приймати лише додатні значення, тому отримуємо щільність
    \begin{gather*}
        f_\eta(y) = \begin{cases}
            \frac{1}{\sqrt{2\pi}\sigma y} e^{-\frac{(\ln y-a)^2}{2\sigma^2}}, & y > 0 \\
            0, & y \leq 0
        \end{cases}
    \end{gather*}
    Розподіл $\eta$ називається \emph{логнормальним} (оскільки $\ln \eta$ має нормальний розподіл)
    та позначається $\eta \sim \mathrm{LogN}(a, \sigma^2)$.
\end{example}
\begin{exercise}
    Перевірити, що $\E\eta = e^{a + \frac{\sigma^2}{2}}$ і
    $\D\eta = \left(e^{\sigma^2} - 1\right) e^{2a + \sigma^2}$.
\end{exercise}

\begin{example}
    Нехай $\xi$ --- довільна НВВ, $F_\xi (x)$ --- її функція розподілу. 
    Знайти закон розподілу $\eta = F_\xi (\xi)$.

    \noindent $F_\eta (y) = \P\{\eta < y\} = \P\{F_\xi(\xi) < y\} = \P\{\xi < F_\xi^{-1}(y)\} = F_\xi(F_\xi^{-1}(y)) = \begin{cases}
        0, & y \leq 0 \\
        y, & 0 < y \leq 1 \\
        1, & y > 1
    \end{cases}$.

    \noindent Отже, $F_\xi (\xi) \sim U\left<0;1\right>$. Ще раз зауважимо, що результат має місце для довільної НВВ $\xi$.
\end{example}

\begin{remark}
    Якщо $\varphi$ не є неперервною, то $\eta = \varphi(\xi)$ може не бути НВВ. Наведемо декілька прикладів.
    \begin{enumerate}
        \item Нехай $\xi \sim \mathrm{Exp}(1)$. Знайдемо розподіл $\eta = \left[ \xi\right]$, де $\left[ \cdot\right]$ --- ціла частина.
        Оскільки ціла частина приймає значення $0, 1, 2, ...$, знайдемо відповідні ймовірності. 
        
        $\P\left\{ \left[ \xi\right] = n\right\} = \P\left\{ \xi \in [n; n+1)\right\} = \int\limits_n^{n+1} e^{-x} dx = e^{-n} - e^{-(n+1)} = e^{-n}(1 - e^{-1})$, $n = 0, 1, 2, ...$

        Отже, $\eta = \left[ \xi\right] \sim \mathrm{Geom}(1-e^{-1})$.

        За цим прикладом можна зробити висновок: якщо $\varphi$ --- кусково стала, а $\xi$ --- НВВ, то $\varphi(\xi)$ --- ДВВ.
        \item Нехай $\xi \sim \mathrm{Exp}(1)$. Знайдемо розподіл $\eta = \left\{ \xi\right\} = \xi - \left[ \xi\right]$. 
        В цьому випадку $\varphi$ є кусково неперервною функцією та приймає значення з інтервалу $[0; 1)$.
        
        Для $0< y < 1$ $F_\eta (y) = \P\left\{\eta < y\right\} = \P\left\{\eta \in [0; y)\right\} = \P\left\{\xi \in [n; n+y), n = 0, 1, 2, ...\right\}$.

        Для кожного $n$ $\P\left\{\xi \in [n; n+y)\right\} = \int\limits_n^{n+y} e^{-x} dx = e^{-n}(1 - e^{-y})$. 
        Отже, $\P\left\{\eta \in [0; y)\right\} = \sum\limits_{n=0}^{\infty} \P\left\{\xi \in [n; n+y)\right\} = (1 - e^{-y}) \sum\limits_{n=0}^{\infty} e^{-n} = \frac{1 - e^{-y}}{1 - e^{-1}}$. Отримали функцію розподілу $\eta$:

        $F_\eta (y) = \begin{cases}
            0, & y \leq 0 \\
            \frac{1 - e^{-y}}{1 - e^{-1}}, & 0 < y \leq 1 \\
            1, & y > 1
        \end{cases}$. Таким чином, $\eta$ --- НВВ.
    \end{enumerate}
\end{remark}

Узагальнимо формулу \eqref{eq:phi_xi} на випадок немонотонної функції $\varphi$.
Нехай $\xi$ --- НВВ, $f_\xi(x)$ --- її щільність розподілу, а $\varphi$ --- неперервна немонотонна числова функція.
Тоді область можливих значень $\xi$ можна розбити на проміжки, на яких $\varphi$ буде монотонною. 
Тоді, скориставшись результатом для монотонної функції, отримаємо формулу для щільності $\eta = \varphi(\xi)$:
\begin{equation}
    f_\eta (y) = \sum\limits_{k=1}^m f_\xi\left(\varphi_k^{-1} (y)\right) \cdot \left|\left(\varphi_k^{-1} (y) \right)^{\prime}\right|
\end{equation}
де $m$ --- кількість проміжків монотонності, а $\varphi_k^{-1}$ --- відповідні обернені функції, які мають бути диференційовними майже скрізь.

\begin{example}
    \begin{enumerate}
        \item $\xi \sim \mathrm{U}(-\frac{\pi}{2}, \frac{\pi}{2})$, знайти розподіл $\eta = \cos\xi$.

        На проміжках $(-\frac{\pi}{2}; 0]$ та $[0; \frac{\pi}{2})$ $\cos x$ є монотонною функцією, 
        відповідні обернені --- $\varphi_1^{-1} (y) = -\arccos y$, $\varphi_2^{-1} (y) = \arccos y$.

        Тоді $f_\eta (y) = \frac{1}{\pi} \cdot \left| - \frac{1}{\sqrt{1-y^2}}\right| + \frac{1}{\pi} \cdot \left|\frac{1}{\sqrt{1-y^2}}\right| = \begin{cases}
            \frac{2}{\pi} \frac{1}{\sqrt{1-y^2}}, & y \in [0; 1) \\
            0, & y \notin [0; 1)
        \end{cases}$.
        \item\label{ex:squared_norm_distr} $\xi \sim \mathrm{N}(a, \sigma)$, знайти розподіл $\eta = \xi^2$.

        На проміжках $(-\infty; 0]$ та $[0; +\infty)$ $x^2$ є монотонною функцією, 
        відповідні обернені --- $\varphi_1^{-1} (y) = -\sqrt{y}$, $\varphi_2^{-1} (y) = \sqrt{y}$.
        Щільність розподілу $\xi$ --- $f_\xi (x) = \frac{1}{\sqrt{2\pi}\sigma} e^{-\frac{(x-a)^2}{2\sigma^2}}$.

        Тоді $f_\eta (y) = \begin{cases}
            \frac{1}{\sqrt{2\pi}\sigma} e^{-\frac{(-\sqrt{y}-a)^2}{2\sigma^2}} \cdot \frac{1}{2\sqrt{y}} + 
        \frac{1}{\sqrt{2\pi}\sigma} e^{-\frac{(\sqrt{y}-a)^2}{2\sigma^2}} \cdot \frac{1}{2\sqrt{y}}, & y > 0 \\
        0, & y \leq 0
        \end{cases}$.

        При $\xi \sim \mathrm{N}(0, 1)$ $f_\eta(y) = \begin{cases}
            \frac{1}{\sqrt{2\pi y}} e^{-y/2}, y > 0 \\
            0, y \leq 0
        \end{cases}$, що означає $\xi^2 \sim \Gamma(\frac{1}{2}, \frac{1}{2})$.
    \end{enumerate}
\end{example}

\subsection{Числові характеристики функції неперервного випадкового аргументу}
Для функцій від ДВВ вже було отримано, що
$\E(\varphi(\xi))^k = \sum\limits_{i=1}^{n(\infty)} \varphi^k(x_i) \P\{\xi = x_i\}$.
Доведемо аналогічне твердження для функцій від НВВ.

\begin{proposition*}
    $\E (\varphi(\xi))^k = \int\limits_{-\infty}^{+\infty} \varphi^k(x) f_\xi(x) dx$ для всіх $k \in \mathbb{N}$.
\end{proposition*}
\begin{proof}
    Достатньо довести для монотонно зростаючої $\varphi$.

    \noindent$\eta = \varphi(\xi)$, $\E\eta^k = \int\limits_{-\infty}^{+\infty} y^k f_\eta(y) dy = \int\limits_{-\infty}^{+\infty} y^k f_\xi (\varphi^{-1} (y)) (\varphi^{-1} (y))^\prime dy =
    \int\limits_{-\infty}^{+\infty} y^k f_\xi(\varphi^{-1} (y)) d \varphi^{-1} (y) = \left[ \varphi^{-1} (y) = x, y = \varphi(x)\right] = \int\limits_{-\infty}^{+\infty} \varphi^k(x) f_\xi(x) dx$.
    У випадку монотонно спадної $\varphi$ під інтегралом отримаємо $-(\varphi^{-1} (y))^\prime$, а інтеграл після заміни змінної буде від $+\infty$ до $-\infty$.
    У випадку немонотонної $\varphi$ треба буде скористатися адитивністю інтеграла.
\end{proof}

\begin{remark}
    Для знаходження числових характеристик функції неперервного випадкового аргументу розподіл самої функції знаходити не потрібно.
\end{remark}

\begin{exercise}
    Дослідити, за яких умов на $\varphi$ ця формула має місце.
\end{exercise}

У випадках, коли $\varphi(x) = \cos (s x)$ або $\sin (s x)$ ($s\in \mathbb{R}$),
для знаходження числових характеристик можна скористатися характеристичною функцією $\xi$,
оскільки $\chi_{\xi}(s) = \E e^{i s \xi} = \E \cos(s\xi) + i \cdot\E\sin(s\xi)$,
тому $\E\cos(s\xi) = \mathrm{Re} \chi_{\xi}(s)$ і
$\E\sin(s\xi) = \mathrm{Im} \chi_{\xi}(s)$.
Для знаходження моментів вищих порядків можна скористатися формулами пониження степеня.
        % !TEX root = ../main.tex

\section{Функції кількох випадкових аргументів}
\subsection{Випадок довільної функції}
Нехай $\varphi : \mathbb{R}^n \to \mathbb{R}$ --- задана вимірна числова функція.

Якщо $\vec{\xi} = \left(\xi_1, ..., \xi_n\right)^T$ --- дискретний випадковий вектор, тоді $\eta = \varphi(\vec{\xi})$ --- ДВВ.
Побудову закону розподілу $\eta$ доцільно розглянути на прикладі.
\begin{example}
    $\vec{\xi} = \left( \xi_1, \xi_2\right)$ задано таблицею розподілу:
    \begin{center}
        \begin{tabular}{|c|c|c|c|}
            \hline
            \diagbox{$\xi_2$}{$\xi_1$} & $0$ & $1$ & $2$ \\
            \hline
            $-1$ & $0.1$ & $0.2$ & $0.3$ \\
            \hline
            $1$ & $0.2$ & $0.1$ & $0.1$ \\
            \hline
        \end{tabular}    
    \end{center}

    \noindentЗнайдемо закони розподілу $\eta_1 = \xi_1 \xi_2$ та $\eta_2 = \xi_1 + \xi_2$.
    Для цього треба визначити значення, які приймають ці величини, та обчислити відповідні ймовірності.
    
    $\eta_1$ приймає значення $-2$ (коли $\xi_1 = 2$, $\xi_2 = -1$), 
    $-1$ (коли $\xi_1 = 1$, $\xi_2 = -1$), 
    $0$ (коли $\xi_1 = 0$, $\xi_2 = -1$ або $\xi_2 = 1$), 
    $1$ (коли $\xi_1 = 1$, $\xi_2 = 1$), 
    $2$ (коли $\xi_1 = 2$, $\xi_2 = 1$).

    $\eta_2$ приймає значення $-1$ (коли $\xi_1 = 0$, $\xi_2 = -1$), 
    $0$ (коли $\xi_1 = 1$, $\xi_2 = -1$), 
    $1$ (коли $\xi_1 = 0$, $\xi_2 = 1$ або $\xi_1 = 2$, $\xi_2 = -1$), 
    $2$ (коли $\xi_1 = 1$, $\xi_2 = 1$), 
    $3$ (коли $\xi_1 = 2$, $\xi_2 = 1$).
    Відповідні сумісні ймовірності отримуємо з таблиці розподілу $\vec{\xi}$. Отже,
    $\eta_1$ та $\eta_2$ мають наступні ряди розподілу:
    \begin{center}
        \begin{tabular}{|c|c|c|c|c|c|}
            \hline
            $\eta_1$ & $-2$ & $-1$ & $0$ & $1$ & $2$ \\
            \hline
            $p$ & $0.3$ & $0.2$ & $0.3$ & $0.1$ & $0.1$ \\
            \hline
        \end{tabular}
        \begin{tabular}{|c|c|c|c|c|c|}
            \hline
            $\eta_2$ & $-1$ & $0$ & $1$ & $2$ & $3$ \\
            \hline
            $p$ & $0.1$ & $0.2$ & $0.5$ & $0.1$ & $0.1$ \\
            \hline
        \end{tabular}
    \end{center}
\end{example}

Якщо $\vec{\xi} = \left(\xi_1, ..., \xi_n\right)^T$ --- неперервний випадковий вектор
із щільністю $f_{\vec{\xi}} (\vec{x})$, то можна знайти функцію розподілу $\eta = \varphi(\vec{\xi})$:
$$F_\eta (y) = \P \left\{ \eta < y\right\} = \P \left\{ \vec{\xi} \in D_y\right\} = \underset{D_y}{\int ... \int} f_{\vec{\xi}} (\vec{x}) d \vec{x}, \text{ де }D_y = \left\{\vec{x} \in \mathbb{R}^n : \varphi(\vec{x}) < y\right\}$$

Розглянемо тепер взаємно однозначне гладке перетворення $\vec{\psi} : \mathbb{R}^n \to \mathbb{R}^n$ та
знайдемо щільність розподілу $\vec{\eta} = \vec{\psi} (\vec{\xi})$. Для множини $D \subset \mathbb{R}^n$
$\P\left\{ \vec{\psi} (\vec{\xi}) \in D\right\} = \P\left\{ \vec{\xi} \in \vec{\psi}^{-1}(D)\right\} = \int_{\vec{\psi}^{-1}(D)} f_{\vec{\xi}} (\vec{x}) d\vec{x} = 
\left[ \text{заміна }\vec{y} = \vec{\psi}(\vec{x})\right] = \int_D f_{\vec{\xi}} (\vec{\psi}^{-1}(\vec{y})) \left| \mathcal{J}^{-1} (\vec{\psi}^{-1}(\vec{y})\right| d\vec{y}$,
де $\mathcal{J} (\vec{x})$ --- якобіан $\vec{\psi}$. Отже,
\begin{equation}
    f_{\vec{\eta}} (\vec{y}) = f_{\vec{\xi}} (\vec{\psi}^{-1}(\vec{y})) \left| \mathcal{J}^{-1} (\vec{\psi}^{-1}(\vec{y})\right|
\end{equation}
Як в одновимірному випадку, треба вказувати, які значення може приймати $\vec{\eta}$.

\begin{example}
    Нехай $A$ --- невироджена матриця розміру $n \times n$, $\vec{b} \in \mathbb{R}^n$ --- деякий вектор, $\vec{\xi}$ --- неперервний випадковий вектор.
    Знайти щільність розподілу $\vec{\eta} = A \vec{\xi} + \vec{b}$.

    В цьому випадку $\vec{y} = \vec{\psi}(\vec{x}) = A \vec{x} + \vec{b}$, тому $\vec{\psi}^{-1} (\vec{y}) = A^{-1} (\vec{y} - \vec{b})$. Якобіан $\vec{\psi}$ рівний $\left| \det A\right|$. 
    Отже,
    $f_{\vec{\eta}}(\vec{y}) = \frac{1}{\left| \det A\right|} f_{\vec{\xi}}\left(A^{-1} (\vec{y} - \vec{b})\right)$
\end{example}

Далі розглянемо конкретні приклади числових функцій від випадкових векторів.

\subsection{Закон розподілу мінімуму та максимуму}

Нехай випадкові величини $\xi_1, ..., \xi_n$ незалежні та однаково розподілені, як деяка випадкова величина $\xi$.
Знайдемо розподіл їх мінімуму та максимуму.
\begin{enumerate}
    \item Нехай $\mu_1 = \min\left\{\xi_1, \xi_2, ..., \xi_n\right\}$.
    \begin{gather*}
        F_{\mu_1} (x) = \P \left\{ \min\{\xi_1, \xi_2, ..., \xi_n\} < x \right\} =
        1 - \P \left\{ \min\{\xi_1, \xi_2, ..., \xi_n\} \geq x \right\} = \\ =
        1 - \P \left\{ \xi_1 \geq x, \xi_2 \geq x, ..., \xi_n \geq x \right\} = 
        1 - \P\left\{ \xi_1 \geq x\right\} \cdot \P\left\{ \xi_2 \geq x\right\} \cdot ... \cdot \P\left\{ \xi_n \geq x\right\} = \\
        = 1 - (\P\left\{ \xi \geq x\right\})^n = 1 - (1- F_\xi (x))^n
    \end{gather*}

    Якщо $\xi_1, ..., \xi_n$ --- неперервні, то $f_{\mu_1} (x) = \left( F_{\mu_1} (x)\right)^\prime = n (1- F_\xi (x))^{n-1} f_\xi(x)$.
    \item Нехай $\mu_2 = \max\left\{\xi_1, \xi_2, ..., \xi_n\right\}$.
    \begin{gather*}
        F_{\mu_2} (x) = \P \left\{ \max\{\xi_1, \xi_2, ..., \xi_n\} < x \right\} =
        \P \left\{ \xi_1 < x, \xi_2 < x, ..., \xi_n < x \right\} = \\ =
        \P\left\{ \xi_1 < x\right\} \cdot \P\left\{ \xi_2 < x\right\} \cdot ... \cdot \P\left\{ \xi_n < x\right\} =
        (\P\left\{ \xi < x\right\})^n = (F_\xi (x))^n
    \end{gather*}

    Якщо $\xi_1, ..., \xi_n$ --- неперервні, то $f_{\mu_2} (x) = \left( F_{\mu_2} (x)\right)^\prime = n (F_\xi (x))^{n-1} f_\xi(x)$.
\end{enumerate}

\begin{example}
    Нехай $\xi_1, ..., \xi_n$ незалежні та мають розподіл $\mathrm{Exp}(\lambda)$. Знайти розподіл їх мінімуму.
    За отриманими формулами запишемо щільність:

    $f_{\min}(x) = n (1-(1-e^{-\lambda x}))^{n-1} \lambda e^{-\lambda x} = n \lambda e^{-n\lambda x}$ при $x \geq 0$ та $0$ при $x < 0$.
    
    \noindentОтже, $\min\{\xi_1, ..., \xi_n\} \sim \mathrm{Exp} (n \lambda)$.
\end{example}

\subsection{Закон розподілу добутку двох НВВ}
Нехай $\vec{\xi} = (\xi_1, \xi_2)^T$ --- неперервний випадковий вектор, щільність
$f_{\vec{\xi}}(x, y)$ відома.

\noindent\textbf{Задача}: знайти розподіл $\eta = \xi_1\xi_2$.

Знайдемо функцію розподілу за вже відомою схемою:

$$F_\eta(z) = \iint\limits_{D_z}f_{\vec{\xi}}(x, y)dx dy, D_z = \left\{(x, y) \in 
\mathbb{R}^2 : xy < z\right\}$$

Розглянемо два варіанти:
\begin{enumerate}
    \item 
\begin{tabular}{c p{8.8cm}}
    \begin{tikzpicture}[baseline={(current bounding box.north)} ,scale = 0.4]
        \draw [domain=0.2:5, smooth, variable = \x, ultra thick] plot ({\x}, 
        {
            1/\x
        });
        \fill [lightgray, domain=0.2:5, smooth, variable = \x] plot ({\x}, 
        {
            1/\x
        }) -- (5, -5) -- (0, -5) -- (0, 5) -- (0.2, 5);
        \draw [domain=-5:-0.2, smooth, variable = \x, ultra thick] plot ({\x}, 
        {
            1/\x
        });
        \fill [lightgray, domain=-5:-0.2, smooth, variable = \x] plot ({\x}, 
        {
            1/\x
        }) -- (0, -5) -- (0, 5) -- (-5, 5) -- (-5, -0.2);
        \draw [->] (-5, 0) -- (5, 0);
        \draw [->] (0, -5) -- (0, 5);
        \node [below left] at (5, 0) {$x$};
        \node [below left] at (0, 5) {$y$};
        \node [above left] at (5, -5) {$D_z$};
        \node [above right] at (1, 1) {$y = \frac{z}{x}$};
    \end{tikzpicture} &
    $z > 0$, тоді $D_z = 
    \left\{x<0, y>\frac{z}{x}\right\} \cup 
    \left\{x>0, y<\frac{z}{x}\right\}$.

    $F_\eta (z) = \int\limits^{0}_{-\infty}dx\int\limits_{\frac{z}{x}}^{+\infty}
    f_{\vec{\xi}}(x, y) dy + \int\limits_0^{+\infty}dx\int\limits_{-\infty}^{\frac{z}{x}} 
    f_{\vec{\xi}}(x, y) dy$.

    Продиференціюємо обидві частини по $z$:

    $f_\eta(z) = -\int\limits_{-\infty}^0 \frac{1}{x}f_{\vec{\xi}}\left(x, \frac{z}{x}\right)dx + 
    \int\limits_0^{+\infty}\frac{1}{x}f_{\vec{\xi}}\left(x, \frac{z}{x}\right)dx$.
\end{tabular}

\item
\begin{tabular}{c p{8.8cm}}
    \begin{tikzpicture}[baseline={(current bounding box.north)} ,scale = 0.4]
        \draw [domain=0.2:5, smooth, variable = \x, ultra thick] plot ({\x}, 
        {
            -1/\x
        });
        \fill [lightgray, domain=0.2:5, smooth, variable = \x] plot ({\x}, 
        {
            -1/\x
        }) -- (5, -5) -- (0.2, -5);
        \draw [domain=-5:-0.2, smooth, variable = \x, ultra thick] plot ({\x}, 
        {
            -1/\x
        });
        \fill [lightgray, domain=-5:-0.2, smooth, variable = \x] plot ({\x}, 
        {
            -1/\x
        }) -- (-5, 5) -- (-5, 0.2);
        \draw [->] (-5, 0) -- (5, 0);
        \draw [->] (0, -5) -- (0, 5);
        \node [above left] at (5, 0) {$x$};
        \node [below right] at (0, 5) {$y$};
        \node [above left] at (5, -5) {$D_z$};
        \node [above right] at (1, 1) {$y = \frac{z}{x}$};
    \end{tikzpicture}
    &
    $z < 0$, тоді $D_z = 
    \left\{x<0, y>\frac{z}{x}\right\} \cup 
    \left\{x>0, y<\frac{z}{x}\right\}$.

    $F_\eta(z) = \int\limits_{-\infty}^0 dx \int\limits_{\frac{z}{x}}^{+\infty} 
    f_{\vec{\xi}}(x, y)dy + \int\limits_0^{+\infty}dx\int\limits_{-\infty}^{\frac{z}{x}} 
    f_{\vec{\xi}}(x, y) dy$
    
    Продиференціюємо обидві частини по $z$:

    $f_\eta(z) = - \int\limits_{-\infty}^0 \frac{1}{x} f_{\vec{\xi}}\left(x, \frac{z}{x}\right)dx 
    + \int\limits_0^{+\infty}\frac{1}{x}f_{\vec{\xi}}\left(x, \frac{z}{x}\right)dx$

\end{tabular}
\end{enumerate}

\vspace{0.5em}
\noindentОтже, отримали щільність добутку $\xi_1 \xi_2$:
\begin{equation}\label{eq:distr_product}
    f_{\xi_1 \xi_2}(z) = -\int\limits_{-\infty}^0 \frac{1}{x}f_{\vec{\xi}}\left(x, \frac{z}{x}\right)dx + 
    \int\limits_0^{+\infty}\frac{1}{x}f_{\vec{\xi}}\left(x, \frac{z}{x}\right)dx
\end{equation}

\begin{remark}
    Якщо $\xi_1$ та $\xi_2$ незалежні, то $f_{\vec{\xi}}\left(x, \frac{z}{x}\right) = 
    f_{\xi_1}(x)f_{\xi_2}\left(\frac{z}{x}\right)$.
\end{remark}

\begin{example}
    Нехай $\vec{\xi} = (\xi_1, \xi_2)^T$ рівномірно розподілений в квадраті $K = [0; 2]\times[0; 2]$.
    Знайти закон розподілу $\eta = \xi_1 \xi_2$.

    \begin{center}
        \begin{tabular}{c p{4.5cm}}
            \begin{tikzpicture}[baseline={(current bounding box.center)} ,scale = 1]
                \draw [black, ultra thick] 
                (0, 0) -- (2, 0) -- (2, 2) -- (0, 2) -- (0, 0);
                \fill [lightgray] 
                (0, 0) -- (2, 0) -- (2, 2) -- (0, 2) -- (0, 0);
                \draw [->] (-0.5, 0) -- (3, 0);
                \draw [->] (0, -0.5) -- (0, 3);
                \node [below left] at (3, 0) {$x$};
                \node [below left] at (0, 3) {$y$};
                \node [above right] at (0, 0) {$K$};
                \node [left] at (0, 2) {$2$};
                \node [below] at (2, 0) {$2$};
                \node [below left] at (0, 0) {$0$};
            \end{tikzpicture} 
            &
            $f_{\vec{\xi}}(x, y) = 
            \begin{cases}
                \frac{1}{4}, & (x, y) \in K \\
                0, & (x, y) \notin K
            \end{cases}
            $
        \end{tabular}
    \end{center}
    
    В цьому випадку простіше знайти функцію розподілу з геометричних міркувань, а не застосовувати формулу \eqref{eq:distr_product}:
    $F_{\eta}(z) = \P\{\xi_1\xi_2 < z\} = \iint\limits_{\{xy < z\}} f_{\vec{\xi}}(x, y)dxdy$. При $z<0$ $F_{\eta}(z) = 0$,
    а при $z \geq 4$ $F_{\eta}(z) = 1$.

        \begin{tabular}{c p{8.5 cm}}
            \begin{tikzpicture}[baseline={(current bounding box.north)}, scale = 0.9]
                \fill [lightgray, domain=0.5:2, smooth, variable = \x] plot ({\x}, 
                {
                    1/\x
                }) -- (2, 0) -- (0, 0) -- (0, 2) -- (0.5, 2);
                \draw [domain=0.333:3, smooth, variable = \x, thick] plot ({\x}, 
                {
                    1/\x
                });
                \draw [black, thick] 
                (0, 0) -- (2, 0) -- (2, 2) -- (0, 2) -- (0, 0);
                
                \draw [black, dashed] (0.5, 0) -- (0.5, 2);
                \draw [->] (-0.5, 0) -- (3, 0);
                \draw [->] (0, -0.5) -- (0, 3);
                \node [below left] at (3, 0) {$x$};
                \node [below left] at (0, 3) {$y$};
                \node [below] at (0.5, 0) {$\frac{z}{2}$};
                \node [left] at (0, 2) {$2$};
                \node [below] at (2, 0) {$2$};
                \node [below left] at (0, 0) {$0$};
                \node [above right] at (0.5, 2) {$y = \frac{z}{x}$};
                \node [above left] at (2, 0) {$D_z \cap K$};
            \end{tikzpicture} 
            &
            $F_{\eta}(z) = \frac{1}{4}S_{D_z \cap K} = \frac{1}{4} 
            \left(4 - \int\limits_{\frac{z}{2}}^2dx
            \int\limits_{\frac{z}{x}}^2dy\right) = $

            $=\frac{1}{4}\left(4 - \int\limits_{\frac{z}{2}}^2 
            \left(2 -\frac{z}{x}\right)dx\right) $

            \vspace{0.2em}
            $= 1 - \frac{1}{4}\left.(2x - z \ln x)\right|_{\frac{z}{2}}^2 = $
            
            \vspace{0.2em}
            $= 1 - \frac{1}{4}(4 - z\ln2 - z + z\ln\frac{z}{2}) =$
            
            \vspace{0.2em}
            $= \frac{1}{4}(z + z\ln2 - z\ln z + z\ln2) = $
            
            \vspace{0.2em}
            $= \frac{1}{4}(z + 2z\ln2 - z\ln z)$ 
            при $z\in [0; 4)$. \\
            Отже, $f_{\eta}(z) = 
            \begin{cases}
                0, & z \notin (0, 4] \\
                \frac{1}{4}(2\ln2 - \ln z), & z \in (0, 4]
            \end{cases}$.
        \end{tabular}
\end{example}

\subsection{Закон розподілу частки двох НВВ}

Нехай $\vec{\xi} = (\xi_1, \xi_2)^T$ --- неперервний випадковий вектор, щільність
$f_{\vec{\xi}}(x, y)$ відома.

\noindent\textbf{Задача}: знайти розподіл $\eta = \frac{\xi_2}{\xi_1}$.

Знайдемо функцію розподілу за вже відомою схемою:

$$F_\eta(z) = \iint\limits_{D_z}f_{\vec{\xi}}(x, y)dx dy, D_z = \left\{(x, y) \in 
\mathbb{R}^2 : \frac{y}{x} < z\right\}$$

\begin{enumerate}
    \item 
\begin{tabular}{c p{8.8cm}}
    \begin{tikzpicture}[baseline={(current bounding box.north)}, scale = 0.4]
        \fill [lightgray] (0, 5) -- (-5, 5) -- (-5, -5) -- (0, 0);
        \fill [lightgray] (5, -5) -- (0, -5) -- (0, 0) -- (5, 5);
        \draw [domain=-5:5, smooth, variable = \x, thick] plot ({\x}, 
        {
            \x
        });
        \draw [->] (-5, 0) -- (5, 0);
        \draw [->] (0, -5) -- (0, 5);
        \node [below left] at (5, 0) {$x$};
        \node [below left] at (0, 5) {$y$};
        \node [above left] at (5, -5) {$D_z$};
        \node [above left] at (3.7, 3) {$y = zx$};
    \end{tikzpicture} &
    $z > 0$, тоді $D_z = 
    \left\{x<0, y>z x\right\} \cup 
    \left\{x>0, y<z x\right\}$

    $F_\eta(z) = \int\limits_{-\infty}^0 dx \int\limits_{zx}^{+\infty}f_{\vec{\xi}}(x, y)dy 
    + \int\limits_0^{+\infty}dx\int\limits_{-\infty}^{zx}f_{\vec{\xi}}(x, y)dy$

    Продиференціюємо обидві частини по $z$:

    $f_\eta(z) = -\int\limits_{-\infty}^0 x f_{\vec{\xi}}(x, zx) dx + \int\limits_0^{+\infty}
    xf_{\vec{\xi}}(x, zx)dx$
\end{tabular}

\item 
\begin{tabular}{c p{8.8cm}}
    \begin{tikzpicture}[baseline={(current bounding box.north)} ,scale = 0.4]
        \fill [lightgray] (0, 0) -- (0, 5) -- (-5, 5);
        \fill [lightgray] (0, 0) -- (0, -5) -- (5, -5);
        \draw [domain=-5:5, smooth, variable = \x, thick] plot ({\x}, 
        {
            -\x
        });
        \draw [->] (-5, 0) -- (5, 0);
        \draw [->] (0, -5) -- (0, 5);
        \node [below left] at (5, 0) {$x$};
        \node [below left] at (0, 5) {$y$};
        \node [above left] at (4, -5) {$D_z$};
        \node [above right] at (2.6, -3) {$y = zx$};
    \end{tikzpicture} &
    $z < 0$, тоді $D_z = 
    \left\{x<0, y>z x\right\} \cup 
    \left\{x>0, y<z x\right\}$

    $F_\eta(z) = \int\limits_{-\infty}^0 dx \int\limits_{zx}^{+\infty}f_{\vec{\xi}}(x, y) dy +
    \int\limits_0^{+\infty}dx\int\limits_{-\infty}^{zx}f_{\vec{\xi}}(x, y)dy $

    Продиференціюємо обидві частини по $z$:

    $f_\eta(z) = -\int\limits_{-\infty}^0 x f_{\vec{\xi}}(x, zx)dx + \int\limits_0^{+\infty}
    xf_{\vec{\xi}}(x, zx)dx$
\end{tabular}
\end{enumerate}

\noindentОтже, отримали щільність частки $\frac{\xi_2}{\xi_1}$ 
\begin{equation}\label{eq:distr_frac}
    f_{\frac{\xi_2}{\xi_1}} (z)= -\int\limits_{-\infty}^0 xf_{\vec{\xi}}(x, zx)dx + 
    \int\limits_0^{+\infty}xf_{\vec{\xi}}(x, zx)dx    
\end{equation}

\begin{remark}
    Якщо $\xi_1$ та $\xi_2$ незалежні, то $f_{\vec{\xi}}(x, zx) = 
    f_{\xi_1}(x)f_{\xi_2}(zx)$.
\end{remark}

\begin{exercise}
    Перевірити, що щільність розподілу $\frac{\xi_1}{\xi_2}$ має вигляд
    \begin{equation}
        f_{\frac{\xi_1}{\xi_2}}(z) = -\int\limits_{-\infty}^0 y f_{\vec {\xi}}(zy, y) dy + \int\limits_0^{+\infty}
        yf_{\vec{\xi}}(zy, y)dy
    \end{equation}
\end{exercise}

\begin{example}
    Нехай $\xi_1$ та $\xi_2$ --- незалежні випадкові величини, що мають розподіл $\mathrm{N}(0, \sigma^2)$. Знайти закон розподілу $\eta = \frac{\xi_2}{\xi_1}$.

    \noindentСкористаємося знайденою формулою \eqref{eq:distr_frac}, враховуючи незалежність:
    \begin{gather*}
        f_\eta(z) = - \int\limits_{-\infty}^0 x f_{\xi_1}(x) f_{\xi_2}(zx) dx + 
        \int\limits_{0}^{+\infty} x f_{\xi_1}(x) f_{\xi_2}(zx) dx = 
        \\
        =
        -\frac{1}{2\pi\sigma^2}\int\limits_{-\infty}^0 x e^{-\frac{x^2}{2\sigma^2}\left(1+z^2\right)} dx
        +
        \frac{1}{2\pi\sigma^2}\int\limits_0^{+\infty} x e^{-\frac{x^2}{2\sigma^2}\left(1+z^2\right)}dx = 
        \\
        = \left[
        \begin{array}{c}
            \frac{x^2}{2\sigma^2}\left(1+z^2\right) = t \\
            \frac{x}{\sigma^2}\left(1+z^2\right)dx = dt 
        \end{array}
        \right]
        =
        \frac{-\sigma^2}{2\pi\sigma^2\left(1+z^2\right)}\int\limits_{+\infty}^0 e^{-t} dt
        +
        \frac{\sigma^2}{2\pi\sigma^2\left(1+z^2\right)}\int\limits^{+\infty}_0 e^{-t} dt = 
        \\
        = \frac{1}{\pi(1+z^2)}\int\limits_0^{+\infty}e^{-t}dt = \frac{1}{\pi(1+z^2)}, z \in \mathbb{R}
        \text{ --- це щільність розподілу Коші.} 
    \end{gather*}\index{розподіл!Коші}
\end{example}

\subsection{Закон розподілу суми двох НВВ}

Нехай $\vec{\xi} = (\xi_1, \xi_2)^T$ --- неперервний випадковий вектор, щільність
$f_{\vec{\xi}}(x, y)$ відома.

\noindent\textbf{Задача}: знайти розподіл $\eta = \xi_1 + \xi_2$.

Знайдемо функцію розподілу за вже відомою схемою:

$$F_\eta(z) = \iint\limits_{D_z}f_{\vec{\xi}}(x, y)dx dy, D_z = \left\{(x, y) \in 
\mathbb{R}^2 : x + y < z\right\}$$

\begin{tabular}{c p{8.8cm}}
    \begin{tikzpicture}[baseline={(current bounding box.north)} ,scale = 0.4]
        \fill [lightgray, domain=-5:5, smooth, variable = \x] 
        (-4, 5) -- (5, -4) -- (5, -5) -- (-5, -5) -- (-5, 5) -- (-4, 5);
        \draw [domain=-4:5, smooth, variable = \x, thick] plot ({\x}, 
        {
            1 - \x
        });
        \draw [->] (-5, 0) -- (5, 0);
        \draw [->] (0, -5) -- (0, 5);
        \node [below left] at (5, 0) {$x$};
        \node [below left] at (0, 5) {$y$};
        \node [above left] at (-1, 0) {$D_z$};
        \node [above right] at (0.5, 0) {$y = z - x$};
    \end{tikzpicture} &

    $F_\eta(z) = \int\limits_{-\infty}^{+\infty} dx \int\limits_{-\infty}^{z-x} 
    f_{\vec{\xi}}(x, y) dy$

    Продиференціюємо обидві частини по $z$:

    $f_\eta(z) = \int\limits_{-\infty}^{+\infty} f_{\vec{\xi}}(x, z-x) dx$
\end{tabular}

\vspace{0.5em}
\noindentОтже, отримали щільність суми $\xi_1 + \xi_2$:
\begin{equation}\label{eq:distr_sum}
    f_{\xi_1 + \xi_2}(z) = \int\limits_{-\infty}^{+\infty} f_{\vec{\xi}}(x, z-x) dx
\end{equation}

\begin{remark}
    Якщо $\xi_1$ та $\xi_2$ незалежні, то $f_{\xi_1 + \xi_2} (z) = 
    \int\limits_{-\infty}^{+\infty}f_{\xi_1}(x) f_{\xi_2}(z-x) dx = 
    (f_{\xi_1} \ast f_{\xi_2})(z)$. Тут $\ast$ позначає операцію 
    \href{https://uk.wikipedia.org/wiki/%D0%97%D0%B3%D0%BE%D1%80%D1%82%D0%BA%D0%B0_(%D0%BC%D0%B0%D1%82%D0%B5%D0%BC%D0%B0%D1%82%D0%B8%D1%87%D0%BD%D0%B8%D0%B9_%D0%B0%D0%BD%D0%B0%D0%BB%D1%96%D0%B7)}{згортки},
    що вводиться в курсі гармонічного аналізу.
\end{remark}

\begin{exercise}
    Перевірити, що щільність розподілу $\eta_1 = \xi_1 - \xi_2$ має вигляд
    $f_{\eta_1}(z) = \int\limits_{-\infty}^{+\infty} f_{\vec{\xi}}(x, x-z) dx$, а 
    $\eta_2 = \xi_2 - \xi_1$ --- $f_{\eta_2}(z) = \int\limits_{-\infty}^{+\infty} f_{\vec{\xi}}(x, x+z) dx$.
\end{exercise}

\begin{example}
    $\xi_1$ та $\xi_2$ --- незалежні випадкові величини, $\xi_i \sim \mathrm{Exp}(\lambda_i)$, 
    $i = 1,2$, $\lambda_1 \neq \lambda_2$.

    \noindentЗнайти закон розподілу $\eta = \xi_1 + \xi_2$.

    \noindentСкористаємося знайденою формулою \eqref{eq:distr_sum}, враховуючи незалежність:
    \begin{equation*}
        f_\eta(z) = \left[z - x > 0 \Rightarrow x < z\right] = 
        \int\limits_0^z \lambda_1 e^{-\lambda_1 x}\lambda_2e^{-\lambda_2(z-x)}dx
        =
        \lambda_1\lambda_2e^{-\lambda_2z}\int\limits_0^z e^{-x(\lambda_1 - \lambda_2)} dx=
    \end{equation*}
    \begin{equation*}
        =\lambda_1\lambda_2e^{-\lambda_2z}
        \left.
        \left(\frac{1}{-(\lambda_1 - \lambda_2)} e^{-x(\lambda_1 - \lambda_2)}\right)
        \right|_0^z = 
        \frac{\lambda_1\lambda_2}{\lambda_2 - \lambda_1}e^{-\lambda_2z}
        (e^{-z(\lambda_1 - \lambda_2)} - 1) =
    \end{equation*}
    \begin{equation*}
        = \frac{\lambda_1\lambda_2}{\lambda_2 - \lambda_1}(e^{-\lambda_1 z} - e^{-\lambda_2 z})
        ,\; z \geq 0 \text{ та } 0 \text{ інакше.}
    \end{equation*}
Це щільність так званого \emph{закону Ерланга 2-го порядку} \index{розподіл!Ерланга}.
\noindentПри $\lambda_1 = \lambda_2 = \lambda$ отримаємо
$f_{\eta}(z) = \lambda^2 z e^{-\lambda z}$ при $z \geq 0$ та $0$ інакше.
Це вже буде щільність гамма-розподілу $\Gamma(2, \lambda)$.
\end{example}

\subsection{Числові характеристики функції багатьох випадкових аргументів}
Розглядаємо неперервний випадковий вектор $\vec{\xi} = (\xi_1, \xi_2, ..., \xi_n)$ та
випадкову величину $\eta = \varphi(\vec{\xi})$.
\begin{exercise}
    Довести, що для НВВ $\eta$, що приймає невід'ємні значення, $\E\eta = \int\limits_0^{+\infty} (1 - F_{\eta}(x)) dx$.
\end{exercise}

\begin{proposition*} 
    $\E\eta^k = \int_{\mathbb{R}^n} \varphi^k(\vec{x}) f_{\vec{\xi}} (\vec{x}) d\vec{x}$ для всіх $k \in \mathbb{N}$.
\end{proposition*}
\begin{proof}
    Доведемо у випадку $n=2$, та $\varphi(x, y) \geq 0$.
    \begin{gather*}
        \E\eta^k = \int\limits_0^{+\infty} \left(1 - F_{\eta^k}(z)\right) dz = 
    \int\limits_0^{+\infty} \left({\iint_{\varphi^k(x,y) \geq z}} f_{\vec{\xi}}(x,y) dx dy\right) dz = \\
    =\iint\limits_{\mathbb{R}^2} \left( \int\limits_0^{\varphi^k(x,y)} dz\right) f_{\vec{\xi}}(x,y) dx dy = 
     \iint\limits_{\mathbb{R}^2} \varphi^k(x,y) f_{\vec{\xi}}(x,y) dx dy
    \end{gather*}
\end{proof}

\begin{exercise}
    Довести твердження в більш загальному вигляді.
\end{exercise}

\begin{example}\label{proof:expectation}
    $\xi_1, \xi_2$ --- НВВ, знайдемо $\E(\xi_1 + \xi_2)$ та $\E\xi_1\xi_2$
    у випадку незалежності цих НВВ.
    $$E(\xi_1 + \xi_2) = \iint\limits_{\mathbb{R}^2} (x+y) f_{\vec{\xi}}(x,y) dx dy = 
    \iint\limits_{\mathbb{R}^2} x f_{\vec{\xi}}(x,y) dx dy + 
    \iint\limits_{\mathbb{R}^2} y f_{\vec{\xi}}(x,y) dx dy = \E\xi_1 + \E\xi_2$$
    Якщо $\xi_1, \xi_2$ незалежні, то $f_{\vec{\xi}}(x,y) = f_{\xi_1}(x) f_{\xi_2}(y)$.
    $$\E\xi_1 \xi_2 =
    \iint\limits_{\mathbb{R}^2} xy f_{\xi_1}(x) f_{\xi_2}(y) dx dy =
    \left( \int\limits_{\mathbb{R}} x f_{\xi_1}(x) dx \right)
    \left( \int\limits_{\mathbb{R}} y f_{\xi_2}(y) dy \right) = \E\xi_1 E\xi_2$$
\end{example}
    \chapter{Основні розподіли математичної статистики}
        % !TEX root = ../main.tex

У цьому розділі буде наведено виведення законів розподілу, що застосовуються
в задачах математичної статистики, та їх числових характеристик.

\section{Розподіл \texorpdfstring{$\chi^2$}{x2} (Пірсона)}
\noindent\textbf{Означення:}
    нехай $\xi_k \sim \mathrm{N}(a_k, \sigma_k)$, $k= 1,..., n$ --- незалежні у сукупності.
    Тоді $\xi = \sum\limits_{k=1}^n \left( \frac{\xi_k - a_k}{\sigma_k}\right)^2$ має
    розподіл \emph{$\chi^2$ (хі-квадрат, Пірсона) з $n$ ступенями вільності}.

    \noindent$\mathring{\xi}_{k} = \frac{\xi_k - a_k}{\sigma_k} \sim \mathrm{N}(0, 1)$, тому можна ще записати
    $\xi = \sum\limits_{k=1}^n (\mathring{\xi}_{k})^2$.

\noindent\textbf{Коротке позначення:} $\xi \sim \chi_n^2$, $n\in\mathbb{N}$ --- кількість ступенів вільності.

\noindent\textbf{Щільність розподілу:}
відомо, що якщо $\eta \sim \mathrm{N}(0, 1)$, то $\eta^2 \sim \Gamma\left(\frac{1}{2}, \frac{1}{2}\right)$.
Гамма-розподіл стійкий при $\beta_1 = \beta_2 = ... = \beta_n$, $\mathring{\xi}_{k}$ незалежні у сукупності,
тому $\xi = \sum\limits_{k=1}^n (\mathring{\xi}_{k})^2 \sim \Gamma\left(\frac{n}{2}, \frac{1}{2}\right) = \chi_n^2$.
\begin{equation*}
    f_{\chi_n^2}(x) = \begin{cases}
        \frac{1}{2^{\frac{n}{2}} \Gamma\left(\frac{n}{2}\right)} x^{\frac{n}{2}-1} e^{-\frac{x}{2}}, & x \geq 0 \\
        0, & x < 0
    \end{cases}
\end{equation*}

\noindent \textbf{Крива розподілу:} графіки для різних значень $n$.
\begin{center}
    \begin{tikzpicture}[yscale = 12, xscale = 0.5, baseline={(current bounding box.center)}]
        \pgfmathsetmacro{\a}{2};
        \pgfmathsetmacro{\b}{0.5};
        \pgfmathsetmacro{\c}{3};
        \pgfmathsetmacro{\d}{4};
    
        \draw [->] (-2, 0) -- (18.3, 0);
        \draw [->] (0, -0.05) -- (0, 0.2);
        \draw [thick] (-2, 0) -- (0, 0);
        \draw [domain=0:18, smooth, variable = \x, thick] plot ({\x}, {\b^(\a)*\x^(\a-1)/factorial(\a-1) * e^(-\x*\b)});
        \draw [domain=0:18, smooth, variable = \x, thick] plot ({\x}, {\b^(\c)*\x^(\c-1)/factorial(\c-1) * e^(-\x*\b)});
        \draw [domain=0:18, smooth, variable = \x, thick] plot ({\x}, {\b^(\d)*\x^(\d-1)/factorial(\d-1) * e^(-\x*\b)});
        \node [below] at (18.2, 0) {$x$};
        \node [left] at (0, 0.2) {$f_{\chi_n^2}(x)$};
        \node [below left] at (0, 0) {$0$};
    \end{tikzpicture}
\end{center}

\noindent\textbf{Числові характеристики:}
\begin{enumerate}
    \item $\E\chi_n^2 = \frac{n/2}{1/2} = n$.
    \item $\D\chi_n^2 = \frac{n/2}{1/4} = 2n$.
\end{enumerate}

\section{Розподіл \texorpdfstring{$\chi$}{x}}
\noindent\textbf{Означення:} нехай випадкова величина $\xi$ має розподіл $\chi_n^2$.
Тоді $\eta = \sqrt{\xi}$ має розподіл
\emph{$\chi$ (хі) з $n$ ступенями вільності}.

\noindent\textbf{Коротке позначення:} $\eta \sim \chi_n$, $n\in\mathbb{N}$ --- кількість ступенів вільності.

\noindent\textbf{Щільність розподілу:} скористаємося формулою для визначення щільності розподілу функції від
випадкової величини. $\eta = \sqrt{\xi}$, тому позначимо $\varphi(x) = \sqrt{x}$, $\varphi^{-1}(y) = y^2$,
$\left( \varphi^{-1}(y) \right)^{\prime} = 2y$.
\begin{equation*}
    f_{\chi_n}(y) = f_{\chi_n^2}\left(\varphi^{-1} (y)\right) \cdot \left|\left(\varphi^{-1} (y) \right)^{\prime}\right| =
f_{\chi_n^2}(y^2) \cdot 2y = 
\begin{cases}
    \frac{1}{2^{\frac{n}{2} - 1} \Gamma\left(\frac{n}{2}\right)} y^{n-1} e^{-\frac{y^2}{2}}, & y \geq 0 \\
    0, & y < 0
\end{cases}
\end{equation*}
\noindent\textbf{Числові характеристики:}
\begin{enumerate}
    \item $\E\chi_n = \sqrt{2} \cdot \frac{\Gamma\left(\frac{n+1}{2}\right)}{\Gamma\left(\frac{n}{2}\right)}$.
    \item $\D\chi_n = n - \left( \E\chi_n \right)^2$.
\end{enumerate}

\begin{remark}
    Нескладно помітити, що $\chi_2$ --- це розподіл Релея.
\end{remark} 

\noindent Знайдемо ще розподіл $\zeta = \frac{\chi_n}{\sqrt{n}}$. $\varphi(y) = \frac{y}{\sqrt{n}}$,
$\varphi^{-1}(z) = z \sqrt{n}$, 
$\left( \varphi^{-1}(z) \right)^{\prime} = \sqrt{n}$.

\begin{equation*}
    f_{\frac{\chi_n}{\sqrt{n}}}(z) = f_{\chi_n}\left(\varphi^{-1} (z)\right) \cdot \left|\left(\varphi^{-1} (y) \right)^{\prime}\right| =
f_{\chi_n}(z \sqrt{n}) \cdot \sqrt{n} = 
\begin{cases}
    \frac{n^{\frac{n}{2}}}{2^{\frac{n}{2} - 1} \Gamma\left(\frac{n}{2}\right)} z^{n-1} e^{-\frac{nz^2}{2}}, & z \geq 0 \\
    0, & z < 0
\end{cases}
\end{equation*}
\begin{exercise}
    Записати числові характеристики випадкової величини, що має розподіл $\frac{\chi_n}{\sqrt{n}}$.
\end{exercise}

\section{Розподіл Стьюдента (\texorpdfstring{$t$}{t}-розподіл)}
\noindent\textbf{Означення:} якщо $\xi \sim \mathrm{N}(0, 1)$ та $\eta \sim \frac{\chi_n}{\sqrt{n}}$ незалежні,
то $\zeta = \frac{\xi}{\eta} = \frac{\xi}{\chi_n /\sqrt{n}}$ має \emph{розподіл Стьюдента з $n$ ступенями вільності}.

\noindent\textbf{Коротке позначення:} $\zeta \sim \mathrm{St}_n$, $n\in\mathbb{N}$ --- кількість ступенів вільності.

\noindent\textbf{Щільність розподілу:} скористаємося формулою для визначення щільності розподілу частки
двох незалежних НВВ.
\begin{gather*}
    f_{\mathrm{St}_n} (z) = \int\limits_0^{+\infty} x f_{\xi}(z x) f_{\frac{\chi_n}{\sqrt{n}}} (x) dx =
    \frac{1}{\sqrt{2\pi}} \cdot \frac{n^{\frac{n}{2}}}{2^{\frac{n}{2} - 1} \Gamma\left(\frac{n}{2}\right)}
    \int\limits_0^{+\infty} x e^{-\frac{z^2 x^2}{2}} x^{n-1} e^{-\frac{nx^2}{2}} dx = \\
    = \frac{n^{\frac{n}{2}}}{\sqrt{\pi} 2^{\frac{n-1}{2}} \Gamma\left(\frac{n}{2}\right)}
    \int\limits_0^{+\infty} x^n e^{-\frac{x^2}{2}(z^2+n)} dx = 
    \left[ \frac{x^2}{2} (z^2 + n) = t, x = \frac{\sqrt{2} \sqrt{t}}{\sqrt{z^2 + n}},
    dx = \frac{1}{\sqrt{2}\sqrt{z^2 + n}} \frac{dt}{\sqrt{t}}\right] = \\
    = \frac{n^{\frac{n}{2}}}{\sqrt{\pi} 2^{\frac{n-1}{2}} \Gamma\left(\frac{n}{2}\right)} \cdot
    \frac{2^{\frac{n}{2}}}{(z^2 + n)^{\frac{n}{2}}} \frac{1}{2^{\frac{1}{2}}(z^2 + n)^{\frac{1}{2}}}
    \int\limits_0^{+\infty} t^{\frac{n-1}{2}} e^{-t} dt = 
    \frac{n^{\frac{n}{2}} \Gamma\left(\frac{n+1}{2}\right)}{\sqrt{\pi} \Gamma\left(\frac{n}{2}\right)} \cdot
    \frac{1}{(z^2 + n)^{\frac{n+1}{2}}}, \; z \in \mathbb{R}
\end{gather*}
\noindent \textbf{Крива розподілу:} графіки для різних значень $n$, називаються \emph{кривими Стьюдента}.
Вони схожі на криву гауссівського розподілу, але повільніше прямують до 0 на нескінченності.
\begin{center}
    \begin{tikzpicture}[yscale = 6, xscale = 1.3, baseline={(current bounding box.center)}]
        \pgfmathsetmacro{\a}{1}; % n = 2
        \pgfmathsetmacro{\b}{12}; % n = 4
        \pgfmathsetmacro{\c}{0.318309886184}; % n = 1
    
        \draw [->] (-5, 0) -- (5, 0);
        \draw [->] (0, -0.05) -- (0, 0.4);
        \draw [domain=-5:5, smooth, variable = \x, thick] plot ({\x}, {\a/((2+(\x)^2)^((2+1)/2))});
        \draw [domain=-5:5, smooth, variable = \x, thick] plot ({\x}, {\b/((4+(\x)^2)^((4+1)/2))});
        \draw [domain=-5:5, smooth, variable = \x, thick] plot ({\x}, {\c/((1+(\x)^2)^((1+1)/2))});
        \node [below] at (5.2, 0) {$x$};
        \node [left] at (0, 0.4) {$f_{\mathrm{St}_n}(x)$};
        \node [below left] at (0, 0) {$0$};
    \end{tikzpicture}
\end{center}

\noindent\textbf{Числові характеристики:}
\begin{enumerate}
    \item $\E\mathrm{St}_n = 0$.
    \item $\D\mathrm{St}_n = \frac{n}{n-2}$, $n>2$.
\end{enumerate}

\begin{remark}
    Нескладно помітити, що $\mathrm{St}_1$ --- це розподіл Коші.
\end{remark}

\section{Розподіл Фішера-Снедекора (\texorpdfstring{$F$}{F}-розподіл)}
\noindent\textbf{Означення:} випадкова величина $\eta = \frac{\chi_{n_1}^2/n_1}{\chi_{n_2}^2/n_2}$, чисельник
та знаменник якої незалежні, має \emph{розподіл Фішера-Снедекора з $n_1$, $n_2$ ступенями вільності}.

\noindent\textbf{Коротке позначення:} $\eta \sim \mathrm{F}(n_1, n_2)$, $n_1, n_2\in\mathbb{N}$ --- кількість ступенів вільності.

\noindent\textbf{Щільність розподілу:} скористаємося формулою для визначення щільності розподілу частки
двох незалежних НВВ. Нагадаємо, що 
$f_{\chi_n^2}(x) = \begin{cases}
    \frac{1}{2^{\frac{n}{2}} \Gamma\left(\frac{n}{2}\right)} x^{\frac{n}{2}-1} e^{-\frac{x}{2}}, & x \geq 0 \\
    0, & x < 0
\end{cases}$.

\noindentТоді $f_{\chi_n^2/n}(y) = f_{\chi_n^2}(n y) \cdot n = 
\begin{cases}
    \frac{n^{\frac{n}{2}}}{2^{\frac{n}{2}} \Gamma\left(\frac{n}{2}\right)} y^{\frac{n}{2}-1} e^{-\frac{ny}{2}}, & y \geq 0 \\
    0, & y < 0
\end{cases}$.

\begin{gather*}
    f_{\mathrm{F}(n_1, n_2)} (z) = \int\limits_0^{+\infty} x f_{\chi_{n_1}^2/n_1}(zx) f_{\chi_{n_2}^2/n_2}(x) dx = \\
    = \frac{n_1^{\frac{n_1}{2}} n_2^{\frac{n_2}{2}}}{2^{\frac{n_1+n_2}{2}} \Gamma\left(\frac{n_1}{2}\right) \Gamma\left(\frac{n_2}{2}\right)}
    \int\limits_0^{+\infty} x z^{\frac{n_1}{2} - 1} x^{\frac{n_1}{2} - 1} e^{-\frac{n_1 zx}{2}} x^{\frac{n_2}{2} - 1} e^{-\frac{n_2 x}{2}} dx = \\
    = \frac{n_1^{\frac{n_1}{2}} n_2^{\frac{n_2}{2}}}{2^{\frac{n_1+n_2}{2}} \Gamma\left(\frac{n_1}{2}\right) \Gamma\left(\frac{n_2}{2}\right)} \cdot
    z^{\frac{n_1}{2} - 1} \int\limits_0^{+\infty} x^{\frac{n_1 + n_2}{2} - 1} e^{-\frac{x}{2}(n_1 z + n_2)} dx =
    \left[ \frac{x}{2}(n_1 z + n_2) = t, x = \frac{2t}{n_1 z + n_2}\right] = \\
    = \frac{n_1^{\frac{n_1}{2}} n_2^{\frac{n_2}{2}}}{2^{\frac{n_1+n_2}{2}} \Gamma\left(\frac{n_1}{2}\right) \Gamma\left(\frac{n_2}{2}\right)} \cdot
    z^{\frac{n_1}{2} - 1} \cdot 2^{\frac{n_1+n_2}{2}} \cdot \frac{1}{(n_1 z + n_2)^{\frac{n_1+n_2}{2}}}
    \int\limits_0^{+\infty} t^{\frac{n_1+n_2}{2} - 1} e^{-t} dt = \\
    = n_1^{\frac{n_1}{2}} n_2^{\frac{n_2}{2}} \frac{\Gamma\left(\frac{n_1+n_2}{2}\right)}{\Gamma\left(\frac{n_1}{2}\right) \Gamma\left(\frac{n_2}{2}\right)} \cdot
    \frac{z^{\frac{n_1}{2} - 1}}{(n_1 z + n_2)^{\frac{n_1+n_2}{2}}}, \; z \geq 0 \text{ та } 0 \text{ інакше}.
\end{gather*}

\noindent \textbf{Крива розподілу:} графіки для різних значень $n_1$, $n_2$, називаються \emph{кривими Фішера}.

\begin{center}
    \begin{tikzpicture}[yscale = 6, xscale = 3, baseline={(current bounding box.center)}]
        \pgfmathsetmacro{\a}{3.30797337253}; % n1 = 3, n2 = 1
        \pgfmathsetmacro{\b}{64}; % n1 = 4, n2 = 2
        \pgfmathsetmacro{\c}{68.7549354157}; % n = 1
    
        \draw [->] (-0.5, 0) -- (4, 0);
        \draw [->] (0, -0.05) -- (0, 0.7);
        \draw [domain=0:4, smooth, variable = \x, thick, samples = 200] plot ({\x}, {\a*((\x)^(3/2 - 1))/(3*\x + 1)^((3+1)/2)});
        \draw [domain=0:4, smooth, variable = \x, thick, samples = 400] plot ({\x}, {\b*((\x)^(4/2 - 1))/(4*\x + 2)^((4+2)/2)});
        \draw [domain=0:4, smooth, variable = \x, thick, samples = 200] plot ({\x}, {\c*((\x)^(3/2 - 1))/(3*\x + 3)^((3+3)/2)});
        \node [below] at (4, 0) {$x$};
        \draw [thick] (-0.5, 0) -- (0, 0);
        \node [left] at (0, 0.7) {$f_{\mathrm{F}(n_1, n_2)}(x)$};
        \node [below left] at (0, 0) {$0$};
    \end{tikzpicture} 
\end{center}

\noindent\textbf{Числові характеристики:}
\begin{enumerate}
    \item $\E\mathrm{F}(n_1, n_2) = \frac{n_2}{n_2 - 2}$, $n_2 > 2$.
    \item $\D\mathrm{F}(n_1, n_2) = \frac{2 n_2^2 (n_1 + n_2 - 2)}{n_1 (n_2 - 2)^2 (n_2 -4)}$, $n_2>4$.
\end{enumerate}

\begin{remark}
    Якщо $\eta \sim \mathrm{F}(n_1, n_2)$, то $\frac{1}{\eta} \sim \mathrm{F}(n_2, n_1)$.
\end{remark}
    %\chapter{Граничні теореми теорії ймовірності}
        %% !TEX root = ../main.tex
\section{Послідовності випадкових величин}
\subsection{Нерівності Маркова та Чебишова}
\begin{theorem*}[нерівність Маркова]
    Нехай модуль випадкової величини $\xi$ має скінченне математичне сподівання: $E|\xi| < +\infty$.
    Тоді 
    \begin{gather}\label{Markov_ineq}
        \forall \; \varepsilon >0 : P\left\{ |\xi| \geq \varepsilon\right\} \leq \frac{E|\xi|}{\varepsilon}
    \end{gather}
\end{theorem*}
\begin{proof}
    Запишемо випадкову величину $|\xi|$ через події-індикатори: 
    $|\xi| = |\xi|\cdot I\left\{|\xi| \geq \varepsilon\right\} + |\xi|\cdot I\left\{|\xi| < \varepsilon\right\} \geq
    \varepsilon\cdot I\left\{|\xi| \geq \varepsilon\right\}$. Звідси $E|\xi| \geq E \left( \varepsilon\cdot I\left\{|\xi| \geq \varepsilon\right\}\right) =
    \varepsilon \cdot P\left\{ |\xi| \geq \varepsilon\right\}$.
\end{proof}
\begin{remark}
    Еквівалентною нерівністю є $P\left\{ |\xi| < \varepsilon\right\} \geq 1 - \frac{E|\xi|}{\varepsilon}$.
\end{remark}

\begin{example}
    Багатьма спостереженнями з'ясовано, що середня кількість
    сонячних дні у Києві складає 220. Оцінити ймовірність того, що
    сонячних днів за рік буде не менше 300.

    \noindentПозначимо $\xi$ кількість сонячних днів. За умовою $\xi$ невід'ємна та $E\xi = 220$,
    тому за нерівністю Маркова $P\left\{ \xi \geq 300\right\} \leq \frac{220}{300} = \frac{11}{15}$.
\end{example}

\begin{theorem*}[нерівність Чебишова]
    Нехай випадкова величина $\xi$ має скінченні математичне сподівання та дисперсію.
Тоді
\begin{gather}\label{Cheb_ineq}
    \forall \; \varepsilon >0 : P\left\{ \left|\xi - E\xi\right| \geq \varepsilon\right\} \leq \frac{D\xi}{\varepsilon^2}
\end{gather}
\end{theorem*} 
\begin{proof}
    $P\left\{ |\xi - E\xi| \geq \varepsilon\right\} = P\left\{ (\xi - E\xi)^2 \geq \varepsilon^2\right\}$.
    Застосуємо нерівність Маркова:
    
    \noindent $P\left\{ (\xi - E\xi)^2 \geq \varepsilon^2\right\} \leq \frac{E(\xi - E\xi)^2}{\varepsilon^2} = \frac{D\xi}{\varepsilon^2}$.
\end{proof}
\begin{remark}
    Еквівалентною нерівністю є $P\left\{ \left|\xi - E\xi\right| < \varepsilon\right\} \geq 1 - \frac{D\xi}{\varepsilon^2}$.
\end{remark}
\begin{example}
    Отримаємо <<правило $3 \sigma$>> для довільної випадкової величини зі скінченними математичним сподівання та дисперсією.
    $P\left\{ |\xi - E\xi| < 3 \sigma\right\} \geq 1 - \frac{D\xi}{9 \sigma^2} = 1 - \frac{\sigma^2}{9 \sigma^2} = \frac{8}{9}$.
\end{example}
Розглянемо застосування \emph{нерівності Чебишова в схемі Бернуллі}. Нехай $\xi \sim \mathrm{Bin}(n, p)$, $E\xi = np$, $D\xi = npq$.
Відношення $\frac{\xi}{n}$ називається відносною частотою появи успіху або частістю. $E\left( \frac{\xi}{n}\right) = p$, 
$D\left( \frac{\xi}{n}\right) = \frac{pq}{n}$. З нерівності Чебишова 
$P\left\{ \left|\frac{\xi}{n} - p\right| \geq \varepsilon\right\} \leq \frac{pq}{n \epsilon^2}$.

\subsection{Послідовності випадкових величин}
Розглядаємо фіксований ймовірнісний простір $\left\{ \Omega, \mathcal{F}, P\right\}$ та
послідовність випадкових величин $\left\{ \xi_n (\omega)\right\}_{n=1}^{\infty}$.
Якщо для кожного $n \in \mathbb{N}$ події $\xi_1, \xi_2, ..., \xi_n$ \emph{незалежні у сукупності},
то послідовність $\left\{ \xi_n (\omega)\right\}_{n=1}^{\infty}$ називається \emph{послідовністю незалежних випадкових величин}.

Послідовності випадкових величин можна задавати різними способами:
\begin{enumerate}
    \item Нехай $\xi$ --- деяка випадкова величина, можна задати $\xi_n = f_n(\xi)$, де $f_n$ --- деяка числова функція. 
    Наприклад: $\xi_n = \xi^n$, $\xi_n = \cos (n\xi)$.
    \item $n$ може входити як параметр закону розподілу $\xi_n$. Наприклад, $\xi_n \sim \mathrm{Exp}(n)$, $\xi_n \sim \mathrm{N}(0, \frac{1}{n})$.
    \item Для послідовностей ДВВ $n$ може входити як в значення, що приймає $\xi_n$, так і у відповідні ймовірності. Наприклад:
    \begin{tabular}{|c|c|c|c|}
        \hline
        $\xi_n$ & $-\sqrt{n}$ & $0$ & $\sqrt{n}$ \\
        \hline
        $p$ & $1/n$ & $1 - 2/n$ & $1/n$ \\
        \hline
    \end{tabular}
\end{enumerate}
В курсі функціонального аналізу вводяться різні види збіжності послідовності вимірних функцій та
зв'язок між цими видами збіжності.

\subsection{Види збіжності послідовності випадкових величин}
Нагадаємо класичне означення границі числової послідовності. Число $a$ називають границею послідовності 
$\left\{ a_n\right\}_{n=1}^{\infty}$, якщо
$$ \forall \; \varepsilon > 0 \; \exists \; N(\varepsilon) \in \mathbb{N}: \forall \; n \geq N(\varepsilon) \; \left| a_n - a\right| < \varepsilon$$
Оскільки послідовність випадкових величин $\left\{ \xi_n\right\}_{n=1}^{\infty}$ є послідовністю функцій з $\Omega$ в $\mathbb{R}$, то це означення
не є застосовним, бо $\left| \xi_n - \xi\right| < \varepsilon$ є випадковою подією, що виконується, взагалі кажучи, не для всіх елементарних подій
$\omega \in \Omega$. ших. Тому ми маємо ввести інше означення границі (та збіжності) послідовності
випадкових величин. Виявляється, що таких означень можна запропонувати декілька, причому вони
не є еквівалентними одне одному.
\vspace{1em}
\begin{enumerate}
    \item \textbf{Збіжність майже напевно (сильна збіжність, збіжність з ймовірністю 1).}
    \noindent$\left\{ \xi_n\right\}_{n=1}^{\infty}$ збігається до $\xi$ \emph{майже напевно}, якщо
    $P\left\{ \omega: \underset{n\to\infty}{\lim} \xi_n(\omega) = \xi(\omega)\right\} = 1$. 
    Це еквівалентно умові $P\left\{ \omega: \underset{n\to\infty}{\lim} \xi_n(\omega) \neq \xi(\omega)\right\} = 0$.

    Позначення: $\xi_n \overset{P1}{\longrightarrow} \xi, n \to \infty$ або $\xi_n \overset{\text{м.н.}}{\longrightarrow} \xi, n \to \infty$.
    \begin{exercise}
        Довести ще одне еквівалентне визначення збіжності майже напевно: $\xi_n \overset{P1}{\longrightarrow} \xi, n \to \infty$, якщо
        $\forall \; \varepsilon > 0: P\left( \bigcap\limits_{n=1}^{\infty} \bigcup\limits_{k = n}^{\infty}
        \left\{ \omega : \left| \xi_k(\omega) - \xi(\omega)\right| > \varepsilon\right\}\right) = 0$.
    \end{exercise}

    Це найбільш природний з інтуїтивної точки зору вид збіжності. Зауважимо, однак, що він має доволі дивні властивості. Наприклад, можна навести приклад
    послідовності $\xi_n$, що не збігається до деякої $\xi$, але будь-яка її підпослідовність $\xi_{n_k}$ містить свою підпослідовність
    $\xi_{n_{k_l}}$, яка все ж таки збігається до $\xi$.
    
    \textbf{Лема Бореля-Кантеллі.} Якщо для послідовності подій $\left\{A_n \right\}_{n=1}^{\infty}$ ряд
    $\sum\limits_{n=1}^{\infty} P(A_n)$ збігається, то $P\left( \bigcap\limits_{n=1}^{\infty} \bigcup\limits_{k = n}^{\infty} A_k\right) = 0$.
    Це означає, що ймовірність того, що відбудеться нескінченна кількість цих подій, є нульовою.
    \begin{proof}
        Послідовність подій $B_n = \bigcup\limits_{k = n}^{\infty} A_k$ монотонно спадна, тому за теоремою неперервності \ref{th:2}
        $P\left( \bigcap\limits_{n=1}^{\infty} B_n\right) = \underset{n\to\infty}{\lim} P(B_n)$.
        $P(B_n) = P\left( \bigcup\limits_{k = n}^{\infty} A_k\right) \leq \sum\limits_{k=n}^{\infty}P(A_n) \to 0$ при $n\to \infty$ 
        зі збіжності ряду. Тому $\underset{n\to\infty}{\lim} P(B_n) = P\left( \bigcap\limits_{n=1}^{\infty} \bigcup\limits_{k = n}^{\infty} A_k\right) = 0$.
    \end{proof}
    Застосуванням цієї леми до послідовності подій $A_n = \left\{ \omega : \left| \xi_n(\omega) - \xi(\omega)\right| > \varepsilon\right\}$ отримаємо зручну для використання
    ознаку збіжності майже напевно: якщо для будь-якого $\varepsilon >0$ ряд
    $\sum\limits_{n=1}^{\infty} P\left\{\left| \xi_n - \xi\right| > \varepsilon\right\}$ збігається, то 
    $\xi_n \overset{P1}{\longrightarrow} \xi, n \to \infty$.
    \begin{example}
        Довести, що послідовність $\xi_n \sim \mathrm{U}(-\frac{1}{n}, \frac{1}{n})$ при $n\to\infty$ збігається до 0 майже напевно.
        
        Візьмемо довільне $\varepsilon > 0$ та знайдемо $P\left\{ |\xi_n| > \varepsilon\right\}$. Для $\varepsilon \geq 1$
        ця ймовірність, очевидно, рівна 0. В іншому випадку, для кожного $\varepsilon \in (0; 1)$ можна знайти такий номер $N$,
        для якого $\varepsilon$ буде більше за $\frac{1}{n}$ при $n \geq N$. Тому для будь-якого $\varepsilon >0$ ймовірності
        $P\left\{ |\xi_n| > \varepsilon\right\}$ рівні 0, починаючи з якогось $n$. Отже, для будь-якого $\varepsilon >0$ ряд 
        $\sum\limits_{n=1}^{\infty} P\left\{\left| \xi_n \right| > \varepsilon\right\}$ збігається і 
        $\xi_n \overset{P1}{\longrightarrow} 0, n \to \infty$.
    \end{example}
    \item \textbf{Збіжність за ймовірністю.}
    \noindent$\left\{ \xi_n\right\}_{n=1}^{\infty}$ збігається до $\xi$ \emph{за ймовірністю}, якщо 
    $\forall \; \varepsilon > 0: \underset{n \to \infty}{\lim} P\left\{|\xi_n - \xi| \geq \varepsilon\right\}= 0$.
    Це еквівалентно умові $\forall \; \varepsilon > 0: \underset{n \to \infty}{\lim} P\left\{|\xi_n - \xi| \leq \varepsilon\right\}= 1$.
    Позначення: $\xi_n \overset{P}{\longrightarrow} \xi, n \to \infty$.
    \begin{example}
        Нехай $\xi_n$ --- послідовність ДВВ: 
        \begin{tabular}{|c|c|c|}
            \hline
            $\xi_n$ & $0$ & $n^7$ \\
            \hline
            $p$ & $1-1/n$ & $1/n$ \\
            \hline
        \end{tabular}.
        Перевірити збіжність $\xi_n \overset{P}{\longrightarrow} 0, n \to \infty$.
        
        \noindentДля $\varepsilon >0$ $P\left\{|\xi_n -0| \geq \varepsilon\right\} = P\left\{ \xi_n \geq \varepsilon\right\}$.
        $\forall \; \varepsilon >0 \; \exists \; N: \forall n\geq N:n^7 > \varepsilon$, тому з якогось номера
        $P\left\{ \xi_n \geq \varepsilon\right\} = P\left\{ \xi_n = n^7 \right\} = \frac{1}{n} \to 0, n\to\infty$, тому $\xi_n \overset{P}{\longrightarrow} 0, n \to \infty$.
    \end{example}
    \item \textbf{Збіжність в середньому.}
    \noindent$\left\{ \xi_n\right\}_{n=1}^{\infty}$ збігається до $\xi$ \emph{в середньому},
    якщо $\underset{n \to \infty}{\lim} E|\xi_n - \xi| = 0$.
    Позначення: $\xi_n \overset{\text{С}}{\longrightarrow} \xi, n \to \infty$.
    \item \textbf{Збіжність в середньому квадратичному.}
    \noindent$\left\{ \xi_n\right\}_{n=1}^{\infty}$ збігається до $\xi$ \emph{в середньому квадратичному},
    якщо $\underset{n \to \infty}{\lim} E(\xi_n - \xi)^2 = 0$.
    
    \noindentПозначення: $\xi_n \overset{\text{СК}}{\longrightarrow} \xi, n \to \infty$.
    \item \textbf{Збіжність за розподілом (слабка збіжність, збіжність в основному).}

    \noindent$\left\{ \xi_n\right\}_{n=1}^{\infty}$ збігається до $\xi$ \emph{за розподілом}, якщо функціональна послідовність
    $\left\{ F_{\xi_n} (x)\right\}_{n=1}^{\infty}$ збігається до $F_{\xi}(x)$ в точках її неперервності.

    Позначення: $\xi_n \overset{F}{\longrightarrow} \xi, n \to \infty$.

    Ця збіжність за характером відрізняється від інших тим, що не враховує залежність або незалежність $\xi_n$.
    Є також \emph{еквівалентне означення збіжності} за розподілом: якщо для будь-якої обмеженої неперервної функції $\varphi : \mathbb{R} \to \mathbb{R}$
    $\underset{n \to \infty}{\lim} E \varphi(\xi_n) = E\varphi(\xi)$.
    Доведення еквівалентності цих двох означень виходить за рамки курсу: його ідея полягає в тому, що $F_{\xi}(x) = P\left\{ \xi < x\right\} = E 1_{(-\infty, x)}(\xi)$,
    де $1_{(-\infty, x)}$ --- індикатор множини $(-\infty, x)$, і такі функції-індикатори можна наблизити з будь-якою заданою точністю неперервними функціями та навпаки.
    \begin{example}
        Нехай $\xi_n \sim \mathrm{Exp}(\frac{1}{n})$. Перевірити $\xi_n \overset{F}{\to} 0, n \to \infty$.
        Тут під 0 розуміється ДВВ, що приймає значення 0 з ймовірністю 1.

        $F_{\xi_n}(x) = \begin{cases}
            1 - e^{\frac{x}{n}}, & x > 0 \\
            0, & x \leq 0
        \end{cases}$. Видно, що при $n \to \infty$ $F_{\xi_n}(x) \to F_0(x) = \begin{cases}
            1, & x > 0 \\
            0, & x \leq 0
        \end{cases}
        $ --- функція розподілу 0.
    \end{example}
\end{enumerate}

\begin{remark}
    Збіжності в середньому та середньому квадратичному є частковими випадками \emph{збіжності порядку $k$}, для якої
$\underset{n \to \infty}{\lim} E|\xi_n - \xi|^k = 0$. Оскільки 
для $0 < s < t$ має місце
$( E \left| \xi\right|^s)^{\frac{1}{s}} \leq ( E \left| \xi\right|^t)^{\frac{1}{t}}$
то збіжність порядку $k$ гарантує збіжність порядків менше $k$.
\end{remark}
\begin{proposition*}
     Нехай послідовність $\left\{ \xi_n\right\}_{n=1}^{\infty}$ збігається за ймовірністю до двох випадкових
    величин $\xi$ та $\eta$. Тоді $P\left\{ \xi = \eta\right\} = 1$.
\end{proposition*}
\begin{proof}
    Для будь-якого $\varepsilon > 0$ : $\left\{ |\xi - \eta| > \varepsilon\right\} \subset 
    \left\{ |\xi - \xi_n| > \frac{\varepsilon}{2}\right\} \cup \left\{ |\eta - \xi_n| > \frac{\varepsilon}{2}\right\}$, тому \\
    $P\left\{ |\xi - \eta| > \varepsilon\right\} \leq P\left\{ |\xi - \xi_n| > \frac{\varepsilon}{2}\right\} + 
    P\left\{ |\eta - \xi_n| > \frac{\varepsilon}{2}\right\} \to 0, n\to\infty$.
    Отже, $P\left\{ |\xi - \eta| > \varepsilon\right\} = 0$, і за довільністю $\varepsilon > 0$ маємо $P\left\{ \xi \neq \eta\right\} = 0$,
    або ж $P\left\{ \xi = \eta\right\} = 1$.
\end{proof}
\begin{exercise}
    Довести, що це твердження виконується для збіжностей з ймовірністю 1, в середньому та середньому квадратичному, але не виконується
    для збіжності за розподілом.
\end{exercise}
Можна довести, що усім збіжностям, крім збіжності за розподілом, притаманні відомі арифметичні властивості: збіжність суми $\xi_n + \eta_n$ та добутку
$\xi_n \eta_n$ до $\xi + \eta$ та $\xi \eta$ відповідно за умови збіжностей $\xi_n$ до $\xi$ та $\eta_n$ до $\eta$. Також, $\varphi(\xi_n)$ збігається
до $\varphi(\xi)$ за умови неперервності $\varphi$.

В курсі функціонального аналізу встановлюється зв'язок між розглянутими видами збіжності:
\vspace{5 pt}
\begin{center}
\begin{tikzpicture}
\label{scheme}
\node[block] (a) {Збіжність з ймовірністю 1 \\ ($P1$)};
\node[block, above right = 0cm and 2cm of a] (b) {Збіжність в середньому\\ квадратичному ($\text{СК}$)};
\node[block, below right = 0cm and 2cm of a]   (c) {Збіжність в середньому \\ ($\text{С}$)};
\node[block, below right = 2cm and -1.5cm of a]   (d) {Збіжність за ймовірністю \\ ($P$)};
\node[block, below = 1cm of d]   (e){Збіжність за розподілом \\ ($F$)};
\node[below left = 1cm and -2.5 cm of a]   (1){\circled 1};
\node[below right = 0.25cm and -1.6 cm of d]   (2){\circled 2};

\draw[line] (a.south) -- ([xshift=-0.5cm]d.north);
\draw[line] (b.south) -- (c.north);
\draw[line] (c.south) -- ([xshift=0.5cm]d.north);
\draw[line] ([xshift=-0.5cm]d.south) -- ([xshift=-0.5cm]e.north);
\draw[line,dashed] ([xshift=-1.5cm]d.north) -- ([xshift=-1cm]a.south);
\draw[line,dashdotted] ([xshift=0.5cm]e.north) -- ([xshift=0.5cm]d.south);
\end{tikzpicture}
\end{center}

Штрихова лінія \circled 1 означає, що у кожній послідовності $\left\{ \xi_n\right\}_{n=1}^{\infty}$, що збігається за ймовірністю, 
міститься підпослідовність $\left\{ \xi_{n_k}\right\}_{k=1}^{\infty}$, що збігається з імовірністю 1. 
Штрих-пунктирна лінія \circled 2 показує, що відповідний перехід справедливий, коли гранична випадкова величина $\xi$ є константою.
Доведемо деякі з цих тверджень.

\begin{proposition*} 
    Зі збіжності в середньому квадратичному випливає збіжність в середньому ($\text{СК} \Rightarrow \text{С}$).
\end{proposition*}
\begin{proof}
    $E|\xi_n -\xi| \leq \sqrt{E(\xi_n - \xi)^2}$, як було сказано вище, тому з $E(\xi_n - \xi)^2 \to 0$
    випливає $E|\xi_n - \xi| \to 0, n\to\infty$.
\end{proof}
\begin{proposition*}
    Зі збіжності в середньому випливає збіжність за ймовірністю.
\end{proposition*}
\begin{proof}
    З нерівності Маркова (\ref{Markov_ineq})
    $\forall \; \varepsilon > 0: P\left\{|\xi_n - \xi| \geq \varepsilon\right\} \leq \frac{E|\xi_n - \xi|}{\varepsilon}$, 
    тому з $E|\xi_n - \xi| \to 0$ маємо $P\left\{|\xi_n - \xi| \geq \varepsilon\right\} \to 0, n\to\infty$.
\end{proof}
\begin{proposition*}
     Зі збіжності за розподілом \emph{до константи} випливає збіжність за ймовірністю.
\end{proposition*}
\begin{proof}
    Нехай $\xi_n \overset{F}{\longrightarrow} c$, $c \in \mathbb{R}$. Візьмемо $\varepsilon > 0$.
    $P\left\{ |\xi_n - c| \leq \varepsilon\right\} = P\left\{ c - \varepsilon \leq \xi_n \leq c + \varepsilon \right\} \geq
    P\left\{ c - \varepsilon \leq \xi_n < c + \varepsilon \right\} = F_{\xi_n}(c+\varepsilon) - F_{\xi_n}(c-\varepsilon)$, тому
    $\underset{n \to \infty}{\lim} P\left\{ |\xi_n - c| \leq \varepsilon\right\} \geq \underset{n \to \infty}{\lim}  F_{\xi_n}(c+\varepsilon) -
    \underset{n \to \infty}{\lim}  F_{\xi_n}(c-\varepsilon) = F_c(c+\varepsilon) - F_c(c-\varepsilon) = 1$, звідки $\xi_n \overset{P}{\longrightarrow} c$.
\end{proof}
\begin{example}
    Нехай $\xi_n$ --- послідовність ДВВ: 
        \begin{tabular}{|c|c|c|c|}
            \hline
            $\xi_n$ & $-n$ & $0$ & $2n$ \\
            \hline
            $p$ & $1/{2n}$ & $1 - 1/n$ & $1/{2n}$ \\
            \hline
        \end{tabular}.

    Зі збільшенням $n$ $\xi_n$ все з більшою ймовірністю набувають значення 1. 
    Водночас, інші два можливі значення ($-n$ та $2n$) розбігаються до нескінченності. 
    Ці два процеси спричиняють протилежні ефекти --- перший наближає $\xi_n$ до 1, а другий --- віддаляє. 
    Тому з точки зору одних видів збіжності, більш чутливих до <<аномальних викидів>>, послідовність $\xi_n$ не буде мати границі, 
    а з точки зору інших, менш чутливих, границя дорівнюватиме 1.
    Дійсно: 
    \begin{gather*}
        E\left|\xi_n -1 \right| = \frac{1}{2n} |-n-1|  + \left( 1 - \frac{1}{n}\right)|1-1| + \frac{1}{2n}|2n-1| \to \frac{3}{2} \\
        E\left(\xi_n -1 \right)^2 = \frac{1}{2n} (-n-1)^2  + \left( 1 - \frac{1}{n}\right)(1-1)^2 + \frac{1}{2n}(2n-1)^2 \to \infty
    \end{gather*}
    Отже, збіжності до 1 ні в середньому, ні в середньому квадратичному немає. Зауважимо, що якби <<викиди>> прямували до нескінченності повільніше або відповідні
    ймовірності прямували до нуля швидше, то ці збіжності могли б бути: наприклад, якщо замість $n$ та $-2n$ $\xi_n$ набували значення $\sqrt{n}$ та $2\sqrt{n}$
    з ймовірностями $\frac{1}{2n^2}$, то послідовність збігалася б до 1 і в середньому, і в середньому квадратичному.
    З іншого боку, для будь-якого $\varepsilon>0$ значення $-n$ та $2n$ рано чи пізно вийдуть за межі відрізку $[1-\varepsilon;1+\varepsilon]$, і тому
    $\xi_n \overset{P}{\longrightarrow} 1$. Тепер стає зрозуміло, чому перевіряли збіжність в середньому та середньому квадратичному лише до 1: якби
    якась з цих границь існувала, то вона була б і границею за ймовірністю.
\end{example}
\begin{remark}Для збіжності $\xi_n \overset{\text{СК}}{\longrightarrow} C$ ($C$ --- стала) необхідно і достатньо, щоб
$\underset{n \to \infty}{\lim} E\xi_n = C$ та $\underset{n \to \infty}{\lim} D\xi_n = 0$.
$\underset{n \to \infty}{\lim} E\xi_n = C \Leftrightarrow \underset{n \to \infty}{\lim} (E\xi_n - EC) = 0$,
$E(\xi_n - C)^2 = D\xi_n + (E(\xi_n -C))^2$,
звідки отримуємо твердження при $n \to \infty$, оскільки $D\xi_n \geq 0$ та $(E(\xi_n -C))^2 \geq 0$.
Зауважимо, що ці дві умови є достатніми для збіжності $\xi_n \overset{P}{\longrightarrow} C$.
\end{remark}
\begin{example}
    Нехай $\xi_n \sim \mathrm{N}\left(5 + \frac{1}{n}, \sigma = \frac{1}{n}\right)$. Перевірити $\xi_n \overset{P}{\longrightarrow} 5$.
    
    \noindent Оскільки $E\xi_n = 5+\frac{1}{n} \to 5$, $D\xi_n = \frac{1}{n^2} \to 0$ при $n \to \infty$,
    то $\xi_n \overset{\text{СК}}{\longrightarrow} 5 \Rightarrow \xi_n \overset{P}{\longrightarrow} 5$.
\end{example}

Наведемо без доведення важливу теорему, що стосується збіжності за розподілом.

\begin{theorem*}[теорема неперевності Леві]
    Збіжність за розподілом послідовності випадкових величин $\left\{ \xi_n\right\}_{n=1}^{\infty}$ еквівалентна поточковій збіжності їх
    характеристичних функцій:
    $$ \xi_n \overset{F}{\longrightarrow} \xi \Longleftrightarrow \chi_{\xi_n}(t) \to \chi_\xi (t) \; \forall \; t\in \mathbb{R}$$
\end{theorem*}
    \chapter*{Таблиця значень функції Лапласа}
    \addcontentsline{toc}{chapter}{Таблиця значень функції Лапласа} 
        % !TEX root = ../main.tex

$$\Phi(x) = \frac{1}{\sqrt{2 \pi}} \int\limits_0^x e^{-\frac{t^2}{2}} dt$$
\noindent \begin{tabular}{|c|c|c|c|c|c|c|c|c|c|c|}
\hline
$\pmb{x}$ & \textbf{ 0.00 } & \textbf{ 0.01 } & \textbf{ 0.02 } & \textbf{ 0.03 } & \textbf{ 0.04 } & \textbf{ 0.05 } & \textbf{ 0.06 } & \textbf{ 0.07 } & \textbf{ 0.08 } & \textbf{ 0.09 } \\ 
 \hline 
\textbf{ 0.00 } & 0.00000 & 0.00399 & 0.00798 & 0.01197 & 0.01595 & 0.01994 & 0.02392 & 0.02790 & 0.03188 & 0.03586\\ 
\hline
\textbf{ 0.10 } & 0.03983 & 0.04380 & 0.04776 & 0.05172 & 0.05567 & 0.05962 & 0.06356 & 0.06749 & 0.07142 & 0.07535\\ 
\hline
\textbf{ 0.20 } & 0.07926 & 0.08317 & 0.08706 & 0.09095 & 0.09483 & 0.09871 & 0.10257 & 0.10642 & 0.11026 & 0.11409\\ 
\hline
\textbf{ 0.30 } & 0.11791 & 0.12172 & 0.12552 & 0.12930 & 0.13307 & 0.13683 & 0.14058 & 0.14431 & 0.14803 & 0.15173\\ 
\hline
\textbf{ 0.40 } & 0.15542 & 0.15910 & 0.16276 & 0.16640 & 0.17003 & 0.17364 & 0.17724 & 0.18082 & 0.18439 & 0.18793\\ 
\hline
\textbf{ 0.50 } & 0.19146 & 0.19497 & 0.19847 & 0.20194 & 0.20540 & 0.20884 & 0.21226 & 0.21566 & 0.21904 & 0.22240\\ 
\hline
\textbf{ 0.60 } & 0.22575 & 0.22907 & 0.23237 & 0.23565 & 0.23891 & 0.24215 & 0.24537 & 0.24857 & 0.25175 & 0.25490\\ 
\hline
\textbf{ 0.70 } & 0.25804 & 0.26115 & 0.26424 & 0.26730 & 0.27035 & 0.27337 & 0.27637 & 0.27935 & 0.28230 & 0.28524\\ 
\hline
\textbf{ 0.80 } & 0.28814 & 0.29103 & 0.29389 & 0.29673 & 0.29955 & 0.30234 & 0.30511 & 0.30785 & 0.31057 & 0.31327\\ 
\hline
\textbf{ 0.90 } & 0.31594 & 0.31859 & 0.32121 & 0.32381 & 0.32639 & 0.32894 & 0.33147 & 0.33398 & 0.33646 & 0.33891\\ 
\hline
\textbf{ 1.00 } & 0.34134 & 0.34375 & 0.34614 & 0.34849 & 0.35083 & 0.35314 & 0.35543 & 0.35769 & 0.35993 & 0.36214\\ 
\hline
\textbf{ 1.10 } & 0.36433 & 0.36650 & 0.36864 & 0.37076 & 0.37286 & 0.37493 & 0.37698 & 0.37900 & 0.38100 & 0.38298\\ 
\hline
\textbf{ 1.20 } & 0.38493 & 0.38686 & 0.38877 & 0.39065 & 0.39251 & 0.39435 & 0.39617 & 0.39796 & 0.39973 & 0.40147\\ 
\hline
\textbf{ 1.30 } & 0.40320 & 0.40490 & 0.40658 & 0.40824 & 0.40988 & 0.41149 & 0.41309 & 0.41466 & 0.41621 & 0.41774\\ 
\hline
\textbf{ 1.40 } & 0.41924 & 0.42073 & 0.42220 & 0.42364 & 0.42507 & 0.42647 & 0.42785 & 0.42922 & 0.43056 & 0.43189\\ 
\hline
\textbf{ 1.50 } & 0.43319 & 0.43448 & 0.43574 & 0.43699 & 0.43822 & 0.43943 & 0.44062 & 0.44179 & 0.44295 & 0.44408\\ 
\hline
\textbf{ 1.60 } & 0.44520 & 0.44630 & 0.44738 & 0.44845 & 0.44950 & 0.45053 & 0.45154 & 0.45254 & 0.45352 & 0.45449\\ 
\hline
\textbf{ 1.70 } & 0.45543 & 0.45637 & 0.45728 & 0.45818 & 0.45907 & 0.45994 & 0.46080 & 0.46164 & 0.46246 & 0.46327\\ 
\hline
\textbf{ 1.80 } & 0.46407 & 0.46485 & 0.46562 & 0.46638 & 0.46712 & 0.46784 & 0.46856 & 0.46926 & 0.46995 & 0.47062\\ 
\hline
\textbf{ 1.90 } & 0.47128 & 0.47193 & 0.47257 & 0.47320 & 0.47381 & 0.47441 & 0.47500 & 0.47558 & 0.47615 & 0.47670\\ 
\hline
\textbf{ 2.00 } & 0.47725 & 0.47778 & 0.47831 & 0.47882 & 0.47932 & 0.47982 & 0.48030 & 0.48077 & 0.48124 & 0.48169\\ 
\hline
\textbf{ 2.10 } & 0.48214 & 0.48257 & 0.48300 & 0.48341 & 0.48382 & 0.48422 & 0.48461 & 0.48500 & 0.48537 & 0.48574\\ 
\hline
\textbf{ 2.20 } & 0.48610 & 0.48645 & 0.48679 & 0.48713 & 0.48745 & 0.48778 & 0.48809 & 0.48840 & 0.48870 & 0.48899\\ 
\hline
\textbf{ 2.30 } & 0.48928 & 0.48956 & 0.48983 & 0.49010 & 0.49036 & 0.49061 & 0.49086 & 0.49111 & 0.49134 & 0.49158\\ 
\hline
\textbf{ 2.40 } & 0.49180 & 0.49202 & 0.49224 & 0.49245 & 0.49266 & 0.49286 & 0.49305 & 0.49324 & 0.49343 & 0.49361\\ 
\hline
\textbf{ 2.50 } & 0.49379 & 0.49396 & 0.49413 & 0.49430 & 0.49446 & 0.49461 & 0.49477 & 0.49492 & 0.49506 & 0.49520\\ 
\hline
\textbf{ 2.60 } & 0.49534 & 0.49547 & 0.49560 & 0.49573 & 0.49585 & 0.49598 & 0.49609 & 0.49621 & 0.49632 & 0.49643\\ 
\hline
\textbf{ 2.70 } & 0.49653 & 0.49664 & 0.49674 & 0.49683 & 0.49693 & 0.49702 & 0.49711 & 0.49720 & 0.49728 & 0.49736\\ 
\hline
\textbf{ 2.80 } & 0.49744 & 0.49752 & 0.49760 & 0.49767 & 0.49774 & 0.49781 & 0.49788 & 0.49795 & 0.49801 & 0.49807\\ 
\hline
\textbf{ 2.90 } & 0.49813 & 0.49819 & 0.49825 & 0.49831 & 0.49836 & 0.49841 & 0.49846 & 0.49851 & 0.49856 & 0.49861\\ 
\hline
\end{tabular}
\newpage
\thispagestyle{empty}
\noindent \begin{tabular}{|c|c|c|c|c|c|c|c|c|c|c|}
\hline
$\pmb{x}$ & \textbf{ 0.00 } & \textbf{ 0.01 } & \textbf{ 0.02 } & \textbf{ 0.03 } & \textbf{ 0.04 } & \textbf{ 0.05 } & \textbf{ 0.06 } & \textbf{ 0.07 } & \textbf{ 0.08 } & \textbf{ 0.09 } \\ 
 \hline 
\textbf{ 3.00 } & 0.49865 & 0.49869 & 0.49874 & 0.49878 & 0.49882 & 0.49886 & 0.49889 & 0.49893 & 0.49896 & 0.49900\\ 
\hline
\textbf{ 3.10 } & 0.49903 & 0.49906 & 0.49910 & 0.49913 & 0.49916 & 0.49918 & 0.49921 & 0.49924 & 0.49926 & 0.49929\\ 
\hline
\textbf{ 3.20 } & 0.49931 & 0.49934 & 0.49936 & 0.49938 & 0.49940 & 0.49942 & 0.49944 & 0.49946 & 0.49948 & 0.49950\\ 
\hline
\textbf{ 3.30 } & 0.49952 & 0.49953 & 0.49955 & 0.49957 & 0.49958 & 0.49960 & 0.49961 & 0.49962 & 0.49964 & 0.49965\\ 
\hline
\textbf{ 3.40 } & 0.49966 & 0.49968 & 0.49969 & 0.49970 & 0.49971 & 0.49972 & 0.49973 & 0.49974 & 0.49975 & 0.49976\\ 
\hline
\textbf{ 3.50 } & 0.49977 & 0.49978 & 0.49978 & 0.49979 & 0.49980 & 0.49981 & 0.49981 & 0.49982 & 0.49983 & 0.49983\\ 
\hline
\textbf{ 3.60 } & 0.49984 & 0.49985 & 0.49985 & 0.49986 & 0.49986 & 0.49987 & 0.49987 & 0.49988 & 0.49988 & 0.49989\\ 
\hline
\textbf{ 3.70 } & 0.49989 & 0.49990 & 0.49990 & 0.49990 & 0.49991 & 0.49991 & 0.49992 & 0.49992 & 0.49992 & 0.49992\\ 
\hline
\textbf{ 3.80 } & 0.49993 & 0.49993 & 0.49993 & 0.49994 & 0.49994 & 0.49994 & 0.49994 & 0.49995 & 0.49995 & 0.49995\\ 
\hline
\textbf{ 3.90 } & 0.49995 & 0.49995 & 0.49996 & 0.49996 & 0.49996 & 0.49996 & 0.49996 & 0.49996 & 0.49997 & 0.49997\\ 
\hline
\textbf{ 4.00 } & 0.49997 & 0.49997 & 0.49997 & 0.49997 & 0.49997 & 0.49997 & 0.49998 & 0.49998 & 0.49998 & 0.49998\\ 
\hline
\textbf{ 4.10 } & 0.49998 & 0.49998 & 0.49998 & 0.49998 & 0.49998 & 0.49998 & 0.49998 & 0.49998 & 0.49999 & 0.49999\\ 
\hline
\textbf{ 4.20 } & 0.49999 & 0.49999 & 0.49999 & 0.49999 & 0.49999 & 0.49999 & 0.49999 & 0.49999 & 0.49999 & 0.49999\\ 
\hline
\textbf{ 4.30 } & 0.49999 & 0.49999 & 0.49999 & 0.49999 & 0.49999 & 0.49999 & 0.49999 & 0.49999 & 0.49999 & 0.49999\\ 
\hline
\textbf{ 4.40 } & 0.49999 & 0.49999 & 0.50000 & 0.50000 & 0.50000 & 0.50000 & 0.50000 & 0.50000 & 0.50000 & 0.50000\\ 
\hline
\textbf{ 4.50 } & 0.50000 & 0.50000 & 0.50000 & 0.50000 & 0.50000 & 0.50000 & 0.50000 & 0.50000 & 0.50000 & 0.50000\\ 
\hline
\textbf{ 4.60 } & 0.50000 & 0.50000 & 0.50000 & 0.50000 & 0.50000 & 0.50000 & 0.50000 & 0.50000 & 0.50000 & 0.50000\\ 
\hline
\textbf{ 4.70 } & 0.50000 & 0.50000 & 0.50000 & 0.50000 & 0.50000 & 0.50000 & 0.50000 & 0.50000 & 0.50000 & 0.50000\\ 
\hline
\textbf{ 4.80 } & 0.50000 & 0.50000 & 0.50000 & 0.50000 & 0.50000 & 0.50000 & 0.50000 & 0.50000 & 0.50000 & 0.50000\\ 
\hline
\textbf{ 4.90 } & 0.50000 & 0.50000 & 0.50000 & 0.50000 & 0.50000 & 0.50000 & 0.50000 & 0.50000 & 0.50000 & 0.50000\\ 
\hline
\end{tabular} 


\end{document}
